\setcounter{section}{10}
\section{Propiedades topológicas de los morfismos planos de presentación finita. Criterios de planitud}
\label{section:IV.11}

\oldpage[IV-3]{116}

Mientras que en la \S~2 consideramos enunciados relativos a la planitud que no dependen
de ninguna hipótesis de finitud, y la \S~6 estudia la noción de planitud en el
marco de los preesquemas localmente noetherianos (pero sin imponer una hipótesis de finitud
sobre los morfismos), la presente sección está dedicada a la noción de
$f$-planitud cuando el morfismo $f:X\to Y$ es \emph{localmente de presentación
finita}. El interés de la noción de morfismo plano de presentación finita
radica en que parece expresar técnicamente, de la manera más adecuada, la noción
intuitiva de una ``familia de preesquemas algebraicos parametrizada por un esquema
$Y$'', cuyo estudio constituye uno de los objetos principales de la
Geometría Algebraica. Además, aun cuando inicialmente solo interesase el caso
de un esquema base noetheriano, resulta indispensable, por ciertas razones
técnicas (por ejemplo, para algunas aplicaciones de la teoría del
``descenso'', que lleva a introducir esquemas no necesariamente
noetherianos), no limitarse a ese caso al tratar problemas esencialmente
relativos ligados a morfismos localmente de presentación finita. Seguiremos
sistemáticamente este principio, ya respaldado por los resultados de las \S\S~8 y 9,
en todo el resto de este Capítulo e incluso en todo el resto de este Tratado,
aunque ello pueda sacrificar ocasionalmente la sencillez de ciertas demostraciones
que las hipótesis noetherianas permiten a veces aligerar.\footnote{Este principio
se inspira asimismo en la necesidad de dar cabida, como ``espacios de parámetros''
para familias de esquemas algebraicos, a espacios anillados arbitrarios (e incluso
a ``topos'' anillados arbitrarios), para los cuales ya no cabe considerar en general
hipótesis noetherianas. Parece bastante claro que esta nueva ampliación de la
Geometría Algebraica no podrá eludirse por mucho más tiempo, y conviene desarrollar
desde ahora las nociones y técnicas de naturaleza ``relativa'' de la teoría de
esquemas de modo que puedan adaptarse, prácticamente sin cambios, a este marco más
general.} En la presente sección, esto nos lleva a retomar, en el contexto de la
``presentación finita'' (en particular en el n$^\circ$~3), ciertos enunciados sobre
planitud ya obtenidos en el contexto noetheriano. La herramienta técnica esencial
para la reducción al caso noetheriano es el teorema de compatibilidad de la
planitud con los límites proyectivos de preesquemas \sref{IV.11.2.6}, que completa
los resultados generales de la \S~8. Demostramos también de paso en \sref{IV.11.3.1}
un resultado utilizado con frecuencia

\oldpage[IV-3]{117}
más adelante, según el cual el conjunto de puntos de planitud de un morfismo
localmente de presentación finita es abierto.

En los n$^\circ$s~4 a 8 estudiamos la cuestión del ``descenso'' de la
planitud, que consiste en hallar condiciones útiles sobre un morfismo de cambio
de base $Y'\to Y$ (no necesariamente plano) bajo las cuales pueda concluirse,
si $X\times_YY'$ es plano sobre $Y'$, que $X$ es plano sobre $Y$. Estos resultados,
más técnicos que los de los n$^\circ$s~1 a 3, se utilizan con menor frecuencia
en lo que sigue; desempeñan, no obstante, un papel importante en las técnicas de
construcción no proyectivas del capítulo siguiente. El único resultado de los
n$^\circ$s~4 a 8 utilizado posteriormente en el Capítulo IV es el criterio valuativo
de planitud \sref{IV.11.8}, que se aplicará en \sref{IV.15.2}.

Por último, los n$^\circ$s~9 y 10 están dedicados al estudio de una noción que
refina, en la teoría de esquemas, la densidad en sentido topológico: la noción de
una familia de subpreesquemas \emph{esquemáticamente densa} en un preesquema dado,
y en particular al estudio del comportamiento de esta noción bajo cambio de base
(plano o arbitrario). Por el momento, esta noción se utiliza sobre todo en el
estudio de los esquemas en grupos.


\subsection{Lugares de planitud (caso noetheriano)}
\label{subsection:IV.11.1}

\begin{theorem}[11.1.1]
\label{IV.11.1.1}
Sea $\mathrm{Y}$ un preesquema localmente noetheriano, $f:\mathrm{X}\to\mathrm{Y}$ un morfismo localmente de tipo finito y $\sh{F}$ un $\sh{O}_\mathrm{X}$-Módulo coherente. Entonces el conjunto $\mathrm{U}$ de los $x\in\mathrm{X}$ tales que $\sh{F}$ es $f$-plano en el punto $x$ es abierto en $\mathrm{X}$.
\end{theorem}

\begin{proof}
Como la cuestión es evidentemente local sobre $\mathrm{X}$ y $\mathrm{Y}$, podemos suponer $\mathrm{Y}=\Spec(\mathrm{A})$, $\mathrm{X}=\Spec(\mathrm{B})$, con $\mathrm{A}$ noetheriano y $\mathrm{B}$ una A-álgebra de tipo finito. Se tiene entonces $\sh{F}=\widetilde{\mathrm{M}}$, donde $\mathrm{M}$ es un B-módulo de tipo finito. Aplicando el criterio $(\mathbf{0}_{\mathrm{III}},~9.2.6)$, basta por tanto demostrar la afirmación siguiente:

\medskip
\phantomsection\label{IV.11.1.1.1}
\noindent\textbf{(11.1.1.1)}
\emph{Sean $\mathrm{A}$ un anillo noetheriano, $\mathrm{B}$ una $\mathrm{A}$-álgebra de tipo finito, $\mathrm{M}$ un $\mathrm{B}$-módulo de tipo finito, $\mathfrak{q}$ un ideal primo de $\mathrm{B}$ y $\mathfrak{p}$ su imagen inversa en $\mathrm{A}$. Supongamos que $\mathrm{M}_\mathfrak{q}$ es un $\mathrm{A}_\mathfrak{p}$-módulo plano. Existe entonces $g\in\mathrm{B}-\mathfrak{q}$ tal que, para todo ideal primo $\mathfrak{q}'\supset\mathfrak{q}$ de $\mathrm{B}$ con $g\notin\mathfrak{q}'$, $\mathrm{M}_{\mathfrak{q}'}$ es un $\mathrm{A}_{\mathfrak{p}'}$-módulo plano, donde $\mathfrak{p}'$ denota la imagen inversa de $\mathfrak{q}'$ en $\mathrm{A}$ (equivale a lo mismo $(\mathbf{0}_\mathrm{I},~6.3.1)$ decir que $\mathrm{M}_{\mathfrak{q}'}$ es un $\mathrm{A}$-módulo plano).}
\medskip

Para ello, consideremos $\mathrm{B}_{\mathfrak{q}'}$ como una $\mathrm{A}$-álgebra; evidentemente
$\mathfrak{p}\mathrm{B}_{\mathfrak{q}'}\subset
\mathfrak{q}'\mathrm{B}_{\mathfrak{q}'}$. Por
$(\mathbf{0}_{\mathrm{III}},~10.2.2)$, el $\mathrm{A}$-módulo
$\mathrm{M}_{\mathfrak{q}'}$ es plano si y solo si
$\mathrm{M}_{\mathfrak{q}'}/\mathfrak{p}\mathrm{M}_{\mathfrak{q}'}$ es un
$(\mathrm{A}/\mathfrak{p})$-módulo plano y
\[
  \operatorname{Tor}_1^\mathrm{A}
  (\mathrm{M}_{\mathfrak{q}'},\mathrm{A}/\mathfrak{p})=0.
\]
Ahora bien, $\mathrm{M}_{\mathfrak{q}'}=
\mathrm{M}\otimes_{\mathrm{B}}\mathrm{B}_{\mathfrak{q}'}$, y
$\mathrm{B}_{\mathfrak{q}'}$ es plano sobre $\mathrm{B}$; por tanto
\[
  \mathrm{M}_{\mathfrak{q}'}/\mathfrak{p}\mathrm{M}_{\mathfrak{q}'}
    =(\mathrm{M}/\mathfrak{p}\mathrm{M})_{\mathfrak{q}'},
  \qquad
  \operatorname{Tor}_1^\mathrm{A}
  (\mathrm{M}_{\mathfrak{q}'},\mathrm{A}/\mathfrak{p})
    =\bigl(\operatorname{Tor}_1^\mathrm{A}
      (\mathrm{M},\mathrm{A}/\mathfrak{p})\bigr)_{\mathfrak{q}'}.
\]
La segunda igualdad se obtiene definiendo Tor a partir de una resolución
proyectiva de $\mathrm{A}/\mathfrak{p}$. Por la misma razón, puesto que exigimos
$g\notin\mathfrak{q}'$, el anillo $\mathrm{B}_{\mathfrak{q}'}$ es plano sobre
$\mathrm{B}_g$, y por consiguiente
\[
  (\mathrm{M}/\mathfrak{p}\mathrm{M})_{\mathfrak{q}'}
    =\bigl((\mathrm{M}/\mathfrak{p}\mathrm{M})_g
      \bigr)_{\mathfrak{q}'\mathrm{B}_g},
  \qquad
  \bigl(\operatorname{Tor}_1^\mathrm{A}
      (\mathrm{M},\mathrm{A}/\mathfrak{p})\bigr)_{\mathfrak{q}'}
    =\bigl(\bigl(\operatorname{Tor}_1^\mathrm{A}
      (\mathrm{M},\mathrm{A}/\mathfrak{p})\bigr)_g
      \bigr)_{\mathfrak{q}'\mathrm{B}_g},
\]
entendiendo que $\mathrm{M}/\mathfrak{p}\mathrm{M}$ y
$\operatorname{Tor}_1^\mathrm{A}
(\mathrm{M},\mathrm{A}/\mathfrak{p})$ se consideran
como $\mathrm{B}$-módulos. En virtud de $(\mathbf{0}_\mathrm{I},~6.3.2)$, quedamos
así reducidos a demostrar el lema siguiente.

\begin{lemma}[11.1.1.2]
\label{IV.11.1.1.2}
Bajo las condiciones de \emph{{\hyperlink{ega4.c04.sec.11-1-1-1}{(11.1.1.1)}}}, existe $g\in\mathrm{B}-\mathfrak{q}$ tal que:
\begin{enumerate}[label=\emph{\roman*)}]
  \item $(\mathrm{M}/\mathfrak{p}\mathrm{M})_g$ es un $(\mathrm{A}/\mathfrak{p})$-módulo plano;
  \item $(\operatorname{Tor}_1^\mathrm{A}(\mathrm{M},\mathrm{A}/\mathfrak{p}))_g=0$.
\end{enumerate}
\end{lemma}

\oldpage[IV-3]{118}

Por el teorema de planitud genérica $(6.9.1)$, aplicado al anillo íntegro
$\mathrm{A}/\mathfrak{p}$, a la
$(\mathrm{A}/\mathfrak{p})$-álgebra de tipo finito
$\mathrm{B}/\mathfrak{p}\mathrm{B}$ y al
$(\mathrm{B}/\mathfrak{p}\mathrm{B})$-módulo de tipo finito
$\mathrm{M}/\mathfrak{p}\mathrm{M}$, existe
$h\in\mathrm{A}-\mathfrak{p}$ tal que, si $\bar h$ denota su imagen canónica
en $\mathrm{A}/\mathfrak{p}$,
$(\mathrm{M}/\mathfrak{p}\mathrm{M})_{\bar h}$ es un
$(\mathrm{A}/\mathfrak{p})$-módulo plano. Por otra parte,
$\mathrm{M}_{\mathfrak q}$ es plano sobre $\mathrm{A}_{\mathfrak p}$ y,
por tanto, en virtud de $(\mathbf{0}_\mathrm{I},~6.3.1)$, es plano sobre $\mathrm{A}$.
En consecuencia,
\[
  \operatorname{Tor}_1^\mathrm{A}
  (\mathrm{M}_{\mathfrak q},\mathrm{A}/\mathfrak p)=0,
  \qquad\text{o, equivalentemente,}\qquad
  \bigl(\operatorname{Tor}_1^\mathrm{A}
  (\mathrm{M},\mathrm{A}/\mathfrak p)\bigr)_{\mathfrak q}=0.
\]
Como $\mathrm{A}$ y $\mathrm{B}$ son noetherianos,
$\operatorname{Tor}_1^\mathrm{A}
(\mathrm{M},\mathrm{A}/\mathfrak p)$ es un $\mathrm{B}$-módulo de tipo finito.
Así, $(\mathbf{0}_\mathrm{I},~5.2.2)$ proporciona un elemento
$g\in\mathrm{B}-\mathfrak q$ tal que
\[
  \bigl(\operatorname{Tor}_1^\mathrm{A}
  (\mathrm{M},\mathrm{A}/\mathfrak p)\bigr)_g=0.
\]
Además,
$(\mathrm{M}/\mathfrak p\mathrm{M})_{\bar h}
=(\mathrm{M}/\mathfrak p\mathrm{M})_h$, donde a la derecha el módulo se
considera como $\mathrm{A}$-módulo. Como
$(\mathrm{M}/\mathfrak p\mathrm{M})_h$ es también un $\mathrm{B}$-módulo,
$(\mathrm{M}/\mathfrak p\mathrm{M})_{hg}$ sigue siendo plano sobre
$\mathrm{A}/\mathfrak p$: es
$(\mathrm{M}/\mathfrak p\mathrm{M})_{\overline{hg}}$, donde
$\overline{hg}$ es la imagen canónica de $hg$ en
$\mathrm{B}/\mathfrak p\mathrm{B}$, y se aplica
$(\mathbf{0}_\mathrm{I},~6.3.2)$. Finalmente,
\[
  \bigl(\operatorname{Tor}_1^\mathrm{A}
  (\mathrm{M},\mathrm{A}/\mathfrak p)\bigr)_{hg}=0,
\]
y $hg\in\mathrm{B}-\mathfrak q$ porque
$h\in\mathrm{A}-\mathfrak p$. Esto demuestra el lema y el teorema.
\end{proof}

\begin{corollary}[11.1.2]
\label{IV.11.1.2}
Sea $\mathrm{Y}$ un preesquema localmente noetheriano, $f:\mathrm{X}\to\mathrm{Y}$ un morfismo localmente de tipo finito, $\sh{F}$ y $\sh{F}'$ dos $\sh{O}_\mathrm{X}$-Módulos coherentes, y $u:\sh{F}'\to\sh{F}$ un homomorfismo de $\sh{O}_\mathrm{X}$-Módulos. Supongamos que $\sh{F}$ es $f$-plano. Entonces el conjunto $\mathrm{U}$ de los $x\in\mathrm{X}$ tales que, poniendo $y=f(x)$, el homomorfismo $u_x\otimes 1:\sh{F}'_x\otimes_{\sh{O}_y}\boldsymbol{k}(y)\to\sh{F}_x\otimes_{\sh{O}_y}\boldsymbol{k}(y)$ es inyectivo, es abierto en $\mathrm{X}$.
\end{corollary}

\begin{proof}
Sean $\sh{N}$ (resp.\ $\sh{P}$) el núcleo (resp.\ el conúcleo) de $u$.
Apliquemos $(\mathbf{0}_{\mathrm{III}},~10.2.4)$\footnote{Damos aquí una demostración de
$(\mathbf{0}_{\mathrm{III}},~10.2.4)$ que no fue publicada en el
\emph{Alg\`ebre commutative} de N.~Bourbaki. En virtud de
$(\mathbf{0}_\mathrm{I},~6.1.2)$, basta demostrar que b) implica a).
Pongamos $\mathrm{P}=\operatorname{Im}(u)$,
$\mathrm{Q}=\operatorname{Coker}(u)$ y
$\mathrm{R}=\operatorname{Ker}(u)$. El compuesto
$\mathrm{M}\otimes k\to\mathrm{P}\otimes k\to\mathrm{N}\otimes k$ es
inyectivo, mientras que $\mathrm{M}\otimes k\to\mathrm{P}\otimes k$ es
sobreyectivo. Por tanto, $\mathrm{P}\otimes k\to\mathrm{N}\otimes k$ es
inyectivo y $\mathrm{M}\otimes k\to\mathrm{P}\otimes k$ es biyectivo.
La sucesión exacta $0\to\mathrm{P}\to\mathrm{N}\to\mathrm{Q}\to0$ proporciona
la sucesión exacta
\[
0=\operatorname{Tor}_1^\mathrm{A}(\mathrm{N},k)\to\operatorname{Tor}_1^\mathrm{A}(\mathrm{Q},k)\to\mathrm{P}\otimes k\to\mathrm{N}\otimes k\to\mathrm{Q}\otimes k\to 0
\]
y, como $\mathrm{P}\otimes k\to\mathrm{N}\otimes k$ es inyectivo,
$\operatorname{Tor}_1^\mathrm{A}(\mathrm{Q},k)=0$. Puesto que $\mathrm{Q}$ es un
$\mathrm{B}$-módulo de tipo finito,
$(\mathbf{0}_{\mathrm{III}},~10.2.2)$ muestra que es plano sobre
$\mathrm{A}$; entonces $\mathrm{P}$ es también plano sobre $\mathrm{A}$ por
$(\mathbf{0}_\mathrm{I},~6.1.2)$. A partir de la sucesión exacta
$0\to\mathrm{R}\to\mathrm{M}\to\mathrm{P}\to0$, el mismo resultado proporciona una
sucesión exacta
$0\to\mathrm{R}\otimes k\to\mathrm{M}\otimes k\to
\mathrm{P}\otimes k\to0$. Como la última aplicación es biyectiva,
$\mathrm{R}\otimes k=0$. Pero $\mathrm{B}$ es noetheriano, de modo que
$\mathrm{R}$ es un $\mathrm{B}$-módulo de tipo finito; el lema de Nakayama
da por tanto $\mathrm{R}=0$.}
a los anillos locales $\sh{O}_y$ y $\sh{O}_x$, y a los
$\sh{O}_x$-módulos $\sh{F}'_x$ y $\sh{F}_x$. La inyectividad de
$u_x\otimes1$ equivale a las dos condiciones
$\sh{N}_x=0$ y que $\sh{P}$ sea $f$-plano en $x$. Como
$\sh{N}$ y $\sh{P}$ son coherentes $(\mathbf{0}_\mathrm{I},~5.3.4)$,
el lugar donde $\sh{N}_x=0$ es abierto
$(\mathbf{0}_\mathrm{I},~5.2.2)$, y el lugar donde $\sh{P}$ es
$f$-plano es abierto por \sref{IV.11.1.1}. De aquí resulta la afirmación.

En particular:
\end{proof}

\begin{corollary}[11.1.3]
\label{IV.11.1.3}
Sea $\mathrm{Y}$ un preesquema localmente noetheriano,
$f:\mathrm{X}\to\mathrm{Y}$ un morfismo plano localmente de tipo finito, y
$g$ una sección de $\sh{O}_\mathrm{X}$ sobre $\mathrm{X}$. El conjunto de los
$x\in\mathrm{X}$ tales que $g_x\otimes 1$ es un no divisor de cero en
$\sh{O}_x\otimes_{\sh{O}_{f(x)}}\boldsymbol{k}(f(x))$ es abierto en
$\mathrm{X}$.
\end{corollary}

\begin{proof}
Basta aplicar \sref{IV.11.1.2} al endomorfismo de $\sh{O}_\mathrm{X}$ definido por $g$.
\end{proof}

\begin{corollary}[11.1.4]
\label{IV.11.1.4}
Sea $\mathrm{Y}$ un preesquema localmente noetheriano, $f:\mathrm{X}\to\mathrm{Y}$ un morfismo localmente de tipo finito, y $\sh{F}$ un $\sh{O}_\mathrm{X}$-Módulo coherente y $f$-plano. Sea $(g_i)_{1\leqslant i\leqslant n}$ una sucesión de secciones de $\sh{O}_\mathrm{X}$ sobre $\mathrm{X}$. El conjunto $\mathrm{U}$ de los $x\in\mathrm{X}$ tales que la sucesión $((g_i)_x\otimes 1)$ es $(\sh{F}_x\otimes_{\sh{O}_{f(x)}}\boldsymbol{k}(f(x)))$-regular es abierto en $\mathrm{X}$.
\end{corollary}

\begin{proof}
\oldpage[IV-3]{119}
Como $\sh{F}$ es $f$-plano, se deduce de
$(\mathbf{0},~15.1.16)$ que $\mathrm{U}$ es también el conjunto de los puntos
$x\in\mathrm{X}$ para los cuales la sucesión $((g_i)_x)$ es
$\sh{F}_x$-regular y el $\sh{O}_x$-módulo
\[
  \sh{G}_x,\qquad
  \sh{G}=\sh{F}\big/\Bigl(\sum_{i=1}^n g_i\sh{F}\Bigr),
\]
es $f$-plano. Los módulos $\sh{F}$ y $\sh{G}$ son coherentes, de modo que el
corolario se deduce de \sref{IV.11.1.1} y de
$(\mathbf{0},~15.2.4)$.
\end{proof}

\begin{corollary}[11.1.5]
\label{IV.11.1.5}
Sean $\mathrm{X}$, $\mathrm{Y}$ y $\mathrm{Z}$ tres preesquemas localmente noetherianos, $f:\mathrm{X}\to\mathrm{Y}$, $g:\mathrm{Y}\to\mathrm{Z}$ dos morfismos de tipo finito, y $\sh{F}$ un $\sh{O}_\mathrm{X}$-Módulo coherente. Entonces el conjunto $\mathrm{U}$ de los $y\in\mathrm{Y}$ tales que, para toda generalización $y'$ de $y$, $\sh{F}$ es $(g\circ f)$-plano en todos los puntos de $f^{-1}(y')$ (es decir,\ tales que $\sh{F}\otimes_\mathrm{Y}\Spec(\sh{O}_{\mathrm{Y},y})$ es plano respecto del morfismo $\mathrm{X}\times_\mathrm{Y}\Spec(\sh{O}_{\mathrm{Y},y})\to\mathrm{Z}$) es abierto en $\mathrm{Y}$.
\end{corollary}

\begin{proof}
Si $\mathrm{V}$ es el conjunto de los puntos $x\in\mathrm{X}$ en los que $\sh{F}$ es
$(g\circ f)$-plano, entonces $\mathrm{U}$ es el conjunto de los $y\in\mathrm{Y}$ tales
que $f^{-1}(y')\subset\mathrm{V}$ para toda generalización $y'$ de $y$.
Ahora bien, $\mathrm{V}$ es abierto por \sref{IV.11.1.1}, y por tanto localmente
construible en $\mathrm{X}$, y el conjunto $\mathrm{W}$ de los
$y\in\mathrm{Y}$ tales que $f^{-1}(y)\subset\mathrm{V}$ es
\[
  \mathrm{W}=\mathrm{Y}-f(\mathrm{X}-\mathrm{V}).
\]
Es, por consiguiente, localmente construible en $\mathrm{Y}$ por el teorema de
Chevalley $(1.8.4)$. Se deduce de
$(\mathbf{0}_{\mathrm{III}},~9.2.5)$ que los puntos de $\mathrm{U}$ son
los puntos interiores de $\mathrm{W}$, y por tanto que $\mathrm{U}$ es abierto.
\end{proof}

\begin{corollary}[11.1.6]
\label{IV.11.1.6}
Sea $\mathrm{A}$ un anillo noetheriano, $\mathrm{B}$ una $\mathrm{A}$-álgebra de tipo finito, $\mathrm{C}$ una $\mathrm{B}$-álgebra de tipo finito y $\mathrm{M}$ un $\mathrm{C}$-módulo de tipo finito. Entonces el conjunto de los $\mathfrak{q}\in\Spec(\mathrm{B})$ tales que $\mathrm{M}_\mathfrak{q}$ es un $\mathrm{A}$-módulo plano es abierto en $\Spec(\mathrm{B})$.
\end{corollary}

\begin{proof}
Teniendo en cuenta $(2.1.2)$, esto es consecuencia de \sref{IV.11.1.5} aplicado a $\mathrm{Z}=\Spec(\mathrm{A})$, $\mathrm{Y}=\Spec(\mathrm{B})$, $\mathrm{X}=\Spec(\mathrm{C})$, $\sh{F}=\widetilde{\mathrm{M}}$.

Los resultados de este número se formularán sin hipótesis noetherianas en \sref{IV.11.3}.
\end{proof}

\subsection{Planitud de un límite proyectivo de preesquemas}
\label{subsection:IV.11.2}

\oldpage[IV-3]{119}

\begin{env}[11.2.1]
\label{IV.11.2.1}
Sea $\mathrm{A}$ un anillo, sean $\mathrm{M}$ y $\mathrm{N}$ dos A-módulos, y sea $\mathrm{A}'$ una A-álgebra; pongamos $\mathrm{M}'=\mathrm{M}\otimes_\mathrm{A}\mathrm{A}'$, $\mathrm{N}'=\mathrm{N}\otimes_\mathrm{A}\mathrm{A}'$. Recordemos $(\textbf{III},~6.3.8)$ que, para todo $i$, se define un homomorfismo canónico de A-módulos
\begin{equation}
  \psi_i:\operatorname{Tor}_i^\mathrm{A}(\mathrm{M},\mathrm{N})\to\operatorname{Tor}_i^{\mathrm{A}'}(\mathrm{M}',\mathrm{N}')
  \label{IV.11.2.1.1}
\tag{11.2.1.1}
\end{equation}
de la manera siguiente: se considera una resolución por la izquierda de $\mathrm{M}$ mediante A-módulos libres
\begin{equation}
  \ldots\longrightarrow\mathrm{L}_{i+1}\xrightarrow{f_i}\mathrm{L}_i\longrightarrow\ldots\longrightarrow\mathrm{L}_0\xrightarrow{\varepsilon}\mathrm{M}\longrightarrow 0
  \label{IV.11.2.1.2}
\tag{11.2.1.2}
\end{equation}
de la cual se deduce, por producto tensorial con $\mathrm{A}'$, un complejo de $\mathrm{A}'$-módulos
\begin{equation}
  \ldots\longrightarrow\mathrm{L}'_{i+1}\xrightarrow{f'_i}\mathrm{L}'_i\longrightarrow\ldots\longrightarrow\mathrm{L}'_0\xrightarrow{\varepsilon'}\mathrm{M}'\longrightarrow 0
  \label{IV.11.2.1.3}
\tag{11.2.1.3}
\end{equation}
donde hemos puesto $\mathrm{L}'_i=\mathrm{L}_i\otimes_\mathrm{A}\mathrm{A}'$, $f'_i=f_i\otimes 1_{\mathrm{A}'}$, $\varepsilon'=\varepsilon\otimes 1_{\mathrm{A}'}$. Consideremos, por otra parte, una resolución por la izquierda de $\mathrm{M}'$ mediante $\mathrm{A}'$-módulos libres
\begin{equation}
  \ldots\longrightarrow\mathrm{L}''_{i+1}\xrightarrow{f''_i}\mathrm{L}''_i\longrightarrow\ldots\longrightarrow\mathrm{L}''_0\xrightarrow{\varepsilon''}\mathrm{M}'\longrightarrow 0
  \label{IV.11.2.1.4}
\tag{11.2.1.4}
\end{equation}
\end{env}

\oldpage[IV-3]{120}

Como los $\mathrm{L}'_i$ son $\mathrm{A}'$-módulos libres, sabemos $(\mathrm{M},~\mathrm{V},~1.1)$ que existen homomorfismos $\mathrm{A}'$-lineales $u_i$ que forman un diagrama conmutativo
\[
\xymatrix{
  \ldots \ar[r] & \mathrm{L}'_{i+1} \ar[r]^{f'_i} \ar[d]_{u_{i+1}} & \mathrm{L}'_i \ar[r] \ar[d]_{u_i} & \ldots \ar[r] & \mathrm{L}'_0 \ar[r]^{\varepsilon'} \ar[d]_{u_0} & \mathrm{M}' \ar[r] \ar[d]_{1_{\mathrm{M}'}} & 0 \\
  \ldots \ar[r] & \mathrm{L}''_{i+1} \ar[r]_{f''_i} & \mathrm{L}''_i \ar[r] & \ldots \ar[r] & \mathrm{L}''_0 \ar[r]_{\varepsilon''} & \mathrm{M}' \ar[r] & 0
}
\]

Al componer el morfismo de complejos
$u_\bullet:\mathrm{L}'_\bullet\to\mathrm{L}''_\bullet$ con el morfismo canónico
$\mathrm{L}_\bullet\to\mathrm{L}'_\bullet$, se obtiene un morfismo de
complejos de $\mathrm{A}$-módulos
$\mathrm{L}_\bullet\to\mathrm{L}''_\bullet$. Como
\[
  (\mathrm{L}_i\otimes_\mathrm{A}\mathrm{N})
    \otimes_\mathrm{A}\mathrm{A}'
  =\mathrm{L}'_i\otimes_{\mathrm{A}'}\mathrm{N}',
\]
obtenemos un morfismo de complejos
\[
  \mathrm{L}_\bullet\otimes_\mathrm{A}\mathrm{N}
  \longrightarrow
  \mathrm{L}''_\bullet\otimes_{\mathrm{A}'}\mathrm{N}'.
\]
Al pasar a la homología se obtienen los homomorfismos canónicos
\eref{IV.11.2.1.1}. Como los $u_i$ están determinados salvo
homotopía $(\mathrm{M},~\mathrm{V},~1.1)$, estos homomorfismos no dependen
ni de la elección de los $u_i$ ni de la elección de las resoluciones libres
$\mathrm{L}_\bullet$ y $\mathrm{L}''_\bullet$.

Como los $\operatorname{Tor}_i^{\mathrm{A}'}(\mathrm{M}',\mathrm{N}')$ son $\mathrm{A}'$-módulos, de \eref{IV.11.2.1.1} se deducen canónicamente homomorfismos $\mathrm{A}'$-lineales
\begin{equation}
  \varphi_i:\operatorname{Tor}_i^\mathrm{A}(\mathrm{M},\mathrm{N})\otimes_\mathrm{A}\mathrm{A}'\to\operatorname{Tor}_i^{\mathrm{A}'}(\mathrm{M}',\mathrm{N}').
  \label{IV.11.2.1.5}
\tag{11.2.1.5}
\end{equation}

Consideremos ahora dos homomorfismos de anillos
\[
  \rho:\mathrm{A}\to\mathrm{A}^{(1)},\quad \sigma:\mathrm{A}^{(1)}\to\mathrm{A}^{(2)},
\]
y su compuesto $\sigma\circ\rho:\mathrm{A}\to\mathrm{A}^{(2)}$; pongamos $\mathrm{M}^{(j)}=\mathrm{M}\otimes_\mathrm{A}\mathrm{A}^{(j)}$, $\mathrm{N}^{(j)}=\mathrm{N}\otimes_\mathrm{A}\mathrm{A}^{(j)}$ para $j=1,2$. Entonces el homomorfismo canónico compuesto
\[
  \operatorname{Tor}_i^\mathrm{A}(\mathrm{M},\mathrm{N})
  \longrightarrow
  \operatorname{Tor}_i^{\mathrm{A}^{(1)}}
    (\mathrm{M}^{(1)},\mathrm{N}^{(1)})
  \longrightarrow
  \operatorname{Tor}_i^{\mathrm{A}^{(2)}}
    (\mathrm{M}^{(2)},\mathrm{N}^{(2)})
\]
coincide con el homomorfismo canónico deducido de $\sigma\circ\rho$; esto resulta de que, si $\mathrm{L}_\bullet^{(j)}$ es una resolución libre de $\mathrm{M}^{(j)}$, el diagrama
\[
\xymatrix{
  \mathrm{L}_\bullet \ar[r] & \mathrm{L}_\bullet\otimes_\mathrm{A}\mathrm{A}^{(1)} \ar[r] \ar[d] & \mathrm{L}_\bullet^{(1)} \ar[d] \\
  & \mathrm{L}_\bullet\otimes_\mathrm{A}\mathrm{A}^{(2)} \ar[r] &
  \mathrm{L}_\bullet^{(1)}\otimes_{\mathrm{A}^{(1)}}\mathrm{A}^{(2)} \ar[r] & \mathrm{L}_\bullet^{(2)}
}
\]
es conmutativo.

\begin{env}[11.2.2]
\label{IV.11.2.2}
Con la notación de \sref{IV.11.2.1}, sea
$(\mathrm{A}_\alpha,g_{\beta\alpha})$ un sistema inductivo filtrante de
$\mathrm{A}$-álgebras y pongamos
\[
  \mathrm{M}_\alpha=\mathrm{M}\otimes_\mathrm{A}\mathrm{A}_\alpha,
  \qquad
  \mathrm{N}_\alpha=\mathrm{N}\otimes_\mathrm{A}\mathrm{A}_\alpha.
\]
Para cada $i$, los módulos
$\operatorname{Tor}_i^{\mathrm{A}_\alpha}
(\mathrm{M}_\alpha,\mathrm{N}_\alpha)$, con las aplicaciones canónicas proporcionadas
por \sref{IV.11.2.1}, forman un sistema inductivo. Pongamos
\[
  \mathrm{A}'=\varinjlim_\alpha\mathrm{A}_\alpha,\qquad
  \mathrm{M}'=\varinjlim_\alpha\mathrm{M}_\alpha
    =\mathrm{M}\otimes_\mathrm{A}\mathrm{A}',\qquad
  \mathrm{N}'=\varinjlim_\alpha\mathrm{N}_\alpha
    =\mathrm{N}\otimes_\mathrm{A}\mathrm{A}'.
\]
Si $g_\alpha:\mathrm{A}_\alpha\to\mathrm{A}'$ denota el homomorfismo
canónico, obtenemos aplicaciones canónicas
\[
  \psi_{i\alpha}:
  \operatorname{Tor}_i^{\mathrm{A}_\alpha}
    (\mathrm{M}_\alpha,\mathrm{N}_\alpha)
  \longrightarrow
  \operatorname{Tor}_i^{\mathrm{A}'}(\mathrm{M}',\mathrm{N}')
\]
que forman también un sistema inductivo. Completamos el resultado de
$(\mathrm{M},~\mathrm{V},~9.4^*)$ demostrando que
\begin{equation}
  \psi_i=\varinjlim\psi_{i\alpha}:\varinjlim_\alpha\operatorname{Tor}_i^{\mathrm{A}_\alpha}(\mathrm{M}_\alpha,\mathrm{N}_\alpha)\to\operatorname{Tor}_i^{\mathrm{A}'}(\mathrm{M}',\mathrm{N}')
  \label{IV.11.2.2.1}
\tag{11.2.2.1}
\end{equation}
son \emph{isomorfismos de $\mathrm{A}'$-módulos}. Para ello procedemos como en $(\mathrm{M},~\mathrm{V},~9.5^*)$, asociando a cada $\mathrm{M}_\alpha$ su \emph{resolución libre canónica}. Teniendo en cuenta la exactitud del funtor $\varinjlim$, esto equivale a demostrar el
\end{env}

\oldpage[IV-3]{121}

\begin{lemma}[11.2.2.2]
\label{IV.11.2.2.2}
Sea $(\mathrm{A}_\alpha,g_{\beta\alpha})$ un sistema inductivo filtrante
de anillos y $(\mathrm{M}_\alpha,h_{\beta\alpha})$ un sistema inductivo
filtrante de conjuntos. Pongamos
$\mathrm{A}'=\varinjlim\mathrm{A}_\alpha$ y
$\mathrm{M}'=\varinjlim\mathrm{M}_\alpha$, y sean
$g_\alpha:\mathrm{A}_\alpha\to\mathrm{A}'$ y
$h_\alpha:\mathrm{M}_\alpha\to\mathrm{M}'$ las aplicaciones canónicas. Para
cada $\alpha$, sea $\mathrm{F}(\mathrm{M}_\alpha)$ el
$\mathrm{A}_\alpha$-módulo de las combinaciones lineales formales de elementos de
$\mathrm{M}_\alpha$; análogamente, sea $\mathrm{F}(\mathrm{M}')$ el
$\mathrm{A}'$-módulo de las combinaciones lineales formales de elementos de
$\mathrm{M}'$. Si
$h_{\beta\alpha}^0:\mathrm{F}(\mathrm{M}_\alpha)\to
\mathrm{F}(\mathrm{M}_\beta)$ (para $\alpha\leqslant\beta$) y
$h'_\alpha:\mathrm{F}(\mathrm{M}_\alpha)\to\mathrm{F}(\mathrm{M}')$
son los homomorfismos $\mathrm{A}_\alpha$-lineales inducidos respectivamente por
$h_{\beta\alpha}$ y $h_\alpha$, entonces
$(\mathrm{F}(\mathrm{M}_\alpha),h^0_{\beta\alpha})$ es un sistema inductivo
de $\mathrm{A}_\alpha$-módulos y $(h'_\alpha)$ es un sistema inductivo
de homomorfismos. Además,
\[
  h''=\varinjlim h'_\alpha:\varinjlim\mathrm{F}(\mathrm{M}_\alpha)\to\mathrm{F}(\mathrm{M}')=\mathrm{F}(\varinjlim\mathrm{M}_\alpha)
\]
\emph{es un isomorfismo.}
\end{lemma}

\begin{proof}
Para la demostración, véase Bourbaki, \emph{Alg.}, cap.~II, 3\textsuperscript{e}~\'{e}d., \textsection\,6, n\textsuperscript{o}\,6, cor.\ de la prop.~10.
\end{proof}

\begin{env}[11.2.3]
\label{IV.11.2.3}
Retomemos la notación de \sref{IV.11.2.1} y consideremos el caso $i=1$.
Pongamos
$\mathrm{T}=\operatorname{Tor}_1^\mathrm{A}(\mathrm{M},\mathrm{N})$,
$\mathrm{T}'=\mathrm{T}\otimes_\mathrm{A}\mathrm{A}'$, y
$\mathrm{T}''=\operatorname{Tor}_1^{\mathrm{A}'}
(\mathrm{M}',\mathrm{N}')$. Entonces $\psi_1$ es el homomorfismo
\[
  \frac{\operatorname{Ker}(f_0\otimes1_\mathrm{N})}
       {\operatorname{Im}(f_1\otimes1_\mathrm{N})}
  \longrightarrow
  \frac{\operatorname{Ker}(f''_0\otimes1_{\mathrm{N}'})}
       {\operatorname{Im}(f''_1\otimes1_{\mathrm{N}'})}
\]
inducido en los cocientes por la restricción
\[
  \operatorname{Ker}(f_0\otimes1_\mathrm{N})
  \longrightarrow
  \operatorname{Ker}(f''_0\otimes1_{\mathrm{N}'})
\]
de
$u_1\otimes1:\mathrm{L}_1\otimes_\mathrm{A}\mathrm{N}\to
\mathrm{L}''_1\otimes_{\mathrm{A}'}\mathrm{N}'$.

Pongamos $\mathrm{R}=\operatorname{Ker}(\varepsilon)$, $\mathrm{R}''=\operatorname{Ker}(\varepsilon'')$, de modo que se tienen las sucesiones exactas
\begin{equation}
  0\longrightarrow\mathrm{R}\xrightarrow{j}\mathrm{L}_0\xrightarrow{\varepsilon}\mathrm{M}\longrightarrow 0
  \quad\text{y}\quad
  0\longrightarrow\mathrm{R}''\xrightarrow{j''}\mathrm{L}''_0\xrightarrow{\varepsilon''}\mathrm{M}'\longrightarrow 0,
\end{equation}
de las cuales se deducen las sucesiones exactas de homología
\begin{equation}
  0=\operatorname{Tor}_1^\mathrm{A}(\mathrm{L}_0,\mathrm{N})\to\mathrm{T}\xrightarrow{\partial}\mathrm{R}\otimes_\mathrm{A}\mathrm{N}\xrightarrow{j\otimes 1}\mathrm{L}_0\otimes_\mathrm{A}\mathrm{N}\xrightarrow{\varepsilon\otimes 1}\mathrm{M}\otimes_\mathrm{A}\mathrm{N}\to 0
  \label{IV.11.2.3.1}
\tag{11.2.3.1}
\end{equation}
\begin{equation}
  0=\operatorname{Tor}_1^{\mathrm{A}'}(\mathrm{L}''_0,\mathrm{N}')\to\mathrm{T}''\xrightarrow{\partial''}\mathrm{R}''\otimes_{\mathrm{A}'}\mathrm{N}'\xrightarrow{j''\otimes 1}\mathrm{L}''_0\otimes_{\mathrm{A}'}\mathrm{N}'\xrightarrow{\varepsilon''\otimes 1}\mathrm{M}'\otimes_{\mathrm{A}'}\mathrm{N}'\to 0
  \label{IV.11.2.3.2}
\tag{11.2.3.2}
\end{equation}

Por otra parte, existe un homomorfismo de $\mathrm{A}'$-módulos
\[
  v:\mathrm{R}'=\operatorname{Ker}(\varepsilon)\otimes_\mathrm{A}\mathrm{A}'\longrightarrow\operatorname{Ker}(\varepsilon')\xrightarrow{u_0}\operatorname{Ker}(\varepsilon'')=\mathrm{R}''.
\]
\end{env}

\oldpage[IV-3]{122}

Demostremos que el diagrama
\begin{equation}
\xymatrix{
  & \mathrm{T}' \ar[r]^-{\partial\otimes 1} \ar[d]_{\varphi_1} & \mathrm{R}'\otimes_{\mathrm{A}'}\mathrm{N}' \ar[r] \ar[d]_{v\otimes 1} & \mathrm{L}'_0\otimes_{\mathrm{A}'}\mathrm{N}' \ar[r] \ar[d]_{u_0\otimes 1} & \mathrm{M}'\otimes_{\mathrm{A}'}\mathrm{N}' \ar[r] \ar@{=}[d] & 0 \\
  0 \ar[r] & \mathrm{T}'' \ar[r]_{\partial''} & \mathrm{R}''\otimes_{\mathrm{A}'}\mathrm{N}' \ar[r] & \mathrm{L}''_0\otimes_{\mathrm{A}'}\mathrm{N}' \ar[r] & \mathrm{M}'\otimes_{\mathrm{A}'}\mathrm{N}' \ar[r] & 0
}
  \label{IV.11.2.3.3}
\tag{11.2.3.3}
\end{equation}
es \emph{conmutativo}. Para ello se comprueba inmediatamente $(\mathrm{M},~\mathrm{IV},~1)$ que el homomorfismo $\partial:\mathrm{T}\to\mathrm{R}\otimes_\mathrm{A}\mathrm{N}$ se obtiene (en el caso presente), por paso al cociente, del homomorfismo $w:\operatorname{Ker}(f_0\otimes 1_\mathrm{N})\to\mathrm{R}\otimes_\mathrm{A}\mathrm{N}$, restricción del homomorfismo $g_0\otimes 1_\mathrm{N}:\mathrm{L}_1\otimes_\mathrm{A}\mathrm{N}\to\mathrm{R}\otimes_\mathrm{A}\mathrm{N}$, donde $g_0:\mathrm{L}_1\to\mathrm{R}$ es tal que $f_0=j\circ g_0$; y análogamente para $\partial''$. Basta, por tanto, ver que el diagrama
\[
\xymatrix{
  \operatorname{Ker}(f_0\otimes 1_\mathrm{N}) \ar[r]^{w} \ar[d] & \mathrm{R}\otimes_\mathrm{A}\mathrm{N} \ar[r] & \mathrm{R}'\otimes_{\mathrm{A}'}\mathrm{N}'=(\mathrm{R}\otimes_\mathrm{A}\mathrm{N})\otimes_\mathrm{A}\mathrm{A}' \ar[d]_{v\otimes 1} \\
  \operatorname{Ker}(f''_0\otimes 1_{\mathrm{N}'}) \ar[rr] & & \mathrm{R}''\otimes_{\mathrm{A}'}\mathrm{N}'
}
\]
es conmutativo, lo que resulta de la conmutatividad del diagrama
\[
\xymatrix{
  \mathrm{L}_1 \ar[r] \ar[d] & \mathrm{R} \ar[d] \\
  \mathrm{L}''_1 \ar[r] & \mathrm{R}''
}
\]

\begin{lemma}[11.2.4]
\label{IV.11.2.4}
Sea $\mathrm{A}$ un anillo, $\mathrm{C}$ una $\mathrm{A}$-álgebra, $\mathrm{M}$ un $\mathrm{A}$-módulo, $\mathrm{N}$ un $\mathrm{C}$-módulo y $\mathrm{A}'$ una $\mathrm{A}$-álgebra. Pongamos $\mathrm{C}'=\mathrm{C}\otimes_\mathrm{A}\mathrm{A}'$, $\mathrm{M}'=\mathrm{M}\otimes_\mathrm{A}\mathrm{A}'$, $\mathrm{N}'=\mathrm{N}\otimes_\mathrm{A}\mathrm{A}'=\mathrm{N}\otimes_\mathrm{C}\mathrm{C}'$. Supongamos que $\mathrm{M}\otimes_\mathrm{A}\mathrm{N}$ es un $\mathrm{C}$-módulo plano. Entonces el homomorfismo canónico
\[
  \varphi_1:\operatorname{Tor}_1^\mathrm{A}(\mathrm{M},\mathrm{N})\otimes_\mathrm{A}\mathrm{A}'\to\operatorname{Tor}_1^{\mathrm{A}'}(\mathrm{M}',\mathrm{N}')
\]
(cf.\ \eref{IV.11.2.1.5}) \emph{es sobreyectivo}.
\end{lemma}

\begin{proof}
Conservemos la notación de \sref{IV.11.2.3}. La exactitud por la derecha del producto
tensorial muestra que
\[
  \mathrm{L}'_1\xrightarrow{f'_0}\mathrm{L}'_0
  \xrightarrow{\varepsilon'}\mathrm{M}'\longrightarrow0
\]
es exacta. Como $\mathrm{L}'_0$ y $\mathrm{L}'_1$ son
$\mathrm{A}'$-módulos libres, podemos suponer que
$\mathrm{L}''_0=\mathrm{L}'_0$ y
$\mathrm{L}''_1=\mathrm{L}'_1$, que $u_0$ y $u_1$ son las aplicaciones
identidad y que $f''_0=f'_0$. Como
$\operatorname{Ker}(\varepsilon')=\operatorname{Im}(f'_0)$ y
$\mathrm{R}''=\operatorname{Im}(f''_0)=\operatorname{Im}(f'_0)$, el
homomorfismo $v$ es sobreyectivo y, por tanto, también lo es $v\otimes1$. La demostración
quedará completa si la primera fila de
\eref{IV.11.2.3.3} es exacta, pues $v\otimes1$ es
sobreyectivo y $u_0\otimes1$ es inyectivo (Bourbaki,
\emph{Alg.\ comm.}, cap.~I, \S\,1, n.º~4, Cor.~2 de la Prop.~2). Para ello
utilizamos

\oldpage[IV-3]{123}

la hipótesis de que $\mathrm{M}\otimes_\mathrm{A}\mathrm{N}$ es un C-módulo plano. Poniendo $\mathrm{Q}=\operatorname{Ker}(\varepsilon\otimes 1)=\operatorname{Im}(j\otimes 1)$, la sucesión exacta \eref{IV.11.2.3.1} da las dos sucesiones exactas $0\to\mathrm{Q}\to\mathrm{L}_0\otimes_\mathrm{A}\mathrm{N}\to\mathrm{M}\otimes_\mathrm{A}\mathrm{N}\to 0$ y $\mathrm{T}\to\mathrm{R}\otimes_\mathrm{A}\mathrm{N}\to\mathrm{Q}\to 0$, cuyos homomorfismos son homomorfismos de \emph{C-módulos}; utilizando la hipótesis de planitud, se deduce de $(\mathbf{0}_\mathrm{I},~6.1.2)$ la sucesión exacta
\[
  0\to\mathrm{Q}\otimes_\mathrm{C}\mathrm{C}'\to(\mathrm{L}_0\otimes_\mathrm{A}\mathrm{N})\otimes_\mathrm{C}\mathrm{C}'\to(\mathrm{M}\otimes_\mathrm{A}\mathrm{N})\otimes_\mathrm{C}\mathrm{C}'\to 0
\]
y, por otra parte, como el producto tensorial es exacto por la derecha, se tiene la sucesión exacta
\[
  \mathrm{T}\otimes_\mathrm{C}\mathrm{C}'\to(\mathrm{R}\otimes_\mathrm{A}\mathrm{N})\otimes_\mathrm{C}\mathrm{C}'\to\mathrm{Q}\otimes_\mathrm{C}\mathrm{C}'\to 0
\]
de donde se obtiene finalmente la sucesión exacta
\[
  \mathrm{T}\otimes_\mathrm{C}\mathrm{C}'\to(\mathrm{R}\otimes_\mathrm{A}\mathrm{N})\otimes_\mathrm{C}\mathrm{C}'\to(\mathrm{L}_0\otimes_\mathrm{A}\mathrm{N})\otimes_\mathrm{C}\mathrm{C}'\to(\mathrm{M}\otimes_\mathrm{A}\mathrm{N})\otimes_\mathrm{C}\mathrm{C}'\to 0
\]
Pero, por definición, para todo C-módulo $\mathrm{P}$, $\mathrm{P}\otimes_\mathrm{C}\mathrm{C}'=\mathrm{P}\otimes_\mathrm{A}\mathrm{A}'$, de donde resulta la conclusión.
\end{proof}

\begin{lemma}[11.2.5]
\label{IV.11.2.5}
Sea $\mathrm{A}$ un anillo, $\mathfrak{J}$ un ideal de $\mathrm{A}$,
$\mathrm{B}$ una $\mathrm{A}$-álgebra, $\mathrm{M}$ un
$\mathrm{B}$-módulo y $\mathrm{A}\to\mathrm{A}'$ un homomorfismo de
anillos. Pongamos
$\mathfrak{J}'=\mathfrak{J}\mathrm{A}'$,
$\mathrm{B}'=\mathrm{B}\otimes_\mathrm{A}\mathrm{A}'$, y
$\mathrm{M}'=\mathrm{M}\otimes_\mathrm{A}\mathrm{A}'
=\mathrm{M}\otimes_\mathrm{B}\mathrm{B}'$. Sea
$\mathfrak{p}'$ un ideal primo de $\mathrm{B}'$ que contiene
$\mathfrak{J}'\mathrm{B}'$. Supongamos que se cumple una de las hipótesis
siguientes:
\begin{enumerate}[label=\emph{\alph*)}]
  \item $\mathfrak{J}$ es nilpotente.
  \item $\mathrm{A}'$ es noetheriano, $\mathrm{B}'$ es una $\mathrm{A}'$-álgebra de tipo finito y $\mathrm{M}'$ un $\mathrm{B}'$-módulo de tipo finito.
\end{enumerate}
Bajo estas condiciones, supongamos que se verifican las dos propiedades siguientes:
\begin{enumerate}[label=\emph{(\roman*)}]
  \item $\mathrm{M}\otimes_\mathrm{A}(\mathrm{A}/\mathfrak{J})$ es un $(\mathrm{A}/\mathfrak{J})$-módulo plano.
  \item El homomorfismo canónico compuesto
\[
  \operatorname{Tor}_1^\mathrm{A}(\mathrm{M},\mathrm{A}/\mathfrak{J})\xrightarrow{\psi_1}\operatorname{Tor}_1^{\mathrm{A}'}(\mathrm{M}',\mathrm{A}'/\mathfrak{J}')\xrightarrow{\theta}(\operatorname{Tor}_1^{\mathrm{A}'}(\mathrm{M}',\mathrm{A}'/\mathfrak{J}'))_{\mathfrak{p}'}
\]
(donde $\psi_1$ es el homomorfismo \eref{IV.11.2.1.1} y $\theta$ el homomorfismo canónico de un $\mathrm{B}'$-módulo en su localizado en $\mathfrak{p}'$) es nulo.
\end{enumerate}

Entonces $\mathrm{M}'_{\mathfrak{p}'}$ es un $\mathrm{A}'$-módulo plano.
\end{lemma}

\begin{proof}
Bajo la hipótesis b), $\mathrm{M}'_{\mathfrak{p}'}$ es un
$\mathrm{B}'_{\mathfrak{p}'}$-módulo de tipo finito,
$\mathrm{B}'_{\mathfrak{p}'}$ es una
$\mathrm{A}'$-álgebra noetheriana y
$\mathfrak{J}'\mathrm{B}'_{\mathfrak{p}'}$ está contenido en el radical
$\mathfrak{p}'\mathrm{B}'_{\mathfrak{p}'}$ de
$\mathrm{B}'_{\mathfrak{p}'}$. Bajo la hipótesis a),
$\mathfrak{J}'$ es nilpotente. Podemos, pues, aplicar el criterio de planitud
$(\mathbf{0}_{\mathrm{III}},~10.2.1)$ o
$(\mathbf{0}_{\mathrm{III}},~10.2.2)$, según se cumpla a) o b).
En primer lugar,
\[
\begin{split}
  \mathrm{M}'_{\mathfrak{p}'}\otimes_{\mathrm{A}'}
    (\mathrm{A}'/\mathfrak{J}')
  &=(\mathrm{M}'\otimes_{\mathrm{A}'}
    (\mathrm{A}'/\mathfrak{J}'))_{\mathfrak{p}'}\\
  &=\bigl((\mathrm{M}\otimes_\mathrm{A}
    (\mathrm{A}/\mathfrak{J}))
    \otimes_{\mathrm{A}/\mathfrak{J}}
    (\mathrm{A}'/\mathfrak{J}')\bigr)_{\mathfrak{p}'}.
\end{split}
\]
La hipótesis (i) implica, por tanto, en virtud de
$(\mathbf{0}_\mathrm{I},~6.2.1\text{ y }6.3.2)$, que este es un
$(\mathrm{A}'/\mathfrak{J}')$-módulo plano. Queda por demostrar que
\[
  \operatorname{Tor}_1^{\mathrm{A}'}
  (\mathrm{M}'_{\mathfrak{p}'},\mathrm{A}'/\mathfrak{J}')=0.
\]
Por la planitud de $\mathrm{B}'_{\mathfrak{p}'}$ sobre $\mathrm{B}'$, este
módulo es
\[
  \bigl(\operatorname{Tor}_1^{\mathrm{A}'}
  (\mathrm{M}',\mathrm{A}'/\mathfrak{J}')\bigr)_{\mathfrak{p}'}.
\]
La hipótesis (ii) dice que $\theta\circ\varphi_1$ es nulo. Por otra
parte, \sref{IV.11.2.4}, aplicado con
$\mathrm{C}=\mathrm{A}/\mathfrak{J}$ y $\mathrm{N}=\mathrm{C}$,
muestra, teniendo en cuenta (i), que $\varphi_1$ es sobreyectivo
(pues $\mathrm{A}'/\mathfrak{J}'=\mathrm{C}\otimes_\mathrm{A}
\mathrm{A}'$). Por consiguiente, $\theta$ es nulo. Como la imagen por $\theta$
de $\operatorname{Tor}_1^{\mathrm{A}'}
(\mathrm{M}',\mathrm{A}'/\mathfrak{J}')$ genera su localizado como
$\mathrm{B}'_{\mathfrak{p}'}$-módulo, dicho localizado es nulo.
\end{proof}

\begin{theorem}[11.2.6]
\label{IV.11.2.6}
Con la notación de $(8.5.1)$ y $(8.8.1)$, supongamos $\mathrm{S}_\alpha$ cuasicompacto, $\mathrm{X}_\alpha$ y $\mathrm{Y}_\alpha$ de presentación finita sobre $\mathrm{S}_\alpha$; sea $f_\alpha:\mathrm{X}_\alpha\to\mathrm{Y}_\alpha$ un $\mathrm{S}_\alpha$-morfismo y $\sh{F}_\alpha$ un $\sh{O}_{\mathrm{X}_\alpha}$-Módulo cuasicoherente de presentación finita.

\oldpage[IV-3]{124}

\begin{enumerate}[label=\emph{(\roman*)}]
  \item \emph{Sea $x$ un punto de $\mathrm{X}$ y $x_\lambda$ su proyección canónica en $\mathrm{X}_\lambda$. Para que $\sh{F}$ sea $f$-plano en el punto $x$, es necesario y suficiente que exista $\lambda\geqslant\alpha$ tal que $\sh{F}_\lambda$ sea $f_\lambda$-plano en el punto $x_\lambda$.}
  \item \emph{Para que $\sh{F}$ sea $f$-plano, es necesario y suficiente que exista $\lambda\geqslant\alpha$ tal que $\sh{F}_\lambda$ sea $f_\lambda$-plano.}
\end{enumerate}
\end{theorem}

\begin{proof}
Podemos suponer que $\mathrm{S}_\alpha=\mathrm{S}_0$; como $\mathrm{Y}_0$ es de presentación finita sobre $\mathrm{S}_0$, $\mathrm{Y}_0$ es cuasicompacto, y $f_0:\mathrm{X}_0\to\mathrm{Y}_0$ es un morfismo de presentación finita $(1.6.2,~(\mathrm{v}))$, por lo que podemos restringirnos también al caso $\mathrm{S}_0=\mathrm{Y}_0$. Además, por la cuasicompacidad de $\mathrm{S}_0$ y porque el conjunto de índices $\mathrm{L}$ es filtrante, el mismo razonamiento permite restringirse al caso en que $\mathrm{S}_0=\Spec(\mathrm{A}_0)$ es afín. Asimismo, $\mathrm{X}_0$ es también cuasicompacto, y el mismo razonamiento permite suponer $\mathrm{X}_0=\Spec(\mathrm{B}_0)$ afín; se tiene entonces $\sh{F}_0=\widetilde{\mathrm{M}}_0$, donde $\mathrm{M}_0$ es un $\mathrm{B}_0$-módulo de presentación finita, y el enunciado \sref{IV.11.2.6} equivale en este caso al siguiente (por $(\mathbf{0}_\mathrm{I},~6.3.1)$):

\begin{corollary}[11.2.6.1]
\label{IV.11.2.6.1}
Sea $\mathrm{A}_0$ un anillo,
$(\mathrm{A}_\lambda)_{\lambda\in\mathrm{L}}$ un sistema inductivo
filtrante de $\mathrm{A}_0$-álgebras, $\mathrm{B}_0$ una
$\mathrm{A}_0$-álgebra de presentación finita y $\mathrm{M}_0$ un
$\mathrm{B}_0$-módulo de presentación finita. Pongamos
\[
\begin{gathered}
  \mathrm{B}_\lambda
    =\mathrm{B}_0\otimes_{\mathrm{A}_0}\mathrm{A}_\lambda,\qquad
  \mathrm{M}_\lambda
    =\mathrm{M}_0\otimes_{\mathrm{A}_0}\mathrm{A}_\lambda
    =\mathrm{M}_0\otimes_{\mathrm{B}_0}\mathrm{B}_\lambda,\\
  \mathrm{A}=\varinjlim_\lambda\mathrm{A}_\lambda,\qquad
  \mathrm{B}=\varinjlim_\lambda\mathrm{B}_\lambda
    =\mathrm{B}_0\otimes_{\mathrm{A}_0}\mathrm{A},\\
  \mathrm{M}=\varinjlim_\lambda\mathrm{M}_\lambda
    =\mathrm{M}_0\otimes_{\mathrm{A}_0}\mathrm{A}
    =\mathrm{M}_0\otimes_{\mathrm{B}_0}\mathrm{B}.
\end{gathered}
\]
\begin{enumerate}[label=\emph{(\roman*)}]
  \item \emph{Sea $\mathfrak{p}$ un ideal primo de $\mathrm{B}$ y, para todo $\lambda$, sea $\mathfrak{p}_\lambda$ su imagen inversa en $\mathrm{B}_\lambda$. Para que $\mathrm{M}_\mathfrak{p}$ sea un $\mathrm{A}$-módulo plano, es necesario y suficiente que exista $\lambda$ tal que $(\mathrm{M}_\lambda)_{\mathfrak{p}_\lambda}$ sea un $\mathrm{A}_\lambda$-módulo plano.}
  \item \emph{Para que $\mathrm{M}$ sea un $\mathrm{A}$-módulo plano, es necesario y suficiente que exista $\lambda$ tal que $\mathrm{M}_\lambda$ sea un $\mathrm{A}_\lambda$-módulo plano.}
\end{enumerate}
\end{corollary}

Solo hay que demostrar que las condiciones son necesarias $(2.1.4)$. Procederemos en varias etapas.

\medskip
\noindent\textbf{I)} \emph{Reducción al caso en que los
$\mathrm{A}_\lambda$ son noetherianos.} --- Por $(8.9.1)$ existen un
subanillo $\mathrm{A}'_0\subset\mathrm{A}_0$ que es una
$\mathbf{Z}$-álgebra de tipo finito, una $\mathrm{A}'_0$-álgebra de tipo finito
$\mathrm{B}'_0$ y un $\mathrm{B}'_0$-módulo de tipo finito
$\mathrm{M}'_0$, tales que
\[
  \mathrm{B}_0=\mathrm{B}'_0\otimes_{\mathrm{A}'_0}\mathrm{A}_0,
  \qquad
  \mathrm{M}_0=\mathrm{M}'_0\otimes_{\mathrm{A}'_0}\mathrm{A}_0.
\]
Como
$\mathrm{B}_\lambda=\mathrm{B}'_0\otimes_{\mathrm{A}'_0}
\mathrm{A}_\lambda$ y
$\mathrm{M}_\lambda=\mathrm{M}'_0\otimes_{\mathrm{A}'_0}
\mathrm{A}_\lambda$, podemos sustituir
$\mathrm{A}_0,\mathrm{B}_0,\mathrm{M}_0$ por
$\mathrm{A}'_0,\mathrm{B}'_0,\mathrm{M}'_0$, considerando los
$\mathrm{A}_\lambda$ como $\mathrm{A}'_0$-álgebras. Podemos, por tanto,
suponer ya que $\mathrm{A}_0$ es noetheriano.

Sea $\mathrm{H}$ el conjunto de los pares $(\lambda,\mathrm{C}_\lambda)$,
donde $\mathrm{C}_\lambda$ es una $\mathrm{A}_0$-subálgebra de tipo finito
de $\mathrm{A}_\lambda$. Ordenemos $\mathrm{H}$ declarando
$(\lambda,\mathrm{C}_\lambda)\leqslant
(\mu,\mathrm{D}_\mu)$ si $\lambda\leqslant\mu$ y
$\varphi_{\mu\lambda}(\mathrm{C}_\lambda)\subset\mathrm{D}_\mu$.
Así $\mathrm{H}$ es filtrante: dos pares están dominados eligiendo
$\nu\geqslant\lambda,\mu$ y tomando como $\mathrm{E}_\nu$ la
$\mathrm{A}_0$-subálgebra engendrada por
$\varphi_{\nu\lambda}(\mathrm{C}_\lambda)$ y
$\varphi_{\nu\mu}(\mathrm{D}_\mu)$. Para
$\xi=(\lambda,\mathrm{C}_\lambda)$ pongamos
$\mathrm{A}_\xi=\mathrm{C}_\lambda$; las aplicaciones de transición restringidas
evidentes forman un sistema inductivo filtrante de $\mathrm{A}_0$-álgebras.
Pongamos
$\mathrm{B}_\xi=\mathrm{B}_0\otimes_{\mathrm{A}_0}\mathrm{A}_\xi$
y
$\mathrm{M}_\xi=\mathrm{M}_0\otimes_{\mathrm{A}_0}\mathrm{A}_\xi$.
Los anillos $\mathrm{A}_\xi$ son noetherianos, y la fórmula del doble límite
inductivo (Bourbaki, \emph{Alg.}, cap.~II, 3.ª ed., \S\,6, n.º~4,
Prop.~7) da
\[
  \varinjlim_{\mathrm H}\mathrm{A}_\xi=\mathrm{A},\qquad
  \varinjlim_{\mathrm H}\mathrm{B}_\xi=\mathrm{B},\qquad
  \varinjlim_{\mathrm H}\mathrm{M}_\xi=\mathrm{M}.
\]
Supongamos demostrado \sref{IV.11.2.6.1} para el sistema
$(\mathrm{A}_\xi)_{\xi\in\mathrm H}$, y sea $\mathfrak p$ un ideal primo
de $\mathrm B$ tal que $\mathrm M_{\mathfrak p}$ es plano sobre
$\mathrm A$. Existe entonces $\xi\in\mathrm H$ tal que
\oldpage[IV-3]{125}
$(\mathrm{M}_\xi)_{\mathfrak{p}_\xi}$ es plano sobre
$\mathrm{A}_\xi$, donde $\mathfrak p_\xi$ es la imagen inversa de
$\mathfrak p$. Escribamos $\xi=(\lambda,\mathrm{C}_\lambda)$. Como
$\mathfrak{p}_\lambda$ es la imagen inversa de $\mathfrak p$ en
$\mathrm{B}_\lambda$, $\mathfrak p_\xi$ es la imagen inversa de
$\mathfrak p_\lambda$ en $\mathrm B_\xi$, y
\[
  (\mathrm{M}_\lambda)_{\mathfrak p_\lambda}
  =(\mathrm{M}_\xi\otimes_{\mathrm{C}_\lambda}
    \mathrm{A}_\lambda)_{\mathfrak p_\lambda}
  =\bigl((\mathrm{M}_\xi)_{\mathfrak p_\xi}
    \otimes_{\mathrm{C}_\lambda}\mathrm{A}_\lambda
    \bigr)_{\mathfrak p_\lambda}.
\]
Así, $(\mathrm{M}_\lambda)_{\mathfrak p_\lambda}$ es plano sobre
$\mathrm{A}_\lambda$ por
$(\mathbf{0}_\mathrm{I},~6.2.1\text{ y }6.3.2)$. El enunciado (ii) se
trata de la misma manera. En lo sucesivo podemos, pues, suponer que todos los
$\mathrm{A}_\lambda$ son noetherianos, aunque el propio $\mathrm A$ no tenga
por qué serlo.

\medskip
\noindent\textbf{II)} \emph{Reducción del enunciado global (ii) al enunciado puntual (i).} --- Supongamos que $\sh{F}$ es $f$-plano. Para todo $\lambda$, sea $\mathrm{U}_\lambda$ el conjunto de los $x_\lambda\in\mathrm{X}_\lambda$ tales que $\sh{F}_\lambda$ es $f_\lambda$-plano en el punto $x_\lambda$; sabemos que $\mathrm{U}_\lambda$ es abierto en $\mathrm{X}_\lambda$, pues $\mathrm{S}_\lambda$ es noetheriano y $f_\lambda$ es de tipo finito \sref{IV.11.1.1}; sea $\mathrm{V}_\lambda=v_\lambda^{-1}(\mathrm{U}_\lambda)$ su imagen inversa en $\mathrm{X}$. Por hipótesis, para todo $x\in\mathrm{X}$, la proyección de $x$ en $\mathrm{X}_\lambda$ pertenece a $\mathrm{U}_\lambda$ para algún $\lambda$, lo que significa que $\mathrm{X}=\bigcup\mathrm{V}_\lambda$. Además, por $(2.1.4)$, si $\lambda\leqslant\mu$, se tiene $\mathrm{V}_\lambda\subset\mathrm{V}_\mu$; por tanto, como $\mathrm{X}$ es cuasicompacto, existe un índice $\mu$ tal que $\mathrm{X}=\mathrm{V}_\mu$. Como los $\mathrm{X}_\lambda$ son cuasicompactos, se deduce de $(8.3.4)$ que existe $\nu\geqslant\mu$ tal que $v_{\mu\nu}^{-1}(\mathrm{U}_\mu)=\mathrm{X}_\nu$; pero, por $(2.1.4)$, esto implica que $\sh{F}_\nu$ es $f_\nu$-plano.

\medskip
\noindent\textbf{III)} \emph{Fin de la demostración.} --- Queda por demostrar
(i), con $\mathrm S_0$ afín y noetheriano. Sea $y_0$ la
proyección de $x$ en $\mathrm S_0$. Por $(2.1.4)$ y
$(\mathbf I,~3.6.5)$ podemos sustituir $\mathrm S_0$ por
$\Spec(\sh O_{\mathrm S_0,y_0})$; podemos, pues, suponer que
$\mathrm A_0$ es un anillo local noetheriano con ideal maximal
$\mathfrak m_0=\mathfrak j_{y_0}$. Por definición,
$\mathfrak m_0$ es la imagen inversa del ideal primo
$\mathfrak p=\mathfrak j_x$ de $\mathrm B$, y
$\mathrm M_{\mathfrak p}$ es plano sobre $\mathrm A$. Por tanto,
\[
  \operatorname{Tor}_1^\mathrm A
  (\mathrm A/\mathfrak m_0\mathrm A,\mathrm M_{\mathfrak p})=0.
\]
Como los módulos Tor considerados son $\mathrm B$-módulos y
$\mathrm B_{\mathfrak p}$ es plano sobre $\mathrm B$, esto equivale
a
\[
  \bigl(\operatorname{Tor}_1^\mathrm A
  (\mathrm A/\mathfrak m_0\mathrm A,\mathrm M)\bigr)_{\mathfrak p}=0.
\]

El $\mathrm B_0$-módulo
$\operatorname{Tor}_1^{\mathrm A_0}
(\mathrm A_0/\mathfrak m_0,\mathrm M_0)$ es de tipo finito. En efecto, puede
calcularse a partir de una resolución de $\mathrm A_0/\mathfrak m_0$ mediante
$\mathrm A_0$-módulos libres finitos; al tomar el producto tensorial con $\mathrm M_0$ se obtiene un
complejo de $\mathrm B_0$-módulos de tipo finito, cuya homología es también
de tipo finito porque $\mathrm B_0$ es noetheriano. Elijamos generadores
$t_i^0$ $(1\leqslant i\leqslant m)$ de este módulo, y sean $t_i$
sus imágenes canónicas en
$\operatorname{Tor}_1^\mathrm A
(\mathrm A/\mathfrak m_0\mathrm A,\mathrm M)$. La igualdad de localización
precedente proporciona un elemento
$h\in\mathrm B-\mathfrak p$ tal que $ht_i=0$ para todo $i$.

Por \eref{IV.11.2.2.1},
\[
  \operatorname{Tor}_1^\mathrm A
  (\mathrm A/\mathfrak m_0\mathrm A,\mathrm M)
  =\varinjlim_\lambda
  \operatorname{Tor}_1^{\mathrm A_\lambda}
  (\mathrm A_\lambda/\mathfrak m_0\mathrm A_\lambda,
   \mathrm M_\lambda).
\]
Por consiguiente, para algún $\lambda$ existe un elemento
$h_\lambda\in\mathrm B_\lambda$ cuya imagen es $h$ tal que
$h_\lambda t_i^\lambda=0$ para todo $i$, donde $t_i^\lambda$ es la
imagen de $t_i^0$. Si $\mathfrak p_\lambda$ es la imagen inversa de
$\mathfrak p$ en $\mathrm B_\lambda$, entonces
$h_\lambda\notin\mathfrak p_\lambda$. Así, las imágenes de todos los
$t_i^0$ en
\[
  \bigl(\operatorname{Tor}_1^{\mathrm A_\lambda}
  (\mathrm A_\lambda/\mathfrak m_0\mathrm A_\lambda,
   \mathrm M_\lambda)\bigr)_{\mathfrak p_\lambda}
\]
son nulas, y la aplicación canónica de
$\operatorname{Tor}_1^{\mathrm A_0}
(\mathrm A_0/\mathfrak m_0,\mathrm M_0)$ a este localizado es nula.
Se satisfacen, por tanto, las hipótesis de \sref{IV.11.2.5}; la condición
(i) es automática porque $\mathrm A_0/\mathfrak m_0$ es un cuerpo.
Se sigue que $(\mathrm M_\lambda)_{\mathfrak p_\lambda}$ es plano
sobre $\mathrm A_\lambda$, lo que concluye la demostración.

\begin{corollary}[11.2.7]
\label{IV.11.2.7}
Sea $\mathrm{S}=\Spec(\mathrm{A})$ un esquema afín, $f:\mathrm{X}\to\mathrm{S}$ un morfismo, $\sh{F}$ un $\sh{O}_\mathrm{X}$-Módulo cuasicoherente y $x$ un punto de $\mathrm{X}$. Las condiciones siguientes son equivalentes:
\begin{enumerate}[label=\emph{\alph*)}]
  \item $f$ es un morfismo de presentación finita, $\sh{F}$ es un $\sh{O}_\mathrm{X}$-Módulo de presentación finita y $\sh{F}$ es $f$-plano en el punto $x$ (resp.\ $f$-plano).
\oldpage[IV-3]{126}
  \item Existe un esquema afín noetheriano $\mathrm{S}_0=\Spec(\mathrm{A}_0)$, un morfismo de tipo finito $f_0:\mathrm{X}_0\to\mathrm{S}_0$, un $\sh{O}_{\mathrm{X}_0}$-Módulo coherente $\sh{F}_0$ y un morfismo $\mathrm{S}\to\mathrm{S}_0$ tales que el S-preesquema $\mathrm{X}_0\otimes_{\mathrm{S}_0}\mathrm{S}$ es S-isomorfo a $\mathrm{X}$, que $\sh{F}$ es isomorfo a $\sh{F}_0\otimes_{\sh{O}_{\mathrm{S}_0}}\sh{O}_\mathrm{S}$ y que $\sh{F}_0$ es $f_0$-plano en el punto $x_0$, proyección de $x$ en $\mathrm{X}_0$ (resp.\ $f_0$-plano).
  \item Se satisfacen las condiciones de $b)$ y, además, $\mathrm{A}_0$ es una sub-$\mathbf{Z}$-álgebra de tipo finito de $\mathrm{A}$, correspondiendo el morfismo $\mathrm{S}\to\mathrm{S}_0$ a la inyección canónica $\mathrm{A}_0\to\mathrm{A}$.
\end{enumerate}
\end{corollary}

\begin{proof}
Es claro que $c)$ implica $b)$ y que $b)$ implica $a)$ por
$(2.1.4)$. Por otra parte, podemos considerar $\mathrm{A}$ como el límite
inductivo filtrante de sus $\mathbf{Z}$-subálgebras de tipo finito, y por $(8.9.1)$
existen una subálgebra de este tipo $\mathrm{A}_0$ y un morfismo de tipo finito
$f_0:\mathrm{X}_0\to\mathrm{S}_0$ tales que $\mathrm{X}$ es
$\mathrm{S}$-isomorfo a
$\mathrm{X}_0\times_{\mathrm{S}_0}\mathrm{S}$ y $\sh{F}$ es isomorfo
a $\sh{F}_0\otimes_{\sh{O}_{\mathrm{S}_0}}\sh{O}_{\mathrm{S}}$. Podemos,
por tanto, escribir
\[
  \mathrm{X}=\varprojlim_\lambda \mathrm{X}_\lambda,\qquad
  \mathrm{X}_\lambda=\mathrm{X}_0\otimes_{\mathrm{A}_0}\mathrm{A}_\lambda,
\]
donde $\mathrm{A}_\lambda$ recorre las
$\mathbf{Z}$-subálgebras de tipo finito de $\mathrm{A}$ que contienen $\mathrm{A}_0$, y
$\sh{F}=v_\lambda^*(\sh{F}_\lambda)$, donde
$v_\lambda:\mathrm{X}\to\mathrm{X}_\lambda$ es la proyección canónica y
$\sh{F}_\lambda=\sh{F}_0\otimes_{\mathrm{A}_0}\mathrm{A}_\lambda$. Se
deduce de \sref{IV.11.2.6} que existe $\lambda$ tal que
$\sh{F}_\lambda$ es $f_\lambda$-plano en
$x_\lambda=v_\lambda(x)$ (resp.\ $f_\lambda$-plano); el subanillo
$\mathrm{A}_\lambda$ de $\mathrm{A}$ satisface entonces la condición $c)$.
\end{proof}

\begin{proposition}[11.2.8]
\label{IV.11.2.8}
Sean $g:\mathrm{X}\to\mathrm{S}$ y $h:\mathrm{Y}\to\mathrm{S}$ dos morfismos de presentación finita, $f:\mathrm{X}\to\mathrm{Y}$ un S-morfismo y $\sh{F}$ un $\sh{O}_\mathrm{X}$-Módulo cuasicoherente de presentación finita; para todo $s\in\mathrm{S}$, sean $\mathrm{X}_s=g^{-1}(s)=\mathrm{X}\times_\mathrm{S}\Spec(\boldsymbol{k}(s))$, $\mathrm{Y}_s=h^{-1}(s)=\mathrm{Y}\times_\mathrm{S}\Spec(\boldsymbol{k}(s))$, $f_s$ el morfismo $f\times 1:\mathrm{X}_s\to\mathrm{Y}_s$ y $\sh{F}_s=\sh{F}\otimes_{\sh{O}_\mathrm{S}}\boldsymbol{k}(s)$. Entonces el conjunto $\mathrm{E}$ de los $s\in\mathrm{S}$ tales que $\sh{F}_s$ es $f_s$-plano es localmente construible.
\end{proposition}

\begin{proof}
La propiedad cuya constructibilidad queremos demostrar satisface la condición
$(9.2.1,~(\mathrm{i}))$, por $(2.2.13)$ y $(2.5.1)$. En virtud de
$(9.2.3)$, podemos, pues, restringirnos al caso en que $\mathrm{S}$ es
afín, noetheriano e íntegro, con punto genérico $\eta$, y queda por
demostrar que $\mathrm{E}$ o $\mathrm{S}-\mathrm{E}$ es un entorno
de $\eta$ en $\mathrm{S}$. Si
$\eta\in\mathrm{S}-\mathrm{E}$, esto resulta inmediatamente del Lema
$(9.4.7.1)$. Si, por el contrario, $\eta\in\mathrm{E}$, se deduce primero
de \sref{IV.11.2.6}, aplicado mediante el procedimiento de $(8.1.2,~a))$, que existe
un entorno abierto $\mathrm{U}$ de $\eta$ en $\mathrm{S}$ tal que
$\sh{F}|_{g^{-1}(\mathrm{U})}$ es plano sobre $h^{-1}(\mathrm{U})$; \emph{a
fortiori}, por $(2.1.4)$, $\sh{F}_s$ es $f_s$-plano para todo
$s\in\mathrm{U}$, lo que concluye la demostración.
\end{proof}

El teorema siguiente generaliza $(11.2.6.1,~(\mathrm{ii}))$.

\begin{theorem}[11.2.9]
\label{IV.11.2.9}
(Raynaud). --- Sea $\mathrm{A}_0$ un anillo, $(\mathrm{A}_\lambda)_{\lambda\in\mathrm{L}}$ un sistema inductivo filtrante de $\mathrm{A}_0$-álgebras, $\mathrm{B}_0$ una $\mathrm{A}_0$-álgebra de presentación finita, $\mathfrak{J}_0$ un ideal de tipo finito de $\mathrm{B}_0$ y $\mathrm{M}_0$ un $\mathrm{B}_0$-módulo de presentación finita. Pongamos $\mathrm{B}_\lambda=\mathrm{B}_0\otimes_{\mathrm{A}_0}\mathrm{A}_\lambda$, $\mathfrak{J}_\lambda=\mathfrak{J}_0\mathrm{B}_\lambda$, $\mathrm{M}_\lambda=\mathrm{M}_0\otimes_{\mathrm{A}_0}\mathrm{A}_\lambda=\mathrm{M}_0\otimes_{\mathrm{B}_0}\mathrm{B}_\lambda$, $\mathrm{A}=\varinjlim\mathrm{A}_\lambda$, $\mathrm{B}=\varinjlim\mathrm{B}_\lambda=\mathrm{B}_0\otimes_{\mathrm{A}_0}\mathrm{A}$, $\mathfrak{J}=\mathfrak{J}_0\mathrm{B}$ y $\mathrm{M}=\varinjlim\mathrm{M}_\lambda=\mathrm{M}_0\otimes_{\mathrm{A}_0}\mathrm{A}=\mathrm{M}_0\otimes_{\mathrm{B}_0}\mathrm{B}$. Para que $\operatorname{gr}^*_\mathfrak{J}(\mathrm{M})$ sea un $\mathrm{A}$-módulo plano, es necesario y suficiente que exista $\lambda$ tal que $\operatorname{gr}^*_{\mathfrak{J}_\lambda}(\mathrm{M}_\lambda)$ sea un $\mathrm{A}_\lambda$-módulo plano; los homomorfismos canónicos
\begin{equation}
  \begin{cases}
    \operatorname{gr}^*_{\mathfrak{J}_\lambda}(\mathrm{M}_\lambda)\otimes_{\mathrm{A}_\lambda}\mathrm{A} \to \operatorname{gr}^*_\mathfrak{J}(\mathrm{M}) \\
    \operatorname{gr}^*_{\mathfrak{J}_\lambda}(\mathrm{M}_\lambda)\otimes_{\mathrm{A}_\lambda}\mathrm{A}_\mu \to \operatorname{gr}^*_{\mathfrak{J}_\mu}(\mathrm{M}_\mu)
  \end{cases}
  \qquad\text{para }\lambda\leqslant\mu
  \label{IV.11.2.9.1}
\tag{11.2.9.1}
\end{equation}
son entonces biyectivos y $\operatorname{gr}^*_{\mathfrak{J}_\mu}(\mathrm{M}_\mu)$ es un $\mathrm{A}_\mu$-módulo plano para $\lambda\leqslant\mu$.
\end{theorem}

\medskip
\noindent\emph{Demostración.}
Demostremos primero que las condiciones son suficientes, lo que equivale a probar
la biyectividad de $(11.2.9.1)$. Esto se deduce del lema siguiente:

\begin{lemma}[11.2.9.2]
\label{IV.11.2.9.2}
Sea $\mathrm{A}$ un anillo, $\mathrm{B}$ una $\mathrm{A}$-álgebra, $\mathrm{M}$ un $\mathrm{B}$-módulo y $\mathfrak{J}$ un ideal de $\mathrm{B}$. Sea $\mathrm{A}'$ una $\mathrm{A}$-álgebra; pongamos $\mathrm{B}'=\mathrm{B}\otimes_\mathrm{A}\mathrm{A}'$, $\mathrm{M}'=\mathrm{M}\otimes_\mathrm{A}\mathrm{A}'=\mathrm{M}\otimes_\mathrm{B}\mathrm{B}'$ y $\mathfrak{J}'=\mathfrak{J}\mathrm{B}'$. Si $\operatorname{gr}^*_\mathfrak{J}(\mathrm{M})$ es un $\mathrm{A}$-módulo plano, el homomorfismo canónico
\[
  (\operatorname{gr}^*_\mathfrak{J}(\mathrm{M}))\otimes_\mathrm{A}\mathrm{A}'\to\operatorname{gr}^*_{\mathfrak{J}'}(\mathrm{M}')
\]
\emph{es bijetivo.}
\end{lemma}

\oldpage[IV-3]{127}

En efecto, por inducción sobre $k$, la hipótesis de que los $\mathfrak{J}^k\mathrm{M}/\mathfrak{J}^{k+1}\mathrm{M}$ son A-módulos planos implica primero, por $(\mathbf{0}_\mathrm{I},~6.1.2)$, que $\mathrm{M}/\mathfrak{J}^{k+1}\mathrm{M}$ es un A-módulo plano; además, por $(\mathbf{0}_\mathrm{I},~6.1.2)$, la sucesión
\[
  0\to(\mathfrak{J}^{k+1}\mathrm{M})\otimes_\mathrm{A}\mathrm{A}'\to\mathrm{M}\otimes_\mathrm{A}\mathrm{A}'\to(\mathrm{M}/\mathfrak{J}^{k+1}\mathrm{M})\otimes_\mathrm{A}\mathrm{A}'\to 0
\]
es entonces exacta; en otras palabras, $(\mathfrak{J}^{k+1}\mathrm{M})\otimes_\mathrm{A}\mathrm{A}'$ se identifica con su imagen canónica $\mathfrak{J}'^{k+1}\mathrm{M}'$ en $\mathrm{M}'=\mathrm{M}\otimes_\mathrm{A}\mathrm{A}'$. Por otra parte, siempre por $(\mathbf{0}_\mathrm{I},~6.1.2)$, la sucesión
\[
  0\to(\mathfrak{J}^{k+1}\mathrm{M})\otimes_\mathrm{A}\mathrm{A}'\to(\mathfrak{J}^k\mathrm{M})\otimes_\mathrm{A}\mathrm{A}'\to(\mathfrak{J}^k\mathrm{M}/\mathfrak{J}^{k+1}\mathrm{M})\otimes_\mathrm{A}\mathrm{A}'\to 0
\]
es exacta, lo que demuestra el lema.

Para demostrar la necesidad de las condiciones de \sref{IV.11.2.9}, procedemos en varios pasos.

\medskip
\noindent\textbf{I)} \emph{Reducción al caso en que los $\mathrm{A}_\lambda$
son noetherianos.} --- Se procede como en la reducción I) de
\sref{IV.11.2.6.1}, conservando la notación; basta con empezar por sustituir
$\mathrm{A}'_0$ por una $\mathrm{A}'_0$-subálgebra de tipo finito
$\mathrm{A}''_0$ de $\mathrm{A}_0$ tal que, si se pone
$\mathrm{B}''_0=\mathrm{B}'_0\otimes_{\mathrm{A}'_0}\mathrm{A}''_0$, hay
en $\mathrm{B}''_0$ un ideal finitamente generado $\mathfrak{J}''_0$ tal que
$\mathfrak{J}''_0\mathrm{B}_0=\mathfrak{J}_0$. Para ello, consideremos las
$\mathrm{A}'_0$-subálgebras de tipo finito $\mathrm{A}'_\alpha$ de
$\mathrm{A}_0$, que forman una familia filtrante; se tiene
$\mathrm{B}_0=\varinjlim\mathrm{B}'_\alpha$, donde
$\mathrm{B}'_\alpha=\mathrm{B}'_0\otimes_{\mathrm{A}'_0}\mathrm{A}'_\alpha$.
Por tanto, hay un índice $\beta$ tal que un conjunto finito de generadores de
$\mathfrak{J}_0$ está formado por las imágenes en $\mathrm{B}_0$ de elementos de
$\mathrm{B}'_\beta$. Tomemos $\mathrm{A}''_0=\mathrm{A}'_\beta$,
$\mathrm{B}''_0=\mathrm{B}'_\beta$, y como $\mathfrak{J}''_0$ el ideal
generado por estos elementos. Podemos, por tanto, suponer que $\mathrm{A}_0$
(y, por consiguiente, también $\mathrm{B}_0$) es noetheriano. Se definen entonces, como en
\emph{loc.\ cit.}, el conjunto filtrado $\mathrm{H}$ y los $\mathrm{A}_\xi$,
$\mathrm{B}_\xi$, $\mathrm{M}_\xi$ para $\xi\in\mathrm{H}$; pongamos también
$\mathfrak{J}_\xi=\mathfrak{J}''_0\mathrm{B}_\xi$ para todo
$\xi\in\mathrm{H}$. Supongamos demostrado que existe
$\xi=(\lambda,\mathrm{C}_\lambda)$ tal que
$\operatorname{gr}^*_{\mathfrak{J}_\xi}(\mathrm{M}_\xi)$ es un
$\mathrm{C}_\lambda$-módulo plano. Como
$\mathrm{M}_\lambda=\mathrm{M}_\xi\otimes_{\mathrm{C}_\lambda}
\mathrm{A}_\lambda$, se deduce de \sref{IV.11.2.9.2} que
$\operatorname{gr}^*_{\mathfrak{J}_\lambda}(\mathrm{M}_\lambda)$ es un
$\mathrm{A}_\lambda$-módulo plano.

\medskip
\noindent\textbf{II)} \emph{Lemas preliminares.}

\begin{lemma}[11.2.9.3]
\label{IV.11.2.9.3}
Sea $\mathrm{A}$ un anillo, sea
$\mathrm{B}=\bigoplus_{i\geqslant 0}\mathrm{B}_i$ una
$\mathrm{A}$-álgebra graduada en grados no negativos y sea
$\mathrm{M}=\bigoplus_{i\in\mathbf{Z}}\mathrm{M}_i$ un
$\mathrm{B}$-módulo graduado. Supongamos que $\mathrm{B}_0$ es un anillo local y
que cada $\mathrm{M}_i$ es un $\mathrm{B}_0$-módulo de tipo finito.
Para que $\mathrm{M}$ sea un $\mathrm{B}$-módulo de tipo finito, es necesario
y suficiente que, si $\mathfrak{q}$ es la imagen inversa en
$\mathrm{A}$ del ideal maximal $\mathfrak{m}$ de $\mathrm{B}_0$,
$\mathrm{M}\otimes_\mathrm{A}\boldsymbol{k}(\mathfrak{q})$ sea un
$(\mathrm{B}\otimes_\mathrm{A}\boldsymbol{k}(\mathfrak{q}))$-módulo
de tipo finito.
\end{lemma}

\oldpage[IV-3]{128}

\begin{proof}
La condición es evidentemente necesaria; demostremos que es
suficiente. Como $\mathrm{B}_0$ (y, por tanto, $\mathrm{B}$) es una
$\mathrm{A}_{\mathfrak{q}}$-álgebra, podemos sustituir $\mathrm{A}$ por
$\mathrm{A}_{\mathfrak{q}}$; en otras palabras, podemos suponer que
$\mathrm{A}$ es también un anillo local con ideal maximal
$\mathfrak{q}$. Por hipótesis, hay un entero $\mathrm{N}$ tal
que el homomorfismo canónico
\[
 \left(\bigoplus_{-\mathrm{N}\leqslant i\leqslant\mathrm{N}}
   \mathrm{M}_i\otimes_{\mathrm{B}_0}\mathrm{B}\right)
   \otimes_\mathrm{A}\boldsymbol{k}(\mathfrak{q})
 \longrightarrow
 \mathrm{M}\otimes_\mathrm{A}\boldsymbol{k}(\mathfrak{q})
\]
es sobreyectivo. Esto significa también que, para todo entero $n$, el
homomorfismo canónico de $\boldsymbol{k}(\mathfrak{q})$-módulos
\[
 \left(\bigoplus_{-\mathrm{N}\leqslant i\leqslant\mathrm{N}}
   \mathrm{M}_i\otimes_{\mathrm{B}_0}\mathrm{B}_{n-i}\right)
   \otimes_\mathrm{A}\boldsymbol{k}(\mathfrak{q})
 \longrightarrow
 \mathrm{M}_n\otimes_\mathrm{A}\boldsymbol{k}(\mathfrak{q})
\]
es sobreyectivo. Ahora bien, $\mathrm{M}_n$ es un
$\mathrm{B}_0$-módulo de tipo finito, $\mathrm{B}_0$ es local y
$\mathfrak{q}\mathrm{B}_0\subset\mathfrak{m}$; el lema de Nakayama
muestra, por tanto, que cada homomorfismo canónico
\[
 \bigoplus_{-\mathrm{N}\leqslant i\leqslant\mathrm{N}}
   \mathrm{M}_i\otimes_{\mathrm{B}_0}\mathrm{B}_{n-i}
 \longrightarrow \mathrm{M}_n
\]
es sobreyectivo, lo que demuestra la afirmación.
\end{proof}

\begin{corollary}[11.2.9.4]
\label{IV.11.2.9.4}
Bajo las hipótesis de \sref{IV.11.2.9.3}, supongamos además que cada $\mathrm{B}_i$ y cada $\mathrm{M}_i$ es un $\mathrm{B}_0$-módulo de presentación finita, y que $\mathrm{M}$ es un $\mathrm{A}$-módulo plano. Para que $\mathrm{M}$ sea un $\mathrm{B}$-módulo de presentación finita, es necesario y suficiente que $\mathrm{M}\otimes_\mathrm{A}\boldsymbol{k}(\mathfrak{q})$ sea un $(\mathrm{B}\otimes_\mathrm{A}\boldsymbol{k}(\mathfrak{q}))$-módulo de presentación finita.
\end{corollary}

\begin{proof}
Por \sref{IV.11.2.9.3}, si se cumple la condición del enunciado,
$\mathrm{M}$ es un $\mathrm{B}$-módulo de tipo finito. Existe, por tanto, un
$\mathrm{B}$-módulo graduado libre de tipo finito $\mathrm{L}$ (que posee,
pues, una base finita de elementos homogéneos) y un homomorfismo graduado
sobreyectivo de grado $0$,
\[
 u:\mathrm{L}\longrightarrow\mathrm{M}.
\]
Sea $\mathrm{R}=\operatorname{Ker}(u)$, que es un
$\mathrm{B}$-módulo graduado. Hay un número finito de enteros $m_j$
($1\leqslant j\leqslant r$) tales que, para todo entero $i$,
$\mathrm{R}_i$ es el núcleo de un homomorfismo sobreyectivo
\[
 \bigoplus_{1\leqslant j\leqslant r}\mathrm{B}_{i+m_j}
 \longrightarrow \mathrm{M}_i.
\]
Las hipótesis sobre los $\mathrm{M}_i$ y los $\mathrm{B}_i$ implican entonces
que $\mathrm{R}_i$ es un $\mathrm{B}_0$-módulo de tipo finito
(Bourbaki, \emph{Alg.\ comm.}, cap.~I,
\textsection\,2, n\textsuperscript{o}\,8, lema~9). Para demostrar que
$\mathrm{R}$ es un $\mathrm{B}$-módulo de tipo finito, observemos que,
por la hipótesis de planitud y $(\mathbf{0}_\mathrm{I},~6.1.2)$, la
sucesión
\[
 0\longrightarrow
 \mathrm{R}\otimes_\mathrm{A}\boldsymbol{k}(\mathfrak{q})
 \longrightarrow
 \mathrm{L}\otimes_\mathrm{A}\boldsymbol{k}(\mathfrak{q})
 \longrightarrow
 \mathrm{M}\otimes_\mathrm{A}\boldsymbol{k}(\mathfrak{q})
 \longrightarrow 0
\]
es exacta. La hipótesis implica, por tanto (Bourbaki,
\emph{loc.\ cit.}) que
$\mathrm{R}\otimes_\mathrm{A}\boldsymbol{k}(\mathfrak{q})$ es un
$(\mathrm{B}\otimes_\mathrm{A}\boldsymbol{k}(\mathfrak{q}))$-módulo
de tipo finito. Basta, por tanto, aplicar \sref{IV.11.2.9.3} a $\mathrm{R}$.
\end{proof}

\medskip
\noindent\textbf{III)} \emph{Reducción al caso en que los homomorfismos de transición $\varphi_{\mu\lambda}:\mathrm{A}_\lambda\to\mathrm{A}_\mu$ (para $\lambda\leqslant\mu$) son inyectivos.} --- Sea $\mathrm{A}'_\lambda$ la imagen de $\mathrm{A}_\lambda$ por el homomorfismo canónico $\varphi_\lambda:\mathrm{A}_\lambda\to\mathrm{A}$; es claro que los $\mathrm{A}'_\lambda$ forman un sistema inductivo de subanillos \emph{noetherianos} de $\mathrm{A}$ y que los homomorfismos de transición $\mathrm{A}'_\lambda\to\mathrm{A}'_\mu$ (para $\lambda\leqslant\mu$) son inyectivos. Pongamos $\mathrm{B}'_\lambda=\mathrm{B}_0\otimes_{\mathrm{A}_0}\mathrm{A}'_\lambda$, $\mathfrak{J}'_\lambda=\mathfrak{J}_0\mathrm{B}'_\lambda$ y $\mathrm{M}'_\lambda=\mathrm{M}_0\otimes_{\mathrm{A}_0}\mathrm{A}'_\lambda=\mathrm{M}_\lambda\otimes_{\mathrm{A}_\lambda}\mathrm{A}'_\lambda$; se sigue teniendo $\mathrm{B}=\varinjlim\mathrm{B}'_\lambda$ y $\mathrm{M}=\varinjlim\mathrm{M}'_\lambda$. Supongamos demostrado que existe $\lambda$ tal que, para $\mu\geqslant\lambda$, $\operatorname{gr}^*_{\mathfrak{J}'_\mu}(\mathrm{M}'_\mu)$ es un $\mathrm{A}'_\mu$-módulo plano. Sea $a_\lambda$ el núcleo del homomorfismo $\varphi_\lambda$, que es un ideal del anillo noetheriano $\mathrm{A}_\lambda$. De la definición del límite inductivo se deduce que existe un índice $\mu\geqslant\lambda$ tal que $\varphi_{\mu\lambda}(a_\lambda)=0$; en otras palabras, $\varphi_{\mu\lambda}$ se factoriza como $\varphi_{\mu\lambda}:\mathrm{A}_\lambda\to\mathrm{A}'_\lambda\to\mathrm{A}_\mu$, y se puede escribir $\mathrm{B}_\mu=\mathrm{B}'_\lambda\otimes_{\mathrm{A}'_\lambda}\mathrm{A}_\mu$, $\mathfrak{J}_\mu=\mathfrak{J}'_\lambda\mathrm{B}_\mu$ y $\mathrm{M}_\mu=\mathrm{M}'_\lambda\otimes_{\mathrm{A}'_\lambda}\mathrm{A}_\mu$; se deduce entonces de \sref{IV.11.2.9.2} que $\operatorname{gr}^*_{\mathfrak{J}_\mu}(\mathrm{M}_\mu)$ es un $\mathrm{A}_\mu$-módulo plano.

\medskip
\noindent\textbf{IV)} \emph{Reducción al caso en que $\mathrm{B}_0=\mathrm{A}_0[\mathrm{T}_1,\ldots,\mathrm{T}_n]$ y $\operatorname{gr}^*_{\mathfrak{J}_0}(\mathrm{B}_0)=\mathrm{B}_0$.} --- En virtud de la hipótesis, existe un sistema $(t_i)_{1\leqslant i\leqslant n}$ de generadores de la $\mathrm{A}_0$-álgebra de tipo finito $\mathrm{B}_0$, tal que los $t_i$ para $i\leqslant m$ generan el ideal $\mathfrak{J}_0$ de $\mathrm{B}_0$. Pongamos $\mathrm{B}'_0=\mathrm{A}_0[\mathrm{T}_1,\ldots,\mathrm{T}_n]$, un álgebra de polinomios (y, por tanto, noetheriana), y sea $\mathfrak{J}'_0$ el ideal de $\mathrm{B}'_0$ generado por los $\mathrm{T}_i$ con $i\leqslant m$; existe entonces un $\mathrm{A}_0$-homomorfismo sobreyectivo $u_0:\mathrm{B}'_0\to\mathrm{B}_0$ que envía cada $\mathrm{T}_i$ a $t_i$ ($1\leqslant i\leqslant n$) y que, por tanto, satisface $u_0(\mathfrak{J}'_0)=\mathfrak{J}_0$. Este homomorfismo permite considerar $\mathrm{M}_0$ como un $\mathrm{B}'_0$-módulo de tipo finito. Pongamos
\[
  \mathrm{B}'_\lambda=\mathrm{B}'_0\otimes_{\mathrm{A}_0}\mathrm{A}_\lambda
  =\mathrm{A}_\lambda[\mathrm{T}_1,\ldots,\mathrm{T}_n],
  \qquad
  \mathfrak{J}'_\lambda=\mathfrak{J}'_0\mathrm{B}'_\lambda,
\]
\oldpage[IV-3]{129}
y también
\[
  \mathrm{B}'=\mathrm{B}'_0\otimes_{\mathrm{A}_0}\mathrm{A}
  =\mathrm{A}[\mathrm{T}_1,\ldots,\mathrm{T}_n],
  \qquad
  \mathfrak{J}'=\mathfrak{J}'_0\mathrm{B}'.
\]
Así, $\mathfrak{J}'$ es el ideal de $\mathrm{B}'$ generado por los $\mathrm{T}_i$ con $i\leqslant m$. Además, para todo entero $k\geqslant 0$ se tiene claramente
\[
  \mathfrak{J}'^k\mathrm{M}=\mathfrak{J}^k\mathrm{M},
  \qquad
  \mathfrak{J}'_\lambda{}^k\mathrm{M}_\lambda
  =\mathfrak{J}_\lambda^k\mathrm{M}_\lambda
\]
para todo $\lambda$. Por tanto,
\[
  \operatorname{gr}^*_{\mathfrak{J}'}(\mathrm{M})
  =\operatorname{gr}^*_{\mathfrak{J}}(\mathrm{M})
\]
como $\mathrm{A}$-módulos, y
\[
  \operatorname{gr}^*_{\mathfrak{J}'_\lambda}(\mathrm{M}_\lambda)
  =\operatorname{gr}^*_{\mathfrak{J}_\lambda}(\mathrm{M}_\lambda)
\]
como $\mathrm{A}_\lambda$-módulos. Podemos, por tanto, sustituir $\mathrm{B}'$, $\mathrm{B}'_\lambda$, $\mathfrak{J}'$ y $\mathfrak{J}'_\lambda$ por $\mathrm{B}$, $\mathrm{B}_\lambda$, $\mathfrak{J}$ y $\mathfrak{J}_\lambda$, respectivamente, en la demostración. Obsérvese también que, por construcción, $\operatorname{gr}^*_{\mathfrak{J}'}(\mathrm{B}')$ se identifica con $\mathrm{B}'$, y $\operatorname{gr}^*_{\mathfrak{J}'_\lambda}(\mathrm{B}'_\lambda)$ con $\mathrm{B}'_\lambda$.

\medskip
\noindent\textbf{V)} \emph{Demostración de que $\operatorname{gr}^*_\mathfrak{J}(\mathrm{M})$ es un $\operatorname{gr}^*_\mathfrak{J}(\mathrm{B})$-módulo de presentación finita.} --- Supondremos, pues, en lo sucesivo, que $\mathrm{A}_0$ y los $\mathrm{A}_\lambda$ son noetherianos, que los homomorfismos de transición $\mathrm{A}_\lambda\to\mathrm{A}_\mu$ son inyectivos, que $\mathrm{B}_0=\mathrm{A}_0[\mathrm{T}_1,\ldots,\mathrm{T}_n]$ y que $\mathfrak{J}_0$ está generado por los $\mathrm{T}_i$ con $i\leqslant m$. El anillo $\mathrm{C}_0=\operatorname{gr}^*_{\mathfrak{J}_0}(\mathrm{B}_0)$ se identifica, por tanto, con $\mathrm{B}_0$, y $\mathrm{C}=\operatorname{gr}^*_\mathfrak{J}(\mathrm{B})$ se identifica con $\mathrm{B}$; al comienzo de lo que sigue solo utilizaremos el hecho de que $\mathrm{C}\cong\mathrm{C}_0\otimes_{\mathrm{A}_0}\mathrm{A}$.

Necesitaremos la siguiente variante de $(6.9.3)$:

\begin{lemma}[11.2.9.5]
\label{IV.11.2.9.5}
Sea $\mathrm{A}_0$ un anillo noetheriano, $\mathrm{B}_0$ una
$\mathrm{A}_0$-álgebra de tipo finito, $\mathfrak{J}_0$ un ideal de
$\mathrm{B}_0$ y $\mathrm{M}_0$ un $\mathrm{B}_0$-módulo de tipo finito.
Existe una sucesión $(\mathrm{S}_{0i})_{1\leqslant i\leqslant m}$ de
subesquemas de $\mathrm{S}_0=\Spec(\mathrm{A}_0)$ con las propiedades
siguientes:
\begin{enumerate}[label=\arabic{enumi}\textsuperscript{o}]
  \item Los espacios subyacentes de los $\mathrm{S}_{0i}$ son disjuntos dos a dos y forman un recubrimiento de $\mathrm{S}_0$.
  \item Para todo $i$, el conjunto $\mathrm{S}_{0i}$ es abierto en $\bigcup_{j\geqslant i}\mathrm{S}_{0j}$.
  \item Cada esquema $\mathrm{S}_{0i}$ es afín.
  \item Si $\mathrm{A}_{0i}$ es el anillo de $\mathrm{S}_{0i}$ y se pone $\mathrm{B}_{0i}=\mathrm{B}_0\otimes_{\mathrm{A}_0}\mathrm{A}_{0i}$, $\mathfrak{J}_{0i}=\mathfrak{J}_0\mathrm{B}_{0i}$, entonces $\operatorname{gr}^*_{\mathfrak{J}_{0i}}(\mathrm{M}_0\otimes_{\mathrm{A}_0}\mathrm{A}_{0i})$ es un $\mathrm{A}_{0i}$-módulo plano.
\end{enumerate}
\end{lemma}

\begin{proof}
Procedamos como en $(6.9.3)$ por inducción noetheriana, suponiendo conocido el lema tras sustituir $\mathrm{A}_0$ por
\[
  \mathrm{A}'_0=\mathrm{A}_0/\mathfrak{a}_0,
\]
donde $\mathfrak{a}_0$ es un ideal tal que $\mathrm{V}(\mathfrak{a}_0)\neq\mathrm{S}_0$, y sustituyendo $\mathrm{B}_0$, $\mathfrak{J}_0$ y $\mathrm{M}_0$ por
\[
  \mathrm{B}'_0=\mathrm{B}_0\otimes_{\mathrm{A}_0}\mathrm{A}'_0,
  \qquad
  \mathfrak{J}'_0=\mathfrak{J}_0\mathrm{B}'_0,
  \qquad
  \mathrm{M}_0\otimes_{\mathrm{A}_0}\mathrm{A}'_0,
\]
respectivamente. Sea $\mathfrak{N}_0$ el nilradical de $\mathrm{A}_0$. Es
evidente que basta demostrar el lema después de sustituir $\mathrm{A}_0$ por
$\mathrm{A}_0/\mathfrak{N}_0$ y $\mathrm{B}_0$,
$\mathfrak{J}_0$ y $\mathrm{M}_0$ por los objetos correspondientes que acabamos
de describir. En otras palabras, podemos suponer que $\mathrm{A}_0$ es
\emph{reducida}. Ahora
$\operatorname{gr}^*_{\mathfrak{J}_0}(\mathrm{B}_0)$ es una
$\mathrm{A}_0$-álgebra de tipo finito y
$\operatorname{gr}^*_{\mathfrak{J}_0}(\mathrm{M}_0)$ es un módulo de tipo finito sobre
$\operatorname{gr}^*_{\mathfrak{J}_0}(\mathrm{B}_0)$. Por el teorema de
planitud genérica $(6.9.1)$, existe un subconjunto abierto
$\mathrm{D}(t_0)\subset\mathrm{S}_0$ tal que, poniendo
$\mathrm{A}_{01}=(\mathrm{A}_0)_{t_0}$,
\[
  \operatorname{gr}^*_{\mathfrak{J}_0}(\mathrm{M}_0)
  \otimes_{\mathrm{A}_0}\mathrm{A}_{01}
\]
es un $\mathrm{A}_{01}$-módulo plano. Como $\mathrm{A}_{01}$ es un $\mathrm{A}_0$-módulo plano, este módulo se identifica con
\[
  \operatorname{gr}^*_{\mathfrak{J}_{01}}
  (\mathrm{M}_0\otimes_{\mathrm{A}_0}\mathrm{A}_{01}),
\]
donde
\[
  \mathrm{B}_{01}=\mathrm{B}_0\otimes_{\mathrm{A}_0}\mathrm{A}_{01},
  \qquad
  \mathfrak{J}_{01}=\mathfrak{J}_0\mathrm{B}_{01}.
\]
El complemento de $\mathrm{D}(t_0)$ en $\mathrm{S}_0$ es entonces de la forma $\mathrm{V}(\mathfrak{a}_0)$, y la conclusión se obtiene aplicando la hipótesis de inducción noetheriana.
\end{proof}

Apliquemos este lema a la situación actual, conservando las mismas notaciones; pongamos $\mathrm{U}_{0i}=\bigcup_{j\leqslant i}\mathrm{S}_{0j}$, que es un abierto cuasicompacto de $\mathrm{S}_0$. Existe entonces una familia $(f_{ik})_{1\leqslant i\leqslant m,\,1\leqslant k\leqslant n}$ de elementos de $\mathrm{A}_0$ tal que, para todo $i$, $\mathrm{U}_{0i}$ es la unión de los $\mathrm{D}(f_{ik})$ ($1\leqslant k\leqslant n$).

\oldpage[IV-3]{130}

El $\mathrm{C}_0$-módulo $\operatorname{gr}^*_{\mathfrak{J}_0}(\mathrm{M}_0)$ está generado por un número finito de elementos homogéneos de grado $0$, por lo que existe un epimorfismo de $\mathrm{C}_0$-módulos graduados
\[
  u_0:\mathrm{C}_0^r\longrightarrow\operatorname{gr}^*_{\mathfrak{J}_0}(\mathrm{M}_0).
\]
Como $\mathrm{C}=\mathrm{C}_0\otimes_{\mathrm{A}_0}\mathrm{A}$ y el homomorfismo canónico $\operatorname{gr}^*_{\mathfrak{J}_0}(\mathrm{M}_0)\otimes_{\mathrm{A}_0}\mathrm{A}\to\operatorname{gr}^*_\mathfrak{J}(\mathrm{M})$ es sobreyectivo, se deduce de $u_0$ un epimorfismo de $\mathrm{C}$-módulos graduados
\[
  u:\mathrm{C}^r\xrightarrow{u_0\otimes 1_{\mathrm{A}}}\operatorname{gr}^*_{\mathfrak{J}_0}(\mathrm{M}_0)\otimes_{\mathrm{A}_0}\mathrm{A}\longrightarrow\operatorname{gr}^*_\mathfrak{J}(\mathrm{M}),
\]
y queda por demostrar que el $\mathrm{C}$-módulo graduado
$\mathrm{Q}=\operatorname{Ker}(u)$ es \emph{de tipo finito}.
\\[0.4\baselineskip]
\begin{lemma}[11.2.9.6]
\label{IV.11.2.9.6}
Bajo las hipótesis precedentes, sea $\mathfrak{R}_0$ un ideal de $\mathrm{A}_0$; pongamos
\[
  \mathrm{A}'_0=\mathrm{A}_0/\mathfrak{R}_0,\qquad
  \mathrm{B}'_0=\mathrm{B}_0\otimes_{\mathrm{A}_0}\mathrm{A}'_0
  =\mathrm{B}_0/\mathfrak{R}_0\mathrm{B}_0,\qquad
  \mathfrak{J}'_0=\mathfrak{J}_0\mathrm{B}'_0,
\]
y $\mathfrak{R}=\mathfrak{R}_0\mathrm{A}$, $\mathrm{A}'=\mathrm{A}/\mathfrak{R}$, y supongamos además que
\[
  \operatorname{gr}^*_{\mathfrak{J}'_0}
  (\mathrm{M}_0\otimes_{\mathrm{A}_0}\mathrm{A}'_0)
\]
es un $\mathrm{A}'_0$-módulo plano. Entonces existe un número finito de elementos $q_h$ ($1\leqslant h\leqslant s$) de $\mathrm{Q}$ tales que, para todo ideal primo $\mathfrak{p}\supset\mathfrak{R}\mathrm{B}$ de $\mathrm{B}$, las imágenes canónicas de los $q_h$ en $\mathrm{Q}_\mathfrak{p}$ generan el $\mathrm{C}_\mathfrak{p}$-módulo $\mathrm{Q}_\mathfrak{p}$.
\end{lemma}

Supongamos demostrado el lema y observemos que sigue siendo aplicable después de sustituir
$\mathrm{A}_0$ por $(\mathrm{A}_0)_{t_0}$, y $\mathrm{B}_0$,
$\mathfrak{J}_0$, $\mathrm{M}_0$, $\mathrm{A}_\lambda$ y $\mathrm{A}$ por
los objetos localizados correspondientes, para cualquier $t_0\in\mathrm{A}_0$, pues
$(\mathrm{A}_0)_{t_0}$ es plano sobre $\mathrm{A}_0$. Apliquemos el lema
sustituyendo sucesivamente $\mathrm{A}_0$ por cada uno de los anillos
$(\mathrm{A}_0)_{f_{ik}}$, y $\mathfrak{R}_0$ por el ideal
$\mathfrak{R}_{ik}$ que define el subesquema cerrado de $\mathrm{D}(f_{ik})$
inducido por $\mathrm{S}_{0i}$ sobre el subconjunto abierto
$\mathrm{S}_{0i}\cap\mathrm{D}(f_{ik})$ de $\mathrm{S}_{0i}$. Como la
hipótesis de planitud de \sref{IV.11.2.9.6} se satisface en cada uno de estos casos
por \sref{IV.11.2.9.5}, para cada par $(i,k)$ se obtiene una familia finita
$(a_{ikh})_{1\leqslant h\leqslant l}$ de elementos de
$\mathrm{Q}_{f_{ik}}$ cuyas imágenes en $\mathrm{Q}_{\mathfrak{p}_{ik}}$
generan este $\mathrm{C}_{\mathfrak{p}_{ik}}$-módulo para todo ideal primo
$\mathfrak{p}_{ik}$ de $\mathrm{B}_{f_{ik}}$ que contiene
$\mathfrak{R}_{ik}\mathrm{B}_{f_{ik}}$. Para un entero adecuado $d>0$, se
puede escribir
\[
  a_{ikh}=\frac{b_{ikh}}{f_{ik}^{d}}
\]
para todos $i,k,h$, con $b_{ikh}\in\mathrm{Q}$. Como los $\mathrm{S}_{0i}$ recubren $\mathrm{S}_0$, todo ideal primo $\mathfrak{p}$ de $\mathrm{B}$ tiene imagen en algún $\mathrm{S}_{0i}\cap\mathrm{D}(f_{ik})$ y, por tanto, su imagen $\mathfrak{p}_{ik}$ en $\mathrm{B}_{f_{ik}}$ contiene $\mathfrak{R}_{ik}\mathrm{B}_{f_{ik}}$. Las imágenes de los $b_{ikh}$ ($1\leqslant h\leqslant l$) en $\mathrm{Q}_{\mathfrak{p}_{ik}}=\mathrm{Q}_{\mathfrak{p}}$ generan, por consiguiente, este $\mathrm{C}_{\mathfrak{p}}$-módulo. Se deduce que los $b_{ikh}$ ($1\leqslant i\leqslant m$, $1\leqslant k\leqslant n$, $1\leqslant h\leqslant l$) generan el $\mathrm{C}$-módulo $\mathrm{Q}$ (Bourbaki, \emph{Alg.\ comm.}, cap.~II, \S3, n.º~3, teorema~1).

Queda por demostrar el Lema \sref{IV.11.2.9.6}. Pongamos
\[
  \mathrm{C}'=\mathrm{C}\otimes_\mathrm{A}\mathrm{A}'=\mathrm{C}/\mathfrak{R}\mathrm{C},\qquad
  \mathrm{B}'=\mathrm{B}\otimes_\mathrm{A}\mathrm{A}',\qquad
  \mathfrak{J}'=\mathfrak{J}\mathrm{B}',\qquad
  \mathrm{M}'=\mathrm{M}\otimes_\mathrm{A}\mathrm{A}'.
\]
Por hipótesis, $\operatorname{gr}^*_\mathfrak{J}(\mathrm{M})$ es un $\mathrm{A}$-módulo plano; por tanto, $\operatorname{gr}^*_{\mathfrak{J}'}(\mathrm{M}')$ se identifica con $\operatorname{gr}^*_\mathfrak{J}(\mathrm{M})\otimes_\mathrm{A}\mathrm{A}'$ \sref{IV.11.2.9.2}, y se tiene la sucesión exacta $(\mathbf{0}_\mathrm{I},~6.1.2)$
\begin{equation}
  0\longrightarrow\mathrm{Q}\otimes_\mathrm{A}\mathrm{A}'
  \longrightarrow(\mathrm{C}')^r
  \xrightarrow{u\otimes 1_{\mathrm{A}'}}
  \operatorname{gr}^*_{\mathfrak{J}'}(\mathrm{M}')
  \longrightarrow 0.
  \label{IV.11.2.9.7}
\tag{11.2.9.7}
\end{equation}

Pongamos
\[
  \mathrm{C}'_0=\mathrm{C}_0\otimes_{\mathrm{A}_0}\mathrm{A}'_0
  =\mathrm{C}_0/\mathfrak{R}_0\mathrm{C}_0
\]
y consideremos el epimorfismo de $\mathrm{C}'_0$-módulos graduados
\[
  u'_0:(\mathrm{C}'_0)^r
  \xrightarrow{u_0\otimes 1_{\mathrm{A}'_0}}
  \operatorname{gr}^*_{\mathfrak{J}_0}(\mathrm{M}_0)\otimes_{\mathrm{A}_0}\mathrm{A}'_0
  \longrightarrow
  \operatorname{gr}^*_{\mathfrak{J}'_0}
  (\mathrm{M}_0\otimes_{\mathrm{A}_0}\mathrm{A}'_0).
\]

Sea $\mathrm{Q}_0=\operatorname{Ker}(u'_0)$. Utilizando el hecho de que
\[
  \operatorname{gr}^*_{\mathfrak{J}'_0}
  (\mathrm{M}_0\otimes_{\mathrm{A}_0}\mathrm{A}'_0)
\]
es un $\mathrm{A}'_0$-módulo plano, se ve que su producto tensorial sobre $\mathrm{A}'_0$ con $\mathrm{A}'$ se identifica con $\operatorname{gr}^*_{\mathfrak{J}'}(\mathrm{M}')$ \sref{IV.11.2.9.2}, y se tiene la sucesión exacta $(\mathbf{0}_\mathrm{I},~6.1.2)$
\begin{equation}
  0\longrightarrow\mathrm{Q}_0\otimes_{\mathrm{A}'_0}\mathrm{A}'
  \longrightarrow(\mathrm{C}')^r
  \xrightarrow{u\otimes 1_{\mathrm{A}'}}
  \operatorname{gr}^*_{\mathfrak{J}'}(\mathrm{M}')
  \longrightarrow 0.
  \label{IV.11.2.9.8}
\tag{11.2.9.8}
\end{equation}

\oldpage[IV-3]{131}

de donde se obtiene, por comparación con {\hyperlink{ega4.c04.thm.11-2-9}{(11.2.9.7)}}, un isomorfismo de $\mathrm{C}'$-módulos graduados
\[
  \mathrm{Q}\otimes_\mathrm{A}\mathrm{A}'
  \xrightarrow{\sim}
  \mathrm{Q}_0\otimes_{\mathrm{A}'_0}\mathrm{A}'.
\]

Como $\mathrm{C}'_0$ es noetheriano, $\mathrm{Q}_0$ es un $\mathrm{C}'_0$-módulo de tipo finito; por tanto, $\mathrm{Q}_0\otimes_{\mathrm{A}'_0}\mathrm{A}'$ es un $\mathrm{C}'$-módulo graduado de tipo finito y, por consiguiente, también lo es $\mathrm{Q}\otimes_\mathrm{A}\mathrm{A}'$. Sean $q_h$ ($1\leqslant h\leqslant s$) elementos homogéneos de $\mathrm{Q}$ cuyas imágenes en $\mathrm{Q}\otimes_\mathrm{A}\mathrm{A}'$ generan este $\mathrm{C}'$-módulo.

Consideremos ahora un ideal primo arbitrario $\mathfrak{p}\supset\mathfrak{R}\mathrm{B}$ de $\mathrm{B}$. Por planitud,
\[
  \mathrm{C}_\mathfrak{p}
  =\operatorname{gr}^*_{\mathfrak{J}_\mathfrak{p}}(\mathrm{B}_\mathfrak{p}),
\]
y $\operatorname{gr}^0_{\mathfrak{J}_\mathfrak{p}}(\mathrm{B}_\mathfrak{p})=\mathrm{B}_\mathfrak{p}/\mathfrak{J}_\mathfrak{p}$ es el anillo cero o un anillo local; para demostrar el lema, basta evidentemente considerar el segundo caso. Cada $(\mathrm{B}_\mathfrak{p}/\mathfrak{J}_\mathfrak{p})$-módulo $\operatorname{gr}^i_{\mathfrak{J}_\mathfrak{p}}(\mathrm{B}_\mathfrak{p})$ es entonces de presentación finita. Por otra parte, para todo $i$,
\[
  \mathrm{M}/\mathfrak{J}^i\mathrm{M}
  \simeq
  (\mathrm{M}_0/\mathfrak{J}_0^i\mathrm{M}_0)
  \otimes_{\mathrm{A}_0}\mathrm{A},
\]
y es, por tanto, un $\mathrm{B}$-módulo de presentación finita. Por inducción sobre $i$, la hipótesis de que $\operatorname{gr}^*_\mathfrak{J}(\mathrm{M})$ es plano sobre $\mathrm{A}$ implica, por $(\mathbf{0}_\mathrm{I},~6.1.2)$, que $\mathrm{M}/\mathfrak{J}^i\mathrm{M}$ es plano sobre $\mathrm{A}$. Aplicando \sref{IV.11.2.6.1}, con $\mathrm{M}_0$ sustituido por $\mathrm{M}_0/\mathfrak{J}_0^k\mathrm{M}_0$ para $k\leqslant i$, se ve que hay un índice $\lambda_i$ tal que, para $\mu\geqslant\lambda_i$, cada $\mathrm{M}_\mu/\mathfrak{J}_\mu^{k+1}\mathrm{M}_\mu$ es plano sobre $\mathrm{A}_\mu$; por consiguiente, de nuevo por $(\mathbf{0}_\mathrm{I},~6.1.2)$, cada $\operatorname{gr}^k_{\mathfrak{J}_\mu}(\mathrm{M}_\mu)$ es plano sobre $\mathrm{A}_\mu$ para $k\leqslant i$. Se deduce de \sref{IV.11.2.9.2} que cada $(\mathrm{B}/\mathfrak{J})$-módulo $\operatorname{gr}^i_\mathfrak{J}(\mathrm{M})$ es de presentación finita y, por tanto, que $\operatorname{gr}^i_{\mathfrak{J}_\mathfrak{p}}(\mathrm{M}_\mathfrak{p})$ es un $(\mathrm{B}_\mathfrak{p}/\mathfrak{J}_\mathfrak{p})$-módulo de presentación finita.

Por otra parte, las imágenes de los $q_h$ en
\[
  \mathrm{Q}_\mathfrak{p}\otimes_{\mathrm{B}_\mathfrak{p}}\boldsymbol{k}(\mathfrak{p})
  =
  (\mathrm{Q}\otimes_\mathrm{A}\mathrm{A}')_\mathfrak{p}
  \otimes_{\mathrm{A}'}\boldsymbol{k}(\mathfrak{p})
\]
generan este módulo sobre $\mathrm{C}_\mathfrak{p}\otimes_{\mathrm{B}_\mathfrak{p}}\boldsymbol{k}(\mathfrak{p})$, pues
\[
  \mathrm{C}_\mathfrak{p}\otimes_{\mathrm{B}_\mathfrak{p}}\boldsymbol{k}(\mathfrak{p})
  =
  \mathrm{C}'_\mathfrak{p}\otimes_{\mathrm{A}'}\boldsymbol{k}(\mathfrak{p}).
\]
Como se tiene una sucesión exacta
\[
  \mathrm{Q}_\mathfrak{p}\otimes_{\mathrm{B}_\mathfrak{p}}\boldsymbol{k}(\mathfrak{p})
  \longrightarrow
  \mathrm{C}_\mathfrak{p}^{\,r}\otimes_{\mathrm{B}_\mathfrak{p}}\boldsymbol{k}(\mathfrak{p})
  \longrightarrow
  \operatorname{gr}^*_{\mathfrak{J}_\mathfrak{p}}(\mathrm{M}_\mathfrak{p})
  \otimes_{\mathrm{B}_\mathfrak{p}}\boldsymbol{k}(\mathfrak{p})
  \longrightarrow 0,
\]
se deduce que $\operatorname{gr}^*_{\mathfrak{J}_\mathfrak{p}}(\mathrm{M}_\mathfrak{p})\otimes_{\mathrm{B}_\mathfrak{p}}\boldsymbol{k}(\mathfrak{p})$ es un módulo de presentación finita sobre $\mathrm{C}_\mathfrak{p}\otimes_{\mathrm{B}_\mathfrak{p}}\boldsymbol{k}(\mathfrak{p})$. Aplicando \sref{IV.11.2.9.4}, con $\mathrm{A}$, $\mathrm{B}$ y $\mathrm{M}$ sustituidos, respectivamente, por $\mathrm{B}_\mathfrak{p}$, $\mathrm{C}_\mathfrak{p}$ y $\operatorname{gr}^*_{\mathfrak{J}_\mathfrak{p}}(\mathrm{M}_\mathfrak{p})$, concluimos que $\mathrm{Q}_\mathfrak{p}$ es un $\mathrm{C}_\mathfrak{p}$-módulo de tipo finito. El lema de Nakayama (y el hecho de que $\mathrm{C}=\mathrm{B}$) muestra entonces que las imágenes de los $q_h$ en $\mathrm{Q}_\mathfrak{p}$ generan este $\mathrm{C}_\mathfrak{p}$-módulo. Esto completa la demostración de \sref{IV.11.2.9.6} y demuestra también que $\operatorname{gr}^*_\mathfrak{J}(\mathrm{M})$ es un $\mathrm{C}$-módulo de presentación finita.

\medskip
\noindent\textbf{VI)} \emph{Fin de la demostración.} --- Para abreviar, pongamos $\mathrm{C}_\lambda=\mathrm{C}_0\otimes_{\mathrm{A}_0}\mathrm{A}_\lambda$ (de hecho, igual a $\mathrm{B}_\lambda$), $\mathrm{N}_\lambda=\operatorname{gr}^*_{\mathfrak{J}_\lambda}(\mathrm{M}_\lambda)$ y $\mathrm{N}=\operatorname{gr}^*_\mathfrak{J}(\mathrm{M})$. Obsérvese primero que, para todo entero $k$, $\mathfrak{J}^k\mathrm{M}$ se identifica con $\varinjlim_\lambda\mathfrak{J}_\lambda^k\mathrm{M}_\lambda$ por la exactitud del funtor $\varinjlim_\lambda$: la imagen de $\mathfrak{J}_\lambda^k\mathrm{M}_\lambda$ en $\mathrm{M}_\mu$ para $\lambda\leqslant\mu$ (respectivamente, en $\mathrm{M}$) genera $\mathfrak{J}_\mu^k\mathrm{M}_\mu$ (respectivamente, $\mathfrak{J}^k\mathrm{M}$). Utilizando una vez más la exactitud de $\varinjlim$, concluimos que $\mathrm{N}$ se identifica canónicamente con $\varinjlim_\lambda\mathrm{N}_\lambda$. Con esta identificación, demostraremos primero:
\begin{equation}
\label{IV.11.2.9.9}
\tag{11.2.9.9}
\text{\emph{Para $\lambda$ suficientemente grande, el homomorfismo canónico $v_\lambda:\mathrm{N}_\lambda\otimes_{\mathrm{A}_\lambda}\mathrm{A}\to\mathrm{N}$ es biyectivo.}}
\end{equation}
Como los $\mathrm{C}_\lambda$ son noetherianos y $\mathrm{N}_\lambda$ es un $\mathrm{C}_\lambda$-módulo de tipo finito, los $\mathrm{C}$-módulos $\mathrm{N}_\lambda\otimes_{\mathrm{A}_\lambda}\mathrm{A}$ son de presentación finita y forman un sistema inductivo filtrante cuyo límite inductivo se identifica canónicamente con $\mathrm{N}$, pues los límites inductivos filtrantes conmutan con los productos tensoriales. Además, los homomorfismos de transición $v_{\mu\lambda}:\mathrm{N}_\lambda\otimes_{\mathrm{A}_\lambda}\mathrm{A}\to\mathrm{N}_\mu\otimes_{\mathrm{A}_\mu}\mathrm{A}$

\oldpage[IV-3]{132}

para $\lambda\leqslant\mu$, así como los homomorfismos $v_\lambda$, son
\emph{sobreyectivos}. Fijemos $\lambda$ y consideremos
\[
  \operatorname{Ker}(v_{\mu\lambda})=\mathrm{N}'_\mu
\]
como un $\mathrm{C}$-submódulo de
$\mathrm{N}_\lambda\otimes_{\mathrm{A}_\lambda}\mathrm{A}$. Por la
definición del límite inductivo, los $\mathrm{N}'_\mu$ forman una
familia filtrante creciente cuya unión es $\operatorname{Ker}(v_\lambda)$.
Demostramos en V) que $\mathrm{N}$ es un $\mathrm{C}$-módulo de presentación finita;
en consecuencia, $\operatorname{Ker}(v_\lambda)$ es un $\mathrm{C}$-módulo de tipo finito
(Bourbaki, \emph{Alg.\ comm.}, cap.~I, \S~2, no.~8, lema~9).
Hay, pues, un índice $\mu\geqslant\lambda$ tal que
$\mathrm{N}'_\mu=\operatorname{Ker}(v_\lambda)$. Como $v_\mu$ es
sobreyectivo, esta igualdad dice que $v_\mu$ es también inyectivo y,
por tanto, demuestra {\hyperlink{ega4.c04.thm.11-2-9}{(11.2.9.9)}}.

Después de sustituir $\mathrm{A}_0$ y $\mathrm{M}_0$ por
$\mathrm{A}_\lambda$ y $\mathrm{M}_\lambda$ para un $\lambda$ suficientemente grande,
podemos suponer, por tanto, que el homomorfismo canónico
$v_\lambda$ es biyectivo para \emph{todo} $\lambda$. Pongamos
\[
  \mathrm{P}_\lambda
  =\mathrm{N}_0\otimes_{\mathrm{C}_0}\mathrm{C}_\lambda
  =\mathrm{N}_0\otimes_{\mathrm{A}_0}\mathrm{A}_\lambda,
\]
de modo que $\mathrm{N}=\varinjlim\mathrm{P}_\lambda$. Como
$\mathrm{C}_0$ es una $\mathrm{A}_0$-álgebra de presentación finita y
$\mathrm{P}_0$ es un $\mathrm{C}_0$-módulo de presentación finita,
\sref{IV.11.2.6.1}, aplicado con
$\mathrm{B}_0=\mathrm{C}_0$ y $\mathrm{M}_0=\mathrm{P}_0$, muestra que
hay un índice $\lambda$ tal que $\mathrm{P}_\mu$ es un
$\mathrm{A}_\mu$-módulo plano para todo $\mu\geqslant\lambda$. Para tal
$\mu$ tenemos un diagrama conmutativo
\[
\xymatrix{
  \mathrm{P}_\mu \ar[r] \ar[d]^{w_\mu} & \mathrm{N} \ar@{=}[d] \\
  \mathrm{N}_\mu \ar[r] & \mathrm{N}
}
\]
y las definiciones muestran que $w_\mu$ es \emph{sobreyectivo}. Por III), el
homomorfismo $\mathrm{A}_\mu\to\mathrm{A}$ es inyectivo; como
$\mathrm{P}_\mu$ es plano sobre $\mathrm{A}_\mu$, el homomorfismo canónico
\[
  \mathrm{P}_\mu\longrightarrow
  \mathrm{N}
  =\mathrm{P}_\mu\otimes_{\mathrm{A}_\mu}\mathrm{A}
\]
es, por tanto, \emph{inyectivo}. El diagrama anterior muestra entonces que
$w_\mu$ es también inyectivo y, por consiguiente, biyectivo. Así,
$\mathrm{N}_\mu$ es un $\mathrm{A}_\mu$-módulo plano para
$\mu\geqslant\lambda$, lo que completa la demostración de
\sref{IV.11.2.9}.
\end{proof}

\begin{remark}[11.2.10]
\label{IV.11.2.10}
No sabemos si es válido el análogo del teorema de Raynaud para la
generalización de \sref{IV.11.2.6},~(i).
\end{remark}

\subsection{Aplicación a la eliminación de hipótesis noetherianas}
\label{subsection:IV.11.3}
\label{IV.11.3}

\oldpage[IV-3]{132}

\begin{theorem}[11.3.1]
\label{IV.11.3.1}
Sea $f:\mathrm{X}\to\mathrm{Y}$ un morfismo localmente de
presentación finita, y sea $\sh{F}$ un $\sh{O}_\mathrm{X}$-módulo
cuasicoherente de presentación finita. Entonces el conjunto
$\mathrm{U}$ de los puntos $x\in\mathrm{X}$ en los que $\sh{F}$ es
$f$-plano es abierto en $\mathrm{X}$. Si, además,
$\operatorname{Supp}(\sh{F})=\mathrm{X}$, la restricción
$f|_\mathrm{U}$ es un morfismo abierto.
\end{theorem}

\begin{proof}
La segunda afirmación ya se ha demostrado en $(2.4.6)$ y se incluye
aquí únicamente como referencia.

Como la cuestión es local en $\mathrm{X}$ y en $\mathrm{Y}$, podemos
suponer que ambos son afines y que $f$ es de presentación finita. Sea
$x\in\mathrm{X}$ un punto en el que $\sh{F}$ es $f$-plano. Aplicando
\sref{IV.11.2.7}, podemos asumir que
\[
  \mathrm{X}=\mathrm{X}_0\times_{\mathrm{Y}_0}\mathrm{Y},
  \qquad
  f=f_0\times 1,
  \qquad
  \sh{F}
  =\sh{F}_0\otimes_{\sh{O}_{\mathrm{Y}_0}}\sh{O}_\mathrm{Y},
\]
en donde $\mathrm{Y}_0$ es noetheriano,
$f_0:\mathrm{X}_0\to\mathrm{Y}_0$ es de tipo finito, y
$\sh{F}_0$ es un $\sh{O}_{\mathrm{X}_0}$-módulo coherente. Además, si
$x_0$ es la proyección canónica de $x$ en $\mathrm{X}_0$, podemos
suponer que $\sh{F}_0$ es $f_0$-plano en $x_0$. Por
\sref{IV.11.1.1}, el conjunto $\mathrm{U}_0$ de puntos de
$\mathrm{X}_0$ en los que $\sh{F}_0$ es $f_0$-plano es, por tanto, una
vecindad de $x_0$. Se deduce de $(2.1.4)$ que $\sh{F}$ es
$f$-plano en cada punto de la imagen inversa de $\mathrm{U}_0$ en
$\mathrm{X}$; así pues, $\mathrm{U}$ es una vecindad de $x$.
\end{proof}

\oldpage[IV-3]{133}

\begin{corollary}[11.3.2]
\label{IV.11.3.2}
Sean $f:\mathrm{X}\to\mathrm{Y}$ y $g:\mathrm{Y}\to\mathrm{Z}$ dos morfismos de presentación finita, y sea $\sh{F}$ un $\sh{O}_\mathrm{X}$-módulo cuasicoherente de presentación finita. Entonces el conjunto $\mathrm{U}$ de los $y\in\mathrm{Y}$ tales que, para toda generización $y'$ de $y$, $\sh{F}$ es $(g\circ f)$-plano en todos los puntos de $f^{-1}(y')$ (es decir, tales que $\sh{F}\otimes_\mathrm{Y}\Spec(\sh{O}_{\mathrm{Y},y})$ es plano respecto del morfismo $\mathrm{X}\times_\mathrm{Y}\Spec(\sh{O}_{\mathrm{Y},y})\to\mathrm{Z}$) es abierto en $\mathrm{Y}$.
\end{corollary}

\begin{proof}
El mismo razonamiento que en \sref{IV.11.1.5} muestra que, si $\mathrm{V}$ es el conjunto de los $x\in\mathrm{X}$ en los que $\sh{F}$ es $(g\circ f)$-plano, $\mathrm{U}$ es el conjunto de los $y\in\mathrm{Y}$ cuyas generizaciones pertenecen todas a $\mathrm{E}=\mathrm{Y}-f(\mathrm{X}-\mathrm{V})$. Como $g\circ f$ es de presentación finita, $\mathrm{V}$ es abierto en $\mathrm{X}$ en virtud de \sref{IV.11.3.1}, y por consiguiente $\mathrm{V}$ es ind-construible $(1.9.6)$ y $\mathrm{X}-\mathrm{V}$ es proconstruible; ahora bien, como $f$ es cuasicompacto, $f(\mathrm{X}-\mathrm{V})$ es proconstruible en $\mathrm{Y}$ $(1.9.5,~(\mathrm{vii}))$, y por tanto $\mathrm{E}$ es ind-construible en $\mathrm{Y}$. Se deduce entonces de $(1.10.1)$ que $\mathrm{U}$ es el interior de $\mathrm{E}$ y, por consiguiente, es abierto en $\mathrm{Y}$.
\end{proof}

\begin{corollary}[11.3.3]
\label{IV.11.3.3}
Sea $\mathrm{A}$ un anillo, $\mathrm{B}$ una $\mathrm{A}$-álgebra de presentación finita, $\mathrm{C}$ una $\mathrm{B}$-álgebra de presentación finita y $\mathrm{M}$ un $\mathrm{C}$-módulo de presentación finita. Entonces el conjunto de los $\mathfrak{q}\in\Spec(\mathrm{B})$ tales que $\mathrm{M}_\mathfrak{q}$ es un $\mathrm{A}$-módulo plano es abierto en $\Spec(\mathrm{B})$.
\end{corollary}

\begin{proposition}[11.3.4]
\label{IV.11.3.4}
Sea $f:\mathrm{X}\to\mathrm{S}$ un morfismo localmente de presentación finita, $\mathscr{J}$ un ideal cuasicoherente de tipo finito de $\sh{O}_\mathrm{X}$, $\mathrm{Y}$ el subpreesquema cerrado de $\mathrm{X}$ definido por $\mathscr{J}$ y $\sh{F}$ un $\sh{O}_\mathrm{X}$-módulo de presentación finita. Supongamos que $\operatorname{gr}^*_\mathscr{J}(\sh{F})$ es $f$-plano. En estas condiciones:
\begin{enumerate}[label=\emph{(\roman*)}]
  \item $\sh{F}$ es $f$-plano en los puntos de $\mathrm{Y}$.
  \item Si $\operatorname{gr}^*_\mathscr{J}(\sh{O}_\mathrm{X})$ es $f$-plano, $\operatorname{gr}^*_\mathscr{J}(\sh{O}_\mathrm{X})$ es una $(\sh{O}_\mathrm{X}/\mathscr{J})$-álgebra de presentación finita y $\operatorname{gr}^*_\mathscr{J}(\sh{F})$ es un $(\operatorname{gr}^*_\mathscr{J}(\sh{O}_\mathrm{X}))$-módulo de presentación finita.
  \item Supongamos que $\operatorname{gr}^*_\mathscr{J}(\sh{O}_\mathrm{X})$ es $f$-plano (lo cual implica que $\sh{O}_\mathrm{X}/\mathscr{J}$ es $f$-plano). Entonces el conjunto $\mathrm{W}$ de los $y\in\mathrm{Y}$ tales que $(\operatorname{gr}^*_\mathscr{J}(\sh{F}))_y$ es un $(\sh{O}_\mathrm{X}/\mathscr{J})_y$-módulo plano es abierto en $\mathrm{Y}$.
  \item Supongamos que $\operatorname{gr}^*_\mathscr{J}(\sh{O}_\mathrm{X})$ es $f$-plano. Sea $\mathrm{S}'\to\mathrm{S}$ un morfismo cualquiera, sea $\mathrm{X}'=\mathrm{X}\times_\mathrm{S}\mathrm{S}'$ y sea $p:\mathrm{X}'\to\mathrm{X}$ la proyección canónica. Pongamos $\mathrm{Y}'=p^{-1}(\mathrm{Y})=\mathrm{Y}\times_\mathrm{S}\mathrm{S}'$, el subpreesquema cerrado de $\mathrm{X}'$ definido por $\mathscr{J}'=p^*(\mathscr{J})\sh{O}_{\mathrm{X}'}$, y pongamos $\sh{F}'=p^*(\sh{F})=\sh{F}\otimes_{\sh{O}_\mathrm{S}}\sh{O}_{\mathrm{S}'}$. Entonces, si $\mathrm{W}'=p^{-1}(\mathrm{W})$, se tiene $\operatorname{gr}^*_{\mathscr{J}'}(\sh{F}')|_{\mathrm{W}'}\cong p^*(\operatorname{gr}^*_\mathscr{J}(\sh{F})|_\mathrm{W})$, y $\operatorname{gr}^*_{\mathscr{J}'}(\sh{F}')$ es un $(\sh{O}_{\mathrm{X}'}/\mathscr{J}')$-módulo plano en los puntos de $\mathrm{W}'$.
\end{enumerate}
\end{proposition}

\begin{proof}
Como las cuestiones son locales en $\mathrm{X}$ y en $\mathrm{S}$, podemos suponer que
$\mathrm{S}=\Spec(\mathrm{A})$ y $\mathrm{X}=\Spec(\mathrm{B})$ son afines,
en donde $\mathrm{B}$ es una $\mathrm{A}$-álgebra de presentación finita, que
$\sh{F}=\widetilde{\mathrm{M}}$, donde $\mathrm{M}$ es una
$\mathrm{B}$-módulo de presentación finita, y que
$\mathscr{J}=\widetilde{\mathfrak{J}}$, donde $\mathfrak{J}$ es un ideal
de tipo finito de $\mathrm{B}$. En virtud de $(8.9.1)$ y $(8.5.11)$,
existen un subanillo noetheriano $\mathrm{A}_0$ de $\mathrm{A}$, una
$\mathrm{A}_0$-álgebra de tipo finito $\mathrm{B}_0$ de tal manera que
$\mathrm{B}=\mathrm{B}_0\otimes_{\mathrm{A}_0}\mathrm{A}$, un ideal
$\mathfrak{J}_0$ de $\mathrm{B}_0$ tal que
$\mathfrak{J}=\mathfrak{J}_0\mathrm{B}$, y un tipo finito
$\mathrm{B}_0$-módulo $\mathrm{M}_0$ tal que
$\mathrm{M}=\mathrm{M}_0\otimes_{\mathrm{A}_0}\mathrm{A}
=\mathrm{M}_0\otimes_{\mathrm{B}_0}\mathrm{B}$. Además,
$\mathrm{A}$ es el límite inductivo de sus
$\mathrm{A}_0$-subálgebras de tipo finito $\mathrm{A}_\lambda$. Pongamos
$\mathrm{B}_\lambda=\mathrm{B}_0\otimes_{\mathrm{A}_0}\mathrm{A}_\lambda$,
$\mathrm{M}_\lambda=\mathrm{M}_0\otimes_{\mathrm{A}_0}\mathrm{A}_\lambda
=\mathrm{M}_0\otimes_{\mathrm{B}_0}\mathrm{B}_\lambda$, y
$\mathfrak{J}_\lambda=\mathfrak{J}_0\mathrm{B}_\lambda$, de modo que
$\mathrm{B}=\varinjlim\mathrm{B}_\lambda$ y
$\mathrm{M}=\varinjlim\mathrm{M}_\lambda$. Por hipótesis,
$\operatorname{gr}^*_\mathfrak{J}(\mathrm{M})$ es un
$\mathrm{A}$-módulo plano; se deduce, pues, de \sref{IV.11.2.9} que
existe $\lambda$ tal que
$\operatorname{gr}^*_{\mathfrak{J}_\lambda}(\mathrm{M}_\lambda)$ es un
$\mathrm{A}_\lambda$-módulo plano. Para probar que $\sh{F}$ es $f$-plano en
los puntos de $\mathrm{Y}$, podemos por tanto limitarnos al caso en que
$\mathrm{S}$ es noetheriano; pero entonces $(\mathbf{0}_\mathrm{I},~6.1.2)$,
aplicado por inducción sobre $n$, demuestra que los
$\sh{F}/\mathscr{J}^{n+1}\sh{F}$ son $f$-planos, y se concluye mediante
$(\mathbf{0}_{\mathrm{III}},~10.2.6)$.
\oldpage[IV-3]{134}

La demostración de (ii) se reduce análogamente al caso en que $\mathrm{S}$ es noetheriano, utilizando \sref{IV.11.2.9}; la conclusión es entonces evidente, pues en este caso $\operatorname{gr}^*_\mathfrak{J}(\mathrm{B})$ es una $(\mathrm{B}/\mathfrak{J})$-álgebra de tipo finito y $\operatorname{gr}^*_\mathfrak{J}(\mathrm{M})$ es un $\operatorname{gr}^*_\mathfrak{J}(\mathrm{B})$-módulo de tipo finito.

Para demostrar (iii), nos reducimos, como en (i), al caso en que $\mathrm{S}$ y $\mathrm{X}$ son afines. Con las mismas notaciones, podemos suponer, en virtud de \sref{IV.11.2.9}, que $\operatorname{gr}^*_{\mathfrak{J}_\lambda}(\mathrm{B}_\lambda)$ y $\operatorname{gr}^*_{\mathfrak{J}_\lambda}(\mathrm{M}_\lambda)$ son $\mathrm{A}_\lambda$-módulos planos y que
\[
\operatorname{gr}^*_\mathfrak{J}(\mathrm{B})
=\operatorname{gr}^*_{\mathfrak{J}_\lambda}(\mathrm{B}_\lambda)\otimes_{\mathrm{A}_\lambda}\mathrm{A},
\qquad
\operatorname{gr}^*_\mathfrak{J}(\mathrm{M})
=\operatorname{gr}^*_{\mathfrak{J}_\lambda}(\mathrm{M}_\lambda)\otimes_{\mathrm{A}_\lambda}\mathrm{A}.
\]
En consecuencia, $\operatorname{gr}^*_\mathfrak{J}(\mathrm{B})$ es una $(\mathrm{B}/\mathfrak{J})$-álgebra de presentación finita y $\operatorname{gr}^*_\mathfrak{J}(\mathrm{M})$ es un $\operatorname{gr}^*_\mathfrak{J}(\mathrm{B})$-módulo de presentación finita. Si $\mathrm{W}^m$ es el conjunto de los $\mathfrak{p}\in\Spec(\mathrm{B}/\mathfrak{J})$ tales que $\operatorname{gr}^m_\mathfrak{J}(\mathrm{M})_{\mathfrak{p}}$ es un $(\mathrm{B}/\mathfrak{J})_{\mathfrak{p}}$-módulo plano, entonces $\mathrm{W}^m$ es abierto en $\Spec(\mathrm{B}/\mathfrak{J})$ por \sref{IV.11.3.1}, y $\mathrm{W}=\bigcap_m\mathrm{W}^m$; por tanto, $\mathrm{W}$ es estable por generización. La aserción (iii) se deduce ahora de \sref{IV.11.3.2}, aplicado con $\mathrm{X}=\Spec(\operatorname{gr}^*_\mathfrak{J}(\mathrm{B}))$, $\mathrm{Y}=\mathrm{Z}=\Spec(\mathrm{B}/\mathfrak{J})$ y $\sh{F}=(\operatorname{gr}^*_\mathfrak{J}(\mathrm{M}))^\sim$.

Por último, (iv) se deduce inmediatamente de \sref{IV.11.2.9.2}.

Generalizando la definición de $(6.10.1)$, sea $\mathrm{Y}$ un subpreesquema cerrado de un preesquema $\mathrm{X}$, definido por un ideal cuasicoherente $\mathscr{J}$ de $\sh{O}_\mathrm{X}$, y sea $\sh{F}$ un $\sh{O}_\mathrm{X}$-módulo cuasicoherente. Se dice que $\sh{F}$ es \emph{normalmente plano a lo largo de $\mathrm{Y}$ en un punto $y$} si $(\operatorname{gr}^*_\mathscr{J}(\sh{F}))_y$ es un $\sh{O}_{\mathrm{Y},y}$-módulo plano. Se dice que $\sh{F}$ es \emph{normalmente plano a lo largo de $\mathrm{Y}$} si lo es en todo punto de $\mathrm{Y}$.
\end{proof}

\begin{corollary}[11.3.5]
\label{IV.11.3.5}
Bajo las hipótesis generales de \sref{IV.11.3.4}, supongamos que $\operatorname{gr}^*_\mathscr{J}(\sh{O}_\mathrm{X})$ y $\operatorname{gr}^*_\mathscr{J}(\sh{F})$ son $f$-planos. Entonces el conjunto $\mathrm{W}$ de los $y\in\mathrm{Y}$ tales que $\sh{F}$ es normalmente plano a lo largo de $\mathrm{Y}$ en el punto $y$ es abierto en $\mathrm{Y}$ y (con las notaciones de \sref{IV.11.3.4},~(iv)) $\sh{F}'$ es normalmente plano a lo largo de $\mathrm{Y}'$ en todos los puntos de $\mathrm{W}'=p^{-1}(\mathrm{W})$.
\end{corollary}

\begin{proposition}[11.3.6]
\label{IV.11.3.6}
Con las notaciones de $(8.5.1)$ y $(8.8.1)$, supongamos que $\mathrm{S}_\alpha$ es cuasicompacto, que $\mathrm{X}_\alpha$ es de presentación finita sobre $\mathrm{S}_\alpha$ y que $\mathrm{Y}_\alpha$ es un subpreesquema cerrado de $\mathrm{X}_\alpha$ definido por un ideal cuasicoherente $\mathscr{J}_\alpha$ de tipo finito de $\sh{O}_{\mathrm{X}_\alpha}$; supongamos que $\operatorname{gr}^*_{\mathscr{J}_\alpha}(\sh{O}_{\mathrm{X}_\alpha})$ es plano sobre $\mathrm{S}_\alpha$; supongamos, por último, que $\sh{F}_\alpha$ es un $\sh{O}_{\mathrm{X}_\alpha}$-módulo de presentación finita. Para que $\sh{F}$ sea normalmente plano a lo largo de $\mathrm{Y}$, es necesario y suficiente que exista $\lambda$ tal que $\sh{F}_\lambda$ sea normalmente plano a lo largo de $\mathrm{Y}_\lambda$.
\end{proposition}

\begin{proof}
Obsérvese que $\mathrm{Y}$ (resp.\ $\mathrm{Y}_\lambda$) es el subpreesquema cerrado de $\mathrm{X}$ (resp.\ $\mathrm{X}_\lambda$) definido por $\mathscr{J}=\mathscr{J}_\alpha\sh{O}_\mathrm{X}$ (resp.\ $\mathscr{J}_\lambda=\mathscr{J}_\alpha\sh{O}_{\mathrm{X}_\lambda}$). Por la hipótesis y \sref{IV.11.2.9.2}, se tiene
\[
\operatorname{gr}^*_\mathscr{J}(\sh{O}_\mathrm{X})
=\operatorname{gr}^*_{\mathscr{J}_\alpha}(\sh{O}_{\mathrm{X}_\alpha})
 \otimes_{\sh{O}_{\mathrm{S}_\alpha}}\sh{O}_\mathrm{S},
\qquad
\operatorname{gr}^*_{\mathscr{J}_\lambda}(\sh{O}_{\mathrm{X}_\lambda})
=\operatorname{gr}^*_{\mathscr{J}_\alpha}(\sh{O}_{\mathrm{X}_\alpha})
 \otimes_{\sh{O}_{\mathrm{S}_\alpha}}\sh{O}_{\mathrm{S}_\lambda}
\quad(\lambda\geqslant\alpha),
\]
y $\operatorname{gr}^*_\mathscr{J}(\sh{O}_\mathrm{X})$ (resp.\ $\operatorname{gr}^*_{\mathscr{J}_\lambda}(\sh{O}_{\mathrm{X}_\lambda})$) es plano sobre $\mathrm{S}$ (resp.\ $\mathrm{S}_\lambda$). Se deduce que $\mathrm{Y}$ es plano sobre $\mathrm{S}$ (resp.\ $\mathrm{Y}_\lambda$ es plano sobre $\mathrm{S}_\lambda$). Si $\sh{F}_\lambda$ es normalmente plano a lo largo de $\mathrm{Y}_\lambda$, entonces $\left.\operatorname{gr}^*_{\mathscr{J}_\lambda}(\sh{F}_\lambda)\right|_{\mathrm{Y}_\lambda}$ es plano sobre $\mathrm{Y}_\lambda$ y, por tanto, también sobre $\mathrm{S}_\lambda$; y como
\[
\left.\operatorname{gr}^n_{\mathscr{J}_\lambda}(\sh{F}_\lambda)\right|_{\mathrm{X}_\lambda\setminus\mathrm{Y}_\lambda}=0
\quad\text{para }n\geqslant0,
\]
$\operatorname{gr}^*_{\mathscr{J}_\lambda}(\sh{F}_\lambda)$ es plano sobre $\mathrm{S}_\lambda$. Por \sref{IV.11.2.9.2},
\[
\operatorname{gr}^*_\mathscr{J}(\sh{F})
=\operatorname{gr}^*_{\mathscr{J}_\lambda}(\sh{F}_\lambda)
 \otimes_{\sh{O}_{\mathrm{S}_\lambda}}\sh{O}_\mathrm{S},
\]
por lo que $\sh{F}$ es normalmente plano a lo largo de $\mathrm{Y}$.

Recíprocamente, podemos suponer que $\mathrm{S}_\alpha$ y $\mathrm{X}_\alpha$ son afines y adoptar las notaciones de \sref{IV.11.2.9}. Puesto que $\operatorname{gr}^*_\mathfrak{J}(\mathrm{B})$ es un $\mathrm{A}$-módulo plano, lo mismo ocurre con $\operatorname{gr}^*_\mathfrak{J}(\mathrm{M})$; por tanto, en virtud de \sref{IV.11.2.9}, existe $\lambda\geqslant\alpha$ tal que $\operatorname{gr}^*_{\mathfrak{J}_\lambda}(\mathrm{M}_\lambda)$ es un $\mathrm{A}_\lambda$-módulo plano y
\[
\operatorname{gr}^*_\mathfrak{J}(\mathrm{M})
=\operatorname{gr}^*_{\mathfrak{J}_\lambda}(\mathrm{M}_\lambda)
 \otimes_{\mathrm{A}_\lambda}\mathrm{A}.
\]
Además, por \sref{IV.11.3.4},~(ii), $\operatorname{gr}^*_{\mathfrak{J}_\lambda}(\mathrm{M}_\lambda)$ es un $\operatorname{gr}^*_{\mathfrak{J}_\lambda}(\mathrm{B}_\lambda)$-módulo de presentación finita y $\operatorname{gr}^*_{\mathfrak{J}_\lambda}(\mathrm{B}_\lambda)$ es una $(\mathrm{B}_\lambda/\mathfrak{J}_\lambda)$-álgebra de presentación finita. La conclusión se deduce ahora de \sref{IV.11.2.6}.
\end{proof}

\oldpage[IV-3]{135}
\begin{proposition}[11.3.7]
\label{IV.11.3.7}
Sea $f:\mathrm{X}\to\mathrm{Y}$ un morfismo localmente de presentación finita, sean $\sh{F}$ y $\sh{F}'$ dos $\sh{O}_\mathrm{X}$-módulos cuasicoherentes de presentación finita, sea $u:\sh{F}'\to\sh{F}$ un homomorfismo de $\sh{O}_\mathrm{X}$-módulos, sea $x$ un punto de $\mathrm{X}$ y sea $y=f(x)$; supongamos que $\sh{F}$ es $f$-plano en el punto $x$. Las condiciones siguientes son equivalentes:
\begin{enumerate}[label=\emph{\alph*)}]
  \item $(\operatorname{Ker} u)_x=0$ y $\operatorname{Coker} u$ es un $\sh{O}_\mathrm{X}$-módulo $f$-plano en el punto $x$.
  \item El homomorfismo $(u\otimes 1)_x:(\sh{F}'\otimes_{\sh{O}_\mathrm{Y}}\boldsymbol{k}(y))_x\to(\sh{F}\otimes_{\sh{O}_\mathrm{Y}}\boldsymbol{k}(y))_x$ es inyectivo.
\end{enumerate}
Si además $\sh{F}$ es $f$-plano, el conjunto de los puntos $x$ que satisfacen las condiciones equivalentes anteriores es abierto en $\mathrm{X}$.
\end{proposition}

\begin{proof}
La condición $a)$ implica $b)$ en virtud de $(\mathbf{0}_\mathrm{I},~6.1.2)$, sin hipótesis alguna sobre $\sh{F}$. Para demostrar la recíproca, podemos limitarnos al caso en que $\mathrm{X}$ y $\mathrm{Y}$ son afines y, a continuación, en virtud de $(8.9.1)$ y \sref{IV.11.2.7}, suponer que
\[
  \mathrm{X}=\mathrm{X}_0\times_{\mathrm{Y}_0}\mathrm{Y},\qquad
  f=f_0\times 1_\mathrm{Y},\qquad
  \sh{F}=\sh{F}_0\otimes_{\sh{O}_{\mathrm{Y}_0}}\sh{O}_\mathrm{Y},\qquad
  \sh{F}'=\sh{F}'_0\otimes_{\sh{O}_{\mathrm{Y}_0}}\sh{O}_\mathrm{Y},\qquad
  u=u_0\otimes 1_\mathrm{Y},
\]
en donde $\mathrm{Y}_0$ es noetheriano, $f_0:\mathrm{X}_0\to\mathrm{Y}_0$ es un morfismo de tipo finito, $\sh{F}_0$ y $\sh{F}'_0$ son dos $\sh{O}_{\mathrm{X}_0}$-módulos coherentes, y $u_0:\sh{F}'_0\to\sh{F}_0$ es un homomorfismo. Además, si $x_0$ es la proyección de $x$ en $\mathrm{X}_0$, podemos suponer que $\sh{F}_0$ es $f_0$-plano en $x_0$. Pongamos $y_0=f_0(x_0)$. Por la transitividad de las fibras $(\mathbf{I},~3.6.4)$, la proyección $f^{-1}(y)\to f_0^{-1}(y_0)$ es fielmente plana $(2.2.13)$; y como
\[
  (u\otimes 1)_x=((u_0\otimes 1)_{x_0})\otimes 1_{\boldsymbol{k}(y)},
\]
la hipótesis de que $(u\otimes 1)_x$ es inyectivo implica que $(u_0\otimes 1)_{x_0}$ es inyectivo $(2.2.7)$. Aplicando $(\mathbf{0}_{\mathrm{III}},~10.2.4)$ a los anillos locales noetherianos $\sh{O}_{\mathrm{Y}_0,y_0}$ y $\sh{O}_{\mathrm{X}_0,x_0}$ y a los $\sh{O}_{\mathrm{X}_0,x_0}$-módulos $(\sh{F}'_0)_{x_0}$ y $(\sh{F}_0)_{x_0}$, el segundo de los cuales es plano sobre $\sh{O}_{\mathrm{Y}_0,y_0}$, obtenemos
\[
  (\operatorname{Ker}u_0)_{x_0}=0
  \quad\text{y}\quad
  \operatorname{Coker}u_0\ \text{es $f_0$-plano en }x_0.
\]
Se deduce, en primer lugar, que $\operatorname{Coker}u$ es $f$-plano en $x$ $(2.1.4)$. Además, por la exactitud por la derecha del producto tensorial,
\[
  \operatorname{Coker}u=(\operatorname{Coker}u_0)
  \otimes_{\sh{O}_{\mathrm{Y}_0}}\sh{O}_\mathrm{Y}.
\]
Aplicando $(\mathbf{0}_\mathrm{I},~6.1.2)$ a la sucesión exacta
\[
  0\longrightarrow(\sh{F}'_0)_{x_0}
  \longrightarrow(\sh{F}_0)_{x_0}
  \longrightarrow(\operatorname{Coker}u_0)_{x_0}
  \longrightarrow0,
\]
concluimos que $u_x$ es inyectivo, lo que demuestra la condición $a)$.

Por último, se deduce de \sref{IV.11.1.2} que el conjunto $\mathrm{U}_0$ de los puntos $z_0\in\mathrm{X}_0$ tales que $(u_0\otimes 1)_{z_0}$ es inyectivo es abierto. Por planitud, para todo $z\in\mathrm{X}$ situado sobre un punto $z_0\in\mathrm{U}_0$, el homomorfismo $(u\otimes 1)_z$ es inyectivo. El conjunto de tales puntos contiene, por tanto, la imagen inversa $\mathrm{U}$ de $\mathrm{U}_0$ en $\mathrm{X}$ y es, en consecuencia, una vecindad de $x$, lo que termina la demostración.
\end{proof}

\begin{theorem}[11.3.8]
\label{IV.11.3.8}
Sea $f:\mathrm{X}\to\mathrm{Y}$ un morfismo localmente de presentación finita, sea $\sh{F}$ un $\sh{O}_\mathrm{X}$-módulo cuasicoherente de presentación finita y sea $(g_i)_{1\leqslant i\leqslant n}$ una sucesión de secciones de $\sh{O}_\mathrm{X}$ sobre $\mathrm{X}$; pongamos $\sh{F}_i=\sh{F}/\sum_{j=1}^{i}g_j\sh{F}$ para $0\leqslant i\leqslant n$ (con $\sh{F}_0=\sh{F}$). Sea $x$ un punto de $\mathrm{X}$, sea $y=f(x)$ y pongamos $\mathrm{X}_y=f^{-1}(y)=\mathrm{X}\times_\mathrm{Y}\Spec(\boldsymbol{k}(y))$, $\sh{F}_y=\sh{F}\otimes_{\sh{O}_y}\boldsymbol{k}(y)$, que es un $\sh{O}_{\mathrm{X}_y}$-módulo. Supongamos que los $(g_i)_x$ pertenecen al ideal maximal $\mathfrak{m}_x$. Las condiciones siguientes son equivalentes:
\begin{enumerate}[label=\emph{\alph*)}]
  \item La sucesión $((g_i)_x)$ es $\sh{F}_x$-regular y los $\sh{F}_i$ ($0\leqslant i\leqslant n$) son $f$-planos en el punto $x$.
  \item La sucesión $((g_i)_x)$ es $\sh{F}_x$-regular y $\sh{F}_n$ es $f$-plano en el punto $x$.
  \item[\emph{b$'$)}] Existe una vecindad $\mathrm{U}$ de $x$ tal que la sucesión $(g_i|\mathrm{U})$ es $(\sh{F}|\mathrm{U})$-cuasirregular, y $\sh{F}_n$ es $f$-plano en el punto $x$.
\oldpage[IV-3]{136}
  \item $\sh{F}$ es $f$-plano en el punto $x$, y la sucesión $((g_i\otimes 1)_x)$ de elementos de $\sh{O}_{\mathrm{X}_y,x}$, imágenes de los $(g_i)_x$, es $(\sh{F}_y)_x$-regular.
  \item $\sh{F}$ es $f$-plano en el punto $x$, y para todo morfismo $\mathrm{Y}'\to\mathrm{Y}$ y todo punto $x'\in\mathrm{X}'=\mathrm{X}\times_\mathrm{Y}\mathrm{Y}'$ situado sobre $x$, si se pone $\sh{F}'=\sh{F}\otimes_\mathrm{Y}\mathrm{Y}'$, la sucesión $(g_i\otimes 1)_{x'}$ de elementos de $\sh{O}_{\mathrm{X}',x'}$ es $\sh{F}'_{x'}$-regular.
\end{enumerate}
Además, el conjunto de los $x\in\mathrm{X}$ que satisfacen estas condiciones es abierto en el conjunto $\mathrm{Z}$ de los $x$ tales que $g_i(x)=0$ para todo $i$.
\end{theorem}

\begin{proof}
El hecho de que $a)$ implique $d)$ se deduce inmediatamente de $(\mathbf{0},~15.1.15)$, y $c)$ es un caso particular de $d)$; además, $a)$ implica trivialmente $b)$. Demostremos que $b)$ o $c)$ implica $a)$. Como los $\sh{F}_i$ son de presentación finita, el hecho de que $c)$ implique $a)$ se deduce inmediatamente de \sref{IV.11.3.7} por inducción sobre $n$, y esto muestra también que el conjunto de los $x\in\mathrm{X}$ que satisfacen $c)$ es abierto en $\mathrm{Z}$.

Para demostrar que $b)$ implica $a)$, la inducción sobre $n$ nos reduce inmediatamente al caso $n=1$; escribimos $g$ en lugar de $g_1$. Como la cuestión es local en $\mathrm{X}$ y en $\mathrm{Y}$, podemos suponer que $\mathrm{Y}=\Spec(\mathrm{A})$ y $\mathrm{X}=\Spec(\mathrm{B})$ son afines, que $\mathrm{B}$ es una $\mathrm{A}$-álgebra de presentación finita y que $\sh{F}=\widetilde{\mathrm{M}}$, donde $\mathrm{M}$ es un $\mathrm{B}$-módulo de presentación finita. Por $(8.9.1)$, podemos escribir entonces
\[
  \mathrm{B}=\mathrm{B}_0\otimes_{\mathrm{A}_0}\mathrm{A},
  \qquad
  \mathrm{M}=\mathrm{M}_0\otimes_{\mathrm{A}_0}\mathrm{A},
\]
en donde $\mathrm{A}_0$ es un subanillo noetheriano de $\mathrm{A}$, $\mathrm{B}_0$ es una $\mathrm{A}_0$-álgebra de tipo finito y $\mathrm{M}_0$ es un $\mathrm{B}_0$-módulo de tipo finito. Además, $\mathrm{A}$ es el límite inductivo filtrante de sus subanillos $\mathrm{A}_\lambda$ que son $\mathrm{A}_0$-álgebras de tipo finito (y, por tanto, noetherianos). Si ponemos
\[
  \mathrm{B}_\lambda=\mathrm{B}_0\otimes_{\mathrm{A}_0}\mathrm{A}_\lambda,
  \qquad
  \mathrm{M}_\lambda=\mathrm{M}_0\otimes_{\mathrm{A}_0}\mathrm{A}_\lambda,
\]
entonces $\mathrm{B}_\lambda$ es una $\mathrm{A}_\lambda$-álgebra de tipo finito, $\mathrm{M}_\lambda$ es un $\mathrm{B}_\lambda$-módulo de tipo finito, y
\[
  \mathrm{B}=\varinjlim_\lambda\mathrm{B}_\lambda,
  \qquad
  \mathrm{M}=\varinjlim_\lambda\mathrm{M}_\lambda.
\]
Existen, por tanto, un índice $\lambda$ y un elemento $g_\lambda\in\mathrm{B}_\lambda$ tales que $g=g_\lambda\otimes 1$. Volviendo a la notación geométrica y poniendo
\[
  \mathrm{Y}_\lambda=\Spec(\mathrm{A}_\lambda),
  \qquad
  \mathrm{X}_\lambda=\Spec(\mathrm{B}_\lambda),
  \qquad
  \sh{F}_\lambda=\widetilde{\mathrm{M}_\lambda},
\]
bastará demostrar que existe $\mu\geqslant\lambda$ tal que, en el punto $x_\mu\in\mathrm{X}_\mu$ que es la proyección de $x$, $(g_\mu)_{x_\mu}$ es $(\sh{F}_\mu)_{x_\mu}$-regular y $\sh{F}_\mu/g_\mu\sh{F}_\mu$ es plano sobre $\mathrm{Y}_\mu$ en $x_\mu$. En efecto, se deducirá entonces de $(\mathbf{0},~15.1.16)$ que $\sh{F}_\mu$ es plano sobre $\mathrm{Y}_\mu$ en $x_\mu$ y, por consiguiente, que $\sh{F}$ es plano sobre $\mathrm{Y}$ en $x$.

El hecho de que $b)$ implique $a)$ se deducirá, por tanto, de la proposición siguiente.

\begin{proposition}[11.3.9]
\label{IV.11.3.9}
Con las notaciones y las hipótesis generales de $(8.5.1)$ y $(8.8.1)$, supongamos que $f_\alpha:\mathrm{X}_\alpha\to\mathrm{Y}_\alpha$ es un morfismo localmente de presentación finita. Sean $\sh{F}_\alpha$ y $\sh{G}_\alpha$ dos $\sh{O}_{\mathrm{X}_\alpha}$-módulos cuasicoherentes de presentación finita, sea $h_\alpha:\sh{F}_\alpha\to\sh{G}_\alpha$ un homomorfismo de $\sh{O}_{\mathrm{X}_\alpha}$-módulos y sea $\sh{H}_\alpha$ su conúcleo. Por último, sea $x$ un punto de $\mathrm{X}$ y sea $x_\lambda$ su proyección en $\mathrm{X}_\lambda$ para $\lambda\geqslant\alpha$. Para que $h_x$ sea inyectivo y $\sh{H}=\operatorname{Coker}h$ sea $f$-plano en $x$, es necesario y suficiente que exista $\lambda\geqslant\alpha$ tal que $(h_\lambda)_{x_\lambda}$ sea inyectivo y $\sh{H}_\lambda=\operatorname{Coker}h_\lambda$ sea $f_\lambda$-plano en $x_\lambda$. Además, el conjunto de los $x\in\mathrm{X}$ que poseen las propiedades anteriores es abierto en $\mathrm{X}$.
\end{proposition}

\begin{proof}
Recordemos que, por la exactitud por la derecha del producto tensorial,
\[
  \sh{H}_\mu=v^*_{\lambda\mu}(\sh{H}_\lambda)
  \quad\text{para }\lambda\leqslant\mu,
  \qquad
  \sh{H}=v^*_\lambda(\sh{H}_\lambda),
\]
lo que justifica las notaciones.

La suficiencia de la condición se deduce de que, si la sucesión
\[
  0\longrightarrow(\sh{F}_\lambda)_{x_\lambda}
  \longrightarrow(\sh{G}_\lambda)_{x_\lambda}
  \longrightarrow(\sh{H}_\lambda)_{x_\lambda}
  \longrightarrow0
\]
es exacta y $(\sh{H}_\lambda)_{x_\lambda}$ es un $\sh{O}_{f_\lambda(x_\lambda)}$-módulo plano, entonces $\sh{H}_x$ es un $\sh{O}_{f(x)}$-módulo plano por cambio de base $(2.1.4)$ y la sucesión
\[
  0\longrightarrow\sh{F}_x
  \longrightarrow\sh{G}_x
  \longrightarrow\sh{H}_x
  \longrightarrow0
\]
es exacta $(2.1.8)$. Para demostrar que

\oldpage[IV-3]{137}

la condición es necesaria, observemos que $\sh{H}$ es de presentación finita. Como la cuestión es local en $\mathrm{X}$ y en $\mathrm{Y}$, podemos suponer que $\mathrm{X}$ y $\mathrm{Y}$ son afines y, por \sref{IV.11.3.1}, que $\sh{H}$ es $f$-plano. Enunciemos ahora el lema siguiente.

\begin{lemma}[11.3.9.1]
\label{IV.11.3.9.1}
Sea $f:\mathrm{X}\to\mathrm{Y}$ un morfismo localmente de presentación finita, sean $\sh{G}$ y $\sh{H}$ dos $\sh{O}_\mathrm{X}$-módulos cuasicoherentes de presentación finita tales que $\sh{H}$ es $f$-plano, y sea $p:\sh{G}\to\sh{H}$ un homomorfismo sobreyectivo. Entonces $\operatorname{Ker}(p)$ es un $\sh{O}_\mathrm{X}$-módulo de presentación finita.
\end{lemma}

\begin{proof}
Podemos, en efecto, suponer que $\mathrm{X}$ y $\mathrm{Y}$ son afines y que existen un morfismo $\mathrm{Y}\to\mathrm{Y}_0$, donde $\mathrm{Y}_0$ es noetheriano; un morfismo $f_0:\mathrm{X}_0\to\mathrm{Y}_0$ de tipo finito tal que $\mathrm{X}=\mathrm{X}_0\times_{\mathrm{Y}_0}\mathrm{Y}$ y $f=(f_0)_{(\mathrm{Y})}$; dos $\sh{O}_{\mathrm{X}_0}$-módulos coherentes $\sh{G}_0$, $\sh{H}_0$; y un homomorfismo $p_0:\sh{G}_0\to\sh{H}_0$ tales que $\sh{G}$, $\sh{H}$ y $p$ se obtienen de $\sh{G}_0$, $\sh{H}_0$ y $p_0$ por cambio de base $(8.9.1\text{ y }8.5.2)$. Podemos suponer además que $p_0$ es sobreyectivo $(8.5.7)$ y que $\sh{H}_0$ es $f_0$-plano \sref{IV.11.2.7}. Entonces, si $\sh{K}_0=\operatorname{Ker}(p_0)$, $\sh{K}_0$ es un $\sh{O}_{\mathrm{X}_0}$-módulo coherente $(\mathbf{0}_\mathrm{I},~5.3.4)$ y, por $(\mathbf{0}_\mathrm{I},~6.1.2)$, $\sh{K}=\operatorname{Ker}(p)$ se obtiene de $\sh{K}_0$ por cambio de base y es, por tanto, de presentación finita.
\end{proof}

Demostrado este lema, pongamos $\sh{K}=\operatorname{Ker}(\sh{G}\to\sh{H})$, que es, por tanto, de presentación finita. Existe un homomorfismo canónico $q:\sh{F}\to\sh{K}$ y, por hipótesis, $q_x:\sh{F}_x\to\sh{K}_x$ es un isomorfismo. Por consiguiente, según $(\mathbf{0}_\mathrm{I},~5.2.7)$, existe una vecindad $\mathrm{U}$ de $x$ tal que $q|_{\mathrm{U}}$ es un isomorfismo; tras restringir $\mathrm{X}$, podemos suponer que la sucesión
\[
  0\to\sh{F}\xrightarrow{h}\sh{G}\to\sh{H}\to 0
\]
es \emph{exacta}. En estas condiciones, se deduce de \sref{IV.11.2.6} que, para todo $\lambda$ suficientemente grande, $\sh{H}_\lambda$ es $f_\lambda$-plano. Pongamos $\sh{K}_\lambda=\operatorname{Ker}(\sh{G}_\lambda\to\sh{H}_\lambda)$ y $\sh{K}_\mu=v^*_{\lambda\mu}(\sh{K}_\lambda)$ para $\mu\geqslant\lambda$. Se deduce de $(\mathbf{0}_\mathrm{I},~6.1.2)$ que
\[
  \sh{K}_\mu=\operatorname{Ker}(\sh{G}_\mu\to\sh{H}_\mu)
  \quad\text{y}\quad
  \sh{K}=v^*_\lambda(\sh{K}_\lambda)
  =\operatorname{Ker}(\sh{G}\to\sh{H})=\sh{F}.
\]
Por otra parte, para todo $\mu\geqslant\lambda$ existe un homomorfismo canónico $q_\mu:\sh{F}_\mu\to\sh{K}_\mu$ tal que $q_\mu=v^*_{\lambda\mu}(q_\lambda)$ y $q=v^*_\mu(q_\mu)$. Como $q$ es un isomorfismo, se deduce de $(8.5.2.4)$ y \sref{IV.11.3.9.1} que existe $\mu\geqslant\lambda$ para el cual $q_\mu$ es un isomorfismo y, por consiguiente, $h_\mu:\sh{F}_\mu\to\sh{G}_\mu$ es inyectivo.
\end{proof}

Volvamos a la demostración de \sref{IV.11.3.8}.

Como el conjunto de los $x\in\mathrm{X}$ que satisfacen $b)$ es abierto en $\mathrm{X}$, está claro que $b)$ implica $b')$. Demostremos, por último, que $b')$ implica $c)$. En primer lugar, $\sh{F}_n$ es $f$-plano en una vecindad de $x$ \sref{IV.11.3.1}; podemos, por tanto, reducirnos al caso en que $\mathrm{U}=\mathrm{X}$ y $\sh{F}_n$ es $f$-plano. Puesto que, por definición, $\operatorname{gr}^*_\mathscr{J}(\sh{F})$ es isomorfo a $\sh{F}_n\otimes_{\sh{O}_\mathrm{Y}}\sh{O}_\mathrm{Y}[\mathrm{T}_1,\ldots,\mathrm{T}_n]$ $(\mathbf{0},~15.2.2)$, es $f$-plano, y se deduce de \sref{IV.11.3.4},~(i), que el propio $\sh{F}$ es $f$-plano en una vecindad de $x$. Por otra parte, si $(\mathscr{J}_y)_x$ es el ideal de $\sh{O}_{\mathrm{X}_y,x}$ generado por los $(g_i\otimes 1)_x$, se deduce de \sref{IV.11.2.9.2} que, en el diagrama
\[
\xymatrix{
  ((\operatorname{gr}^0_{\mathscr{J}_x}(\sh{F}_x))[\mathrm{T}_1,\ldots,\mathrm{T}_n])\otimes_{\sh{O}_y}\boldsymbol{k}(y) \ar[r] \ar[d] & \operatorname{gr}^*_{\mathscr{J}_x}(\sh{F}_x)\otimes_{\sh{O}_y}\boldsymbol{k}(y) \ar[d] \\
  (\operatorname{gr}^0_{(\mathscr{J}_y)_x}((\sh{F}_y)_x))[\mathrm{T}_1,\ldots,\mathrm{T}_n] \ar[r] & \operatorname{gr}^*_{(\mathscr{J}_y)_x}((\sh{F}_y)_x)
}
\]
\oldpage[IV-3]{138}
las flechas verticales son isomorfismos. Como la primera flecha horizontal es un isomorfismo, también lo es la segunda; por consiguiente, la sucesión $((g_i\otimes 1)_x)$ es $(\sh{F}_y)_x$-cuasirregular y, por tanto, regular, ya que $\mathrm{X}_y$ es localmente de tipo finito sobre $\boldsymbol{k}(y)$.
\end{proof}

\begin{theorem}[11.3.10]
\label{IV.11.3.10}
(Criterio de planitud por las fibras). --- Sea $\mathrm{S}$ un preesquema, sean $g:\mathrm{X}\to\mathrm{S}$ y $h:\mathrm{Y}\to\mathrm{S}$ dos morfismos, sea $f:\mathrm{X}\to\mathrm{Y}$ un $\mathrm{S}$-morfismo, sea $\sh{F}$ un $\sh{O}_\mathrm{X}$-módulo cuasicoherente, sea $x$ un punto de $\mathrm{X}$, sea $y=f(x)$ y sea $s=h(y)=g(x)$. Supongamos que se cumple una de las dos hipótesis siguientes:
\begin{enumerate}[label=\arabic{enumi}\textsuperscript{o}]
  \item $\mathrm{S}$, $\mathrm{X}$ y $\mathrm{Y}$ son localmente noetherianos y $\sh{F}$ es coherente.
  \item $g$ y $h$ son localmente de presentación finita y $\sh{F}$ es de presentación finita.
\end{enumerate}
Entonces, con las notaciones de $(9.4.1)$, si $\sh{F}_x\neq 0$, las condiciones siguientes son equivalentes:
\begin{enumerate}[label=\emph{\alph*)}]
  \item $\sh{F}$ es $g$-plano en el punto $x$ y $\sh{F}_s$ es $f_s$-plano en el punto $x$.
  \item $h$ es plano en el punto $y$ y $\sh{F}$ es $f$-plano en el punto $x$.
\end{enumerate}
Además, bajo la hipótesis 2\textsuperscript{o}, el conjunto de los $x\in\mathrm{X}$ que satisfacen la condición $b)$ es abierto en $\mathrm{X}$.
\end{theorem}

\begin{proof}
La última afirmación resulta de aplicar \sref{IV.11.3.1}, por una parte, a $\sh{O}_\mathrm{Y}$ y al morfismo $h$ y, por otra, a $\sh{F}$ y al morfismo $f$ (que es localmente de presentación finita).

Como $g=h\circ f$, está claro que $b)$ implica $a)$ sin suponer 1\textsuperscript{o} ni 2\textsuperscript{o} $(2.1.6$ y $2.1.4)$. Para demostrar que $a)$ implica $b)$, podemos reducirnos al caso en que $\mathrm{S}$, $\mathrm{X}$ y $\mathrm{Y}$ son afines; bajo la hipótesis 2\textsuperscript{o}, aplicando \sref{IV.11.2.7}, nos reducimos además al caso en que $\mathrm{S}$ es \emph{noetheriano}, es decir, al caso en que se cumple la hipótesis 1\textsuperscript{o}. La afirmación equivale entonces al lema siguiente, que mejora $(\mathbf{0}_{\mathrm{III}},~10.2.5)$:

\begin{lemma}[11.3.10.1]
\label{IV.11.3.10.1}
Sean $\rho:\mathrm{A}\to\mathrm{B}$ y $\sigma:\mathrm{B}\to\mathrm{C}$ homomorfismos de anillos locales noetherianos, sea $k$ el cuerpo residual de $\mathrm{A}$ y sea $\mathrm{M}$ un $\mathrm{C}$-módulo $\neq 0$ de tipo finito. Entonces las condiciones siguientes son equivalentes:
\begin{enumerate}[label=\emph{\alph*)}]
  \item $\mathrm{M}$ es un $\mathrm{A}$-módulo plano y $\mathrm{M}\otimes_\mathrm{A} k$ es un $(\mathrm{B}\otimes_\mathrm{A} k)$-módulo plano.
  \item $\mathrm{B}$ es un $\mathrm{A}$-módulo plano y $\mathrm{M}$ es un $\mathrm{B}$-módulo plano.
\end{enumerate}
\end{lemma}

Estableceremos primero el lema más general siguiente:

\begin{lemma}[11.3.10.2]
\label{IV.11.3.10.2}
Sea $\mathrm{A}$ un anillo conmutativo, sea $\mathrm{B}$ una $\mathrm{A}$-álgebra conmutativa, sea $\mathfrak{J}$ un ideal de $\mathrm{A}$ y sea $\mathrm{M}$ un $\mathrm{B}$-módulo. Consideremos primero las condiciones siguientes:
\begin{enumerate}[label=\emph{(\roman*)}]
  \item $\mathfrak{J}$ es nilpotente.
  \item $\mathrm{B}$ es noetheriano y $\mathrm{M}$ está idealmente separado para la topología $\mathfrak{J}\mathrm{B}$-preádica $(\mathbf{0}_{\mathrm{III}},~10.2.1)$.
  \item $\mathfrak{J}\mathrm{B}$ está contenido en el radical de $\mathrm{B}$.
\end{enumerate}
Consideremos, por otra parte, las cuatro propiedades siguientes:
\begin{enumerate}[label=\emph{\alph*)}]
  \item $\mathrm{M}$ es un $\mathrm{B}$-módulo plano.
  \item $\mathrm{B}$ es un $\mathrm{A}$-módulo plano.
  \item $\mathrm{M}$ es un $\mathrm{A}$-módulo plano y $\mathrm{M}/\mathfrak{J}\mathrm{M}$ es un $(\mathrm{B}/\mathfrak{J}\mathrm{B})$-módulo plano.
  \item $\mathrm{B}/\mathfrak{J}\mathrm{B}$ es un $(\mathrm{A}/\mathfrak{J})$-módulo plano y, para todo ideal maximal $\mathfrak{m}\supset\mathfrak{J}\mathrm{B}$ de $\mathrm{B}$, se tiene $\mathfrak{m}\mathrm{M}\neq\mathrm{M}$.
\end{enumerate}
Entonces:
\begin{enumerate}[label=\arabic{enumi}\textsuperscript{o}]
  \item Si se cumple una de las condiciones (i) o (ii), la conjunción de $a)$ y $b)$ implica $c)$, y $c)$ implica $a)$.
  \item \oldpage[IV-3]{139} Si se cumple la condición (i) o la conjunción de (ii) y (iii), la conjunción de $c)$ y $d)$ implica la conjunción de $a)$ y $b)$.
\end{enumerate}
\end{lemma}

\begin{proof}
1\textsuperscript{o} La primera afirmación es inmediata $(\mathbf{0}_\mathrm{I},~6.2.1)$. Supongamos, pues, que se cumple $c)$ y demostremos $a)$. Consideremos los anillos graduados $\operatorname{gr}^*_\mathfrak{J}(\mathrm{A})$, $\operatorname{gr}^*_\mathfrak{J}(\mathrm{B})$ y el módulo graduado $\operatorname{gr}^*_\mathfrak{J}(\mathrm{M})$ (tanto sobre $\operatorname{gr}^*_\mathfrak{J}(\mathrm{A})$ como sobre $\operatorname{gr}^*_\mathfrak{J}(\mathrm{B})$), relativos a las filtraciones $\mathfrak{J}$-preádicas, así como las aplicaciones canónicas sobreyectivas $(\mathbf{0}_{\mathrm{III}},~10.1.1.2)$
\begin{align*}
  u &: \operatorname{gr}^0(\mathrm{B})\otimes_{\operatorname{gr}^0(\mathrm{A})}\operatorname{gr}^*(\mathrm{A})\to\operatorname{gr}^*(\mathrm{B}) \\
  v &: \operatorname{gr}^0(\mathrm{M})\otimes_{\operatorname{gr}^0(\mathrm{A})}\operatorname{gr}^*(\mathrm{A})\to\operatorname{gr}^*(\mathrm{M}) \\
  w &: \operatorname{gr}^0(\mathrm{M})\otimes_{\operatorname{gr}^0(\mathrm{B})}\operatorname{gr}^*(\mathrm{B})\to\operatorname{gr}^*(\mathrm{M})
\end{align*}
Es claro que se tiene un diagrama conmutativo
\begin{equation}
\xymatrix{
  \operatorname{gr}^0(\mathrm{M})\otimes_{\operatorname{gr}^0(\mathrm{A})}\operatorname{gr}^*(\mathrm{A}) \ar[rr]^{v} \ar[dr]_{1\otimes u} & & \operatorname{gr}^*(\mathrm{M}) \\
  & \operatorname{gr}^0(\mathrm{M})\otimes_{\operatorname{gr}^0(\mathrm{B})}\operatorname{gr}^*(\mathrm{B}) \ar[ur]_{w} &
}
  \label{IV.11.3.10.3}
\tag{11.3.10.3}
\end{equation}

La hipótesis $c)$ implica que $v$ es biyectiva $(\mathbf{0}_{\mathrm{III}},~10.2.1)$; como las otras dos aplicaciones del diagrama son sobreyectivas, también ellas son biyectivas. Puesto que, por la hipótesis $c)$, $\mathrm{M}/\mathfrak{J}\mathrm{M}$ es un $(\mathrm{B}/\mathfrak{J}\mathrm{B})$-módulo plano, se deduce de $(\mathbf{0}_{\mathrm{III}},~10.2.1)$ que $\mathrm{M}$ es un $\mathrm{B}$-módulo plano.

2\textsuperscript{o} Cada una de las condiciones (i), (iii) implica que todo ideal maximal de $\mathrm{B}$ contiene $\mathfrak{J}\mathrm{B}$. Se deduce, por tanto, de 1\textsuperscript{o} y de la conjunción de $c)$ y $d)$ que $\mathrm{M}$ es un $\mathrm{B}$-módulo \emph{fielmente plano} y, en consecuencia, que $\operatorname{gr}^0(\mathrm{M})$ es un $\operatorname{gr}^0(\mathrm{B})$-módulo \emph{fielmente plano} $(\mathbf{0}_{\mathrm{I}},~6.2.1)$. Hemos visto en 1\textsuperscript{o} que la hipótesis $c)$ implica que las tres aplicaciones del diagrama \textup{(11.3.10.3)} son biyectivas; la planitud fiel de $\operatorname{gr}^0(\mathrm{M})$ sobre $\operatorname{gr}^0(\mathrm{B})$ implica, por tanto, que $u$ es biyectiva $(\mathbf{0}_{\mathrm{I}},~6.4.1)$. Por otra parte, las condiciones (ii) y (iii) implican que $\mathrm{B}$, como $\mathrm{A}$-módulo, está idealmente separado para la filtración $\mathfrak{J}$-preádica (Bourbaki, \emph{Alg.\ comm.}, cap.~III, \textsection\,5, n\textsuperscript{o}\,4, prop.~2). Se deduce nuevamente de $(\mathbf{0}_{\mathrm{III}},~10.2.1)$ que, si se cumple la condición (i), o la conjunción de (ii) y (iii), $\mathrm{B}$ es un $\mathrm{A}$-módulo plano.
\end{proof}

\phantomsection\label{IV.11.3.10.4}
\textbf{(11.3.10.4)} Una vez establecido el lema \sref{IV.11.3.10.2}, se deduce \sref{IV.11.3.10.1} tomando como $\mathfrak{J}$ el ideal maximal de $\mathrm{A}$ y observando que se cumplen entonces las condiciones (ii) y (iii) de \sref{IV.11.3.10.2} (Bourbaki, \emph{Alg.\ comm.}, cap.~III, \textsection\,5, n\textsuperscript{o}\,4, prop.~2). Esto termina también la demostración de \sref{IV.11.3.10}.
\end{proof}

\begin{corollary}[11.3.11]
\label{IV.11.3.11}
Sean $g:\mathrm{X}\to\mathrm{S}$ y $h:\mathrm{Y}\to\mathrm{S}$ dos morfismos localmente de presentación finita, y sea $f:\mathrm{X}\to\mathrm{Y}$ un $\mathrm{S}$-morfismo. Las condiciones siguientes son equivalentes:
\begin{enumerate}[label=\emph{\alph*)}]
  \item $g$ es plano y $f_s:\mathrm{X}_s\to\mathrm{Y}_s$ es plano para todo $s\in\mathrm{S}$.
  \item $h$ es plano en todos los puntos de $f(\mathrm{X})$ y $f$ es plano.
\end{enumerate}
\end{corollary}

\begin{proof}
Basta aplicar \sref{IV.11.3.10} con $\sh{F}=\sh{O}_\mathrm{X}$.
\end{proof}

\begin{remark}[11.3.12]
\label{IV.11.3.12}
Sería interesante poder dar demostraciones de \sref{IV.11.3.2} y \sref{IV.11.3.10} que no utilizaran el paso al límite inductivo; para ello bastaría demostrar el criterio siguiente:

Sea $\mathrm{A}$ un anillo, sea $\mathrm{B}$ una $\mathrm{A}$-álgebra de presentación finita, sea $\mathrm{M}$ un $\mathrm{B}$-módulo de presentación finita, sea $\mathfrak{J}$ un \oldpage[IV-3]{140} ideal de $\mathrm{A}$ y sea $\mathfrak{p}$ un ideal primo de $\mathrm{B}$ que contiene $\mathfrak{J}\mathrm{B}$. Entonces las dos condiciones siguientes son equivalentes:
\begin{enumerate}[label=\emph{\alph*)}]
  \item $\mathrm{M}_\mathfrak{p}$ es un $\mathrm{A}$-módulo plano.
  \item $\mathrm{M}_\mathfrak{p}/\mathfrak{J}\mathrm{M}_\mathfrak{p}$ es un $(\mathrm{A}/\mathfrak{J})$-módulo plano y $\operatorname{Tor}_1^\mathrm{A}(\mathrm{A}/\mathfrak{J},\mathrm{M}_\mathfrak{p})=0$.
\end{enumerate}
Cuando $\mathrm{A}$ es \emph{noetheriano}, esto se deduce de $(\mathbf{0}_{\mathrm{III}},~10.2.2)$ aplicado a la $\mathrm{A}$-álgebra noetheriana $\mathrm{B}_\mathfrak{p}$; ¿puede deducirse de aquí el enunciado general mediante el paso al límite inductivo?

Conviene señalar, a este respecto, que tal generalización no es ciertamente posible si se sustituye la condición $b)$ anterior por una cualquiera de las condiciones $c)$, $d)$ de $(\mathbf{0}_{\mathrm{III}},~10.2.1)$, reemplazando $\mathrm{M}$ por $\mathrm{M}_\mathfrak{p}$. Por ejemplo, sea $\mathrm{A}$ un anillo local cuyo ideal maximal $\mathfrak{J}$ es principal y tal que
\[
  \mathfrak{R}=\bigcap_{n\geqslant 0}\mathfrak{J}^{n+1}
\]
no es nulo (por ejemplo, el anillo de gérmenes de funciones reales indefinidamente diferenciables en una vecindad de $0$ en $\mathbf{R}$). Tomemos $\mathrm{B}=\mathrm{A}$, $\mathfrak{p}=\mathfrak{J}$ y $\mathrm{M}=\mathrm{A}/\mathfrak{N}$, donde $\mathfrak{N}$ es un submódulo cíclico no nulo de $\mathfrak{R}$; $\mathrm{M}$ es, por tanto, de presentación finita. Es claro que $\mathrm{M}$ no es un $\mathrm{A}$-módulo plano, pues, al ser de presentación finita, sería entonces libre (Bourbaki, \emph{Alg.\ comm.}, cap.~II, \textsection\,3, n\textsuperscript{o}\,2, cor.~2 a prop.~5), lo cual es absurdo, ya que $\mathfrak{N}\neq 0$. Sin embargo,
\[
  \mathfrak{J}^k\mathrm{M}/\mathfrak{J}^{k+n}\mathrm{M}\simeq \mathrm{A}/\mathfrak{J}^n
\]
para todos los enteros positivos $k$ y $n$; por tanto, $\mathrm{M}$ sí satisface las condiciones $c)$ y $d)$ de $(\mathbf{0}_{\mathrm{III}},~10.2.1)$, ya que $\mathrm{A}/\mathfrak{J}$ es un cuerpo.
\end{remark}

\begin{proposition}[11.3.13]
\label{IV.11.3.13}
Sea $f:\mathrm{X}\to\mathrm{Y}$ un morfismo de preesquemas y sea $x$ un punto de $\mathrm{X}$ tal que $f$ es plano en $x$; pongamos $y=f(x)$.
\begin{enumerate}[label=\emph{(\roman*)}]
  \item Si $\mathrm{X}$ es reducido (resp.\ normal) en el punto $x$, entonces $\mathrm{Y}$ es reducido (resp.\ normal) en el punto $y$.
  \item Supongamos además que $f$ es de presentación finita en el punto $x$. Entonces, si $\mathrm{Y}$ es reducido (resp.\ normal) en el punto $y$ y si el morfismo $f$ es reducido (resp.\ normal) en el punto $x$ $(6.8.1)$, $\mathrm{X}$ es reducido (resp.\ normal) en el punto $x$.
\end{enumerate}
\end{proposition}

\begin{proof}
La primera afirmación se recuerda aquí únicamente por comodidad, pues ya se demostró en $(2.1.13)$. Para demostrar (ii), podemos limitarnos al caso en que $\mathrm{X}=\Spec(\mathrm{B})$ y $\mathrm{Y}=\Spec(\mathrm{A})$, donde $\mathrm{A}$ es un anillo local reducido (resp.\ un anillo local íntegro e íntegramente cerrado) y $\mathrm{B}$ es una $\mathrm{A}$-álgebra de presentación finita. Podemos escribir entonces $(8.9.1)$
\[
  \mathrm{B}=\mathrm{B}_0\otimes_{\mathrm{A}_0}\mathrm{A},
\]
en donde $\mathrm{A}_0$ es una $\mathbf{Z}$-subálgebra de tipo finito de $\mathrm{A}$ y $\mathrm{B}_0$ es una $\mathrm{A}_0$-álgebra de tipo finito. Sea $(\mathrm{C}_\alpha)$ la familia creciente filtrante de las $\mathrm{A}_0$-subálgebras de tipo finito de $\mathrm{A}$, que son, por tanto, $\mathbf{Z}$-álgebras de tipo finito; se tiene $\mathrm{A}=\varinjlim\mathrm{C}_\alpha$. Distinguimos los dos casos:

1\textsuperscript{o} Supongamos que $\mathrm{A}$ es reducido y que $f$ es reducido en $x$. Si $\mathfrak{m}$ es el ideal maximal de $\mathrm{A}$, pongamos $\mathfrak{p}_\alpha=\mathfrak{m}\cap\mathrm{C}_\alpha$ y $\mathrm{A}'_\alpha=(\mathrm{C}_\alpha)_{\mathfrak{p}_\alpha}$, de modo que también se tiene $\mathrm{A}=\varinjlim\mathrm{A}'_\alpha$ $(5.13.3)$. Pongamos $\mathrm{Y}_\alpha=\Spec(\mathrm{A}'_\alpha)$ y $\mathrm{X}_\alpha=\Spec(\mathrm{B}_0\otimes_{\mathrm{A}_0}\mathrm{A}'_\alpha)$, y sean $f_\alpha:\mathrm{X}_\alpha\to\mathrm{Y}_\alpha$, $x_\alpha$ (resp.\ $y_\alpha$) la proyección de $x$ (resp.\ $y$) en $\mathrm{X}_\alpha$ (resp.\ $\mathrm{Y}_\alpha$). Como $f$ es reducido en $x$, existe $\alpha_0$ tal que $f_\alpha$ es reducido en $x_\alpha$ para $\alpha\geqslant\alpha_0$ $((6.7.8)$ y \sref{IV.11.2.6}); como $\mathrm{Y}_\alpha$ y $\mathrm{X}_\alpha$ son noetherianos y $\mathrm{A}'_\alpha$ es reducido (pues $\mathrm{C}_\alpha$ lo es), se deduce de $(3.3.5)$ que $\mathrm{X}_\alpha$ es reducido en $x_\alpha$. Pero
\[
  \mathrm{B}=\varinjlim(\mathrm{B}_0\otimes_{\mathrm{A}_0}\mathrm{A}'_\alpha)
\]
y, en consecuencia, $\sh{O}_{\mathrm{X},x}=\varinjlim\sh{O}_{\mathrm{X}_\alpha,x_\alpha}$ $(5.13.3)$; por tanto, $\sh{O}_{\mathrm{X},x}$ es reducido $(5.13.2)$.

2\textsuperscript{o} Supongamos que $\mathrm{A}$ es íntegramente cerrado y que $f$ es normal en $x$. Como $\mathrm{C}_\alpha$ es universalmente japonés $(7.7.4)$, su clausura integral $\mathrm{C}'_\alpha$ es una $\mathrm{C}_\alpha$-álgebra finita, evidentemente contenida en $\mathrm{A}$. Pongamos $\mathfrak{p}'_\alpha=\mathfrak{m}\cap\mathrm{C}'_\alpha$ y $\mathrm{A}''_\alpha=(\mathrm{C}'_\alpha)_{\mathfrak{p}'_\alpha}$, de modo que $\mathrm{A}''_\alpha$ es un anillo local noetheriano íntegro e íntegramente cerrado, y $\mathrm{A}=\varinjlim\mathrm{A}''_\alpha$ $(5.13.3)$. Pongamos $\mathrm{Y}'_\alpha=\Spec(\mathrm{A}''_\alpha)$ y $\mathrm{X}'_\alpha=\Spec(\mathrm{B}_0\otimes_{\mathrm{A}_0}\mathrm{A}''_\alpha)$, y sean $f'_\alpha:\mathrm{X}'_\alpha\to\mathrm{Y}'_\alpha$, $x'_\alpha$ (resp.\ $y'_\alpha$) la proyección de $x$ (resp.\ $y$) en $\mathrm{X}'_\alpha$ (resp.\ $\mathrm{Y}'_\alpha$). Como $f$ es normal en $x$, existe $\alpha_0$ tal que, para $\alpha\geqslant\alpha_0$, $f'_\alpha$ es normal en $x'_\alpha$ $((6.7.8)$ y \sref{IV.11.2.6}); \oldpage[IV-3]{141} como $\mathrm{X}'_\alpha$ y $\mathrm{Y}'_\alpha$ son noetherianos y $\mathrm{A}''_\alpha$ es íntegramente cerrado, se deduce de $(6.5.4)$ que $\mathrm{X}'_\alpha$ es normal en $x'_\alpha$. Pero los morfismos
\[
  \Spec(\mathrm{A}''_\beta)\longrightarrow\Spec(\mathrm{A}''_\alpha)
  \qquad(\alpha\leqslant\beta)
\]
son dominantes; por tanto, puesto que en virtud de \sref{IV.11.3.1} podemos suponer que los $f'_\alpha$ son planos para $\alpha\geqslant\alpha_0$, se deduce de $(2.3.7)$ que toda componente irreducible de $\mathrm{X}'_\beta$ domina una componente irreducible de $\mathrm{X}'_\alpha$. Concluimos entonces de $(5.13.4)$, aplicado a los esquemas
\[
  \Spec\!\left(\sh{O}_{\mathrm{X}_0,x_0}\otimes_{\mathrm{A}_0}\mathrm{A}''_\alpha\right),
\]
que $\mathrm{X}$ es normal en $x$.
\end{proof}

\begin{corollary}[11.3.14]
\label{IV.11.3.14}
Sea $f:\mathrm{X}\to\mathrm{Y}$ un morfismo localmente de presentación finita, sea $x$ un punto de $\mathrm{X}$ y sea $y=f(x)$. Si $\mathrm{Y}$ es geométricamente unirramificado en $y$ $(6.15.1)$ y si $f$ es normal en $x$ $(6.8.1)$, entonces $\mathrm{X}$ es geométricamente unirramificado en $x$.
\end{corollary}

\begin{proof}
Podemos reducirnos evidentemente al caso en que $\mathrm{Y}=\Spec(\mathrm{A})$, con $\mathrm{A}$ un dominio local íntegro geométricamente unirramificado y $y$ el punto cerrado de $\mathrm{Y}$. Sea $\mathrm{A}'$ la clausura integral de $\mathrm{A}$; pongamos $\mathrm{Y}'=\Spec(\mathrm{A}')$ y sea $y'$ el punto cerrado de $\mathrm{Y}'$. Así, el morfismo $g:\mathrm{Y}'\to\mathrm{Y}$ es radicial en $y$ $(6.15.3)$, integral y birracional. Pongamos $\mathrm{X}'=\mathrm{X}\times_\mathrm{Y}\mathrm{Y}'$, y sean $f':\mathrm{X}'\to\mathrm{Y}'$ y $g':\mathrm{X}'\to\mathrm{X}$ las proyecciones. El morfismo $g'$ es integral y radicial en $x$ $(6.15.3.1)$. Sea, pues, $x'$ el único punto de $g'^{-1}(x)$ situado sobre $y'$. El morfismo $f'$ es de presentación finita y normal en $x'$ $(6.7.8)$; en consecuencia, el anillo local $\sh{O}_{\mathrm{X}',x'}$ es un dominio íntegro e íntegramente cerrado \sref{IV.11.3.13}. Por otra parte, podemos suponer que $f$ es plano \sref{IV.11.3.1}, y entonces $g'$ es birracional $(6.15.4.1)$. Se deduce que $\Spec(\sh{O}_{\mathrm{X},x})$ es irreducible; como $\sh{O}_{\mathrm{X},x}$ es reducido \sref{IV.11.3.13}, es un dominio íntegro y es geométricamente unirramificado por $(6.15.5)$.
\end{proof}

\begin{proposition}[11.3.15]
\label{IV.11.3.15}
Sea $\mathrm{A}$ un anillo, sea $\mathrm{B}$ una $\mathrm{A}$-álgebra de presentación finita y sea $\mathrm{M}$ un $\mathrm{B}$-módulo de presentación finita que es plano como $\mathrm{A}$-módulo. Entonces existe una sucesión finita $(f_i)_{1\leqslant i\leqslant n}$ de elementos de $\mathrm{A}$ tal que el ideal generado por los $f_i$ es $\mathrm{A}$ y, para $0\leqslant i<n$,
\[
  \frac{\mathrm{M}_{f_{i+1}}}
       {\sum_{j=1}^{i}f_j\mathrm{M}_{f_{i+1}}}
\]
es un módulo libre sobre
\[
  \frac{\mathrm{A}_{f_{i+1}}}
       {\sum_{j=1}^{i}f_j\mathrm{A}_{f_{i+1}}}.
\]
\end{proposition}

Equivalentemente, los $\mathrm{D}(f_i)$ forman un recubrimiento abierto de $\mathrm{X}=\Spec(\mathrm{A})$ y, si se pone $\mathrm{Z}_1=\mathrm{D}(f_1)$ y luego, por inducción,
\[
  \mathrm{Z}_{i+1}
  =\mathrm{D}(f_{i+1})
   \cap\mathrm{V}(f_1\mathrm{A}+\cdots+f_i\mathrm{A}),
\]
los $\mathrm{Z}_i$ forman una partición de $\mathrm{X}$ en subconjuntos localmente cerrados, mientras que
\[
  \frac{\mathrm{A}_{f_{i+1}}}
       {\sum_{j=1}^{i}f_j\mathrm{A}_{f_{i+1}}}
\]
es el anillo de un subpreesquema afín de $\mathrm{X}$ cuyo espacio subyacente es $\mathrm{Z}_{i+1}$.

\begin{proof}
En virtud de \sref{IV.11.2.7}, podemos suponer primero que existen un subanillo noetheriano $\mathrm{A}_0$ de $\mathrm{A}$, una $\mathrm{A}_0$-álgebra de tipo finito $\mathrm{B}_0$ y un $\mathrm{B}_0$-módulo de tipo finito $\mathrm{M}_0$, plano sobre $\mathrm{A}_0$, tales que
\[
  \mathrm{B}=\mathrm{B}_0\otimes_{\mathrm{A}_0}\mathrm{A},
  \qquad
  \mathrm{M}=\mathrm{M}_0\otimes_{\mathrm{A}_0}\mathrm{A}.
\]
Es claro que basta demostrar la proposición para $\mathrm{A}_0$, $\mathrm{B}_0$ y $\mathrm{M}_0$: si los elementos $f_i\in\mathrm{A}_0$ satisfacen en ese caso las condiciones del enunciado, también las satisfacen para $\mathrm{A}$, $\mathrm{B}$ y $\mathrm{M}$, pues
\[
  \frac{\mathrm{M}_{f_{i+1}}}
       {\sum_{j=1}^{i}f_j\mathrm{M}_{f_{i+1}}}
  =
  \left(
    \frac{(\mathrm{M}_0)_{f_{i+1}}}
         {\sum_{j=1}^{i}f_j(\mathrm{M}_0)_{f_{i+1}}}
  \right)\otimes_{\mathrm{A}_0}\mathrm{A}
\]
(Bourbaki, \emph{Alg.\ comm.}, cap.~II, \textsection\,2, n\textsuperscript{o}\,7, prop.~18). Podemos, por tanto, limitarnos en lo sucesivo al caso en que $\mathrm{A}$ es noetheriano.

Observemos ahora que, si $\mathrm{C}$ es un anillo noetheriano, $\mathfrak{N}$ su nilradical y $\mathrm{P}$ un $\mathrm{C}$-módulo plano, se deduce de $(\mathbf{0}_{\mathrm{III}},10.1.2)$ que, para que $\mathrm{P}$ sea un $\mathrm{C}$-módulo libre, \oldpage[IV-3]{142} es necesario y suficiente que $\mathrm{P}\otimes_\mathrm{C}(\mathrm{C}/\mathfrak{N})$ sea un $(\mathrm{C}/\mathfrak{N})$-módulo libre. Obsérvese también que, si $\mathrm{C}$ es un anillo noetheriano reducido, existe $g\in\mathrm{C}$ tal que $\mathrm{C}_g$ es un dominio íntegro.

Aplicamos ahora el lema $(6.9.2)$. Puede definirse por inducción una sucesión $(f_i)_{i\geqslant1}$ de elementos de $\mathrm{A}$ como sigue:
\begin{enumerate}
\item[$1^\circ$] se elige $f_1$ de modo que
\[
  \mathrm{A}_1=(\mathrm{A}_{\mathrm{red}})_{f_1}
\]
sea un dominio íntegro y $\mathrm{M}\otimes_\mathrm{A}\mathrm{A}_1$ sea un $\mathrm{A}_1$-módulo libre;
\item[$2^\circ$] si el ideal $\mathfrak{J}_i$ generado por $f_1,\ldots,f_i$ es igual a $\mathrm{A}$, se pone $f_{i+1}=f_i$;
\item[$3^\circ$] si, por el contrario, $\mathfrak{J}_i\neq\mathrm{A}$, se elige $f_{i+1}\notin\mathfrak{J}_i$ de modo que
\[
  \mathrm{A}_{i+1}
  =\bigl((\mathrm{A}/\mathfrak{J}_i)_{\mathrm{red}}\bigr)_{f_{i+1}}
\]
sea un dominio íntegro y $\mathrm{M}\otimes_\mathrm{A}\mathrm{A}_{i+1}$ sea un $\mathrm{A}_{i+1}$-módulo libre.
\end{enumerate}
Como $\mathrm{A}$ es noetheriano, la sucesión creciente $(\mathfrak{J}_i)$ es estacionaria. Existe, por tanto, $n$ tal que $f_1,\ldots,f_n$ generan el ideal unidad de $\mathrm{A}$, y estos elementos satisfacen las condiciones requeridas porque
\[
  \bigl((\mathrm{A}/\mathfrak{J}_i)_{\mathrm{red}}\bigr)_{f_{i+1}}
  =
  \bigl((\mathrm{A}/\mathfrak{J}_i)_{f_{i+1}}\bigr)_{\mathrm{red}}.
\]
\end{proof}

\begin{proposition}[11.3.16]
\label{IV.11.3.16}
Sea $f:\mathrm{X}\to\mathrm{Y}$ un morfismo fielmente plano de presentación finita, y sea $g:\mathrm{Y}\to\mathrm{Z}$ un morfismo tal que $g\circ f:\mathrm{X}\to\mathrm{Z}$ es de tipo finito (resp.\ de presentación finita). Entonces $g$ es de tipo finito (resp.\ de presentación finita).
\end{proposition}

\begin{proof}
Como $f$ es sobreyectivo y $g\circ f$ es cuasicompacto, $g$ es cuasicompacto $(1.1.3)$. Demostremos primero que, si $g\circ f$ es de presentación finita, entonces $g$ es cuasiseparado. Sea $\mathrm{W}\subset\mathrm{Z}$ un abierto afín y sea $(\mathrm{V}_i)$ un recubrimiento abierto afín finito de $g^{-1}(\mathrm{W})$; basta demostrar que los $\mathrm{V}_i\cap\mathrm{V}_j$ son cuasicompactos $(1.2.6,\ 1.2.7)$. Para cada $i$, sea $(\mathrm{U}_{ih})$ un recubrimiento abierto afín finito de $f^{-1}(\mathrm{V}_i)$. Como $g\circ f$ es de presentación finita, todas las intersecciones $\mathrm{U}_{ih}\cap\mathrm{U}_{jk}$ son cuasicompactas. Puesto que $f$ es sobreyectivo, $\mathrm{V}_i\cap\mathrm{V}_j$ es la unión, cuando varían $h$ y $k$, de las imágenes $f(\mathrm{U}_{ih}\cap\mathrm{U}_{jk})$; es, por tanto, cuasicompacto, pues $f$ es continuo.

La cuestión es, pues, local en $\mathrm{Y}$, y podemos suponer que $\mathrm{Y}=\Spec(\mathrm{B})$ y $\mathrm{Z}=\Spec(\mathrm{A})$ son afines, mientras que $\mathrm{X}$ es la unión de un número finito de abiertos afines $\mathrm{X}_i=\Spec(\mathrm{C}_i)$, donde los $\mathrm{C}_i$ son $\mathrm{B}$-álgebras de tipo finito (resp.\ de presentación finita). Sea $\mathrm{X}'$ la unión disjunta de los $\mathrm{X}_i$, y sea $p:\mathrm{X}'\to\mathrm{X}$ el morfismo cuya restricción a cada $\mathrm{X}_i$ es la inmersión abierta canónica. El morfismo $p$ es fielmente plano y de presentación finita $(1.6.5)$; por consiguiente, $g\circ f\circ p$ es de tipo finito (resp.\ de presentación finita), y $f\circ p$ es fielmente plano y de presentación finita. Podemos suponer además, por tanto, que $\mathrm{X}=\Spec(\mathrm{C})$ es afín, y la proposición se reduce al corolario siguiente.
\end{proof}

\begin{corollary}[11.3.17]
\label{IV.11.3.17}
Sea $\mathrm{A}$ un anillo, sea $\mathrm{B}$ una $\mathrm{A}$-álgebra y sea $\mathrm{C}$ una $\mathrm{B}$-álgebra de presentación finita que es fielmente plana como $\mathrm{B}$-módulo. Si $\mathrm{C}$ es una $\mathrm{A}$-álgebra de tipo finito (resp.\ de presentación finita), entonces $\mathrm{B}$ es una $\mathrm{A}$-álgebra de tipo finito (resp.\ de presentación finita).
\end{corollary}

\begin{proof}
Supongamos primero que $\mathrm{C}$ es de tipo finito sobre $\mathrm{A}$. Sea $(\mathrm{B}_\lambda)$ la familia creciente filtrante de las $\mathrm{A}$-subálgebras de tipo finito de $\mathrm{B}$. Por $(8.8.2)$, existe un índice $\alpha$ tal que
\[
  \mathrm{C}=\mathrm{C}_\alpha\otimes_{\mathrm{B}_\alpha}\mathrm{B},
\]
en donde $\mathrm{C}_\alpha$ es una $\mathrm{B}_\alpha$-álgebra de presentación finita. Además, por $(8.10.5)$ y \sref{IV.11.2.6}, podemos suponer que $\mathrm{C}_\alpha$ es fielmente plana como $\mathrm{B}_\alpha$-módulo. Para $\lambda\geqslant\alpha$, pongamos
\[
  \mathrm{C}_\lambda
  =\mathrm{C}_\alpha\otimes_{\mathrm{B}_\alpha}\mathrm{B}_\lambda.
\]
Entonces $\mathrm{C}=\mathrm{C}_\lambda\otimes_{\mathrm{B}_\lambda}\mathrm{B}$ y $\mathrm{C}_\lambda$ es fielmente plana como $\mathrm{B}_\lambda$-módulo. Como $\mathrm{B}_\lambda\to\mathrm{B}$ es inyectivo, también lo es $\mathrm{C}_\lambda\to\mathrm{C}$. Además, $\mathrm{C}=\varinjlim\mathrm{C}_\lambda$, de modo que todo elemento de $\mathrm{C}$ pertenece a algún $\mathrm{C}_\lambda$; como $\mathrm{C}$ es de tipo finito sobre $\mathrm{A}$, existe $\lambda$ para el cual $\mathrm{C}_\lambda\to\mathrm{C}$ \oldpage[IV-3]{143} es biyectivo. La planitud fiel implica entonces que $\mathrm{B}_\lambda\to\mathrm{B}$ es biyectivo. Así, $\mathrm{B}=\mathrm{B}_\lambda$ es de tipo finito sobre $\mathrm{A}$.

Supongamos ahora que $\mathrm{C}$ es de presentación finita sobre $\mathrm{A}$. Por lo que acabamos de demostrar, $\mathrm{B}$ es de tipo finito sobre $\mathrm{A}$; escribamos
\[
  \mathrm{B}
  =\mathrm{A}[\mathrm{T}_1,\ldots,\mathrm{T}_n]/\widetilde{\mathfrak{J}}
\]
para un ideal $\widetilde{\mathfrak{J}}$. Sea $(\widetilde{\mathfrak{J}}_\lambda)$ la familia filtrante de los ideales de tipo finito de $\mathrm{A}[\mathrm{T}_1,\ldots,\mathrm{T}_n]$ contenidos en $\widetilde{\mathfrak{J}}$. Entonces
\[
  \widetilde{\mathfrak{J}}
  =\varinjlim\widetilde{\mathfrak{J}}_\lambda,
  \qquad
  \mathrm{B}
  =\varinjlim\mathrm{B}_\lambda,
  \qquad
  \mathrm{B}_\lambda
  =\mathrm{A}[\mathrm{T}_1,\ldots,\mathrm{T}_n]/
   \widetilde{\mathfrak{J}}_\lambda.
\]
Aplicando $(8.8.2)$, $(8.10.5)$ y \sref{IV.11.2.6} como antes, se tiene
\[
  \mathrm{C}
  =\mathrm{C}_\alpha\otimes_{\mathrm{B}_\alpha}\mathrm{B},
\]
en donde $\mathrm{C}_\alpha$ es una $\mathrm{B}_\alpha$-álgebra de presentación finita y un $\mathrm{B}_\alpha$-módulo fielmente plano. Para $\lambda\geqslant\alpha$, pongamos
\[
  \mathrm{C}_\lambda
  =\mathrm{C}_\alpha\otimes_{\mathrm{B}_\alpha}\mathrm{B}_\lambda.
\]
Por planitud,
\[
  \mathrm{C}_\lambda
  =
  \mathrm{C}_\alpha\Big/
  \Bigl(
    \mathrm{C}_\alpha\otimes_{\mathrm{B}_\alpha}
    (\widetilde{\mathfrak{J}}_\lambda/
     \widetilde{\mathfrak{J}}_\alpha)
  \Bigr),
\]
y, análogamente,
\[
  \mathrm{C}
  =
  \mathrm{C}_\alpha\Big/
  \Bigl(
    \mathrm{C}_\alpha\otimes_{\mathrm{B}_\alpha}
    (\widetilde{\mathfrak{J}}/
     \widetilde{\mathfrak{J}}_\alpha)
  \Bigr).
\]
Como $\mathrm{C}$ y $\mathrm{C}_\alpha$ son ambas $\mathrm{A}$-álgebras de presentación finita, el ideal
\[
  \mathrm{C}_\alpha\otimes_{\mathrm{B}_\alpha}
  (\widetilde{\mathfrak{J}}/
   \widetilde{\mathfrak{J}}_\alpha)
\]
es de tipo finito en $\mathrm{C}_\alpha$ $(1.4.4)$. Es la unión de los ideales de tipo finito
\[
  \mathrm{C}_\alpha\otimes_{\mathrm{B}_\alpha}
  (\widetilde{\mathfrak{J}}_\lambda/
   \widetilde{\mathfrak{J}}_\alpha).
\]
Por consiguiente, existe $\lambda\geqslant\alpha$ tal que
\[
  \mathrm{C}_\alpha\otimes_{\mathrm{B}_\alpha}
  (\widetilde{\mathfrak{J}}/
   \widetilde{\mathfrak{J}}_\alpha)
  =
  \mathrm{C}_\alpha\otimes_{\mathrm{B}_\alpha}
  (\widetilde{\mathfrak{J}}_\lambda/
   \widetilde{\mathfrak{J}}_\alpha).
\]
La planitud fiel da entonces
$\widetilde{\mathfrak{J}}=\widetilde{\mathfrak{J}}_\lambda$, lo que demuestra que $\mathrm{B}$ es de presentación finita sobre $\mathrm{A}$.
\end{proof}

\subsection{Descenso de la planitud por morfismos cualesquiera: caso de un preesquema base artiniano}
\label{subsection:IV.11.4}

\begin{theorem}[11.4.1]
\label{IV.11.4.1}
Sea $\mathrm{A}$ un anillo, sea $\mathfrak{J}$ un ideal nilpotente de $\mathrm{A}$ y sea $u_\lambda:\mathrm{A}\to\mathrm{B}_\lambda$ $(\lambda\in\mathrm{L})$
una familia de homomorfismos de anillos tal que la intersección de los núcleos de los $u_\lambda$ es $0$.
Sea $\mathrm{M}$ un $\mathrm{A}$-módulo tal que, para todo $\lambda\in\mathrm{L}$, $\mathrm{M}\otimes_\mathrm{A}\mathrm{B}_\lambda$ es un $\mathrm{B}_\lambda$-módulo libre y que $\mathrm{M}\otimes_\mathrm{A}(\mathrm{A}/\mathfrak{J})$
es un $(\mathrm{A}/\mathfrak{J})$-módulo libre. Entonces $\mathrm{M}$ es un $\mathrm{A}$-módulo libre. Si, además, el conjunto de índices $\mathrm{L}$ es
finito, en el enunciado anterior puede sustituirse en todas partes ``libre'' por ``plano''.
\end{theorem}

\begin{proof}
En ambos casos bastará demostrar que $\mathrm{M}$ es un $\mathrm{A}$-módulo \emph{plano}: en efecto,
cuando $\mathrm{M}\otimes_\mathrm{A}(\mathrm{A}/\mathfrak{J})$ es un $(\mathrm{A}/\mathfrak{J})$-módulo libre, se deduce que $\mathrm{M}$ es un $\mathrm{A}$-módulo
libre en virtud de $(\mathbf{0}_{\mathrm{III}},~10.1.2)$.

Utilizaremos el lema siguiente, que generaliza $(\mathbf{0}_{\mathrm{III}},~10.2.1)$:

\begin{lemma}[11.4.1.1]
\label{IV.11.4.1.1}
Sea $\mathrm{A}$ un anillo provisto de una filtración finita $(\mathfrak{J}_n)_{0\leq n\leq\mathrm{N}+1}$ con
$\mathfrak{J}_0=\mathrm{A}$, $\mathfrak{J}_{\mathrm{N}+1}=0$. Sea $\mathrm{M}$ un $\mathrm{A}$-módulo provisto de la filtración $(\mathfrak{J}_n\mathrm{M})_{0\leq n\leq\mathrm{N}+1}$, y denotemos
por $\operatorname{gr}(\mathrm{A})$ y $\operatorname{gr}(\mathrm{M})$, respectivamente, el anillo y el módulo graduados. Supongamos que $\mathrm{M}\otimes_\mathrm{A}(\mathrm{A}/\mathfrak{J}_1)$
es un $(\mathrm{A}/\mathfrak{J}_1)$-módulo plano y que el homomorfismo canónico
\[
\label{IV.11.4.1.2}
\operatorname{gr}_0(\mathrm{M})\otimes_{\operatorname{gr}_0(\mathrm{A})}\operatorname{gr}(\mathrm{A}) \to \operatorname{gr}(\mathrm{M})
\tag{11.4.1.2}
\]
es inyectivo. Entonces $\mathrm{M}$ es un $\mathrm{A}$-módulo plano.
\end{lemma}

El homomorfismo canónico {\hyperlink{ega4.c04.lem.11-4-1-2}{(11.4.1.2)}} se define del mismo modo que
el de $(\mathbf{0}_{\mathrm{III}},~10.1.1)$, siendo en grado $n$ el homomorfismo
\[
(\mathrm{M}/\mathfrak{J}_1\mathrm{M})\otimes(\mathfrak{J}_n/\mathfrak{J}_{n+1})=(\mathrm{M}\otimes\mathfrak{J}_n)/(\operatorname{Im}(\mathrm{M}\otimes\mathfrak{J}_{n+1})+\operatorname{Im}(\mathfrak{J}_1\mathrm{M}\otimes\mathfrak{J}_n))\to\mathfrak{J}_n\mathrm{M}/\mathfrak{J}_{n+1}\mathrm{M}
\]
inducido por el homomorfismo canónico $\mathrm{M}\otimes\mathfrak{J}_n\to\mathfrak{J}_n\mathrm{M}$ al pasar a los cocientes.
El lema se demuestra por inducción sobre $\mathrm{N}$, pues no hay nada que demostrar para $\mathrm{N}=0$.
\oldpage[IV-3]{144}
Las condiciones impuestas a $\mathrm{M}$ implican, por la hipótesis de inducción, que $\mathrm{M}\otimes_\mathrm{A}(\mathrm{A}/\mathfrak{J}_\mathrm{N})$
es un $(\mathrm{A}/\mathfrak{J}_\mathrm{N})$-módulo plano. Observemos ahora que $\mathfrak{J}_\mathrm{N}^2\subset\mathfrak{J}_{\mathrm{N}+1}=0$ si $\mathrm{N}\geq 1$, y
$(\mathrm{M}/\mathfrak{J}_1\mathrm{M})\otimes(\mathfrak{J}_\mathrm{N}/\mathfrak{J}_{\mathrm{N}+1})=(\mathrm{M}/\mathfrak{J}_\mathrm{N}\mathrm{M})\otimes(\mathfrak{J}_\mathrm{N}/\mathfrak{J}_{\mathrm{N}+1})=(\mathrm{M}/\mathfrak{J}_\mathrm{N}\mathrm{M})\otimes\mathfrak{J}_\mathrm{N}$; de ahí el homomorfismo canónico
\[
(\mathrm{M}/\mathfrak{J}_\mathrm{N}\mathrm{M})\otimes\mathfrak{J}_\mathrm{N} \to \mathfrak{J}_\mathrm{N}\mathrm{M}=\mathfrak{J}_\mathrm{N}\mathrm{M}/\mathfrak{J}_\mathrm{N}^2\mathrm{M}
\]
es inyectivo. Aplicando $(\mathbf{0}_{\mathrm{III}},~10.2.1)$ a la filtración $\mathfrak{J}_\mathrm{N}$-preádica, concluimos que
$\mathrm{M}$ es un $\mathrm{A}$-módulo plano.

Para aplicar este lema a \sref{IV.11.4.1}, denotemos por $\mathfrak{J}_n$ el ideal de $\mathrm{A}$ que es la intersección
de las imágenes inversas $u_\lambda^{-1}(\mathfrak{J}^n\mathrm{B}_\lambda)$: es inmediato que $\mathfrak{J}_0=\mathrm{A}$, $\mathfrak{J}_m\mathfrak{J}_n\subset\mathfrak{J}_{m+n}$
para $m\geq 0$, $n\geq 0$; además, si $\mathfrak{J}^{\mathrm{N}+1}=0$, se tiene también $\mathfrak{J}_{\mathrm{N}+1}=0$, ya que la intersección de los
núcleos de los $u_\lambda$ es $0$.

Provistamos $\mathrm{A}$ de la filtración $(\mathfrak{J}_n)$ y $\mathrm{M}$ de la filtración $(\mathfrak{J}_n\mathrm{M})$ y, para cada $\lambda$,
provistamos $\mathrm{B}_\lambda$ y $\mathrm{N}_\lambda=\mathrm{M}\otimes_\mathrm{A}\mathrm{B}_\lambda$ de las filtraciones $\mathfrak{J}$-preádicas; consideremos para cada $\lambda$ el
diagrama conmutativo
\[
\xymatrix{
\operatorname{gr}^0(\mathrm{M})\otimes_{\operatorname{gr}^0(\mathrm{A})}\operatorname{gr}(\mathrm{A}) \ar[r]^-{\varphi} \ar[d]_{f_\lambda} & \operatorname{gr}(\mathrm{M}) \ar[d]^{g_\lambda} \\
\operatorname{gr}^0_\mathfrak{J}(\mathrm{N}_\lambda)\otimes_{\operatorname{gr}^0_\mathfrak{J}(\mathrm{B}_\lambda)}\operatorname{gr}_\mathfrak{J}(\mathrm{B}_\lambda) \ar[r]_-{\varphi_\lambda} & \operatorname{gr}_\mathfrak{J}(\mathrm{N}_\lambda)
}
\]
en donde las flechas horizontales son los homomorfismos canónicos {\hyperlink{ega4.c04.lem.11-4-1-2}{(11.4.1.2)}} y las flechas
verticales se deducen de los homomorfismos canónicos $\mathrm{A}\to\mathrm{B}_\lambda$ y $\mathrm{M}\to\mathrm{N}_\lambda$. La hipótesis
de que $\mathrm{N}_\lambda$ es un $\mathrm{B}_\lambda$-módulo plano implica que $\varphi_\lambda$ es inyectivo $(\mathbf{0}_{\mathrm{III}},~10.2.1)$ y, por tanto,
$\operatorname{Ker}(\varphi)\subset\operatorname{Ker}(f_\lambda)$. Pongamos $\mathrm{A}_0=\mathrm{A}/\mathfrak{J}_1$, $\mathrm{M}_0=\mathrm{M}\otimes_\mathrm{A}\mathrm{A}_0$ y observemos que
\[
\operatorname{gr}^0_\mathfrak{J}(\mathrm{N}_\lambda)=\mathrm{N}_\lambda/\mathfrak{J}\mathrm{N}_\lambda=(\mathrm{M}/\mathfrak{J}\mathrm{M})\otimes_\mathrm{A}\mathrm{B}_\lambda=(\mathrm{M}/\mathfrak{J}_1\mathrm{M})\otimes_{\mathrm{A}}(\mathrm{B}_\lambda/\mathfrak{J}\mathrm{B}_\lambda)
\]
y, por tanto,
\[
\operatorname{gr}^0_\mathfrak{J}(\mathrm{N}_\lambda)\otimes_{\operatorname{gr}^0_\mathfrak{J}(\mathrm{B}_\lambda)}\operatorname{gr}_\mathfrak{J}(\mathrm{B}_\lambda)=\mathrm{M}_0\otimes_{\mathrm{A}_0}\operatorname{gr}_\mathfrak{J}(\mathrm{B}_\lambda).
\]
\oldpage[IV-3]{145}
El homomorfismo $f_\lambda$ puede escribirse, por tanto,
\[
1\otimes\operatorname{gr}(u_\lambda) : \mathrm{M}_0\otimes_{\mathrm{A}_0}\operatorname{gr}(\mathrm{A}) \to \mathrm{M}_0\otimes_{\mathrm{A}_0}\operatorname{gr}_\mathfrak{J}(\mathrm{B}_\lambda),
\]
y, como por hipótesis $\mathrm{M}_0$ es un $\mathrm{A}_0$-módulo \emph{plano}, el núcleo de $f_\lambda$ es igual a $\mathrm{M}_0\otimes_{\mathrm{A}_0}\mathrm{R}_\lambda$,
en donde $\mathrm{R}_\lambda$ es el núcleo de $\operatorname{gr}(u_\lambda):\operatorname{gr}(\mathrm{A})\to\operatorname{gr}_\mathfrak{J}(\mathrm{B}_\lambda)$. Basta, por tanto, demostrar que
$\bigcap_{\lambda\in\mathrm{L}}(\mathrm{M}_0\otimes_{\mathrm{A}_0}\mathrm{R}_\lambda)=0$. Ahora bien, por definición de los $\mathfrak{J}_n$, la intersección de los núcleos de los homomorfismos
$\mathfrak{J}_n\mathrm{A}/\mathfrak{J}_{n+1}\mathrm{A}\to\mathfrak{J}^n\mathrm{B}_\lambda/\mathfrak{J}^{n+1}\mathrm{B}_\lambda$, cuando $\lambda$ recorre $\mathrm{L}$, es $0$; en otras
palabras, $\bigcap_{\lambda\in\mathrm{L}}\mathrm{R}_\lambda=0$. Supongamos primero que $\mathrm{M}_0$ es un $\mathrm{A}_0$-módulo libre; tomando
una base de $\mathrm{M}_0$, vemos inmediatamente que $\mathrm{M}_0\otimes_{\mathrm{A}_0}(\bigcap_{\lambda\in\mathrm{L}}\mathrm{R}_\lambda)=\bigcap_{\lambda\in\mathrm{L}}(\mathrm{M}_0\otimes_{\mathrm{A}_0}\mathrm{R}_\lambda)$, de donde se obtiene la
proposición en este caso. Cuando $\mathrm{L}$ es finito, la fórmula anterior sigue siendo cierta bajo
la sola hipótesis de que $\mathrm{M}_0$ sea un $\mathrm{A}_0$-módulo plano $(\mathbf{0}_\mathrm{I},~6.1.3)$, lo que termina la
demostración.
\end{proof}

\begin{remark}[11.4.2]
\label{IV.11.4.2}
La conclusión de \sref{IV.11.4.1} puede ser falsa si $\mathrm{L}$ es infinito
y solo se supone que $\mathrm{M}\otimes_\mathrm{A}(\mathrm{A}/\mathfrak{J})$ es un $(\mathrm{A}/\mathfrak{J})$-módulo plano. Por ejemplo,
sea $\mathrm{V}$ un anillo de valoración discreta, sea $\mathrm{K}$ su cuerpo de fracciones y sea $\mathrm{A}=\mathrm{V}[\mathrm{T}]/(\mathrm{T}^2)$
($\mathrm{T}$ una indeterminada); tomemos como $\mathfrak{J}$ la imagen de $(\mathrm{T})$ en $\mathrm{A}$, de modo que $\mathrm{A}/\mathfrak{J}=\mathrm{V}$, y
tomemos $\mathrm{M}=\mathrm{K}$, que es un $(\mathrm{A}/\mathfrak{J})$-módulo y, por tanto, es igual a $\mathrm{M}\otimes_\mathrm{A}(\mathrm{A}/\mathfrak{J})$; además, $\mathrm{M}$ es
un $(\mathrm{A}/\mathfrak{J})$-módulo plano, pero no un $\mathrm{A}$-módulo plano: la inclusión $\mathfrak{J}\hookrightarrow\mathrm{A}$ induce, al tensorizar con $\mathrm{K}$,
la aplicación nula $\mathfrak{J}\otimes_{\mathrm{A}}\mathrm{K}\simeq\mathrm{K}\longrightarrow\mathrm{K}$,
que no es inyectiva, pues $\mathrm{K}\neq0$. Por otra parte, consideremos el ideal maximal $\mathfrak{m}$ del anillo local noetheriano
$\mathrm{A}$; se tiene $\mathfrak{m}\mathrm{K}=\mathrm{K}$ y, por tanto, $\mathrm{M}\otimes_\mathrm{A}(\mathrm{A}/\mathfrak{m}^n)=0$ para todo entero $n$; los $(\mathrm{A}/\mathfrak{m}^n)$-módulos
$\mathrm{M}\otimes_\mathrm{A}(\mathrm{A}/\mathfrak{m}^n)$ son, pues, planos para todo $n$, y la intersección de los $\mathfrak{m}^n$ es $0$.
\end{remark}

\begin{corollary}[11.4.3]
\label{IV.11.4.3}
Sea $\mathrm{A}$ un anillo semilocal cuyo radical $\mathfrak{J}$ es nilpotente (por
ejemplo, un anillo artiniano), y sea $u_\lambda:\mathrm{A}\to\mathrm{B}_\lambda$ $(\lambda\in\mathrm{L})$ una familia de homomorfismos de anillos tal
que la intersección de los núcleos de los $u_\lambda$ es $0$. Para que un $\mathrm{A}$-módulo $\mathrm{M}$ sea plano, es necesario y
suficiente que, para todo $\lambda\in\mathrm{L}$, $\mathrm{M}\otimes_\mathrm{A}\mathrm{B}_\lambda$ sea un $\mathrm{B}_\lambda$-módulo plano.
\end{corollary}

\begin{proof}
Como $\mathrm{A}/\mathfrak{J}$ es un producto directo de un número finito de cuerpos (Bourbaki, \emph{Alg. comm.},
cap.~II, \S~3, n\textsuperscript{o}~5, prop.~16) y $\mathfrak{J}$ es nilpotente, $\mathrm{A}$ es un producto directo de un
número finito de anillos \emph{locales} $\mathrm{A}_i$ cuyo radical es nilpotente (\emph{loc.\ cit.}, \S~4, n\textsuperscript{o}~3, cor.~1
de la prop.~15), y $\mathrm{M}$ es, por tanto, una suma directa de $\mathrm{A}_i$-módulos $\mathrm{M}_i$, y cada $\mathrm{M}_i$ es
anulado por los $\mathrm{A}_j$ de índice $j\neq i$; para que $\mathrm{M}$ sea un $\mathrm{A}$-módulo plano, es necesario y
suficiente que cada $\mathrm{M}_i$ sea un $\mathrm{A}_i$-módulo plano; además, la intersección de los núcleos
de los homomorfismos $\mathrm{A}_i\longrightarrow\mathrm{A}\xrightarrow{u_\lambda}\mathrm{B}_\lambda$ es $0$, y $\mathrm{M}_i\otimes_{\mathrm{A}_i}\mathrm{B}_\lambda$ es un sumando directo
de $\mathrm{M}\otimes_\mathrm{A}\mathrm{B}_\lambda$. Podemos, por tanto, limitarnos al caso en que $\mathrm{A}$ es también \emph{local}. Entonces $\mathrm{A}/\mathfrak{J}$ es un
cuerpo, de modo que $\mathrm{M}\otimes_\mathrm{A}(\mathrm{A}/\mathfrak{J})$ es un $(\mathrm{A}/\mathfrak{J})$-módulo libre, y basta ver que, para
todo $\lambda$, $\mathrm{M}\otimes_\mathrm{A}\mathrm{B}_\lambda$ es un $\mathrm{B}_\lambda$-módulo \emph{libre}, en virtud de \sref{IV.11.4.1}. Ahora bien, si se pone $\mathfrak{J}_\lambda=\mathfrak{J}\mathrm{B}_\lambda$,
$(\mathrm{M}\otimes_\mathrm{A}\mathrm{B}_\lambda)\otimes_{\mathrm{B}_\lambda}(\mathrm{B}_\lambda/\mathfrak{J}_\lambda)=(\mathrm{M}\otimes_\mathrm{A}(\mathrm{A}/\mathfrak{J}))\otimes_{\mathrm{A}/\mathfrak{J}}(\mathrm{B}_\lambda/\mathfrak{J}_\lambda)$ es un $(\mathrm{B}_\lambda/\mathfrak{J}_\lambda)$-módulo libre, y $\mathfrak{J}_\lambda$ es
nilpotente. La conclusión se deduce, pues, de la hipótesis de que $\mathrm{M}\otimes_\mathrm{A}\mathrm{B}_\lambda$ es un $\mathrm{B}_\lambda$-módulo
plano y de $(\mathbf{0}_{\mathrm{III}},~10.1.2)$.
\end{proof}

\begin{corollary}[11.4.4]
\label{IV.11.4.4}
Sea $\mathrm{A}$ un anillo y sea $\mathrm{M}$ un $\mathrm{A}$-módulo; supongamos que existe
un ideal nilpotente $\mathfrak{J}$ de $\mathrm{A}$ tal que $\mathrm{M}\otimes_\mathrm{A}(\mathrm{A}/\mathfrak{J})$ es un $(\mathrm{A}/\mathfrak{J})$-módulo libre. Entonces el conjunto de los
ideales $\mathfrak{K}$ de $\mathrm{A}$ tales que $\mathrm{M}\otimes_\mathrm{A}(\mathrm{A}/\mathfrak{K})$ es un $(\mathrm{A}/\mathfrak{K})$-módulo libre admite un elemento mínimo $\mathfrak{K}_0$
\oldpage[IV-3]{146}
(que es también el mínimo de los ideales $\mathfrak{K}$ tales que $\mathrm{M}\otimes_\mathrm{A}(\mathrm{A}/\mathfrak{K})$ es un $(\mathrm{A}/\mathfrak{K})$-módulo \emph{plano}).
\end{corollary}

\begin{proof}
Para todo homomorfismo $u:\mathrm{A}\to\mathrm{A}'$ tal que $\mathrm{M}\otimes_\mathrm{A}\mathrm{A}'$ es un $\mathrm{A}'$-módulo libre (resp.\ un
$\mathrm{A}'$-módulo plano), es necesario y suficiente que $u$ se factorice a través de $\mathrm{A}\to\mathrm{A}/\mathfrak{K}_0\to\mathrm{A}'$ (o, equivalentemente, que $\mathfrak{K}_0\mathrm{A}'=0$).

El hecho de que la intersección $\mathfrak{K}_0$ de los ideales $\mathfrak{K}$ para los que $\mathrm{M}\otimes_\mathrm{A}(\mathrm{A}/\mathfrak{K})$ es un
$(\mathrm{A}/\mathfrak{K})$-módulo libre sea ella misma el mínimo de estos ideales se deduce de \sref{IV.11.4.1}, aplicado al
anillo $\mathrm{A}/\mathfrak{K}_0$, al ideal nilpotente $\mathfrak{J}/\mathfrak{K}_0$ y a los homomorfismos $\mathrm{A}/\mathfrak{K}_0\to\mathrm{A}/\mathfrak{K}$,
cuyos núcleos tienen intersección $0$. Si $\mathrm{A}'$ es una $(\mathrm{A}/\mathfrak{K}_0)$-álgebra, entonces
$\mathrm{M}\otimes_\mathrm{A}\mathrm{A}'=(\mathrm{M}\otimes_\mathrm{A}(\mathrm{A}/\mathfrak{K}_0))\otimes_{\mathrm{A}/\mathfrak{K}_0}\mathrm{A}'$ y, por tanto, $\mathrm{M}\otimes_\mathrm{A}\mathrm{A}'$ es un $\mathrm{A}'$-módulo libre. Recíprocamente, si $\mathrm{A}'$ es una $\mathrm{A}$-álgebra tal que $\mathrm{M}\otimes_\mathrm{A}\mathrm{A}'$ es un $\mathrm{A}'$-módulo libre y si $\mathfrak{K}$ es el núcleo
del homomorfismo $\mathrm{A}\to\mathrm{A}'$, se deduce de \sref{IV.11.4.1}, aplicado al anillo $\mathrm{A}/\mathfrak{K}$, al
$(\mathrm{A}/\mathfrak{K})$-módulo $\mathrm{M}\otimes_\mathrm{A}(\mathrm{A}/\mathfrak{K})$, al ideal nilpotente $(\mathfrak{J}+\mathfrak{K})/\mathfrak{K}$ y al homomorfismo \emph{inyectivo}
$\mathrm{A}/\mathfrak{K}\to\mathrm{A}'$, que $\mathrm{M}\otimes_\mathrm{A}(\mathrm{A}/\mathfrak{K})$ es un $(\mathrm{A}/\mathfrak{K})$-módulo libre, y por tanto $\mathfrak{K}\supset\mathfrak{K}_0$. El hecho de que
pueda sustituirse ``libre'' por ``plano'' en lo anterior (conservando, naturalmente,
la hipótesis de que $\mathrm{M}\otimes_\mathrm{A}(\mathrm{A}/\mathfrak{J})$ es un $(\mathrm{A}/\mathfrak{J})$-módulo \emph{libre}) se deduce de $(\mathbf{0}_{\mathrm{III}},~10.1.2)$,
como vimos al comienzo de la demostración de \sref{IV.11.4.1}.
\end{proof}

\begin{proposition}[11.4.5]
\label{IV.11.4.5}
Sea $\mathrm{Y}$ un preesquema irreducible, sea $f:\mathrm{X}\to\mathrm{Y}$ un morfismo de
presentación finita y sea $\sh{F}$ un $\sh{O}_\mathrm{X}$-módulo de presentación finita. Existe un abierto no vacío $\mathrm{U}$
de $\mathrm{Y}$ y un subpreesquema cerrado $\mathrm{Z}$ de $\mathrm{U}$, de presentación finita sobre $\mathrm{U}$, con la propiedad siguiente:
para todo cambio de base $\mathrm{U}'\to\mathrm{U}$, si se pone $\mathrm{X}'=\mathrm{X}\times_\mathrm{Y}\mathrm{U}'$, $f'=f_{(\mathrm{U}')}$ y $\sh{F}'=\sh{F}\otimes_{\sh{O}_\mathrm{Y}}\sh{O}_{\mathrm{U}'}$,
entonces, para que $\sh{F}'$ sea $f'$-plano, es necesario y suficiente que el morfismo $\mathrm{U}'\to\mathrm{U}$ se factorice
a través de $\mathrm{U}'\to\mathrm{Z}\to\mathrm{U}$. Tal esquema $\mathrm{Z}$ tiene el mismo espacio subyacente que $\mathrm{U}$. Supongamos además que $\mathrm{Y}$ es afín
y sea $(\mathrm{W}_i)$ un recubrimiento abierto afín finito de $\mathrm{X}$; entonces puede suponerse elegido $\mathrm{U}$
de modo que, si $\mathrm{U}'=\operatorname{Spec}(\mathrm{A}')$ es un esquema afín y $\mathrm{U}'\to\mathrm{U}$ un morfismo que se factoriza
a través de $\mathrm{U}'\to\mathrm{Z}\to\mathrm{U}$, cada uno de los $\Gamma(\mathrm{W}_i\times_\mathrm{Y}\mathrm{U}',\sh{F}')$ sea un $\mathrm{A}'$-módulo libre.
\end{proposition}

\begin{proof}
Podemos limitarnos evidentemente al caso en que $\mathrm{Y}=\operatorname{Spec}(\mathrm{A})$ es afín. Utilizando
\sref{IV.8.9.1}, existen un subanillo noetheriano $\mathrm{A}_1$ de $\mathrm{A}$, un morfismo de tipo finito
$f_1:\mathrm{X}_1\to\mathrm{Y}_1=\operatorname{Spec}(\mathrm{A}_1)$ y un $\sh{O}_{\mathrm{X}_1}$-módulo coherente $\sh{F}_1$ tales que $f=f_{(\mathrm{Y})}$ y $\sh{F}=\sh{F}_1\otimes_{\sh{O}_{\mathrm{Y}_1}}\sh{O}_\mathrm{Y}$;
podemos suponer además que los $\mathrm{W}_i$ son imágenes inversas de abiertos afines de $\mathrm{X}_1$.
Observemos asimismo que $\mathrm{Y}_1$ es irreducible, pues el morfismo $\mathrm{Y}\to\mathrm{Y}_1$ es dominante (\textbf{I},~$1.2.7$).
Supongamos demostrada la proposición para $\mathrm{Y}_1$, $\mathrm{X}_1$ y $\sh{F}_1$, y sean $\mathrm{U}_1$ y $\mathrm{Z}_1$ el abierto de $\mathrm{Y}_1$
y el subpreesquema cerrado de $\mathrm{U}_1$ con las propiedades requeridas. Pongamos $\mathrm{U}=\mathrm{U}_1\times_{\mathrm{Y}_1}\mathrm{Y}$ y $\mathrm{Z}=\mathrm{Z}_1\times_{\mathrm{Y}_1}\mathrm{Y}$, sus imágenes inversas. Entonces, si $\mathrm{U}'\to\mathrm{U}$ es un cambio de base
tal que $\sh{F}'=\sh{F}\otimes_{\sh{O}_\mathrm{Y}}\sh{O}_{\mathrm{U}'}$ es $f'$-plano, el morfismo $\mathrm{U}'\to\mathrm{U}_1$ se factoriza
a través de $\mathrm{U}'\to\mathrm{Z}_1\to\mathrm{U}_1$; pero, como $\mathrm{U}'\to\mathrm{U}\to\mathrm{U}_1$ se factoriza también a través de $\mathrm{U}'\to\mathrm{U}\to\mathrm{U}_1$, la definición
del producto de preesquemas muestra que $\mathrm{U}'\to\mathrm{U}$ se factoriza a través de $\mathrm{U}'\to\mathrm{Z}\to\mathrm{U}$.

Podemos, por tanto, limitarnos al caso en que $\mathrm{A}$ es \emph{noetheriano}; sea $\mathfrak{N}$ su nilradical, que es
nilpotente, y pongamos $\mathrm{A}_0=\mathrm{A}_{\mathrm{red}}=\mathrm{A}/\mathfrak{N}$, $\mathrm{Y}_0=\operatorname{Spec}(\mathrm{A}_0)=\mathrm{Y}_{\mathrm{red}}$, $\mathrm{X}_0=\mathrm{X}\times_\mathrm{Y}\mathrm{Y}_0$, $\sh{F}_0=\sh{F}\otimes_{\sh{O}_\mathrm{Y}}\sh{O}_{\mathrm{Y}_0}$,
$\mathrm{W}_{i0}=\mathrm{W}_i\times_\mathrm{Y}\mathrm{Y}_0$; si $\mathrm{M}_i=\Gamma(\mathrm{W}_i,\sh{F})$, entonces $\mathrm{M}_{i0}=\Gamma(\mathrm{W}_{i0},\sh{F}_0)$ es igual a $\mathrm{M}_i\otimes_\mathrm{A}\mathrm{A}_0$.
Como $\mathrm{A}_0$ es \emph{íntegro}, podemos, en virtud del teorema de planitud genérica $(6.9.2)$,
sustituyendo si es preciso $\mathrm{Y}$ por un abierto no vacío de $\mathrm{Y}$, suponer que los $\mathrm{M}_{i0}$ son
$\mathrm{A}_0$-módulos libres. En virtud de \sref{IV.11.4.4}, existe, por tanto, para cada $i$ un ideal $\mathfrak{J}_i\subset\mathfrak{N}$ tal
\oldpage[IV-3]{147}
que las $\mathrm{A}$-álgebras $\mathrm{A}'$ para las que los $\mathrm{M}_i\otimes_\mathrm{A}\mathrm{A}'$ son $\mathrm{A}'$-módulos libres (o planos)
son exactamente aquellas para las que $\mathfrak{J}_i\mathrm{A}'=0$. Sea $\mathfrak{J}=\sum_i\mathfrak{J}_i$, que es un ideal contenido
en $\mathfrak{N}$; decir que $\sh{F}\otimes_{\mathrm{A}}\mathrm{A}'$ es $\mathrm{A}'$-plano equivale a decir que todos los $\mathrm{M}_i\otimes_\mathrm{A}\mathrm{A}'$ son
$\mathrm{A}'$-módulos planos, es decir, que $\mathfrak{J}\mathrm{A}'=0$. Se deduce que, tomando $\mathrm{Z}=\mathrm{V}(\mathfrak{J})$, se
resuelve la cuestión, pues, para que un morfismo $\mathrm{U}'\to\mathrm{U}$ tenga la propiedad del enunciado,
es necesario y suficiente que, para todo abierto afín $\mathrm{V}'\subset\mathrm{U}'$, el morfismo $\mathrm{V}'\to\mathrm{U}$
tenga la misma propiedad.
\end{proof}

\begin{corollary}[11.4.6]
\label{IV.11.4.6}
Sea $\mathrm{A}$ un anillo tal que $\operatorname{Spec}(\mathrm{A})$ es irreducible, sea $\mathfrak{p}$ su único
ideal primo minimal, sea $\mathrm{B}$ una $\mathrm{A}$-álgebra de presentación finita y sea $\mathrm{M}$ un $\mathrm{B}$-módulo de presentación finita.
Supongamos que $\mathrm{M}_\mathfrak{p}$ es un $\mathrm{A}_\mathfrak{p}$-módulo plano; entonces existe $t\in\mathrm{A}-\mathfrak{p}$ tal que $\mathrm{M}_t$ es un $\mathrm{A}_t$-módulo libre.
\end{corollary}

\begin{proof}
Aplicando \sref{IV.11.4.5} con $\mathrm{Y}=\operatorname{Spec}(\mathrm{A})$, $\mathrm{X}=\operatorname{Spec}(\mathrm{B})$, $\sh{F}=\widetilde{\mathrm{M}}$, podemos (sustituyendo
si es preciso $\mathrm{A}$ por $\mathrm{A}_h$, donde $h\in\mathrm{A}-\mathfrak{p}$) suponer que existe un ideal de tipo finito $\mathfrak{J}$ de $\mathrm{A}$
tal que las $\mathrm{A}$-álgebras $\mathrm{A}'$ para las que $\mathrm{M}\otimes_\mathrm{A}\mathrm{A}'$ es un $\mathrm{A}'$-módulo libre (o plano) son
exactamente aquellas tales que $\mathfrak{J}\mathrm{A}'=0$. En particular, la hipótesis implica que $\mathrm{M}_\mathfrak{p}$ es un
$\mathrm{A}_\mathfrak{p}$-módulo plano, y por tanto $\mathfrak{J}\mathrm{A}_\mathfrak{p}=0$, o $\mathfrak{J}_\mathfrak{p}=0$. Ahora bien, como $\mathfrak{J}$ es de tipo finito,
existe $t\in\mathrm{A}-\mathfrak{p}$ tal que $\mathfrak{J}_t=0$, o equivalentemente $\mathfrak{J}\mathrm{A}_t=0$, y por tanto $\mathrm{M}_t$ es un $\mathrm{A}_t$-módulo libre.
\end{proof}

\begin{proposition}[11.4.7]
\label{IV.11.4.7}
Sea $\mathrm{A}$ un anillo noetheriano, sea $\mathfrak{J}$ un ideal de $\mathrm{A}$ y sea $\mathrm{M}$ un $\mathrm{A}$-módulo
idealmente separado $(\mathbf{0}_{\mathrm{III}},~10.2.1)$ para la topología $\mathfrak{J}$-preádica. Sea $(\mathfrak{p}_i)_{1\leq i\leq r}$ una familia finita
de ideales primos de $\mathrm{A}$ que contienen $\mathfrak{J}$; para todo entero $n\geq 0$, sea $\mathfrak{p}_i^{(n)}$ la $n$-ésima
potencia simbólica (es decir, el núcleo del homomorfismo $\mathrm{A}\to\mathrm{A}_{\mathfrak{p}_i}/\mathfrak{p}_i^n\mathrm{A}_{\mathfrak{p}_i}$); pongamos $\mathfrak{J}_n=\bigcap_{i=1}^{r}\mathfrak{p}_i^{(n)}$, de modo que $\mathfrak{J}_n\supset\mathfrak{J}^n$,
y supongamos que la topología definida por la filtración $(\mathfrak{J}_n)$ es idéntica a la topología $\mathfrak{J}$-preádica
(en otras palabras, para todo $n$, existe $m$ tal que $\mathfrak{J}^n\supset\mathfrak{J}_m$). Para que $\mathrm{M}$ sea un $\mathrm{A}$-módulo plano,
es necesario y suficiente que $\mathrm{M}\otimes_\mathrm{A}(\mathrm{A}/\mathfrak{J})$ sea un $(\mathrm{A}/\mathfrak{J})$-módulo plano y que, para todo $i$, $\mathrm{M}\otimes_\mathrm{A}\mathrm{A}_{\mathfrak{p}_i}=\mathrm{M}_{\mathfrak{p}_i}$
sea un $\mathrm{A}_{\mathfrak{p}_i}$-módulo plano.
\end{proposition}

\begin{proof}
Como $\mathrm{M}$ está idealmente separado, basta, en virtud de $(\mathbf{0}_{\mathrm{III}},~10.2.1)$, demostrar
que, para todo $n\geq 0$, $\mathrm{M}\otimes_\mathrm{A}(\mathrm{A}/\mathfrak{J}^n)$ es un $(\mathrm{A}/\mathfrak{J}^n)$-módulo plano; como cada $\mathfrak{J}^n$ contiene
un $\mathfrak{J}_m$, esto equivale a demostrar que, para todo $n$, $\mathrm{M}\otimes_\mathrm{A}(\mathrm{A}/\mathfrak{J}_n)$ es un $(\mathrm{A}/\mathfrak{J}_n)$-módulo
plano. Ahora bien, como $\mathfrak{J}_1\supset\mathfrak{J}$, $\mathrm{M}\otimes_\mathrm{A}(\mathrm{A}/\mathfrak{J}_1)$ es un $(\mathrm{A}/\mathfrak{J}_1)$-módulo plano; en el anillo $\mathrm{A}/\mathfrak{J}_n$,
el ideal $\mathfrak{J}_1/\mathfrak{J}_n$ es nilpotente y, por último, la intersección de los núcleos de los homomorfismos
$\mathrm{A}/\mathfrak{J}_n\to\mathrm{A}_{\mathfrak{p}_i}/\mathfrak{p}_i^n\mathrm{A}_{\mathfrak{p}_i}$ es cero en $\mathrm{A}/\mathfrak{J}_n$, por definición de $\mathfrak{J}_n$. Basta, pues, en virtud de \sref{IV.11.4.1},
verificar que $\mathrm{M}\otimes_\mathrm{A}(\mathrm{A}_{\mathfrak{p}_i}/\mathfrak{p}_i^n\mathrm{A}_{\mathfrak{p}_i})$ es un $(\mathrm{A}_{\mathfrak{p}_i}/\mathfrak{p}_i^n\mathrm{A}_{\mathfrak{p}_i})$-módulo plano, lo cual se deduce de
la hipótesis de que $\mathrm{M}\otimes_\mathrm{A}\mathrm{A}_{\mathfrak{p}_i}$ es un $\mathrm{A}_{\mathfrak{p}_i}$-módulo plano.
\end{proof}

\begin{remark}[11.4.8]
\label{IV.11.4.8}
La hipótesis formulada en \sref{IV.11.4.7} sobre la topología definida
por los $\mathfrak{J}_n$ se cumple si, para todo $n$ suficientemente grande, $\operatorname{Ass}(\mathrm{A}/\mathfrak{J}^n)$ \emph{está contenido en el conjunto
de los $\mathfrak{p}_i$}. En efecto, $\mathfrak{J}^n$ es entonces una intersección de ideales primarios correspondientes a los $\mathfrak{p}_i$, cada uno de los cuales
contiene una potencia simbólica de $\mathfrak{p}_i$, de donde se sigue la conclusión. En particular:
\end{remark}

\begin{corollary}[11.4.9]
\label{IV.11.4.9}
Sea $\mathrm{A}$ un anillo noetheriano, sea $\mathfrak{J}$ un ideal nilpotente de $\mathrm{A}$ y sea $\mathrm{M}$ un
$\mathrm{A}$-módulo. Para que $\mathrm{M}$ sea un $\mathrm{A}$-módulo plano (resp.\ libre), es necesario y suficiente que $\mathrm{M}\otimes_\mathrm{A}(\mathrm{A}/\mathfrak{J})$
sea un $(\mathrm{A}/\mathfrak{J})$-módulo plano (resp.\ libre) y que, para todo ideal primo $\mathfrak{p}\in\operatorname{Ass}(\mathrm{A})$, $\mathrm{M}_\mathfrak{p}$ sea un
$\mathrm{A}_\mathfrak{p}$-módulo plano.
\end{corollary}

La afirmación relativa al caso en que $\mathrm{M}\otimes_\mathrm{A}(\mathrm{A}/\mathfrak{J})$ es libre se deduce de nuevo de la afirmación relativa al caso en que $\mathrm{M}\otimes_\mathrm{A}(\mathrm{A}/\mathfrak{J})$ es plano mediante $(\mathbf{0}_{\mathrm{III}},10.1.2)$.

\begin{corollary}[11.4.10]
\label{IV.11.4.10}
Sea $\mathrm{A}$ un anillo noetheriano, sea $\mathfrak{J}$ un ideal de $\mathrm{A}$ y sea $\mathrm{M}$ un $\mathrm{A}$-módulo.
Supongamos que $\mathrm{M}$ está idealmente separado para la topología $\mathfrak{J}$-preádica y que $\operatorname{gr}_\mathfrak{J}(\mathrm{A})$ es un
$(\mathrm{A}/\mathfrak{J})$-módulo plano. Para que $\mathrm{M}$ sea un $\mathrm{A}$-módulo plano, es necesario y suficiente que $\mathrm{M}\otimes_\mathrm{A}(\mathrm{A}/\mathfrak{J})$ sea
un $(\mathrm{A}/\mathfrak{J})$-módulo plano y que, para todo $\mathfrak{p}\in\operatorname{Ass}(\mathrm{A}/\mathfrak{J})$, $\mathrm{M}_\mathfrak{p}$ sea un $\mathrm{A}_\mathfrak{p}$-módulo plano.
\end{corollary}

\begin{proof}
\oldpage[IV-3]{148}
Teniendo en cuenta \sref{IV.11.4.8}, basta demostrar que $\operatorname{Ass}(\mathrm{A}/\mathfrak{J}^n)$ está contenido en
$\operatorname{Ass}(\mathrm{A}/\mathfrak{J})$ para todo $n$. Ahora bien, si $a\in\mathrm{A}$ no pertenece a ningún $\mathfrak{p}\in\operatorname{Ass}(\mathrm{A}/\mathfrak{J})$, la homotecia
de razón $a$ es inyectiva en $\mathrm{A}/\mathfrak{J}$; como cada $\mathfrak{J}^k/\mathfrak{J}^{k+1}$ es un $(\mathrm{A}/\mathfrak{J})$-módulo
plano, $a$ es también un elemento $(\mathfrak{J}^k/\mathfrak{J}^{k+1})$-regular y, por tanto, $a$ es un elemento $(\mathrm{A}/\mathfrak{J}^n)$-regular para todo $n$
(Bourbaki, \emph{Alg.\ comm.}, cap.~III, \S~2, n\textsuperscript{o}~8, cor.~1 del teorema~1) y, por consiguiente, no pertenece
a ningún ideal primo asociado a $\mathrm{A}/\mathfrak{J}^n$, de donde se sigue el corolario (Bourbaki, \emph{Alg.\ comm.}, cap.~II,
\S~1, n\textsuperscript{o}~1, prop.~2).
\end{proof}

\begin{proposition}[11.4.11]
\label{IV.11.4.11}
Sea $\mathrm{A}$ un anillo local artiniano, sea $\mathfrak{m}$ su ideal maximal,
sea $\mathrm{Y}=\operatorname{Spec}(\mathrm{A})$, sea $y$ el único punto de $\mathrm{Y}$, sea $f:\mathrm{X}\to\mathrm{Y}$ un morfismo localmente de tipo finito y sea $\sh{F}$ un
$\sh{O}_\mathrm{X}$-módulo coherente. Sea $(\mathrm{A}'_\alpha)$ una familia de anillos locales y, para cada $\alpha$, sea $u_\alpha:\mathrm{A}\to\mathrm{A}'_\alpha$
un homomorfismo de anillos (necesariamente local). Pongamos $\mathrm{Y}'_\alpha=\operatorname{Spec}(\mathrm{A}'_\alpha)$, $\mathrm{X}'_\alpha=\mathrm{X}\times_\mathrm{Y}\mathrm{Y}'_\alpha$,
$\sh{F}'_\alpha=\sh{F}\otimes_\mathrm{A}\mathrm{A}'_\alpha$. Sea $x$ un punto
de $\mathrm{X}$, y supongamos que se cumplen las siguientes condiciones:
\begin{enumerate}[label=(\roman*)]
  \item La intersección de los núcleos de los $u_\alpha$ es $0$.
  \item La extensión $k(x)$ del cuerpo residual $k$ de $\mathrm{A}$ es primaria $(4.3.1)$ (condición que se cumple automáticamente
  si $k$ es separablemente cerrado).
\end{enumerate}
\begin{enumerate}[label=(\roman*)]
  \setcounter{enumi}{2}
  \item Para todo $\alpha$, existe un punto $x'_\alpha\in\mathrm{X}'_\alpha$ cuyas proyecciones en $\mathrm{X}$ y $\mathrm{Y}'_\alpha$
  son $x$ y el punto cerrado $y'_\alpha$ de $\mathrm{Y}'_\alpha$, y tal que $\sh{F}'_\alpha$ es $\mathrm{Y}'_\alpha$-plano en el punto $x'_\alpha$.
\end{enumerate}

En estas condiciones, $\sh{F}$ es $f$-plano en el punto $x$.
\end{proposition}

\begin{proof}
Como la cuestión es local en $\mathrm{X}$, podemos limitarnos evidentemente al caso en que $f$ es
un morfismo de tipo finito y suponer $\sh{F}_x\neq 0$. Procederemos en varias etapas.

\medskip
\noindent\textbf{I)} \emph{Reducción al caso en que $\mathrm{A}'_\alpha$ es un anillo local de un $\mathrm{Y}$-preesquema de tipo finito y en que el
cuerpo residual $k'_\alpha$ de $\mathrm{A}'_\alpha$ es una extensión finita de $k$.}

\medskip
Como la reducción se efectúa por separado para cada $\mathrm{A}'_\alpha$, podemos omitir el índice $\alpha$ en
esta parte. Sea $\mathfrak{m}'$ el ideal maximal de $\mathrm{A}'$. Consideremos $\mathrm{A}'$ como
límite inductivo de sus $\mathrm{A}$-subálgebras de tipo finito $\mathrm{B}_\lambda$, y pongamos $\mathfrak{n}_\lambda=\mathfrak{m}'\cap\mathrm{B}_\lambda$; $\mathrm{A}'$ es también
límite inductivo de sus subanillos locales $\mathrm{B}'_\lambda=(\mathrm{B}_\lambda)_{\mathfrak{n}_\lambda}$ $(5.13.3)$, y se tiene evidentemente
$\mathfrak{m}'\cap\mathrm{B}'_\lambda=\mathfrak{n}'_\lambda$, ideal maximal de $\mathrm{B}'_\lambda$. Se deduce, pues, de \sref{IV.11.2.6} que, poniendo
$\mathrm{Z}'_\lambda=\operatorname{Spec}(\mathrm{B}'_\lambda)$, $\sh{F}\otimes_\mathrm{A}\mathrm{B}'_\lambda$ es $\mathrm{Z}'_\lambda$-plano en el punto $x'_\lambda$, proyección de $x'$, siendo la proyección
de $x'_\lambda$ en $\mathrm{Z}'_\lambda$ el punto cerrado $z'_\lambda$ de $\mathrm{Z}'_\lambda$.

Podemos, por tanto, suponer que existen un esquema afín $\mathrm{Z}'$ de tipo finito sobre $\mathrm{Y}$ y un
punto $z'$ de $\mathrm{Z}'$ tales que $\mathrm{A}'=\sh{O}_{\mathrm{Z}',z'}$ y que, si se pone $\mathrm{T}'=\mathrm{X}\times_\mathrm{Y}\mathrm{Z}'$, existe un punto
$t'\in\mathrm{T}'$ cuyas proyecciones en $\mathrm{Z}'$ y $\mathrm{X}$ son $z'$ y $x$, y tal que $\sh{F}\otimes_\mathrm{Y}\mathrm{Z}'$ es $\mathrm{Z}'$-plano en el
punto $t'$. Sea $\mathrm{Z}'_1$ (resp.\ $\mathrm{X}_1$) un subpreesquema cerrado de $\mathrm{Z}'$ (resp.\ $\mathrm{X}$) cuyo espacio subyacente sea
la adherencia de $\{z'\}$ (resp.\ $\{x\}$), y pongamos $\mathrm{T}'_1=\mathrm{X}_1\times_\mathrm{Y}\mathrm{Z}'_1$. El conjunto $\mathrm{U}$ de los
puntos de $\mathrm{T}'$ en los que $\sh{F}\otimes_\mathrm{Y}\mathrm{Z}'$ es $\mathrm{Z}'$-plano es abierto en $\mathrm{T}'$ \sref{IV.11.1.1} y contiene $t'$; por tanto,
\oldpage[IV-3]{149}
$\mathrm{V}=\mathrm{U}\cap\mathrm{T}'_1$ es un abierto no vacío de $\mathrm{T}'_1$. Como el anillo $\mathrm{A}$ es artiniano, es un anillo de Jacobson;
por consiguiente, según $(10.4.6)$, $\mathrm{Z}'_1$ y $\mathrm{T}'_1$ son preesquemas de Jacobson. Existe, pues,
un punto $t'_0$ de $\mathrm{V}$ cerrado en $\mathrm{V}$, y su imagen $z'_0$ en $\mathrm{Z}'_1$ es un punto \emph{cerrado} de $\mathrm{Z}'_1$ $(10.4.7)$.
Sea $f_1$ la restricción de $f$ a $\mathrm{X}_1$, y sea $f'_1:\mathrm{T}'_1\to\mathrm{Z}'_1$ la proyección canónica.
$\mathrm{V}\cap f_1'^{-1}(z'_0)$ es un subconjunto abierto no vacío de
$f_1^{-1}(y)\otimes_{k(y)}k(z'_0)$; como este último preesquema es plano sobre $f_1^{-1}(y)$, un punto maximal $t'_1$ de
$\mathrm{V}\cap f_1'^{-1}(z'_0)$ se encuentra necesariamente sobre el único punto maximal $x$ de $\mathrm{X}_1$ $(2.3.4)$.
Por último, $k(z'_0)$ es una extensión \emph{finita} de $k$ (\textbf{I},~$6.4.2$), y el homomorfismo
$\mathrm{A}=\sh{O}_{\mathrm{Y},y}\to\mathrm{A}'=\sh{O}_{\mathrm{Z}',z'}$ se factoriza como
$\sh{O}_{\mathrm{Y},y}\to\sh{O}_{\mathrm{Z}',z'_0}\to\sh{O}_{\mathrm{Z}',z'}$; por lo tanto, su núcleo está contenido en el de
$\sh{O}_{\mathrm{Y},y}\to\sh{O}_{\mathrm{Z}',z'_0}$. Esto completa la reducción anunciada.

\medskip
\noindent\textbf{II)} \emph{Reducción al caso en que solo hay un número finito de $\mathrm{A}'_\alpha$ y estos son $\mathrm{A}$-álgebras finitas.} --- Sea $\mathfrak{m}'_\alpha$
el ideal maximal de $\mathrm{A}'_\alpha$. Como $\mathrm{A}'_\alpha$ es un anillo local noetheriano, la intersección de los $\mathfrak{m}'^n_\alpha$
es $0$ $(\mathbf{0}_\mathrm{I},~7.3.5)$. La intersección de los $u_\alpha^{-1}(\mathfrak{m}'^n_\alpha)$, para todos los índices $n$ y $\alpha$, es, por tanto,
igual a la intersección de los núcleos de los $u_\alpha$ y, en consecuencia, es $0$ por la hipótesis (i). Como $\mathrm{A}$
es artiniano, un número finito de estos ideales tiene ya intersección $0$; denotémoslos por
$u_i^{-1}(\mathfrak{m}'^{n_i}_i)$ $(1\leq i\leq r)$. Como los $\mathrm{A}'_i$ son noetherianos, los $\mathfrak{m}'^h_i/\mathfrak{m}'^{h+1}_i$ son
$(\mathrm{A}'_i/\mathfrak{m}'_i)$-módulos de longitud finita y, como $\mathrm{A}'_i/\mathfrak{m}'_i$ es un $k$-espacio vectorial de dimensión finita,
$\mathrm{A}''_i=\mathrm{A}'_i/\mathfrak{m}'^{n_i}_i$ es una $\mathrm{A}$-álgebra finita y un anillo local artiniano. La reducción anunciada
se deduce, pues, de $(2.1.4)$.

\medskip
\noindent\textbf{III)} \emph{Fin de la demostración.} --- Supongamos ahora que solo hay un número finito de $\mathrm{A}'_i$ $(1\leq i\leq r)$ y que
son $\mathrm{A}$-álgebras finitas. Para todo $i$, el cuerpo residual $k_i$ de $\mathrm{A}'_i$ es una extensión finita de $k$; utilizando
la hipótesis (ii), se concluye que la imagen inversa de $x$ en $\mathrm{X}'_i$ consta del \emph{único} punto $x'_i$ $(4.3.2)$.
Sea $\mathrm{Y}'$ la unión disjunta de los $\mathrm{Y}'_i$, sea $\mathrm{X}'=\mathrm{X}\times_\mathrm{Y}\mathrm{Y}'$ la unión disjunta de los
$\mathrm{X}'_i$ y pongamos $\sh{F}'=\sh{F}\otimes_\mathrm{Y}\mathrm{Y}'$. La hipótesis implica que $\sh{F}'$ es $\mathrm{Y}'$-plano en todo
punto de la imagen inversa de $x$ por la proyección $p:\mathrm{X}'\to\mathrm{X}$. Como $p$ es de tipo finito, existe, por tanto, un
abierto $\mathrm{U}'\supset p^{-1}(x)$ tal que $\sh{F}'$ es $\mathrm{Y}'$-plano en todo punto de $\mathrm{U}'$ \sref{IV.11.1.1}.
Además, el morfismo $\mathrm{Y}'\to\mathrm{Y}$ es finito porque los $\mathrm{A}'_i$ son $\mathrm{A}$-álgebras finitas. Así, $p$ es un
morfismo finito (\textbf{II},~$6.1.5$) y, por tanto, es cerrado (\textbf{II},~$6.1.10$); en consecuencia, existe una vecindad afín $\mathrm{U}$ de
$x$ en $\mathrm{X}$ tal que $p^{-1}(\mathrm{U})\subset\mathrm{U}'$. Sean $\mathrm{B}$ y $\mathrm{A}'$ los anillos de los esquemas
$\mathrm{U}$ y $\mathrm{Y}'$, respectivamente (así, $\mathrm{A}'$ es el producto directo de los $\mathrm{A}'_i$). Sustituyendo $\mathrm{X}$ por
$\mathrm{U}$, se tiene $\sh{F}=\widetilde{\mathrm{M}}$, donde $\mathrm{M}$ es un $\mathrm{B}$-módulo, y por hipótesis
$\mathrm{M}'=\mathrm{M}\otimes_\mathrm{A}\mathrm{A}'$ es un $\mathrm{A}'$-módulo plano $(2.1.2)$. Como el homomorfismo
$\mathrm{A}\to\mathrm{A}'$ es inyectivo por construcción, se puede aplicar \sref{IV.11.4.3}, que demuestra que $\mathrm{M}$ es un
$\mathrm{A}$-módulo plano.
\end{proof}

\begin{proposition}[11.4.12]
\label{IV.11.4.12}
Sea $\mathrm{A}$ un anillo local artiniano con cuerpo residual $k$,
sea $\mathrm{Y}=\operatorname{Spec}(\mathrm{A})$, sea $f:\mathrm{X}\to\mathrm{Y}$ un morfismo localmente de tipo finito y sea $\sh{F}$ un $\sh{O}_\mathrm{X}$-módulo coherente.
Sea $(\mathrm{A}'_\alpha)$ una familia de anillos locales y, para cada $\alpha$, sea $u_\alpha:\mathrm{A}\to\mathrm{A}'_\alpha$ un homomorfismo
de anillos. Pongamos $\mathrm{Y}'_\alpha=\operatorname{Spec}(\mathrm{A}'_\alpha)$, $\mathrm{X}'_\alpha=\mathrm{X}\times_\mathrm{Y}\mathrm{Y}'_\alpha$,
$f'_\alpha=f_{(\mathrm{Y}'_\alpha)}$ y $\sh{F}'_\alpha=\sh{F}\otimes_\mathrm{A}\mathrm{A}'_\alpha$. Sea $x$ un punto
de $\mathrm{X}$, y supongamos que se cumplen las siguientes condiciones:
\begin{enumerate}[label=(\roman*)]
  \item La intersección de los núcleos de los $u_\alpha$ es $0$.
  \item Para todo $\alpha$, $\sh{F}'_\alpha$ es $f'_\alpha$-plano en todo punto $x'_\alpha\in\mathrm{X}'_\alpha$ cuyas proyecciones en $\mathrm{X}$ y $\mathrm{Y}'_\alpha$
  son $x$ y el punto cerrado $y'_\alpha$ de $\mathrm{Y}'_\alpha$.
\end{enumerate}

Entonces $\sh{F}$ es $f$-plano en el punto $x$.
\end{proposition}

\begin{proof}
\oldpage[IV-3]{150}
Por hipótesis, la intersección de los núcleos de los $u_\alpha$ es $0$; como $\mathrm{A}$ es artiniano,
existe ya un número \emph{finito} de estos núcleos cuya intersección es $0$, de modo que podemos
limitarnos al caso en que la familia $(\mathrm{A}'_\alpha)$ es \emph{finita}. Sea $k'$ una clausura algebraica de $k$; se sabe
por $(\mathbf{0}_{\mathrm{III}},~10.3.1)$ que existe un homomorfismo local $\mathrm{A}\to\mathrm{B}$ que hace de $\mathrm{B}$ un $\mathrm{A}$-módulo plano,
tal que $\mathrm{B}$ sea un anillo local artiniano, integral sobre $\mathrm{A}$, y que $\mathrm{B}\otimes_\mathrm{A} k$ es isomorfo a $k'$.
Por planitud, los núcleos de los homomorfismos $\mathrm{B}\to\mathrm{B}'_\alpha=\mathrm{B}\otimes_\mathrm{A}\mathrm{A}'_\alpha$ se deducen de los de
$u_\alpha$ por tensorización con $\mathrm{B}$ y, como son un número finito, su intersección es
$0$ $(\mathbf{0}_\mathrm{I},~6.1.3)$. Consideremos los anillos $\mathrm{B}'_{\alpha\beta}$, localizaciones de $\mathrm{B}'_\alpha$ en sus ideales \emph{maximales};
se sabe (Bourbaki, \emph{Alg.\ comm.}, cap.~II, \S~3, n\textsuperscript{o}~3, cor.~2 del teorema~1) que la intersección
de los núcleos de los homomorfismos $\mathrm{B}'_\alpha\to\mathrm{B}'_{\alpha\beta}$ (para un $\alpha$ dado) es $0$; se concluye, por tanto,
que la intersección de los núcleos de los homomorfismos compuestos $v_{\alpha\beta}:\mathrm{B}\to\mathrm{B}'_\alpha\to\mathrm{B}'_{\alpha\beta}$ (con $\alpha$ y
$\beta$ variables) es $0$. Por otra parte, como $\mathrm{B}$ es integral sobre $\mathrm{A}$, $\mathrm{B}'_\alpha$ es integral sobre $\mathrm{A}'_\alpha$,
de modo que sus ideales maximales se encuentran sobre el ideal maximal de $\mathrm{A}'_\alpha$. Si se pone
$\mathrm{Z}'_{\alpha\beta}=\operatorname{Spec}(\mathrm{B}'_{\alpha\beta})$,
$\mathrm{T}'_{\alpha\beta}=\mathrm{X}'_\alpha\times_{\mathrm{Y}'_\alpha}\mathrm{Z}'_{\alpha\beta}$,
$f'_{\alpha\beta}=f_{(\mathrm{Z}'_{\alpha\beta})}$,
$\sh{G}=\sh{F}\otimes_\mathrm{A}\mathrm{B}$, y
$\sh{G}'_{\alpha\beta}=\sh{G}\otimes_\mathrm{B}\mathrm{B}'_{\alpha\beta}$, se ve
que las hipótesis (i) y (ii) siguen cumpliéndose cuando se sustituyen
$\mathrm{A}$, $\mathrm{A}'_\alpha$, $u_\alpha$, $\sh{F}$ y $x$ por $\mathrm{B}$, $\mathrm{B}'_{\alpha\beta}$, $v_{\alpha\beta}$, $\sh{G}$ y un punto $t$ de $\mathrm{T}=\mathrm{X}\otimes_\mathrm{A}\mathrm{B}$ situado sobre $x$.
Como el cuerpo residual de $\mathrm{B}$ es separablemente cerrado, se deduce de \sref{IV.11.4.11} que
$\sh{G}$ es plano sobre $\mathrm{B}$ en el punto $t$. Pero, como $\mathrm{B}$ es un $\mathrm{A}$-módulo fielmente plano, se deduce de $(2.1.4)$
que $\sh{F}$ es plano sobre $\mathrm{A}$ en el punto $x$, lo que demuestra \sref{IV.11.4.12}.
\end{proof}

\begin{corollary}[11.4.13]
\label{IV.11.4.13}
Sea $\mathrm{A}$ un anillo local artiniano, sea $\mathrm{Y}=\operatorname{Spec}(\mathrm{A})$, sea $f:\mathrm{X}\to\mathrm{Y}$ un
morfismo localmente de tipo finito y sea $\sh{F}$ un $\sh{O}_\mathrm{X}$-módulo coherente. Sea $(\mathrm{Y}'_\alpha)$ una familia de
$\mathrm{Y}$-preesquemas y, para todo $\alpha$, pongamos $\mathrm{X}'_\alpha=\mathrm{X}\times_\mathrm{Y}\mathrm{Y}'_\alpha$, $f'_\alpha=f_{(\mathrm{Y}'_\alpha)}$ y $\sh{F}'_\alpha=\sh{F}\otimes_\mathrm{Y}\mathrm{Y}'_\alpha$. Sea $x$ un punto
de $\mathrm{X}$ y supongamos que se cumplen las siguientes hipótesis:
\begin{enumerate}[label=(\roman*)]
  \item La intersección de los núcleos de los homomorfismos $\mathrm{A}\to\Gamma(\mathrm{Y}'_\alpha,\sh{O}_{\mathrm{Y}'_\alpha})$ correspondientes a los
  morfismos estructurales es $0$.
  \item Para todo $\alpha$, $\sh{F}'_\alpha$ es $f'_\alpha$-plano en todo punto $x'_\alpha\in\mathrm{X}'_\alpha$ cuya proyección en $\mathrm{X}$ es $x$.
\end{enumerate}

Entonces $\sh{F}$ es $f$-plano en el punto $x$.
\end{corollary}

\begin{proof}
En efecto, para todo $y'_\alpha\in\mathrm{Y}'_\alpha$, consideremos el esquema local
$\mathrm{Y}''_{\alpha,y'_\alpha}=\operatorname{Spec}(\sh{O}_{\mathrm{Y}'_\alpha,y'_\alpha})$; en
virtud de $(2.1.4)$, $\sh{F}'_\alpha\otimes_{\mathrm{Y}'_\alpha}\mathrm{Y}''_{\alpha,y'_\alpha}$ es
$\mathrm{Y}''_{\alpha,y'_\alpha}$-plano en los puntos cuyas proyecciones en $\mathrm{X}$ y
$\mathrm{Y}''_{\alpha,y'_\alpha}$ son $x$ y el punto cerrado de $\mathrm{Y}''_{\alpha,y'_\alpha}$.
Además, el núcleo del homomorfismo
$\mathrm{A}\to\Gamma(\mathrm{Y}'_\alpha,\sh{O}_{\mathrm{Y}'_\alpha})$ es la intersección de los núcleos de los homomorfismos
$\mathrm{A}\to\Gamma(\mathrm{Y}''_{\alpha,y'_\alpha},\sh{O}_{\mathrm{Y}''_{\alpha,y'_\alpha}})$:
en efecto, nos reducimos inmediatamente al caso en que $\mathrm{Y}'_\alpha$ es afín, y basta entonces aplicar Bourbaki, \emph{Alg.\ comm.},
cap.~II, \S~3, n\textsuperscript{o}~3, cor.~2 del teorema~1. Sustituyendo la familia $(\mathrm{Y}'_\alpha)$ por la familia de los $\mathrm{Y}''_{\alpha,y'_\alpha}$,
nos reducimos, pues, a \sref{IV.11.4.12}.
\end{proof}

\subsection{Descenso de la planitud por morfismos arbitrarios: el caso general}
\label{subsection:IV.11.5}

\begin{theorem}[11.5.1]
\label{IV.11.5.1}
Sean $\mathrm{Y}$ un preesquema localmente noetheriano, $f:\mathrm{X}\to\mathrm{Y}$ un
morfismo localmente de tipo finito, $\sh{F}$ un $\sh{O}_\mathrm{X}$-módulo coherente, $x$ un punto de $\mathrm{X}$ y $y=f(x)$. Sea $(\mathrm{Y}'_\alpha)$
una familia de $\mathrm{Y}$-preesquemas locales $\mathrm{Y}'_\alpha=\operatorname{Spec}(\mathrm{A}'_\alpha)$ tales que las imágenes de los puntos cerrados $y'_\alpha$
de $\mathrm{Y}'_\alpha$ sean todas iguales a $y$. Para todo $\alpha$, sean $\mathfrak{m}'_\alpha$ el ideal maximal de $\mathrm{A}'_\alpha$ y $u_\alpha:\sh{O}_y\to\mathrm{A}'_\alpha$
\oldpage[IV-3]{151}
el homomorfismo canónico (\textbf{I},~$2.4.4$); supongamos que las intersecciones finitas de los ideales
$u_\alpha^{-1}(\mathfrak{m}'_\alpha{}^{n_\alpha})$
formen un sistema fundamental de entornos de $0$ en $\sh{O}_y$. Pongamos $\mathrm{X}'_\alpha=\mathrm{X}\times_\mathrm{Y}\mathrm{Y}'_\alpha$, $f'_\alpha=f_{(\mathrm{Y}'_\alpha)}$,
$\sh{F}'_\alpha=\sh{F}\otimes_\mathrm{Y}\mathrm{Y}'_\alpha$, y supongamos que se cumple una de las hipótesis siguientes:
\begin{enumerate}[label=(\roman*)]
  \item Para todo $\alpha$, $\sh{F}'_\alpha$ es $f'_\alpha$-plano en todo punto cuya proyección en $\mathrm{X}$ es igual
  a $x$ y cuya proyección en $\mathrm{Y}'_\alpha$ es igual a $y'_\alpha$.
  \item Para todo $\alpha$, existe un punto $x'_\alpha\in\mathrm{X}'_\alpha$ cuya proyección en $\mathrm{X}$ es $x$ y cuya proyección en $\mathrm{Y}'_\alpha$
  es igual a $y'_\alpha$, tal que $\sh{F}'_\alpha$ es $f'_\alpha$-plano en el punto $x'_\alpha$ y $k(x)$ es una extensión primaria de $k(y)$.
\end{enumerate}

En estas condiciones, $\sh{F}$ es $f$-plano en el punto $x$.
\end{theorem}

\begin{proof}
Sea $\mathfrak{m}_y$ el ideal maximal de $\sh{O}_y$; como $\sh{O}_y$ y $\sh{O}_x$ son noetherianos y $\sh{O}_y\to\sh{O}_x$ es un
homomorfismo de anillos locales, basta, en virtud de $(\mathbf{0}_{\mathrm{III}},~10.2.2)$, demostrar que, para
todo $n>0$, $\sh{F}_x/\mathfrak{m}_y^n\sh{F}_x$ es un $(\sh{O}_y/\mathfrak{m}_y^n)$-módulo plano. Denotemos por $(\mathfrak{J}_\lambda)$ la familia de
intersecciones finitas de los $u_\alpha^{-1}(\mathfrak{m}'_\alpha{}^{n_\alpha})$; por hipótesis, existe $\lambda$ tal que
$\mathfrak{J}_\lambda\subset\mathfrak{m}_y^n$, y como
$\sh{F}_x\otimes_{\sh{O}_y}(\sh{O}_y/\mathfrak{m}_y^n)=(\sh{F}_x\otimes_{\sh{O}_y}(\sh{O}_y/\mathfrak{J}_\lambda))\otimes_{\sh{O}_y/\mathfrak{J}_\lambda}(\sh{O}_y/\mathfrak{m}_y^n)$, bastará con probar que $\sh{F}_x/\mathfrak{J}_\lambda\sh{F}_x$ es
un $(\sh{O}_y/\mathfrak{J}_\lambda)$-módulo plano. Ahora bien, para cada $\alpha$ tal que
$\mathfrak{J}_\lambda\subset u_\alpha^{-1}(\mathfrak{m}'_\alpha{}^{n_\alpha})$, el paso a los cocientes da
un homomorfismo $u'_\alpha:\sh{O}_y/\mathfrak{J}_\lambda\to\mathrm{A}'_\alpha/\mathfrak{m}'_\alpha{}^{n_\alpha}$, y la intersección de los núcleos
de los $u'_\alpha$ se reduce a $0$. Teniendo en cuenta (\textbf{I},~$3.6.1$), vemos que nos reducimos a \sref{IV.11.4.11}
en el caso (ii), y a \sref{IV.11.4.12} en el caso (i).
\end{proof}

\begin{corollary}[11.5.2]
\label{IV.11.5.2}
Sean $\mathrm{Y}$ un preesquema localmente noetheriano, $f:\mathrm{X}\to\mathrm{Y}$ un
morfismo localmente de tipo finito, $\sh{F}$ un $\sh{O}_\mathrm{X}$-módulo coherente, $x$ un punto de $\mathrm{X}$ y $y=f(x)$; pongamos
$\mathrm{A}=\sh{O}_y$. Sea $u:\mathrm{A}\to\mathrm{A}'$ un homomorfismo de $\mathrm{A}$ en un anillo de Zariski $\mathrm{A}'$, tal que
la imagen inversa $u^{-1}(\mathfrak{r}')$ del radical $\mathfrak{r}'$ de $\mathrm{A}'$ sea el ideal maximal $\mathfrak{m}$ de $\mathrm{A}$; supongamos además
que el homomorfismo $\hat{u}:\hat{\mathrm{A}}\to\hat{\mathrm{A}}'$ es inyectivo. Pongamos $\mathrm{Y}'=\operatorname{Spec}(\mathrm{A}')$, $\mathrm{X}'=\mathrm{X}\times_\mathrm{Y}\mathrm{Y}'$, $f'=f_{(\mathrm{Y}')}$,
$\sh{F}'=\sh{F}\otimes_\mathrm{Y}\mathrm{Y}'$. Para que $\sh{F}$ sea $f$-plano en el punto $x$, es necesario y suficiente que $\sh{F}'$ sea $f'$-plano en todo
punto cuya proyección en $\mathrm{X}$ es igual a $x$ y cuya proyección en $\mathrm{Y}'$ es igual a un punto cerrado
de $\mathrm{Y}'$.

Si además $\mathrm{A}'$ es un $\mathrm{A}$-álgebra finita, en lo anterior se puede sustituir la hipótesis de que $\hat{u}$
es inyectivo por la hipótesis de que $u$ es inyectivo.
\end{corollary}

\begin{proof}
Como $\mathrm{A}$ (resp.\ $\mathrm{A}'$) se identifica con un subanillo de $\hat{\mathrm{A}}$ (resp.\ $\hat{\mathrm{A}}'$) (Bourbaki,
\emph{Alg.\ comm.}, cap.~III, \S~3, n\textsuperscript{o}~3, prop.~6), vemos primero que $u$ es inyectivo
y que $\hat{u}$ es su extensión por continuidad a $\hat{\mathrm{A}}$.

Sea $(\mathfrak{m}'_\alpha)$ la familia de ideales maximales de $\mathrm{A}'$; como se tiene
\[
\mathfrak{m}'_\alpha{}^n\hat{\mathrm{A}}'=\hat{\mathfrak{m}}'_\alpha{}^n, \qquad \text{y} \qquad \hat{\mathfrak{m}}'_\alpha{}^n\cap\mathrm{A}'=\mathfrak{m}'_\alpha{}^n,
\]
y como los $\mathfrak{m}'_\alpha{}^n$ son abiertos en $\mathrm{A}'$, se tiene
$u^{-1}(\mathfrak{m}'_\alpha{}^n)=\mathrm{A}\cap\hat{u}^{-1}(\hat{\mathfrak{m}}'_\alpha{}^n)$, y será suficiente probar que en
$\hat{\mathrm{A}}$ las intersecciones finitas de los $\hat{u}^{-1}(\hat{\mathfrak{m}}'_\alpha{}^{n_\alpha})$ forman un sistema fundamental de entornos
de $0$. Esto permitirá la aplicación de \sref{IV.11.5.1} a los homomorfismos compuestos
$v_\alpha\circ u:\mathrm{A}\to\mathrm{A}'_{\mathfrak{m}'_\alpha}$, donde
$v_\alpha:\mathrm{A}'\to\mathrm{A}'_{\mathfrak{m}'_\alpha}$ es el homomorfismo canónico. Como $\hat{\mathrm{A}}$ es completo, bastará
probar que la intersección de los $\hat{u}^{-1}(\hat{\mathfrak{m}}'_\alpha{}^{n_\alpha})$ se reduce a $0$ (Bourbaki,
\emph{Alg.\ comm.}, cap.~III, \S~2, n\textsuperscript{o}~7, prop.~8, donde en la demostración la sucesión decreciente puede sustituirse por un conjunto
filtrante arbitrario). Ahora bien, para cada $\alpha$
\oldpage[IV-3]{152}
fijo, la intersección de los ideales $\hat{\mathfrak{m}}'_\alpha{}^n\hat{\mathrm{A}}'_{\hat{\mathfrak{m}}'_\alpha}$ para $n>0$ se reduce a $0$ en el
anillo local noetheriano $\hat{\mathrm{A}}'_{\hat{\mathfrak{m}}'_\alpha}$.
Por otra parte, los $\hat{\mathfrak{m}}'_\alpha$ son los ideales maximales de $\hat{\mathrm{A}}'$; por tanto, el homomorfismo canónico
$\hat{\mathrm{A}}'\to\prod_\alpha\hat{\mathrm{A}}'_{\hat{\mathfrak{m}}'_\alpha}$ es inyector (Bourbaki, \emph{Alg.\ comm.}, cap.~II, \S~3, n\textsuperscript{o}~3, cor.~2 de th.~1), y
y como, por hipótesis, $\hat{u}:\hat{\mathrm{A}}\to\hat{\mathrm{A}}'$ también es inyectivo, esto concluye la demostración en
el caso general. La última afirmación se deduce de que $\hat{\mathrm{A}}$ es un $\mathrm{A}$-módulo fielmente
plano $(\mathbf{0}_\mathrm{I},~7.3.5)$ y $\hat{\mathrm{A}}'=\mathrm{A}'\otimes_\mathrm{A}\hat{\mathrm{A}}$, pues $\mathrm{A}'$ es un $\mathrm{A}$-módulo de tipo finito
(Bourbaki, \emph{Alg.\ comm.}, cap.~III, \S~3, n\textsuperscript{o}~4, th.~3 y cap.~IV, \S~2, n\textsuperscript{o}~5, cor.~3 de prop.~9).
\end{proof}

\begin{proposition}[11.5.3]
\label{IV.11.5.3}
Sean $f:\mathrm{X}\to\mathrm{Y}$ un morfismo localmente de presentación finita,
$\sh{F}$ un $\sh{O}_\mathrm{X}$-módulo cuasicoherente de presentación finita, $x$ un punto de $\mathrm{X}$, $y=f(x)$ y $g=(\psi,\theta):\mathrm{Y}'\to\mathrm{Y}$
un morfismo propio de presentación finita. Supongamos que:
\begin{enumerate}[label=(\roman*)]
  \item El homomorfismo $\theta_y:\sh{O}_y\to(g_*(\sh{O}_{\mathrm{Y}'}))_y$ es inyector.
  \item Para todo $x'\in\mathrm{X}'=\mathrm{X}\times_\mathrm{Y}\mathrm{Y}'$ cuya proyección en $\mathrm{X}$ es $x$, $\sh{F}'=\sh{F}\otimes_\mathrm{Y}\mathrm{Y}'$ es
  plano sobre $\mathrm{Y}'$ en el punto $x'$.
\end{enumerate}

Entonces $\sh{F}$ es $f$-plano en el punto $x$.
\end{proposition}

\begin{proof}
Como la cuestión es local sobre $\mathrm{X}$, podemos suponer que $f$ es de presentación finita. Sea
$p:\mathrm{X}'\to\mathrm{X}$ la proyección canónica. Como
$f'=f_{|\mathrm{Y}'}:\mathrm{X}'\to\mathrm{Y}'$ es de presentación finita
$(1.6.2)$ y $\sh{F}'$ es un $\sh{O}_{\mathrm{X}'}$-módulo de presentación finita $(\mathbf{0}_\mathrm{I},~5.2.5)$, se deduce
de \sref{IV.11.3.1} que el conjunto $\mathrm{U}'$ de puntos de $\mathrm{X}'$ en los que $\sh{F}'$ es $f'$-plano es abierto. Como $\mathrm{U}'$
contiene $p^{-1}(x)$ por hipótesis y $p$ es propio, y por tanto \emph{cerrado}, $\mathrm{U}'$ contiene un conjunto de la forma
$p^{-1}(\mathrm{U})$, donde $\mathrm{U}$ es un entorno de $x$. Sustituyendo $\mathrm{X}$ por $\mathrm{U}$, podemos, por tanto,
suponer que $\sh{F}'$ es $f'$-plano. Por otra parte, teniendo en cuenta (\textbf{I},~$3.6.5$), (\textbf{II},~$5.4.2$)
y $(2.1.4)$, también se puede reemplazar $\mathrm{Y}$ por $\operatorname{Spec}(\sh{O}_y)$, es decir, supongamos que $\mathrm{Y}=\operatorname{Spec}(\mathrm{A})$,
donde $\mathrm{A}$ es un anillo local. En estas condiciones demostraremos que $\sh{F}$ es $f$-plano. En virtud
de $(5.13.3)$, $\mathrm{A}$ es el límite inductivo de subanillos locales noetherianos $\mathrm{A}_\lambda$ tales que
la inyección canónica $\mathrm{A}_\lambda\to\mathrm{A}$ es un homomorfismo \emph{local}. En virtud de $(8.9.1)$, podemos
puede suponer que $\mathrm{X}=\mathrm{X}_\lambda\otimes_{\mathrm{A}_\lambda}\mathrm{A}$, $f=f_\lambda\otimes 1_\mathrm{A}$, $\sh{F}=\sh{F}_\lambda\otimes_{\mathrm{A}_\lambda}\mathrm{A}$,
para un $\lambda$ adecuado, donde $f_\lambda:\mathrm{X}_\lambda\to\mathrm{Y}_\lambda=\operatorname{Spec}(\mathrm{A}_\lambda)$ es un morfismo de tipo finito y
$\sh{F}_\lambda$ es un $\sh{O}_{\mathrm{X}_\lambda}$-módulo coherente.
Del mismo modo, podemos suponer que $\mathrm{Y}'=\mathrm{Y}'_\lambda\otimes_{\mathrm{A}_\lambda}\mathrm{A}$ y $g=g_\lambda\otimes 1_\mathrm{A}$, donde $g_\lambda:\mathrm{Y}'_\lambda\to\mathrm{Y}_\lambda$ es un
morfismo de presentación finita; además, en virtud de $(8.10.5,~(\text{xii}))$, podemos suponer que $g_\lambda$ es
\emph{propio}. Como, por hipótesis, el homomorfismo $\mathrm{A}\to\Gamma(\mathrm{Y}',\sh{O}_{\mathrm{Y}'})$ es inyectivo y $\mathrm{A}_\lambda$
es un subanillo de $\mathrm{A}$, el homomorfismo $\mathrm{A}_\lambda\to\Gamma(\mathrm{Y}',\sh{O}_{\mathrm{Y}'})$ también es inyectivo. Finalmente,
en virtud de \sref{IV.11.2.6}, podemos tomar $\lambda$ de manera que $\sh{F}'_\lambda=\sh{F}_\lambda\otimes_{\mathrm{Y}_\lambda}\mathrm{Y}'_\lambda$ sea $f'_\lambda$-plano,
pues $\sh{F}'=\sh{F}'_\lambda\otimes_{\mathrm{Y}'_\lambda}\mathrm{Y}'$. Estas observaciones muestran que, en adelante, podemos suponer que
el anillo $\mathrm{A}$ es \emph{noetheriano}, manteniéndose las demás hipótesis de \sref{IV.11.5.3}. Pongamos entonces
$\mathrm{B}=\hat{\mathrm{A}}$ y $\mathrm{Z}=\operatorname{Spec}(\mathrm{B})$; como $\mathrm{B}$ es un $\mathrm{A}$-módulo fielmente plano $(\mathbf{0}_\mathrm{I},~7.3.5)$, equivale
a decir que $\sh{F}$ es $f$-plano o que $\sh{F}\otimes_\mathrm{Y}\mathrm{Z}$ es plano sobre $\mathrm{Z}$ $(2.1.4)$; análogamente, si se
pone $\mathrm{Z}'=\mathrm{Y}'\times_\mathrm{Y}\mathrm{Z}$, el morfismo $\mathrm{Z}'\to\mathrm{Y}'$ es fielmente plano $(2.2.13)$, por lo que
equivale a decir que $\sh{F}'$ es $f'$-plano o que $\sh{F}'\otimes_{\mathrm{Y}'}\mathrm{Z}'$ es plano sobre $\mathrm{Z}'$; finalmente, $h=g_{(\mathfrak{m})}:\mathrm{Z}'\to\mathrm{Z}$
es propio (\textbf{II},~$5.4.2$) y de tipo finito $(1.5.2)$, $\hat{\mathrm{A}}$ es noetheriano y, si $z$ es su
\oldpage[IV-3]{153}
único punto cerrado, el homomorfismo $\sh{O}_z\to(h_*(\sh{O}_{\mathrm{Z}'}))_z$ es \emph{inyectivo}, pues se deduce de $(2.3.1)$ que
$h_*(\sh{O}_{\mathrm{Z}'})=g_*(\sh{O}_{\mathrm{Y}'})\otimes_\mathrm{Y}\mathrm{Z}$, y nuestra afirmación se deriva de la hipótesis (i) y de la definición de
módulos planos $(\mathbf{0}_\mathrm{I},~6.1.1)$.

Por tanto, podemos suponer ahora que el anillo local noetheriano $\mathrm{A}$ es \emph{completo}; la demostración
se terminará si se demuestra que la intersección de los núcleos de los homomorfismos
$\mathrm{A}=\sh{O}_y\to\sh{O}_{y'}$ (donde $y'$ recorre $g^{-1}(y)$) se reduce a $0$. En efecto, los $\sh{O}_{y'}$ son anillos locales
noetherianos; por tanto, para cada $y'$ la intersección de los $\mathfrak{m}_{y'}^n$ $(n\geq 0)$ se reduce a $0$;
si $\mathfrak{a}_{n,y'}$ es la imagen inversa de $\mathfrak{m}_{y'}^n$ en $\mathrm{A}$, las intersecciones de finitamente muchos de los $\mathfrak{a}_{n,y'}$ son
entornos de $0$ en $\mathrm{A}$ y la intersección de todos los $\mathfrak{a}_{n,y'}$ se reduce a $0$; las intersecciones de
un número finito de los $\mathfrak{a}_{n,y'}$ formarán, por tanto, un sistema fundamental de entornos de $0$ en $\mathrm{A}$
(Bourbaki, \emph{Alg.\ comm.}, cap.~III, \S~2, n\textsuperscript{o}~7, prop.~8, donde uno puede en la prueba
reemplazar la sucesión decreciente por un conjunto filtrante arbitrario); podremos entonces
aplicar \sref{IV.11.5.1}. Sea ahora $s\in\mathrm{A}$ un elemento perteneciente al núcleo de cada uno de
los homomorfismos $\mathrm{A}\to\sh{O}_{y'}$; la imagen $s'$ de $s$ en $\Gamma(\mathrm{Y}',\sh{O}_{\mathrm{Y}'})$ es, por lo tanto, una sección
de $\sh{O}_{\mathrm{Y}'}$ sobre $\mathrm{Y}'$ tal que $s'_{y'}=0$ para todo $y'\in g^{-1}(y)$; existe entonces un entorno $\mathrm{V}'$
de $g^{-1}(y)$ en $\mathrm{Y}'$ tal que $s'|\mathrm{V}'=0$. Pero, como $g$ es \emph{cerrado}, $\mathrm{V}'$ contiene un conjunto
de la forma $g^{-1}(\mathrm{V})$, donde $\mathrm{V}$ es un entorno abierto de $y$ en $\mathrm{Y}$; ahora bien, $\mathrm{Y}$ es un preesquema local
y $y$ es su punto cerrado, de modo que $\mathrm{V}=\mathrm{Y}$ y, en consecuencia, $\mathrm{V}'=\mathrm{Y}'$,
$s'=0$ y, como $\mathrm{A}\to\Gamma(\mathrm{Y}',\sh{O}_{\mathrm{Y}'})$ es inyectivo por hipótesis, $s=0$.
\end{proof}

El siguiente caso particular de \sref{IV.11.5.3} será útil en el cap.~V:

\begin{corollary}[11.5.4]
\label{IV.11.5.4}
Sea $f=(\psi,\theta):\mathrm{X}\to\mathrm{Y}$ un morfismo propio de presentación finita,
y sea $p:\mathrm{X}\times_\mathrm{Y}\mathrm{X}\to\mathrm{X}$ la primera proyección. Supongamos que $\theta:\sh{O}_\mathrm{Y}\to f_*(\sh{O}_\mathrm{X})$ es inyectivo.
Para que $f$ sea plano, es necesario y suficiente que $p$ sea plano.
\end{corollary}

\begin{proposition}[11.5.5]
\label{IV.11.5.5}
Sean $\mathrm{A}$ un anillo, $\mathrm{Y}=\operatorname{Spec}(\mathrm{A})$, $f:\mathrm{X}\to\mathrm{Y}$ un morfismo
localmente de presentación finita, $\sh{F}$ un cuasi-coherente $\sh{O}_\mathrm{X}$-módulo de presentación finita, $x$ un punto
de $\mathrm{X}$. Sea $\mathrm{A}\to\mathrm{A}'$ un homomorfismo inyectivo que hace de $\mathrm{A}'$ un álgebra integral sobre $\mathrm{A}$. Pongamos
$\mathrm{Y}'=\operatorname{Spec}(\mathrm{A}')$, $\mathrm{X}'=\mathrm{X}\times_\mathrm{Y}\mathrm{Y}'$, $f'=f\times 1$, $\sh{F}'=\sh{F}\otimes_\mathrm{Y}\mathrm{Y}'$. Entonces, si $\sh{F}'$ es $f'$- plano en cada punto
de $\mathrm{X}'$ cuya proyección en $\mathrm{X}$ es $x$, $\sh{F}$ es $f$-plano en el punto $x$.
\end{proposition}

\begin{proof}
Supongamos primero que $\mathrm{A}'$ es un $\mathrm{A}$-álgebra \emph{finita} de \emph{presentación finita}; entonces, el
morfismo $\mathrm{Y}'\to\mathrm{Y}$ es propio (\textbf{II},~$6.1.11$) y de presentación finita, por lo que se satisfacen las hipótesis
de \sref{IV.11.5.3}, de donde resulta la conclusión. En el caso general, la proposición se deduce
de este caso particular, del hecho de que $\mathrm{A}'$ es el límite inductivo de sus sub-$\mathrm{A}$-álgebras \emph{finitas}
$\mathrm{A}'_\lambda$ y de los dos lemas siguientes:
\end{proof}

\begin{lemma}[11.5.5.1]
\label{IV.11.5.5.1}
Toda $\mathrm{A}$-álgebra finita $\mathrm{A}'$ es el límite inductivo de $\mathrm{A}$-álgebras $\mathrm{A}'_\lambda$
que son finitas y de presentación finita.
\end{lemma}

\begin{proof}
Se razona como en $(1.9.3.1)$. En efecto, $\mathrm{A}'=\mathrm{B}/\mathfrak{J}$, donde $\mathrm{B}$ es una $\mathrm{A}$-álgebra finita
que es libre como $\mathrm{A}$-módulo, y $\mathfrak{J}$ es un ideal de $\mathrm{B}$ $(1.4.7.1)$. Ahora bien, $\mathfrak{J}$ es el
límite inductivo de ideales $\mathfrak{J}_\lambda\subset\mathfrak{J}$ de $\mathrm{B}$ que son de tipo finito (y, por tanto, \emph{a fortiori}, son $\mathrm{A}$-módulos \emph{de tipo finito}),
por lo tanto,
por la exactitud del funtor $\varinjlim$, $\mathrm{A}'$ es el límite inductivo de las $\mathrm{A}$-álgebras $\mathrm{B}/\mathfrak{J}_\lambda$; ahora bien, $\mathrm{B}/\mathfrak{J}_\lambda$
es, por definición, un $\mathrm{A}$-módulo de presentación finita y, por tanto, también $(1.4.7)$ una $\mathrm{A}$-álgebra finita
de presentación finita.
\end{proof}

\oldpage[IV-3]{154}

Para aplicar este lema a la situación de \sref{IV.11.5.5}, observemos además
que si el homomorfismo $\mathrm{A}\to\mathrm{A}'$ es inyector, entonces
\emph{a fortiori}, también lo es $\mathrm{A}\to\mathrm{A}'_\lambda$, para todo
$\lambda$.

\begin{lemma}[11.5.5.2]
\label{IV.11.5.5.2}
Sean $\mathrm{A}$ un anillo, $\mathrm{A}'$ una $\mathrm{A}$-álgebra y $(\mathrm{A}'_\lambda)$ un sistema inductivo
de $\mathrm{A}$-álgebras tal que $\mathrm{A}'=\varinjlim\mathrm{A}'_\lambda$; pongamos $\mathrm{Y}=\operatorname{Spec}(\mathrm{A})$, $\mathrm{Y}'=\operatorname{Spec}(\mathrm{A}')$ y $\mathrm{Y}'_\lambda=\operatorname{Spec}(\mathrm{A}'_\lambda)$.
Sea $f:\mathrm{X}\to\mathrm{Y}$ un morfismo de presentación finita, y sea $\sh{F}$ un $\sh{O}_\mathrm{X}$-módulo cuasicoherente de
presentación finita; pongamos $\mathrm{X}'=\mathrm{X}\times_\mathrm{Y}\mathrm{Y}'$, $f'=f_{|\mathrm{Y}'}$, $\sh{F}'=\sh{F}\otimes_\mathrm{Y}\mathrm{Y}'$, $\mathrm{X}'_\lambda=\mathrm{X}\times_\mathrm{Y}\mathrm{Y}'_\lambda$, $f'_\lambda=f_{|\mathrm{Y}'_\lambda}$,
$\sh{F}'_\lambda=\sh{F}\otimes_\mathrm{Y}\mathrm{Y}'_\lambda$. Sea $x$ un punto de $\mathrm{X}$ tal que $\sh{F}'$ sea $f'$-plano en todo punto $x'\in\mathrm{X}'$ situado sobre $x$;
entonces existe $\lambda$ tal que $\sh{F}'_\lambda$ es $f'_\lambda$-plano en todo punto $x'\in\mathrm{X}'_\lambda$ situado sobre $x$.
\end{lemma}

\begin{proof}
Sea $\mathrm{U}'$ el conjunto de los $x'\in\mathrm{X}'$ tales que $\sh{F}'$ es $f'$-plano en el punto $x'$; se sabe por \sref{IV.11.3.1}
que $\mathrm{U}'$ es abierto en $\mathrm{X}'$, pues $f'$ es de presentación finita $(1.6.2)$; análogamente, el
conjunto $\mathrm{U}'_\lambda$ de puntos de $\mathrm{X}'_\lambda$ en los que $\sh{F}'_\lambda$ es $f'_\lambda$-plano es abierto en $\mathrm{X}'_\lambda$, y se sabe además por
\sref{IV.11.2.6} que $\mathrm{U}'$ es la unión de los $v_\lambda^{-1}(\mathrm{U}'_\lambda)$, donde $v_\lambda:\mathrm{X}'\to\mathrm{X}'_\lambda$ es la proyección canónica.
Consideremos el esquema $\mathrm{T}=\operatorname{Spec}(k(x))$; pongamos $\mathrm{T}'=\mathrm{T}\times_\mathrm{Y}\mathrm{Y}'$ y $\mathrm{T}'_\lambda=\mathrm{T}\times_\mathrm{Y}\mathrm{Y}'_\lambda$, y sean
$w_\lambda:\mathrm{T}'\to\mathrm{T}'_\lambda$, $g':\mathrm{T}'\to\mathrm{X}'$ y $g'_\lambda:\mathrm{T}'_\lambda\to\mathrm{X}'_\lambda$ las proyecciones canónicas. Pongamos $\mathrm{V}'_\lambda=g_\lambda'^{-1}(\mathrm{U}'_\lambda)$,
$\mathrm{V}'=g'^{-1}(\mathrm{U}')=\bigcup w_\lambda^{-1}(\mathrm{V}'_\lambda)$. Por hipótesis (y teniendo en cuenta (\textbf{I},~$3.6.1$)), $\mathrm{V}'=\mathrm{T}'$;
como $\mathrm{T}'$ es cuasicompacto, existe $\lambda$ tal que $w_\lambda^{-1}(\mathrm{V}'_\lambda)=\mathrm{T}'$. Se deduce entonces de $(8.3.3)$,
aplicado a los subconjuntos cerrados cuasicompactos $\mathrm{T}'_\lambda-\mathrm{V}'_\lambda$ de $\mathrm{T}'_\lambda$, que existe $\mu\geq\lambda$ tal
que $\mathrm{T}'_\mu=\mathrm{V}'_\mu$; esto significa que $\sh{F}'_\mu$ es $f'_\mu$-plano en cada punto de $\mathrm{X}'_\mu$ cuya proyección
en $\mathrm{X}$ es $x$.
\end{proof}

\subsection{Descenso de la planitud por morfismos arbitrarios: el caso de un preesquema base unirramificado}
\label{subsection:IV.11.6}

\begin{theorem}[11.6.1]
\label{IV.11.6.1}
Sea $\mathrm{A}$ un anillo local integral geométricamente unirramificado
$(\mathbf{0},~23.2.1)$, $\mathrm{Y}=\operatorname{Spec}(\mathrm{A})$, $f:\mathrm{X}\to\mathrm{Y}$ un morfismo localmente de presentación finita y $\sh{F}$ un
$\sh{O}_\mathrm{X}$-módulo cuasicoherente de presentación finita. Sea $\mathrm{A}'$ un anillo local y $u:\mathrm{A}\to\mathrm{A}'$ un
homomorfismo local inyectivo; pongamos $\mathrm{Y}'=\operatorname{Spec}(\mathrm{A}')$, $\mathrm{X}'=\mathrm{X}\times_\mathrm{Y}\mathrm{Y}'$, $f'=f_{|\mathrm{Y}'}$ y $\sh{F}'=\sh{F}\otimes_\mathrm{Y}\mathrm{Y}'$. Sea $x$
un punto de $\mathrm{X}$ cuya proyección $f(x)=y$ es el punto cerrado de $\mathrm{Y}$, y sea $x'$ un punto de $\mathrm{X}'$ cuyas proyecciones
en $\mathrm{X}$ y $\mathrm{Y}'$ son, respectivamente, $x$ y el punto cerrado $y'$ de $\mathrm{Y}'$. Entonces, si $\sh{F}'$ es $f'$-plano en el punto $x'$,
$\sh{F}$ es $f$-plano en el punto $x$.
\end{theorem}

\begin{proof}
Podemos limitarnos al caso en que $f$ es de presentación finita, puesto que la cuestión es local
en $\mathrm{X}$. Procederemos en varias etapas.

\medskip
\noindent\textbf{I)} \emph{Reducción al caso en que $\mathrm{A}$ y $\mathrm{A}'$ son anillos locales íntegramente cerrados.}

\medskip
Como $u$ es inyectivo y $\mathrm{A}$ es integral, existe un ideal primo $\mathfrak{p}'$ de $\mathrm{A}'$ tal que
$u^{-1}(\mathfrak{p}')=(0)$ $(\mathbf{0}_\mathrm{I},~1.5.8)$; por tanto, el homomorfismo compuesto $\mathrm{A}\to\mathrm{A}'\to\mathrm{A}''=\mathrm{A}'/\mathfrak{p}'$ es inyectivo
y local, y, si $\mathrm{Y}''=\operatorname{Spec}(\mathrm{A}'')$, $\mathrm{X}''=\mathrm{X}'\otimes_{\mathrm{A}'}\mathrm{A}''=\mathrm{X}\times_\mathrm{Y}\mathrm{Y}''$ y $\sh{F}''=\sh{F}'\otimes_{\mathrm{A}'}\mathrm{A}''$, entonces $\sh{F}''$ es $\mathrm{Y}''$-plano
en los puntos de $\mathrm{X}''$ situados sobre $x'$ $(2.1.4)$; sustituyendo, si es necesario, $\mathrm{A}'$ por $\mathrm{A}''$ y teniendo
en cuenta (\textbf{I},~$3.4.7$), podemos suponer en primer lugar que $\mathrm{A}'$ es integral. Si $\mathrm{K}'$ es el cuerpo de
fracciones de $\mathrm{A}'$, existe un anillo de valoración $\mathrm{B}'$ de $\mathrm{K}'$ que domina a $\mathrm{A}'$; como el homomorfismo compuesto
$\mathrm{A}\xrightarrow{u}\mathrm{A}'\to\mathrm{B}'$ es inyectivo y local, el mismo razonamiento que en
\oldpage[IV-3]{155}
el argumento precedente permite sustituir $\mathrm{A}'$ por $\mathrm{B}'$; podemos, pues, suponer que el anillo local $\mathrm{A}'$ está íntegramente cerrado,
siendo $\mathrm{A}$ un subanillo local de $\mathrm{A}'$ dominado por $\mathrm{A}'$. Sea $\mathrm{A}_1$ el cierre integral de $\mathrm{A}$; es claro que
$\mathrm{A}\subset\mathrm{A}_1\subset\mathrm{A}'$, y, por hipótesis, $\mathrm{A}_1$ es un anillo \emph{local};
si $\mathfrak{m}$, $\mathfrak{m}_1$ y $\mathfrak{m}'$ son los ideales máximos de $\mathrm{A}$, $\mathrm{A}_1$ y $\mathrm{A}'$, se tiene $\mathfrak{m}\subset\mathfrak{m}_1\subset\mathfrak{m}'$; en efecto, $\mathfrak{m}_1$ es
el \emph{único} ideal primo de $\mathrm{A}_1$ situado sobre $\mathfrak{m}$, puesto que $\mathrm{A}_1$ es un anillo local (Bourbaki,
\emph{Alg.\ comm.}, cap.~V, \S~2, n\textsuperscript{o}~1, prop.~1); como $\mathfrak{m}'\cap\mathrm{A}=\mathfrak{m}$, se tiene $\mathfrak{m}'\cap\mathrm{A}_1\cap\mathrm{A}=\mathfrak{m}$,
y, por consiguiente, $\mathfrak{m}'\cap\mathrm{A}_1=\mathfrak{m}_1$. Pongamos $\mathrm{Y}_1=\operatorname{Spec}(\mathrm{A}_1)$, $\mathrm{X}_1=\mathrm{X}\times_\mathrm{Y}\mathrm{Y}_1$, $f_1=f_{(\mathrm{Y}_1)}$ y $\sh{F}_1=\sh{F}\otimes_\mathrm{Y}\mathrm{Y}_1$, y
sea $x_1$ la proyección de $x'$ en $\mathrm{X}_1$; denotemos además por $y_1$ el único punto cerrado
de $\mathrm{Y}_1$, de modo que $y_1=f_1(x_1)$. Por hipótesis, el morfismo $\operatorname{Spec}(k(y_1))\to\operatorname{Spec}(k(y))$
es \emph{radicial}, de donde se deduce, por la transitividad de las fibras (\textbf{I},~$3.6.4$) y (\textbf{I},~$3.5.7$), que
el morfismo $f_1^{-1}(y_1)\to f^{-1}(y)$ es radicial y, en particular, que $x_1$ es el \emph{único} punto de $\mathrm{X}_1$
cuyas proyecciones en $\mathrm{X}$ y $\mathrm{Y}_1$ son, respectivamente, $x$ y $y_1$; además, como se ha visto, $y_1$
es el \emph{único} punto de $\mathrm{Y}_1$ cuya proyección en $\mathrm{Y}$ es $y$, por lo que $x_1$ es el \emph{único} punto de $\mathrm{X}_1$
cuya proyección en $\mathrm{X}$ es $x$. Si demostramos que $\sh{F}_1$ es $f_1$-plano en el punto $x_1$, podremos
aplicar \sref{IV.11.5.5}, de donde se obtiene la conclusión. Quedamos así reducidos al caso en que $\mathrm{A}$
es también \emph{íntegramente cerrado}.

\medskip
\noindent\textbf{II)} \emph{Reducción al caso en que $\mathrm{A}$ y $\mathrm{A}'$ son anillos locales íntegramente cerrados de $\mathbf{Z}$-álgebras de tipo
finito.}

\medskip
Podemos considerar $\mathrm{A}'$ como el límite inductivo filtrante de sus sub-$\mathbf{Z}$-álgebras de tipo finito
$\mathrm{B}'_\lambda$ que son dominios de integridad; además, como $\mathrm{A}'$ está íntegramente cerrado y el cierre
integral de una $\mathbf{Z}$-álgebra de tipo finito es también una $\mathbf{Z}$-álgebra de tipo finito $(7.8.3)$, se
ve que $\mathrm{A}'$ es el límite inductivo de sus sub-$\mathbf{Z}$-álgebras de tipo finito $\mathrm{B}'_\lambda$ \emph{íntegramente cerradas};
si $\mathfrak{m}'_\lambda=\mathfrak{m}'\cap\mathrm{B}'_\lambda$, $\mathrm{A}'$ es también el límite inductivo de los subanillos locales $(\mathrm{B}'_\lambda)_{\mathfrak{m}'_\lambda}=\mathrm{A}'_\lambda$
dominados por $\mathrm{A}'$ $(5.13.3)$. Para todo $\lambda$, $\mathrm{B}'_\lambda\cap\mathrm{A}$ es asimismo el límite inductivo de sus
sub-$\mathbf{Z}$-álgebras de tipo finito $\mathrm{B}_{\alpha\lambda}$, y, por tanto, $\mathrm{A}=\varinjlim_{\alpha,\lambda}\mathrm{B}_{\alpha\lambda}$; como antes, podemos sustituir $\mathrm{B}_{\alpha\lambda}$
en esta fórmula por su cierre integral (que, por hipótesis, está contenido en $\mathrm{B}'_\lambda\cap\mathrm{A}$), y después por el anillo local
$\mathrm{A}_{\alpha\lambda}=(\mathrm{B}_{\alpha\lambda})_{\mathfrak{m}_{\alpha\lambda}}$, donde
$\mathfrak{m}_{\alpha\lambda}=\mathfrak{m}\cap\mathrm{B}_{\alpha\lambda}=\mathfrak{m}'_\lambda\cap\mathrm{B}_{\alpha\lambda}$, de modo que $\mathrm{A}_{\alpha\lambda}$ está dominado por $\mathrm{A}'_\lambda$ y por $\mathrm{A}$.
Pongamos $\mathrm{Y}_{\alpha\lambda}=\operatorname{Spec}(\mathrm{A}_{\alpha\lambda})$; de $(8.9.1)$ se deduce que existen un par suficientemente grande
$(\alpha,\lambda)$, un morfismo $f_{\alpha\lambda}:\mathrm{X}_{\alpha\lambda}\to\mathrm{Y}_{\alpha\lambda}$ de tipo finito y un
$\sh{O}_{\mathrm{X}_{\alpha\lambda}}$-módulo coherente $\sh{F}_{\alpha\lambda}$ tales que
$\mathrm{X}=\mathrm{X}_{\alpha\lambda}\times_{\mathrm{Y}_{\alpha\lambda}}\mathrm{Y}$,
$f=f_{\alpha\lambda}\times 1_\mathrm{Y}$ y $\sh{F}=\sh{F}_{\alpha\lambda}\otimes_{\mathrm{Y}_{\alpha\lambda}}\mathrm{Y}$; si $x_{\alpha\lambda}$
es la proyección de $x$ en $\mathrm{X}_{\alpha\lambda}$, bastará demostrar que $\sh{F}_{\alpha\lambda}$ es $f_{\alpha\lambda}$-plano en el punto $x_{\alpha\lambda}$.
Como $\mathrm{A}'$ es el límite inductivo de los $\mathrm{A}'_\mu$ para $\mu\geq\lambda$, $\mathrm{X}'$ es el límite proyectivo de los
$\mathrm{X}_{\alpha\lambda}\times_{\mathrm{Y}_{\alpha\lambda}}\mathrm{Y}'_\mu=\mathrm{X}'_\mu$, y se tiene asimismo
$\sh{F}'=\sh{F}'_{\alpha\mu}\otimes_{\mathrm{Y}'_\mu}\mathrm{Y}'$, donde
$\sh{F}'_{\alpha\mu}=\sh{F}_{\alpha\lambda}\otimes_{\mathrm{Y}_{\alpha\lambda}}\mathrm{Y}'_\mu$. Aplicando
\sref{IV.11.2.6}, se ve que podemos tomar $\mu$ suficientemente grande para que $\sh{F}'_{\alpha\mu}$ sea $\mathrm{Y}'_\mu$-plano en el
punto $x'_\mu$, proyección de $x'$ en $\mathrm{X}'_\mu$; además, por la construcción de los $\mathrm{A}'_\mu$, la proyección
de $x'_\mu$ en $\mathrm{Y}'_\mu$ es el punto cerrado de $\mathrm{Y}'_\mu$.

\medskip
\noindent\textbf{III)} \emph{Reducción al caso en que el cuerpo residual $k'$ de $\mathrm{A}'$ es una extensión finita y radicial del
cuerpo residual $k$ de $\mathrm{A}$.}

\medskip
Podemos comenzar repitiendo el razonamiento de la parte I) de la
demostración de \sref{IV.11.4.11}, teniendo en cuenta que $\mathbf{Z}$ es un
\oldpage[IV-3]{156}
anillo de Jacobson; quedamos así reducidos al caso en que $k'$ es una extensión finita de $k$, como supondremos en lo sucesivo.
Sea $k''$ la mayor extensión separable de $k$ contenida en $k'$, y sea $k_1$ una extensión de Galois finita de $k$ que contiene a $k''$, de modo que $k''\otimes_k k_1$ es un producto directo de cuerpos isomorfos a $k_1$.
Como $k'$ es una extensión radicial de $k''$, $k'\otimes_k k_1$ es, por tanto, un producto directo de extensiones radiciales de $k_1$.

Existe un anillo local $\mathrm{A}_1$ que es una $\mathrm{A}$-álgebra y un $\mathrm{A}$-módulo libre de tipo finito, tal que $\mathrm{A}_1/\mathfrak{m}\mathrm{A}_1$ es $k$-isomorfo a $k_1$ $(\mathbf{0}_{\mathrm{III}},~10.3.1.2)$.
Más precisamente, escribamos $k_1=k[\mathrm{T}]/(r)$, donde $r$ es un polinomio mónico irreducible separable de grado $n$, y sea $\mathrm{R}\in\mathrm{A}[\mathrm{T}]$ un polinomio mónico cuya imagen canónica es $r$; podemos tomar $\mathrm{A}_1=\mathrm{A}[\mathrm{T}]/(\mathrm{R})$.
Si $\mathrm{K}$ es el cuerpo de fracciones de $\mathrm{A}$, es claro que $\mathrm{R}$ es un polinomio irreducible separable de $\mathrm{K}[\mathrm{T}]$; se deduce, en primer lugar, que $\mathrm{A}_1$ es un anillo integral cuyo cuerpo de fracciones $\mathrm{K}_1=\mathrm{K}[\mathrm{T}]/(\mathrm{R})$ es una extensión separable de $\mathrm{K}$.
Además, si $t$ es la imagen canónica de $\mathrm{T}$ en $\mathrm{A}_1$, los $t^j$ $(0\leq j<n)$ forman una base del $\mathrm{A}$-módulo $\mathrm{A}_1$, y sus imágenes en $k_1$ forman una base sobre $k$; por tanto, $d=\det\bigl(\operatorname{Tr}_{\mathrm{A}_1/\mathrm{A}}(t^{i+j})\bigr)$ es un elemento de $\mathrm{A}$ cuya clase en $k$ es no nula y que, por consiguiente, es invertible.
El mismo razonamiento que en \sref{IV.6.12.4.1},~I) demuestra entonces que el morfismo $\operatorname{Spec}(\mathrm{A}_1)\to\operatorname{Spec}(\mathrm{A})$ es plano y tiene fibras regulares; de \sref{IV.6.5.4},~(ii), se deduce que $\mathrm{A}_1$ es \emph{íntegramente cerrado}.

Consideremos ahora el anillo $\mathrm{A}'_1=\mathrm{A}'\otimes_\mathrm{A}\mathrm{A}_1$.
Es un $\mathrm{A}'$-módulo libre de tipo finito y, por tanto, un anillo semilocal (Bourbaki, \emph{Alg.\ comm.}, cap.~IV, \S~2, n\textsuperscript{o}~5, cor.~3 de prop.~9); además, todos los ideales máximos de esta $\mathrm{A}'$-álgebra finita están sobre el ideal máximo $\mathfrak{m}'$ de $\mathrm{A}'$ y, \emph{a fortiori}, contienen a $\mathfrak{m}\mathrm{A}'_1$.
Pero
\[
\mathrm{A}'_1/\mathfrak{m}\mathrm{A}'_1=(\mathrm{A}'/\mathfrak{m}\mathrm{A}')\otimes_k k_1,
\]
y, como $k_1$ es una extensión finita separable de $k$, el radical de $\mathrm{A}'_1/\mathfrak{m}\mathrm{A}'_1$ es igual a $(\mathfrak{m}'/\mathfrak{m}\mathrm{A}')\otimes_k k_1$ (Bourbaki, \emph{Alg.}, cap.~VIII, \S~7, n\textsuperscript{o}~2, cor.~2 de prop.~3).
Si $\mathfrak{n}_i$ $(1\leq i\leq r)$ son los ideales máximos de $\mathrm{A}'_1$, los cuerpos $\mathrm{A}'_1/\mathfrak{n}_i$ son, por tanto, los cuerpos que aparecen como factores del álgebra $k'\otimes_k k_1$; en otras palabras, son \emph{extensiones radiciales finitas de $k_1$}.

Como $\mathrm{A}\to\mathrm{A}'$ es inyectivo y $\mathrm{A}_1$ es un $\mathrm{A}$-módulo plano, el homomorfismo $\mathrm{A}_1\to\mathrm{A}'_1$ es inyectivo.
El homomorfismo canónico
\[
\mathrm{A}'_1\longrightarrow\prod_{i=1}^{r}(\mathrm{A}'_1)_{\mathfrak{n}_i}
\]
es también inyectivo (Bourbaki, \emph{Alg.\ comm.}, cap.~II, \S~3, n\textsuperscript{o}~3, th.~1), y, por consiguiente, también lo es el homomorfismo compuesto de $\mathrm{A}_1$ en este producto.
Ahora bien, $\mathrm{A}_1$ es integral y solo hay un número finito de núcleos de los homomorfismos $\mathrm{A}_1\to(\mathrm{A}'_1)_{\mathfrak{n}_i}$; como su intersección es cero, uno de ellos es ya cero.
En otras palabras, para algún $i$ existe un anillo $\mathrm{B}_1=(\mathrm{A}'_1)_{\mathfrak{n}_i}$ tal que el homomorfismo $\mathrm{A}_1\to\mathrm{B}_1$ es \emph{inyectivo} y \emph{local}.

Pongamos $\mathrm{Y}'_1=\operatorname{Spec}(\mathrm{B}_1)$, $\mathrm{X}'_1=\mathrm{X}'\times_{\mathrm{Y}'}\mathrm{Y}'_1$ y $\sh{F}'_1=\sh{F}'\otimes_{\mathrm{Y}'}\mathrm{Y}'_1$.
El módulo $\sh{F}'_1$ es $\mathrm{Y}'_1$-plano en todos los puntos de $\mathrm{X}'_1$ situados sobre $x'$; además, el ideal máximo de $\mathrm{B}_1$ es el único situado sobre $\mathfrak{m}'$, por lo que todos estos puntos se proyectan en el punto cerrado $y'_1$ de $\mathrm{Y}'_1$.
Sea $x'_1$ uno de estos puntos.
Por otra parte, pongamos $\mathrm{Y}_1=\operatorname{Spec}(\mathrm{A}_1)$, $\mathrm{X}_1=\mathrm{X}\times_\mathrm{Y}\mathrm{Y}_1$ y $\sh{F}_1=\sh{F}\otimes_\mathrm{Y}\mathrm{Y}_1$; si $x_1$ es la proyección de $x'_1$ en $\mathrm{X}_1$, sus proyecciones en $\mathrm{X}$ y $\mathrm{Y}_1$ son, respectivamente, $x$ y el punto cerrado $y_1$.
Si demostramos que $(\sh{F}_1)_{x_1}$ es un $\sh{O}_{y_1}$-módulo plano, se deduce que $\sh{F}_x$ es un $\sh{O}_y$-módulo plano: en efecto, $\sh{O}_{y_1}$ es un $\sh{O}_y$-módulo plano, de modo que $(\sh{F}_1)_{x_1}$ es un $\sh{O}_y$-módulo plano $(\mathbf{0}_\mathrm{I},~6.2.1)$.

\oldpage[IV-3]{157}
Pero $(\sh{F}_1)_{x_1}=\sh{F}_x\otimes_{\sh{O}_x}\sh{O}_{x_1}$, y $\sh{O}_{x_1}$ es un $\sh{O}_x$-módulo fielmente plano; por \sref{IV.2.2.11},~(iii), $\sh{F}_x$ es, por tanto, un $\sh{O}_y$-módulo plano.
Como $\mathrm{X}'_1=\mathrm{X}_1\times_{\mathrm{Y}_1}\mathrm{Y}'_1$ y $\sh{F}'_1=\sh{F}_1\otimes_{\mathrm{Y}_1}\mathrm{Y}'_1$, hemos quedado efectivamente reducidos a la situación de \sref{IV.11.6.1}, sustituyendo $\mathrm{A}$ y $\mathrm{A}'$, respectivamente, por $\mathrm{A}_1$ y $\mathrm{B}_1$.

\medskip
\noindent\textbf{IV)} \emph{Fin de la demostración.} --- Finalmente, quedamos reducidos a demostrar \sref{IV.11.6.1} bajo las siguientes hipótesis adicionales:
\begin{enumerate}[label=(\roman*)]
  \item $\mathrm{A}$ y $\mathrm{A}'$ son anillos locales de $\mathbf{Z}$-álgebras de tipo finito (y, por tanto, anillos \emph{excelentes} \sref{IV.7.8.3});
  \item $\mathrm{A}$ está íntegramente cerrado;
  \item el cuerpo residual $k'$ de $\mathrm{A}'$ es una extensión radicial finita del cuerpo residual $k$ de $\mathrm{A}$.
\end{enumerate}
Se sabe entonces (\sref{IV.7.8.3} y \sref{IV.2.3.8}) que, bajo las condiciones~(i) y~(ii), la topología $\mathfrak{m}$-ádica de $\mathrm{A}$ es inducida por la topología $\mathfrak{m}'$-ádica de $\mathrm{A}'$.
Por tanto, el homomorfismo entre los completados $\widehat{u}:\widehat{\mathrm{A}}\to\widehat{\mathrm{A}'}$ es inyectivo.
Por otra parte, como el morfismo $\operatorname{Spec}(k')\to\operatorname{Spec}(k)$ es radicial, lo mismo ocurre con $f'^{-1}(y')\to f^{-1}(y)$ \sref[I]{I.3.5.7}; existe, pues, un único punto $x'$ de $\mathrm{X}'$ cuyas proyecciones en $\mathrm{X}$ y $\mathrm{Y}'$ son, respectivamente, $x$ y $y'$.
Podemos, por tanto, aplicar \sref{IV.11.5.2}.
\end{proof}

\begin{corollary}[11.6.2]
\label{IV.11.6.2}
Sea $\mathrm{A}$ un anillo local integral unirramificado, $\mathrm{Y}=\operatorname{Spec}(\mathrm{A})$, $f:\mathrm{X}\to\mathrm{Y}$ un morfismo localmente de presentación finita y $\sh{F}$ un $\sh{O}_\mathrm{X}$-módulo cuasicoherente de presentación finita.
Sea $\mathrm{A}'$ un anillo local y $\mathrm{A}\to\mathrm{A}'$ un homomorfismo local inyectivo; pongamos $\mathrm{Y}'=\operatorname{Spec}(\mathrm{A}')$, $\mathrm{X}'=\mathrm{X}\times_\mathrm{Y}\mathrm{Y}'$, $f'=f_{(\mathrm{Y}')}$ y $\sh{F}'=\sh{F}\otimes_\mathrm{Y}\mathrm{Y}'$.
Sea $x$ un punto de $\mathrm{X}$ cuya proyección $f(x)=y$ es el punto cerrado de $\mathrm{Y}$, y supongamos que $\sh{F}'$ es $f'$-plano en todos los puntos $x'$ de $\mathrm{X}'$ cuyas proyecciones en $\mathrm{X}$ y $\mathrm{Y}'$ son, respectivamente, $x$ y el punto cerrado $y'$ de $\mathrm{Y}'$.
Entonces $\sh{F}$ es $f$-plano en $x$.
\end{corollary}

\begin{proof}
Podemos repetir la parte~I) de la demostración de \sref{IV.11.6.1}, que muestra (con la misma notación) que, si $\sh{F}_1$ es $f_1$-plano en \emph{todos} los puntos $x_1$ de $\mathrm{X}_1$ cuyas proyecciones en $\mathrm{X}$ y $\mathrm{Y}_1$ son, respectivamente, $x$ y $y_1$, entonces $\sh{F}$ es $f$-plano en $x$.
Quedamos así reducidos al caso en que $\mathrm{A}$ está íntegramente cerrado y, por tanto, es geométricamente unirramificado; la conclusión se sigue de \sref{IV.11.6.1}.
\end{proof}

\subsection{Contraejemplos}
\label{subsection:IV.11.7}

\begin{env}[11.7.1]
\label{IV.11.7.1}
Consideremos primero el caso en que $\mathrm{A}$ es un anillo local \emph{artiniano} y en que se satisfacen las hipótesis de \sref{IV.11.4.11}, \emph{salvo} la condición~(ii) relativa al cuerpo residual $k$ de $\mathrm{A}$. Veremos que, en tal caso, puede fallar la conclusión de \sref{IV.11.4.11}.
Sea $k$ un cuerpo que admite una extensión de Galois $k'$ de grado $[k':k]>1$, y sea $\Gamma$ su grupo de Galois.
Sea $\mathrm{A}$ una $k$-álgebra con una base de tres elementos $1,a,b$ y tabla de multiplicación $a^2=ab=ba=b^2=0$, de modo que $\mathrm{A}$ es un anillo local artiniano cuyo ideal máximo $\mathfrak{m}=ka+kb$ tiene cuadrado cero.
Sea $\mathrm{A}'=\mathrm{A}\otimes_k k'$, que es una $k'$-álgebra con base $1,a,b$ y un anillo local artiniano cuyo ideal máximo de cuadrado cero es $\mathfrak{m}'=k'a+k'b=\mathfrak{m}\mathrm{A}'$; el anillo $\mathrm{A}$ se identifica canónicamente con un subanillo de $\mathrm{A}'$.
Sea $\mathfrak{J}$ el subespacio vectorial sobre $k'$ de $\mathfrak{m}'$ generado por $a+\gamma b$, donde $\gamma\in k'$ no pertenece a $k$; es claro que $\mathfrak{J}$ es un ideal de $\mathrm{A}'$.
Pongamos $\mathrm{B}=\mathrm{A}'/\mathfrak{J}$; este es un anillo artiniano que no es plano como $\mathrm{A}$-módulo. De lo contrario (Bourbaki, \emph{Alg.\ comm.}, cap.~II, \S~3, n\textsuperscript{o}~2, cor.~2 de prop.~5), $\mathrm{B}$ sería un $\mathrm{A}$-módulo libre. Como $\mathrm{A}'$ es asimismo un $\mathrm{A}$-módulo libre y el homomorfismo canónico $\mathrm{A}'/\mathfrak{m}\mathrm{A}'\to\mathrm{B}/\mathfrak{m}\mathrm{B}$ es biyectivo, el homomorfismo canónico $\mathrm{A}'\to\mathrm{B}=\mathrm{A}'/\mathfrak{J}$ también sería
\oldpage[IV-3]{158}
biyectivo (loc.~cit., n\textsuperscript{o}~2, cor.~de prop.~6), lo cual es absurdo. En otras palabras, si $\mathrm{Y}=\Spec(\mathrm{A})$ y $\mathrm{X}=\Spec(\mathrm{B})$, entonces $\sh{O}_\mathrm{X}$ no es $\mathrm{Y}$-plano en el único punto $x$ de $\mathrm{X}$.
Pongamos ahora $\mathrm{A}_1=\mathrm{B}$ y $\mathrm{Y}_1=\Spec(\mathrm{A}_1)$, y consideremos el homomorfismo canónico $\mathrm{A}\to\mathrm{A}_1$, que es local e inyectivo porque $\mathfrak{J}\cap\mathrm{A}=0$. Si $\mathrm{B}_1=\mathrm{B}\otimes_\mathrm{A}\mathrm{A}_1=\mathrm{B}\otimes_\mathrm{A}\mathrm{B}$, veremos que existe un punto de $\mathrm{X}_1=\mathrm{X}\times_\mathrm{Y}\mathrm{Y}_1=\Spec(\mathrm{B}_1)$ en el que $\sh{O}_{\mathrm{X}_1}$ es $\mathrm{Y}_1$-plano.
Para ello, observemos que
\[
\mathrm{B}\otimes_\mathrm{A}\mathrm{B}
=\frac{\mathrm{A}'\otimes_\mathrm{A}\mathrm{A}'}
{\operatorname{Im}(\mathfrak{J}\otimes_\mathrm{A}\mathrm{A}')+\operatorname{Im}(\mathrm{A}'\otimes_\mathrm{A}\mathfrak{J})}.
\]
La estructura de $\mathrm{C}'=\mathrm{A}'\otimes_\mathrm{A}\mathrm{A}'$ se obtiene fácilmente. Consideremos la $\mathrm{A}$-álgebra producto $\mathrm{C}=\prod_{\sigma\in\Gamma}\mathrm{A}'_\sigma$, donde todos los $\mathrm{A}'_\sigma$ son iguales a $\mathrm{A}'$, y la aplicación canónica $\varphi:\mathrm{A}'\otimes_\mathrm{A}\mathrm{A}'\to\mathrm{C}$ definida por $\varphi(x\otimes y)=(\sigma(x)y)_{\sigma\in\Gamma}$, donde $\Gamma$ actúa canónicamente sobre $\mathrm{A}'$. Al pasar a los cocientes se obtiene un homomorfismo $\mathrm{C}'/\mathfrak{m}\mathrm{C}'\to\mathrm{C}/\mathfrak{m}\mathrm{C}$ que no es sino el homomorfismo canónico $k'\otimes_k k'\to\prod_\sigma k'_\sigma$, con $k'_\sigma=k'$ para todo $\sigma\in\Gamma$. Este último homomorfismo es biyectivo (Bourbaki, \emph{Alg.}, cap.~VIII, \S~8, prop.~4), y, por tanto, también lo es $\varphi$, pues $\mathrm{C}'$ y $\mathrm{C}$ son $\mathrm{A}$-módulos libres (Bourbaki, \emph{Alg.\ comm.}, cap.~II, \S~3, n\textsuperscript{o}~2, cor.~de prop.~6).
Se deduce que $\mathrm{B}\otimes_\mathrm{A}\mathrm{B}$ es una $\mathrm{A}$-álgebra semilocal que es producto directo de las $\mathrm{A}$-álgebras locales $\mathrm{A}'/(\sigma(\mathfrak{J})+\mathfrak{J})$. El factor correspondiente al elemento neutro de $\Gamma$ es isomorfo a $\mathrm{A}'/\mathfrak{J}$ y es, por tanto, un $\mathrm{A}_1$-módulo plano. Sin embargo, como $\gamma\notin k$, existe al menos un $\sigma\in\Gamma$ tal que $\sigma(\mathfrak{J})\ne\mathfrak{J}$; entonces $\sigma(\mathfrak{J})+\mathfrak{J}=\mathfrak{m}'$, y $\mathrm{A}'/\mathfrak{m}'$ no es un $(\mathrm{A}'/\mathfrak{J})$-módulo plano, pues es el cociente de $\mathrm{A}'/\mathfrak{J}$ por un ideal no nulo.
\end{env}

\begin{env}[11.7.2]
\label{IV.11.7.2}
Mostraremos ahora que el resultado de \sref{IV.11.5.4} deja de ser válido cuando ya no se supone que $f$ es propio (y, \emph{a fortiori}, \sref{IV.11.5.3} deja de ser exacto cuando ya no se supone que $g$ es propio).
Sea $k$ un cuerpo, sea $\mathrm{A}_0=k[\mathrm{S},\mathrm{T}]$ y sea $\mathrm{A}=\mathrm{A}_0/\mathrm{A}_0\mathrm{S}\mathrm{T}$. Así, $\mathrm{Y}=\Spec(\mathrm{A})$ es la curva reducible formada por los dos ``ejes coordenados'' del plano afín $\mathrm{V}_k^2=\Spec(\mathrm{A}_0)$.
El anillo $\mathrm{A}$ tiene dos ideales primos mínimos $\mathfrak{p}_1=\mathrm{A}_0\mathrm{S}/\mathrm{A}_0\mathrm{S}\mathrm{T}$ y $\mathfrak{p}_2=\mathrm{A}_0\mathrm{T}/\mathrm{A}_0\mathrm{S}\mathrm{T}$. Como $\mathrm{A}$ es reducido, se inyecta canónicamente en $\mathrm{B}=\mathrm{B}_1\oplus\mathrm{B}_2$, donde $\mathrm{B}_1=\mathrm{A}/\mathfrak{p}_1$ y $\mathrm{B}_2=\mathrm{A}/\mathfrak{p}_2$. Además, $\mathrm{B}_1$ se identifica con $k[\mathrm{T}]$ y $\mathrm{B}_2$ con $k[\mathrm{S}]$; estos son anillos integrales íntegramente cerrados y, por consiguiente, $\mathrm{Z}=\Spec(\mathrm{B})$ es precisamente la normalización de $\mathrm{Y}$ con respecto a $\mathrm{R}(\mathrm{Y})$ \sref[II]{II.6.3.8}, suma de los dos esquemas $\mathrm{Z}_1=\Spec(\mathrm{B}_1)$ y $\mathrm{Z}_2=\Spec(\mathrm{B}_2)$.
Denotemos por $y$ el ``punto doble'' de $\mathrm{Y}$ correspondiente al ideal máximo $\mathfrak{p}_1+\mathfrak{p}_2=\mathfrak{m}$ de $\mathrm{A}$, y por $z_1,z_2$ los puntos de $\mathrm{Z}$ situados sobre $y$, correspondientes respectivamente a los ideales máximos $\mathfrak{m}_1=(\mathrm{T})$ y $\mathfrak{m}_2=(\mathrm{S})$ de $\mathrm{B}_1$ y $\mathrm{B}_2$.
Sea $\mathrm{X}$ el subpreesquema de $\mathrm{Z}$ inducido en el complementario de $z_2$. Así, $\mathrm{X}=\Spec(\mathrm{B}_1\oplus(\mathrm{B}_2)_\mathrm{S})$. Es inmediato que el homomorfismo $\mathrm{A}\to\mathrm{A}'=\mathrm{B}_1\oplus(\mathrm{B}_2)_\mathrm{S}$ es inyectivo, pero el morfismo correspondiente $f:\mathrm{X}\to\mathrm{Y}$ no es cerrado, pues $f(\mathrm{Z}_2-\{z_2\})$ no es cerrado en $\mathrm{Y}$ aunque $\mathrm{Z}_2-\{z_2\}$ sea cerrado en $\mathrm{X}$; \emph{a fortiori}, $f$ no es propio.
Veamos ahora que $f$ no es plano en $z_1$. Basta demostrar \sref[0]{0.6.6.2} que $\sh{O}_{z_1}$ no es un $\sh{O}_y$-módulo fielmente plano, y para ello basta observar que el homomorfismo canónico $\sh{O}_y\to\sh{O}_{z_1}$ no es inyectivo. Esto es inmediato, pues $\sh{O}_{z_1}$ es un anillo integral, mientras que $\sh{O}_y$ tiene divisores de cero.
Sin embargo, la primera proyección $p:\mathrm{X}\times_\mathrm{Y}\mathrm{X}\to\mathrm{X}$ es un isomorfismo. En efecto,
\[
\mathrm{B}_1\otimes_\mathrm{A}\mathrm{B}_1
=(\mathrm{A}/\mathfrak{p}_1)\otimes_\mathrm{A}(\mathrm{A}/\mathfrak{p}_1)
=\mathrm{A}/\mathfrak{p}_1,
\]
mientras que $(\mathrm{B}_2)_\mathrm{S}\otimes_\mathrm{A}(\mathrm{B}_2)_\mathrm{S}=(\mathrm{B}_2\otimes_\mathrm{A}\mathrm{B}_2)_\mathrm{S}=(\mathrm{B}_2)_\mathrm{S}$ por la misma razón; finalmente, $\mathrm{B}_1\otimes_\mathrm{A}(\mathrm{B}_2)_\mathrm{S}=0$, pues la imagen canónica de $\mathrm{S}$ en $\mathrm{B}_1$ es cero.
\end{env}

\begin{env}[11.7.3]
\label{IV.11.7.3}
El ejemplo precedente puede generalizarse. Sobre un cuerpo $k$, consideremos una curva algebraica reducida $\mathrm{Y}$ que admite un único ``punto doble ordinario'' $y$ (noción que se definirá más adelante en general), junto con su normalización $\mathrm{Z}$, de modo que $g:\mathrm{Z}\to\mathrm{Y}$ sea finito, la restricción de $g$ a $\mathrm{Z}-g^{-1}(y)$ sea un isomorfismo sobre $\mathrm{Y}-\{y\}$ y $g^{-1}(y)$ conste de dos puntos ``simples'' $z_1,z_2$. Además, el preesquema $g^{-1}(y)$ es la suma de los dos preesquemas $\Spec(k(z_1))$ y $\Spec(k(z_2))$, cada uno de ellos canónicamente isomorfo a $\Spec(k(y))$.
Sea $\mathrm{X}$ el subpreesquema de $\mathrm{Z}$ inducido en el abierto $\mathrm{Z}-\{z_2\}$. El morfismo $f:\mathrm{X}\to\mathrm{Y}$ obtenido restringiendo $g$ no es propio; de lo contrario, \sref[II]{II.5.4.3} implicaría lo mismo para la inyección canónica $j:\mathrm{X}\to\mathrm{Z}$, que no es cerrada.
El morfismo $f$ es radicial: para todo $y'\in\mathrm{Y}$, la fibra $f^{-1}(y')$ consta de un único punto $x'$, y $k(y')\to k(x')$ es un isomorfismo. Se deduce primero que la diagonal $\Delta_f:\mathrm{X}\to\mathrm{X}\times_\mathrm{Y}\mathrm{X}$ es biyectiva \sref{IV.1.7.7.1}. Por otra parte, como $f$ es no ramificado \sref{IV.17.4.2}[(d')], $\Delta_f$ es una inmersión abierta \sref{IV.17.4.2}[(b')]. Por consiguiente, $\Delta_f$ es un isomorfismo y la primera proyección $p:\mathrm{X}\times_\mathrm{Y}\mathrm{X}\to\mathrm{X}$ es su inversa.
Sin embargo, $f$ no es plano en $z_1$. De lo contrario, $\sh{O}_{z_1}$ sería un $\sh{O}_y$-módulo fielmente plano \sref[0]{0.6.6.2}; como $\sh{O}_y$ contiene dos ideales primos mínimos distintos $\mathfrak{p}_1,\mathfrak{p}_2$, correspondientes a las dos ramas de $\mathrm{Y}$ en $y$, existirían entonces dos ideales primos distintos de $\sh{O}_{z_1}$ cuyas imágenes inversas por $u:\sh{O}_y\to\sh{O}_{z_1}$ serían $\mathfrak{p}_1$ y $\mathfrak{p}_2$ \sref[0]{0.6.5.1}. Esto es absurdo, pues $\sh{O}_{z_1}$ solo tiene dos ideales primos distintos, $0$ y su ideal máximo $\mathfrak{m}_1$, mientras que $u^{-1}(\mathfrak{m}_1)$ es el ideal máximo $\mathfrak{m}$ de $\sh{O}_y$.
\end{env}

\begin{env}[11.7.4]
\label{IV.11.7.4}
Obsérvese que, en el ejemplo precedente, el homomorfismo $u$ es \emph{inyectivo} cuando $\mathrm{Y}$ es \emph{irreducible} (se puede tomar, por ejemplo, $\mathrm{Y}=\operatorname{Spec}(k[\mathrm{S},\mathrm{T}]/(\mathrm{S}(\mathrm{S}^2+\mathrm{T}^2)-(\mathrm{S}^2-\mathrm{T}^2)))$, una ``cúbica con un punto doble''). En efecto, tras sustituir $\mathrm{Y}$ por un entorno afín de $y$, podemos suponer que $\mathrm{Y}=\operatorname{Spec}(\mathrm{A})$, donde $\mathrm{A}$ es integral, y, por tanto, que $\mathrm{Z}=\operatorname{Spec}(\mathrm{B})$, donde $\mathrm{B}$ es el cierre integral de $\mathrm{A}$. Como $\mathrm{B}$, $\sh{O}_y$ y $\sh{O}_{z_1}$ se identifican entonces con subanillos del
\oldpage[IV-3]{159}
cuerpo de fracciones de $\mathrm{A}$, $u$ es evidentemente inyectivo.
Por otra parte, el homomorfismo $\widehat{u}:\widehat{\sh{O}}_y\to\widehat{\sh{O}}_{z_1}$ no es inyectivo: $\widehat{\sh{O}}_{z_1}$ es un anillo local integral, pues $z_1$ es un punto simple, mientras que $\widehat{\sh{O}}_y$ tiene dos ideales primos mínimos distintos correspondientes a las dos ramas de $\mathrm{Y}$ y, por tanto, tiene divisores de cero.
Esto proporciona un ejemplo que muestra que, en el enunciado de \sref{IV.11.5.2}, la hipótesis de que $\widehat{u}$ es inyectivo no puede sustituirse por la hipótesis de que el propio $u$ es inyectivo, ni siquiera cuando $\mathrm{A}'$ es un anillo local.
En efecto, con la notación precedente, basta tomar $\mathrm{A}=\sh{O}_y$, $\mathrm{A}'=\sh{O}_{z_1}$, $\mathrm{Y}=\operatorname{Spec}(\mathrm{A})$ y $\mathrm{X}=\operatorname{Spec}(\mathrm{A}')$; el razonamiento de \sref{IV.11.7.3} sigue demostrando que la primera proyección $p:\mathrm{X}\times_\mathrm{Y}\mathrm{X}\to\mathrm{X}$ es un isomorfismo, aunque $f$ no es plano en $z_1$.
\end{env}

\begin{env}[11.7.5]
\label{IV.11.7.5}
Los ejemplos de \sref{IV.11.7.2} y \sref{IV.11.7.3} explican la restricción a los anillos locales \emph{unirramificados} en \sref{IV.11.6.1} y \sref{IV.11.6.2}.
Veremos ahora que, en \sref{IV.11.6.1}, no se puede debilitar la hipótesis sobre $\mathrm{A}$ suponiendo únicamente que $\mathrm{A}$ es unirramificado.
Consideremos el anillo local integral completo $\mathrm{A}=\mathbf{R}[[\mathrm{U},\mathrm{V}]]/(\mathrm{U}^2+\mathrm{V}^2)$, que es unirramificado, pero no geométricamente unirramificado \sref{IV.6.15.11}.
Se sabe (\emph{loc.\ cit.}) que, si $u$ y $v$ son las imágenes de $\mathrm{U}$ y $\mathrm{V}$ en $\mathrm{A}$, el cierre integral de $\mathrm{A}$ es $\bar{\mathrm{A}}=\mathrm{A}[t]$, donde $t=v/u$ y $t^2=-1$, de modo que $\bar{\mathrm{A}}$ es isomorfo a $\mathbf{C}[[\mathrm{U}]]$.
Pongamos $\mathrm{Y}=\operatorname{Spec}(\mathrm{A})$ y $\mathrm{X}=\operatorname{Spec}(\bar{\mathrm{A}})$, la normalización de $\mathrm{Y}$ \sref[II]{II.6.3.8}, y sean $y$ y $x$ los puntos cerrados de $\mathrm{Y}$ y $\mathrm{X}$, respectivamente.
Mostraremos que, para una $\mathrm{A}$-álgebra local adecuada $\mathrm{A}'$, si se pone $\mathrm{Y}'=\operatorname{Spec}(\mathrm{A}')$ y $\mathrm{X}'=\mathrm{X}\times_\mathrm{Y}\mathrm{Y}'$, y si $y'$ denota el punto cerrado de $\mathrm{Y}'$, entonces $\sh{O}_{\mathrm{X}'}$ es $\mathrm{Y}'$-plano en un punto de $\mathrm{X}'$ cuyas proyecciones en $\mathrm{X}$ y $\mathrm{Y}'$ son, respectivamente, $x$ y $y'$, pero no es $\mathrm{Y}'$-plano en todos los puntos que tienen estas proyecciones.
De \sref{IV.2.1.4} se seguirá entonces que $\sh{O}_\mathrm{X}$ no es $\mathrm{Y}$-plano en $x$ (lo cual, por lo demás, es trivial \emph{a priori}, pues $\bar{\mathrm{A}}$ no es un $\mathrm{A}$-módulo libre).

Sea $\mathrm{B}=\mathrm{A}\otimes_\mathbf{R}\mathbf{C}$, que es isomorfo a $\mathbf{C}[[\mathrm{U},\mathrm{V}]]/((\mathrm{U}+i\mathrm{V})(\mathrm{U}-i\mathrm{V}))$.
El anillo $\mathrm{B}$ tiene dos ideales primos mínimos $\mathfrak{p}'$ y $\mathfrak{p}''$, generados, respectivamente, por $u+iv$ y $u-iv$, y $\mathfrak{n}=\mathfrak{p}'+\mathfrak{p}''$ es el ideal máximo del anillo local completo $\mathrm{B}$.
Sea $\bar{\mathrm{B}}=\bar{\mathrm{A}}\otimes_\mathbf{R}\mathbf{C}$; es el producto directo de dos álgebras isomorfas a $\mathbf{C}[[\mathrm{U}]]$, generado por los idempotentes $e'=(1\otimes 1+t\otimes i)/2$ y $e''=(1\otimes 1-t\otimes i)/2$.
Como el homomorfismo $\mathrm{A}\to\bar{\mathrm{A}}$ es inyectivo, también lo es $\mathrm{B}\to\bar{\mathrm{B}}$; las imágenes de $u+iv$ y $u-iv$ por esta inyección son, respectivamente, $2ue'$ y $2ue''$.
Se deduce inmediatamente que $\bar{\mathrm{B}}$ se identifica canónicamente con $(\mathrm{B}/\mathfrak{p}')\oplus(\mathrm{B}/\mathfrak{p}'')$.
Tomemos ahora como $\mathrm{A}'$ el anillo local $\mathrm{B}/\mathfrak{p}'$; entonces $\bar{\mathrm{A}}\otimes_\mathrm{A}\mathrm{A}'$ se identifica con $\bar{\mathrm{B}}\otimes_\mathrm{B}\mathrm{A}'$.
Pero $(\mathrm{B}/\mathfrak{p}')\otimes_\mathrm{B}(\mathrm{B}/\mathfrak{p}')=\mathrm{B}/\mathfrak{p}'$ y $(\mathrm{B}/\mathfrak{p}')\otimes_\mathrm{B}(\mathrm{B}/\mathfrak{p}'')=\mathrm{B}/\mathfrak{n}$, de modo que $\bar{\mathrm{A}}\otimes_\mathrm{A}\mathrm{A}'$ es isomorfo a $\mathrm{A}'\oplus(\mathrm{B}/\mathfrak{n})$.
Esto demuestra nuestra afirmación, pues $\mathrm{B}/\mathfrak{n}=\mathrm{A}'/(\mathfrak{n}/\mathfrak{p}')$ no es un $\mathrm{A}'$-módulo plano; de lo contrario, sería un $\mathrm{A}'$-módulo libre (Bourbaki, \emph{Alg.\ comm.}, cap.~II, \S~3, no.~2, cor.~2 de prop.~5), lo cual es absurdo.
\end{env}

\subsection{Un criterio valuativo de planitud}
\label{subsection:IV.11.8}

\begin{theorem}[11.8.1]
\label{IV.11.8.1}
Sean $f:\mathrm{X}\to\mathrm{Y}$ un morfismo localmente de presentación finita, $\mathscr{F}$ un $\sh{O}_\mathrm{X}$-módulo cuasicoherente de presentación finita, $x$ un punto de $\mathrm{X}$ y $y=f(x)$.
Supongamos que el anillo local $\sh{O}_y$ es íntegro (resp.\ reducido y noetheriano).
Para que $\mathscr{F}$ sea $f$-plano en el punto $x$, es necesario y suficiente que, para todo anillo de valoración (resp.\ todo anillo de valoración discreta) $\mathrm{A}'$ y todo homomorfismo local $\sh{O}_y\to\mathrm{A}'$, se satisfaga la condición siguiente: poniendo $\mathrm{Y}'=\operatorname{Spec}(\mathrm{A}')$, $\mathrm{X}'=\mathrm{X}\times_\mathrm{Y}\mathrm{Y}'$, $f'=f_{(\mathrm{Y}')}$, el $\sh{O}_{\mathrm{X}'}$-módulo $\mathscr{F}'=\mathscr{F}\otimes_\mathrm{Y}\mathrm{Y}'$ es $f'$-plano en todos los puntos $x'$ de $\mathrm{X}'$ cuyas proyecciones en $\mathrm{X}$ y $\mathrm{Y}'$ son, respectivamente, $x$ y el punto cerrado $y'$ de $\mathrm{Y}'$.
\end{theorem}

\begin{proof}
La condición es evidentemente necesaria \sref{IV.2.1.4}; queda por demostrar que es suficiente.
Podemos limitarnos evidentemente (\sref[I]{I.3.6.5} y \sref[I]{I.2.4.4}) al caso en que $\mathrm{Y}=\operatorname{Spec}(\mathrm{A})$ es el espectro de un anillo \emph{local} $\mathrm{A}$ y $y$ es el punto cerrado de $\mathrm{Y}$.
\medskip

(i) \emph{Caso en que $\mathrm{A}$ es íntegro.} --- Sean $\mathrm{K}$ el cuerpo de fracciones de $\mathrm{A}$ y $\mathrm{A}_1$ el cierre integral de $\mathrm{A}$; si ponemos $\mathrm{Y}_1=\operatorname{Spec}(\mathrm{A}_1)$, $\mathrm{X}_1=\mathrm{X}\times_\mathrm{Y}\mathrm{Y}_1$, $f_1=f_{(\mathrm{Y}_1)}$, basta, en virtud de \sref{IV.11.5.5}, demostrar que $\mathscr{F}_1=\mathscr{F}\otimes_\mathrm{Y}\mathrm{Y}_1$ es $f_1$-plano en \emph{todos} los puntos $x_1$ de $\mathrm{X}_1$ cuya proyección es el punto $x$ de $\mathrm{X}$.
Ahora bien, si $f_1(x_1)=y_1$, entonces $\mathfrak{i}_{y_1}=\mathfrak{m}_1$, donde $\mathfrak{m}_1$ es un ideal primo de $\mathrm{A}_1$ cuya contracción a $\mathrm{A}$ es $\mathfrak{m}=\mathfrak{i}_y$ (y $\mathfrak{m}_1$ es, además, necesariamente un ideal maximal). Sea $\mathrm{A}'$ un anillo de valoración de $\mathrm{K}$ que domina $\sh{O}_{y_1}=(\mathrm{A}_1)_{\mathfrak{m}_1}$; como el homomorfismo $\mathrm{A}\to\sh{O}_{y_1}$ es local, también lo es $\mathrm{A}\to\mathrm{A}'$.
Existe entonces al menos un
\oldpage[IV-3]{160}
punto $x'\in\mathrm{X}'$ cuyas proyecciones en $\mathrm{X}_1$ y $\mathrm{Y}'$ son, respectivamente, $x_1$ e $y'$ \sref[I]{I.3.4.7}; como la hipótesis afirma que $\mathscr{F}'$ es $f'$-plano en el punto $x'$ y $\sh{O}_{y_1}$ es íntegramente cerrado, podemos aplicar \sref{IV.11.6.1}, y deducimos que $\mathscr{F}_1$ es $f_1$-plano en el punto $x_1$, de donde resulta el teorema en este caso.
\medskip

(ii) \emph{Caso en que $\mathrm{A}$ es reducido y noetheriano.} --- Sean $\mathfrak{p}_i$ ($1\leq i\leq m$) los ideales primos minimales de $\mathrm{A}$; como $\mathrm{A}$ es reducido, se identifica canónicamente con un subanillo del producto de los $\mathrm{A}_i=\mathrm{A}/\mathfrak{p}_i$, que son anillos locales noetherianos; poniendo $\mathrm{Y}_i=\operatorname{Spec}(\mathrm{A}_i)$, $\mathrm{X}_i=\mathrm{X}\times_\mathrm{Y}\mathrm{Y}_i$, $f_i=f_{(\mathrm{Y}_i)}$, se deduce de \sref{IV.11.5.2} que basta demostrar, para cada $i$, que $\mathscr{F}_i=\mathscr{F}\otimes_\mathrm{Y}\mathrm{Y}_i$ es $f_i$-plano en todo punto $x_i$ de $\mathrm{X}_i$ cuyas proyecciones en $\mathrm{X}$ y $\mathrm{Y}_i$ son, respectivamente, $x$ y el punto cerrado $y_i$ de $\mathrm{Y}_i$.
Ahora bien, como $\mathrm{A}_i$ es un anillo local noetheriano íntegro, existe en su cuerpo de fracciones $\mathrm{K}_i$ un subanillo $\mathrm{A}''_i$ que contiene a $\mathrm{A}_i$, que es una $\mathrm{A}_i$-álgebra finita (y, por tanto, un anillo semilocal noetheriano) y cuyos anillos locales son \emph{geométricamente unirramificados} (\sref{IV.6.15.5} y \sref[0]{0.23.2.5}).
Como los ideales maximales de $\mathrm{A}''_i$ están necesariamente sobre el ideal maximal de $\mathrm{A}_i$, se deduce además de \sref[I]{I.3.4.7} y \sref{IV.11.5.2} que basta, poniendo $\mathrm{Y}''_i=\operatorname{Spec}(\mathrm{A}''_i)$, $\mathrm{X}''_i=\mathrm{X}\times_\mathrm{Y}\mathrm{Y}''_i$, $f''_i=f_{(\mathrm{Y}''_i)}$, demostrar que $\mathscr{F}''_i=\mathscr{F}\otimes_\mathrm{Y}\mathrm{Y}''_i$ es $f''_i$-plano en todo punto $x''_i$ de $\mathrm{X}''_i$ cuyas proyecciones en $\mathrm{X}$ y $\mathrm{Y}''_i$ son, respectivamente, $x$ y el punto cerrado $y''_i$ de $\mathrm{Y}''_i$.
Sea ahora $\mathrm{A}'$ un anillo de valoración discreta de $\mathrm{K}_i$ que domina $\mathrm{A}''_i$, y sea $x'$ un punto de $\mathrm{X}'$ cuyas proyecciones en $\mathrm{X}''_i$ y $\mathrm{Y}'$ son, respectivamente, $x''_i$ e $y'$ \sref[I]{I.3.4.7}; como $\sh{O}_{y''_i}$ es geométricamente unirramificado, podemos aplicar de nuevo \sref{IV.11.6.1}, y deducimos que $\mathscr{F}''_i$ es efectivamente $f''_i$-plano en el punto $x''_i$.
\end{proof}

\begin{remarks}[11.8.2]
\label{IV.11.8.2}
\medskip\noindent
\begin{enumerate}[label=(\roman*)]
  \item En el enunciado de \sref{IV.11.8.1}, podemos limitarnos a suponer que la condición sobre $\mathscr{F}'$ se satisface para los anillos de valoración \emph{completos} $\mathrm{A}'$ cuyo cuerpo residual sea \emph{algebraicamente cerrado}.
  En efecto, se sabe que todo anillo de valoración $\mathrm{A}'$ está dominado por un anillo $\mathrm{A}''$ de esta clase \sref[II]{II.7.1.2}, y que, si $\mathrm{A}'$ es un anillo de valoración discreta, lo mismo ocurre con $\mathrm{A}''$ ($\mathbf{0}_{\mathrm{III}}$,~10.3.1).
  \item La demostración de \sref{IV.11.8.1} se simplifica cuando se supone no solo que $\mathrm{A}$ es íntegro y noetheriano, sino también que su completado $\hat{\mathrm{A}}$ es \emph{íntegro}.
  Sustituyendo $\mathrm{X}$ por $\mathrm{X}\otimes_\mathrm{A}\hat{\mathrm{A}}$ y razonando como en la demostración de \sref{IV.11.5.3}, podemos reducirnos, en efecto, a demostrar \sref{IV.11.8.1} cuando $\mathrm{A}=\sh{O}_y$ es íntegro, noetheriano y \emph{completo}.
  Ahora bien, se sabe por \sref[II]{II.7.1.7} que tal anillo $\mathrm{A}$ está dominado por un anillo de valoración discreta completo; la conclusión resulta, por tanto, directamente de \sref{IV.11.5.2}.
\end{enumerate}
\end{remarks}

\subsection{Familias separantes y universalmente separantes de homomorfismos de haces de módulos}
\label{subsection:IV.11.9}

\begin{env}[11.9.1]
\label{IV.11.9.1}
Sea $\mathrm{X}$ un preesquema, sea $(f_\lambda)_{\lambda\in\mathrm{L}}$ una familia de morfismos $f_\lambda:\mathrm{Z}_\lambda\to\mathrm{X}$ y sea $\mathscr{F}$ un $\sh{O}_\mathrm{X}$-módulo cuasicoherente; para cada $\lambda\in\mathrm{L}$, supongamos dados un $\sh{O}_{\mathrm{Z}_\lambda}$-módulo cuasicoherente $\sh{G}_\lambda$ y un homomorfismo
\[
  u_\lambda:\mathscr{F}\to(f_\lambda)_*(\sh{G}_\lambda).
\label{IV.11.9.1.1}
\]
\end{env}

\oldpage[IV-3]{161}
Decimos que la familia $(u_\lambda)$ (o la familia correspondiente de los $u_\lambda^{\#}:(f_\lambda)^*(\mathscr{F})\to\sh{G}_\lambda$) es \emph{separante} si la intersección de los núcleos de los $u_\lambda$ es \emph{cero}.
En otras palabras, esto significa que, para todo abierto $\mathrm{U}$ de $\mathrm{X}$ y toda sección $t$ de $\mathscr{F}$ sobre $\mathrm{U}$ tales que, para \emph{cada} $\lambda$, la sección $u_\lambda(t)$ (que, por definición, es una sección de $\sh{G}_\lambda$ sobre $f_\lambda^{-1}(\mathrm{U})$) sea cero, la propia $t$ es cero.

\begin{env}[11.9.2]
\label{IV.11.9.2}
Con las notaciones de \sref{IV.11.9.1}, sea $\mathrm{M}$ un segundo conjunto de índices; para cada $\lambda\in\mathrm{L}$, sea $(g_{\lambda\mu})_{\mu\in\mathrm{M}}$ una familia de morfismos $g_{\lambda\mu}:\mathrm{Z}'_{\lambda\mu}\to\mathrm{Z}_\lambda$; para cada par $(\lambda,\mu)$, supongamos dados un $\sh{O}_{\mathrm{Z}'_{\lambda\mu}}$-módulo cuasicoherente $\sh{H}_{\lambda\mu}$ y un homomorfismo $v_{\lambda\mu}:\sh{G}_\lambda\to(g_{\lambda\mu})_*(\sh{H}_{\lambda\mu})$; pongamos $h_{\lambda\mu}=f_\lambda\circ g_{\lambda\mu}:\mathrm{Z}'_{\lambda\mu}\to\mathrm{X}$ y consideremos el homomorfismo compuesto
\[
  w_{\lambda\mu}:\mathscr{F}\xrightarrow{u_\lambda}(f_\lambda)_*(\sh{G}_\lambda)\xrightarrow{(f_\lambda)_*(v_{\lambda\mu})}(h_{\lambda\mu})_*(\sh{H}_{\lambda\mu}).
\]
Supongamos que, para cada $\lambda\in\mathrm{L}$, la familia $(v_{\lambda\mu})_{\mu\in\mathrm{M}}$ es separante: entonces lo mismo ocurre con la familia $(f_\lambda)_*(v_{\lambda\mu})$ ($\mu\in\mathrm{M}$), como se ve inmediatamente.
Se concluye que, para que la familia $(u_\lambda)$ sea separante, es necesario y suficiente que lo sea la familia $(w_{\lambda\mu})_{(\lambda,\mu)\in\mathrm{L}\times\mathrm{M}}$.
\end{env}

\begin{env}[11.9.3]
\label{IV.11.9.3}
Para verificar que la familia $(u_\lambda)$ es separante (con las notaciones de \sref{IV.11.9.1}), podemos reducirnos evidentemente primero al caso en que $\mathrm{X}$ es \emph{afín}, pues la propiedad es local sobre $\mathrm{X}$.
Podemos suponer además que $\mathrm{Z}_\lambda=\mathrm{X}$ \emph{para todo} $\lambda\in\mathrm{L}$.
En efecto, sea $\mathrm{M}$ un conjunto de índices, suma de una familia $(\mathrm{M}_\lambda)_{\lambda\in\mathrm{L}}$, y, para cada $\lambda\in\mathrm{L}$, sea $(\mathrm{Y}_{\lambda\mu})_{\mu\in\mathrm{M}_\lambda}$ un recubrimiento abierto afín de $\mathrm{Z}_\lambda$; sea $j_{\lambda\mu}:\mathrm{Y}_{\lambda\mu}\to\mathrm{Z}_\lambda$ la inmersión canónica y pongamos $\sh{H}_{\lambda\mu}=j_{\lambda\mu}^*(\sh{G}_\lambda)=\sh{G}_\lambda|\mathrm{Y}_{\lambda\mu}$.
Si consideramos el homomorfismo canónico $v_{\lambda\mu}:\sh{G}_\lambda\to(j_{\lambda\mu})_*(j_{\lambda\mu}^*(\sh{G}_\lambda))=(j_{\lambda\mu})_*(\sh{H}_{\lambda\mu})$ relativo a $j_{\lambda\mu}$ \sref[0]{0.4.4.3.2}, es inmediato que, para cada $\lambda\in\mathrm{L}$, la familia $(v_{\lambda\mu})_{\mu\in\mathrm{M}_\lambda}$ es separante.
En virtud de \sref{IV.11.9.2}, nos reducimos, por tanto, a demostrar que es separante la familia de homomorfismos compuestos $((f_\lambda)_*(v_{\lambda\mu}))\circ u_\lambda$, donde los $f_\lambda$ son \emph{afines}.
Pero entonces los $(f_\lambda)_*(\sh{G}_\lambda)$ son $\sh{O}_\mathrm{X}$-módulos cuasicoherentes \sref[I]{I.1.6.2} y, en virtud de la definición, podemos sustituir los $\mathrm{Z}_\lambda$ por $\mathrm{X}$ y los $\sh{G}_\lambda$ por $(f_\lambda)_*(\sh{G}_\lambda)$, de donde resulta nuestra afirmación.

Observemos además que, si $\mathrm{L}$ es \emph{finito} y los $f_\lambda$ son \emph{cuasicompactos}, podemos suponer, en la reducción anterior, que los $\mathrm{M}_\lambda$ son también \emph{finitos}; en este caso nos reducimos así a verificar que una familia \emph{finita} de homomorfismos de $\mathscr{F}$ en $\sh{O}_\mathrm{X}$-módulos cuasicoherentes es separante.
\end{env}

\begin{env}[11.9.4]
\label{IV.11.9.4}
Consideremos, pues, el caso en que $\mathrm{Z}_\lambda=\mathrm{X}$ para todo $\lambda$ y en que $\mathrm{X}=\operatorname{Spec}(\mathrm{A})$ es afín; entonces $\mathscr{F}=\widetilde{\mathrm{M}}$ y $\sh{G}_\lambda=\widetilde{\mathrm{N}}_\lambda$, donde $\mathrm{M}$ y los $\mathrm{N}_\lambda$ son $\mathrm{A}$-módulos, y $u_\lambda=\widetilde{\varphi}_\lambda$, donde los $\varphi_\lambda:\mathrm{M}\to\mathrm{N}_\lambda$ son homomorfismos de $\mathrm{A}$-módulos.
Decir que la familia $(u_\lambda)$ es separante significa entonces que, para todo $s\in\mathrm{A}$, la intersección de los núcleos de los $(\varphi_\lambda)_s:\mathrm{M}_s\to(\mathrm{N}_\lambda)_s$ es cero.
Se dice también entonces que la familia $(\varphi_\lambda)$ es \emph{separante}.
Observemos que, si $\mathrm{L}$ es \emph{finito}, esto equivale a decir que la intersección de los núcleos de los $\varphi_\lambda$ es cero, pues entonces
\[
  \bigcap_{\lambda\in\mathrm{L}}\operatorname{Ker}((\varphi_\lambda)_s)
  =\left(\bigcap_{\lambda\in\mathrm{L}}\operatorname{Ker}(\varphi_\lambda)\right)_s
  \quad \sref[0]{0.1.3.2}.
\]
Pero esta relación deja de ser exacta en general
\oldpage[IV-3]{162}
cuando $\mathrm{L}$ es \emph{infinito}, y el hecho de que la intersección de los núcleos de los $\varphi_\lambda$ sea cero \emph{no implica, por tanto},
en general, que la familia $(\varphi_\lambda)$ sea separante.
Por ejemplo, supongamos que $\mathrm{A}$ es un anillo de valoración discreta de ideal maximal $\mathfrak{m}$, y consideremos la familia de homomorfismos $\varphi_k:\mathrm{A}\to\mathrm{A}/\mathfrak{m}^k$; sus núcleos son los ideales $\mathfrak{m}^k$, cuya intersección es cero. Sin embargo, esta familia no es separante, pues las fibras de todos los $\mathfrak{m}^k$ en el punto genérico $x$ de $\mathrm{X}=\operatorname{Spec}(\mathrm{A})$ (que es abierto en $\mathrm{X}$) son iguales a $\sh{O}_x=k(x)$, el cuerpo de fracciones de $\mathrm{A}$, y su intersección no es, por tanto, cero.
\end{env}

\begin{env}[11.9.5]
\label{IV.11.9.5}
En lo que sigue nos ocuparemos principalmente del problema del \emph{cambio de base} para las familias separantes.
Con las notaciones de \sref{IV.11.9.1}, consideremos un morfismo $g:\mathrm{X}'\to\mathrm{X}$ y pongamos $\mathscr{F}'=\mathscr{F}\otimes_{\sh{O}_\mathrm{X}}\sh{O}_{\mathrm{X}'}=g^*(\mathscr{F})$ y, para cada $\lambda$, $\mathrm{Z}'_\lambda=\mathrm{Z}_\lambda\times_\mathrm{X}\mathrm{X}'$, $f'_\lambda=f_\lambda\times 1$, $\sh{G}'_\lambda=\sh{G}_\lambda\otimes_{\sh{O}_{\mathrm{Z}_\lambda}}\sh{O}_{\mathrm{Z}'_\lambda}=(g'_\lambda)^*(\sh{G}_\lambda)$, donde $g'_\lambda:\mathrm{Z}'_\lambda\to\mathrm{Z}_\lambda$ es la proyección.
Para cada $\lambda$, denotamos por $u'_\lambda:\mathscr{F}'\to(f'_\lambda)_*(\sh{G}'_\lambda)$ el homomorfismo que se obtiene de la manera siguiente: sea
\[
  \psi:(f_\lambda)_*(\sh{G}_\lambda)\to(f_\lambda)_*((g'_\lambda)_*((g'_\lambda)^*(\sh{G}_\lambda)))=g_*((f'_\lambda)_*(\sh{G}'_\lambda))
\]
el homomorfismo $(f_\lambda)_*(\rho_{\sh{G}_\lambda})$, donde $\rho_{\sh{G}_\lambda}$ es el homomorfismo canónico \sref[0]{0.4.4.3.2} correspondiente a $g'_\lambda$.
Entonces $u'_\lambda$ se define como el compuesto
\[
  g^*(\mathscr{F})\xrightarrow{g^*(u_\lambda)}g^*((f_\lambda)_*(\sh{G}_\lambda))\xrightarrow{g^*(\psi)}g^*g_*((f'_\lambda)_*(\sh{G}'_\lambda))\xrightarrow{\sigma}(f'_\lambda)_*(\sh{G}'_\lambda)
\]
donde $\sigma=\sigma_{(f'_\lambda)_*(\sh{G}'_\lambda)}$ es el homomorfismo canónico \sref[0]{0.4.4.3.3} correspondiente a $g$.
Diremos, para abreviar, que $u'_\lambda$ se \emph{deduce de $u_\lambda$ por el cambio de base} $g$.
Cuando $f_\lambda:\mathrm{Z}_\lambda\to\mathrm{X}$ es un morfismo \emph{afín}, $(f'_\lambda)_*(\sh{G}'_\lambda)=g^*((f_\lambda)_*(\sh{G}_\lambda))$ y, en este caso, se tiene simplemente $u'_\lambda=g^*(u_\lambda)$.

También podemos interpretar $u'_\lambda$ de la manera siguiente: basta conocer el valor de $u'_\lambda(t')$, donde $t'$ es una sección de $\mathscr{F}'$ sobre un abierto $\mathrm{U}'$ de $\mathrm{X}'$ del tipo particular siguiente: $t'$ es la restricción a $\mathrm{U}'$ de la imagen canónica, por $\Gamma(\mathrm{U},\mathscr{F})\to\Gamma(g^{-1}(\mathrm{U}),\mathscr{F}')$, de una sección $t$ de $\mathscr{F}$ sobre un abierto $\mathrm{U}$ de $\mathrm{X}$ que contiene a $g(\mathrm{U}')$ (estas secciones generan el $\sh{O}_{\mathrm{X}'}$-módulo $\mathscr{F}'$ \sref[0]{0.3.7.1}).
Consideremos la sección $u_\lambda(t)$ de $\sh{G}_\lambda$ sobre $f_\lambda^{-1}(\mathrm{U})$ y su imagen canónica $t''$ por $\Gamma(f_\lambda^{-1}(\mathrm{U}),\sh{G}_\lambda)\to\Gamma(g_\lambda'^{-1}(f_\lambda^{-1}(\mathrm{U})),\sh{G}'_\lambda)$; entonces $u'_\lambda(t')$ es la restricción de $t''$ a $f_\lambda'^{-1}(\mathrm{U}')$.
Consideremos, en particular, el caso en que $\mathrm{Z}_\lambda$ es el subpreesquema inducido sobre un subconjunto abierto de $\mathrm{X}$, en que $\sh{G}_\lambda=j_\lambda^*(\mathscr{F})$, siendo $j_\lambda:\mathrm{Z}_\lambda\to\mathrm{X}$ la inmersión canónica, y en que $u_\lambda$ es el homomorfismo canónico $\rho_\mathscr{F}:\mathscr{F}\to(j_\lambda)_*(j_\lambda^*(\mathscr{F}))=(j_\lambda)_*(\sh{G}_\lambda)$ \sref[0]{0.4.4.3.2} correspondiente a $j_\lambda$.
Entonces $\mathrm{Z}'_\lambda$ está inducido sobre un subconjunto abierto de $\mathrm{X}'$, y la interpretación anterior muestra que $u'_\lambda$ no es sino el homomorfismo canónico $\rho_{\mathscr{F}'}:\mathscr{F}'\to(j'_\lambda)_*((j'_\lambda)^*(\mathscr{F}'))=(j'_\lambda)_*(\sh{G}'_\lambda)$ correspondiente a la inmersión canónica $j'_\lambda:\mathrm{Z}'_\lambda\to\mathrm{X}'$.
\end{env}

\begin{env}[11.9.6]
\label{IV.11.9.6}
En las condiciones de \sref{IV.11.9.5}, supongamos que $\mathrm{X}$ y $\mathrm{X}'$ son ambos afines y que se quiere demostrar que toda sección $t'$ de $\mathscr{F}'$ sobre $\mathrm{X}'$ cuyas imágenes por todos los $u'_\lambda$ sean cero es ella misma cero.
Entonces también podemos limitarnos al caso en que $\mathrm{Z}_\lambda=\mathrm{X}$ para todo $\lambda\in\mathrm{L}$.
En efecto, por \sref{IV.11.9.3} y \sref{IV.11.9.5}, si ponemos $\mathrm{Y}'_{\lambda\mu}=g_\lambda'^{-1}(\mathrm{Y}_{\lambda\mu})$, el homomorfismo $v'_{\lambda\mu}$ deducido de $v_{\lambda\mu}$ por el cambio de base $g$ no es sino el homomorfismo canónico
\oldpage[IV-3]{163}
$\rho_{\sh{G}'_\lambda}:\sh{G}'_\lambda\to(j'_{\lambda\mu})_*(j_{\lambda\mu}'^*(\sh{G}'_\lambda))$ correspondiente a la inmersión canónica $j'_{\lambda\mu}:\mathrm{Y}'_{\lambda\mu}\to\mathrm{Z}'_\lambda$, como vimos en \sref{IV.11.9.5}.
La afirmación se deduce entonces de los razonamientos de \sref{IV.11.9.2} y \sref{IV.11.9.3}, sustituyendo $\mathrm{Y}_{\lambda\mu}$ y $v_{\lambda\mu}$ por $\mathrm{Y}'_{\lambda\mu}$ y $v'_{\lambda\mu}$.

Se tiene una reducción análoga cuando se quiere demostrar que la familia $(u'_\lambda)$ es \emph{separante} (siendo $\mathrm{X}$ y $\mathrm{X}'$ afines): esto se deduce de nuevo de \sref{IV.11.9.2} y \sref{IV.11.9.3}.
\end{env}

\begin{proposition}[11.9.7]
\label{IV.11.9.7}
Con las notaciones de \sref{IV.11.9.1} y \sref{IV.11.9.5}, supongamos que $\mathrm{X}=\operatorname{Spec}(\mathrm{A})$ y $\mathrm{X}'=\operatorname{Spec}(\mathrm{A}')$ son afines y, además, que $\mathrm{A}'$ es un $\mathrm{A}$-módulo proyectivo.
Entonces, si la familia $(u_\lambda)$ es separante, toda sección $t'$ de $\mathscr{F}'$ sobre $\mathrm{X}'$ cuyas imágenes por todos los $u'_\lambda$ sean cero es ella misma cero.
\end{proposition}

\begin{proof}
Hemos visto en \sref{IV.11.9.6} que podemos limitarnos al caso en que todos los $\mathrm{Z}_\lambda$ son \emph{iguales a} $\mathrm{X}$.
La proposición es entonces una consecuencia del siguiente lema:
\end{proof}

\begin{lemma}[11.9.7.1]
\label{IV.11.9.7.1}
Sea $\mathrm{A}$ un anillo, sea $(\mathrm{M}_\lambda)_{\lambda\in\mathrm{L}}$ una familia de $\mathrm{A}$-módulos, sea $\mathrm{M}$ un $\mathrm{A}$-módulo y, para cada $\lambda$, sea $u_\lambda:\mathrm{M}\to\mathrm{M}_\lambda$ un homomorfismo.
Supongamos que la intersección de los núcleos de los $u_\lambda$ se reduce a cero.
Entonces, para todo $\mathrm{A}$-módulo proyectivo $\mathrm{N}$, la intersección de los núcleos de los homomorfismos $u_\lambda\otimes 1:\mathrm{M}\otimes_\mathrm{A}\mathrm{N}\to\mathrm{M}_\lambda\otimes_\mathrm{A}\mathrm{N}$ se reduce a cero.
\end{lemma}

\begin{proof}
En efecto, $\mathrm{N}$ es un factor directo de un $\mathrm{A}$-módulo libre $\mathrm{L}$, y basta evidentemente demostrar que la intersección de los núcleos de los homomorfismos $v_\lambda:\mathrm{M}\otimes_\mathrm{A}\mathrm{L}\to\mathrm{M}_\lambda\otimes_\mathrm{A}\mathrm{L}$ se reduce a cero, pues $u_\lambda\otimes 1$ es la restricción de $v_\lambda$.
Pero la afirmación resulta entonces trivialmente de la hipótesis.
\end{proof}

\begin{remarks}[11.9.8]
\label{IV.11.9.8}
No sabemos si, bajo las hipótesis de la proposición \sref{IV.11.9.7}, la familia $(u'_\lambda)$ es \emph{separante}: en efecto, por \sref{IV.11.9.4} habría que demostrar que toda sección $t'$ de $\mathscr{F}'$ sobre un subconjunto abierto $\mathrm{D}(h')\subset\mathrm{X}'$ (donde $h'\in\mathrm{A}'$) tal que todos los $u'_\lambda(t')$ sean cero es ella misma cero.
Ahora bien, no podemos aplicar la proposición \sref{IV.11.9.7} a $\mathrm{D}(h')=\operatorname{Spec}(\mathrm{A}'_{h'})$, pues el hecho de que $\mathrm{A}'$ sea un $\mathrm{A}$-módulo proyectivo (o incluso libre) no implica que $\mathrm{A}'_{h'}$ sea un $\mathrm{A}$-módulo proyectivo.
Por ejemplo, podemos tomar como $\mathrm{A}$ un anillo de valoración discreta, como $\mathrm{A}'$ un anillo de valoración discreta que sea un $\mathrm{A}$-módulo libre de rango $2$, y como $\mathrm{A}'_{h'}$ el cuerpo de fracciones de $\mathrm{A}'$.
\end{remarks}

Se tiene, no obstante, el resultado siguiente:

\begin{corollary}[11.9.9]
\label{IV.11.9.9}
Sea $\mathrm{X}$ un preesquema artiniano y sea $g:\mathrm{X}'\to\mathrm{X}$ un morfismo plano (observemos que estas dos condiciones se satisfacen si $\mathrm{X}$ es el espectro de un \emph{cuerpo} y $g$ es un morfismo arbitrario).
Entonces, con las notaciones de \sref{IV.11.9.1} y \sref{IV.11.9.5}, si la familia $(u_\lambda)$ es separante, lo mismo ocurre con la familia $(u'_\lambda)$.
\end{corollary}

\begin{proof}
Podemos limitarnos evidentemente al caso en que $\mathrm{X}=\operatorname{Spec}(\mathrm{A})$ es el espectro de un anillo local artiniano $\mathrm{A}$ \sref[I]{I.6.2.2}; observemos entonces que, para todo abierto afín $\mathrm{U}'=\operatorname{Spec}(\mathrm{A}')$ de $\mathrm{X}'$, $\mathrm{A}'$ es un $\mathrm{A}$-módulo plano y, por tanto, \emph{proyectivo} ($\mathbf{0}_{\mathrm{III}}$,~10.1.3).
Basta, por tanto, aplicar \sref{IV.11.9.7} a todo abierto afín de $\mathrm{X}'$ para obtener el corolario.
\end{proof}

\begin{theorem}[11.9.10]
\label{IV.11.9.10}
Sea $\mathrm{X}$ un preesquema, sea $(u_\lambda)$ una familia de homomorfismos \sref{IV.11.9.1.1}, sea $g:\mathrm{X}'\to\mathrm{X}$ un morfismo y sea $(u'_\lambda)$ la familia de homomorfismos deducida de $(u_\lambda)$ por el cambio de base $g$ \sref{IV.11.9.5}.
\begin{enumerate}[label=(\roman*)]
  \item Si $g$ es un morfismo fielmente plano y la familia $(u'_\lambda)$ es separante, entonces la familia $(u_\lambda)$ es separante.
  \item Supongamos que $g$ es un morfismo plano y que, además, se satisface una de las dos condiciones siguientes:
  \begin{enumerate}[label=\emph{\alph*)}]
    \item $\mathrm{L}$ es finito y los $f_\lambda$ son cuasicompactos.
    \item El morfismo $g$ es localmente de presentación finita.
  \end{enumerate}
  Entonces, si la familia $(u_\lambda)$ es separante, lo mismo ocurre con la familia $(u'_\lambda)$.
\end{enumerate}
\end{theorem}

\begin{proof}
\oldpage[IV-3]{164}
\begin{enumerate}[label=(\roman*)]
  \item En virtud de \sref{IV.2.2.8}, basta demostrar que, si una sección $t$ de $\mathscr{F}$ sobre un abierto $\mathrm{U}$ de $\mathrm{X}$ pertenece al núcleo de cada uno de los $u_\lambda$, su imagen $t'$ por el homomorfismo canónico $\Gamma(\rho):\Gamma(\mathrm{U},\mathscr{F})\to\Gamma(g^{-1}(\mathrm{U}),g^*(\mathscr{F}))=\Gamma(\mathrm{U},g_*(g^*(\mathscr{F})))$ es cero.
  Ahora bien, las imágenes de $t'$ por los $g^*(u_\lambda)$ son las imágenes de $u_\lambda(t)$ por el homomorfismo $\Gamma(\mathrm{U},(f_\lambda)_*(\sh{G}_\lambda))\to\Gamma(g^{-1}(\mathrm{U}),g^*((f_\lambda)_*(\sh{G}_\lambda)))$; por tanto, son cero y, \emph{a fortiori}, $u'_\lambda(t')=0$ para todo $\lambda$. Por hipótesis, se sigue que $t'=0$, lo que demuestra~(i).

  \item Como la cuestión es local sobre $\mathrm{X}$ y $\mathrm{X}'$, podemos restringirnos al caso en que $\mathrm{X}=\operatorname{Spec}(\mathrm{A})$ y $\mathrm{X}'=\operatorname{Spec}(\mathrm{A}')$ son afines y $\mathrm{A}'$ es un $\mathrm{A}$-módulo plano; basta probar que toda sección $z'$ de $\mathscr{F}'$ sobre $\mathrm{X}'$ cuyas imágenes por todos los $u'_\lambda$ sean cero es ella misma cero.
  Hemos visto en \sref{IV.11.9.6} que también podemos restringirnos al caso en que $\mathrm{Z}_\lambda=\mathrm{X}$ para todo $\lambda\in\mathrm{L}$.
  Distinguimos ahora los dos casos.
  \medskip

  \emph{a)} Si $\mathrm{L}$ es finito y los $f_\lambda$ son cuasicompactos, hemos visto en \sref{IV.11.9.3} que podemos reducirnos además al caso en que $\mathrm{Z}_\lambda=\mathrm{X}$ para todo $\lambda\in\mathrm{L}$ y en que, asimismo, $\mathrm{L}$ es finito.
  Entonces, decir que la intersección de los núcleos de los $u_\lambda:\mathscr{F}\to\sh{G}_\lambda$ es cero equivale a decir que el homomorfismo $u=(u_\lambda):\mathscr{F}\to\sh{G}=\bigoplus_\lambda\sh{G}_\lambda$ es inyectivo.
  Como $u'=g^*(u)$ es inyectivo porque $g$ es plano, la proposición queda demostrada en este caso.

  \emph{b)} Sean $\mathscr{F}=\widetilde{\mathrm{M}}$ y $\sh{G}_\lambda=\widetilde{\mathrm{N}}_\lambda$, y pongamos $\mathrm{M}'=\mathrm{M}\otimes_\mathrm{A}\mathrm{A}'$, $\mathrm{N}'_\lambda=\mathrm{N}_\lambda\otimes_\mathrm{A}\mathrm{A}'$, de modo que $\mathscr{F}'=\widetilde{\mathrm{M}'}$ y $\sh{G}'_\lambda=\widetilde{\mathrm{N}'_\lambda}$; por abuso de notación, seguiremos denotando por $u_\lambda$ el homomorfismo $\mathrm{M}\to\mathrm{N}_\lambda$, y por $u'_\lambda$ el homomorfismo $u_\lambda\otimes 1:\mathrm{M}'\to\mathrm{N}'_\lambda$.
  Sea $z'\in\mathrm{M}'$ un elemento tal que $u'_\lambda(z')=0$ para todo $\lambda$; queda por demostrar que $z'=0$.
  Ahora bien, la hipótesis de que $g$ es plano y de presentación finita implica, por \sref{IV.11.3.15}, que existe una sucesión finita $(s_i)_{1\leq i\leq n}$ de elementos de $\mathrm{A}$ tal que, si se pone $\mathfrak{J}_k=s_1\mathrm{A}+\ldots+s_k\mathrm{A}$ y $\mathrm{A}_i=\mathrm{A}_{s_i}/\mathfrak{J}_{i-1}\mathrm{A}_{s_i}$, el anillo $\mathrm{A}'_i=\mathrm{A}'\otimes_\mathrm{A}\mathrm{A}_i$ es \emph{libre} como $\mathrm{A}_i$-módulo para $1\leq i\leq n$, y $\mathfrak{J}_n=\mathrm{A}$.
  La proposición quedará demostrada si probamos la siguiente afirmación para $1\leq i\leq n$:
  \begin{enumerate}[label=$(*_i)$]
    \item \emph{Existe un entero $m_i>0$ tal que $s_j^{m_i}z'=0$ para $j\leq i$.}
  \end{enumerate}
  En efecto, poniendo $k=m_n$ y observando que los $s_i^k$ ($1\leq i\leq n$) generan el ideal unidad de $\mathrm{A}$, la afirmación $(*_n)$ mostrará que $z'$, por ser combinación lineal de los $s_i^k z'$, es cero.

  Demostremos $(*_i)$ por inducción sobre $i$; la afirmación es trivialmente cierta para $i=0$.
  Supongamos $i>0$ y sea $m$ un múltiplo común de los $m_j$ para $j\leq i-1$.
  Observemos (para $1\leq h\leq n$) que, si $\mathfrak{J}'_h$ es el ideal generado por los $s_j^m$ ($1\leq j\leq h$), entonces $\mathfrak{J}_h/\mathfrak{J}'_h$ es nilpotente.
  Sustituir los $s_j$ por $s_j^m$ para $1\leq j\leq n$ equivale, por tanto, a sustituir, para $1\leq h\leq n$, $\mathrm{A}_h$ por
  $\mathrm{B}_h=\mathrm{A}_{s_h}/\mathfrak{J}'_{h-1}\mathrm{A}_{s_h}$, para que
  $\mathrm{A}_h=\mathrm{B}_h/((\mathfrak{J}_{h-1}/\mathfrak{J}'_{h-1})\mathrm{B}_h)$; como $\mathrm{A}'$ es un $\mathrm{A}$-módulo plano, se deduce de ($\mathbf{0}_{\mathrm{III}}$,~10.1.2) que $\mathrm{B}'_h=\mathrm{A}'\otimes_\mathrm{A}\mathrm{B}_h$ vuelve a ser un $\mathrm{B}_h$-módulo libre.
  Por consiguiente, se pueden
  \oldpage[IV-3]{165}
  sustituir todos los $s_j$ ($1\leq j\leq n$) por $s_j^m$ sin alterar las propiedades de $\mathfrak{J}_h$ y $\mathrm{A}'_h$, y suponer que $m=1$.
  Entonces $z'$, al estar aniquilado por $\mathfrak{J}_{i-1}$, se identifica con un elemento de $\operatorname{Hom}_{\mathrm{A}'}(\mathrm{A}'/\mathfrak{J}_{i-1}\mathrm{A}',\mathrm{M}')$; y, como $\mathfrak{J}_{i-1}\mathrm{A}'$ es un ideal de tipo finito de $\mathrm{A}'$, este módulo de homomorfismos se identifica a su vez con $\operatorname{Hom}_\mathrm{A}(\mathrm{A}/\mathfrak{J}_{i-1},\mathrm{M})\otimes_\mathrm{A}\mathrm{A}'$ ($\mathbf{0}_{\mathrm{I}}$,~6.2.2).
  Sea $v_\lambda=\operatorname{Hom}(1,u_\lambda):\operatorname{Hom}_\mathrm{A}(\mathrm{A}/\mathfrak{J}_{i-1},\mathrm{M})\to\operatorname{Hom}_\mathrm{A}(\mathrm{A}/\mathfrak{J}_{i-1},\mathrm{N}_\lambda)$ el homomorfismo deducido de $u_\lambda$; la familia $(v_\lambda)$ también es \emph{separante}.
  En efecto, para todo $t\in\mathrm{A}$,
  $(\operatorname{Hom}_\mathrm{A}(\mathrm{A}/\mathfrak{J}_{i-1},\mathrm{M}))_t
  =\operatorname{Hom}_{\mathrm{A}_t}(\mathrm{A}_t/(\mathfrak{J}_{i-1})_t,\mathrm{M}_t)$,
  y lo mismo vale sustituyendo $\mathrm{M}$ por $\mathrm{N}_\lambda$, pues el ideal $\mathfrak{J}_{i-1}$ es de tipo finito ($\mathbf{0}_{\mathrm{I}}$,~1.3.5); como, por hipótesis, la intersección de los núcleos de los $(u_\lambda)_t$ es cero, lo mismo ocurre con la intersección de los núcleos de los $(v_\lambda)_t=\operatorname{Hom}(1,(u_\lambda)_t)$, de donde nuestra afirmación por \sref{IV.11.9.4}.
  Sustituyendo $\mathrm{A}$ por $\mathrm{A}/\mathfrak{J}_{i-1}$, $\mathrm{M}$ por $\operatorname{Hom}_\mathrm{A}(\mathrm{A}/\mathfrak{J}_{i-1},\mathrm{M})$, $\mathrm{N}_\lambda$ por $\operatorname{Hom}_\mathrm{A}(\mathrm{A}/\mathfrak{J}_{i-1},\mathrm{N}_\lambda)$ (que son $(\mathrm{A}/\mathfrak{J}_{i-1})$-módulos), $u_\lambda$ por $v_\lambda$ y, finalmente, $\mathrm{A}'$ por $\mathrm{A}'/\mathfrak{J}_{i-1}\mathrm{A}'$, vemos que podemos reducirnos al caso en que, en la situación inicial, el elemento $s=s_i\in\mathrm{A}$ es tal que $\mathrm{A}'_s$ es un $\mathrm{A}_s$-módulo libre.
  Ahora bien, la familia $(u_\lambda)_s:\mathrm{M}_s\to(\mathrm{N}_\lambda)_s$ es separante por hipótesis; como $(u'_\lambda)_s(z'/1)=0$, se deduce de \sref{IV.11.9.7} que $z'/1=0$ en $\mathrm{M}'_s$, lo que significa que existe un entero $r$ tal que $s^r z'=0$ en $\mathrm{M}'$. Esto concluye por inducción la demostración de $(*_i)$.
\end{enumerate}
\end{proof}

\begin{remark}[11.9.11]
\label{IV.11.9.11}
Limitémonos, por simplicidad, al caso en que $\mathrm{Z}_\lambda=\mathrm{X}$ para todo $\lambda$.
Conviene observar que, si la familia de homomorfismos $u_\lambda:\mathscr{F}\to\sh{G}_\lambda$ es separante, \emph{no se sigue necesariamente} que, para todo $x\in\mathrm{X}$, la intersección de los núcleos de los homomorfismos $(u_\lambda)_x:\mathscr{F}_x\to(\sh{G}_\lambda)_x$ se reduzca a $0$.
Por ejemplo, sea $\mathrm{X}$ un preesquema localmente noetheriano de Jacobson, de dimensión $\geq 1$, y sea $\mathscr{F}$ un $\sh{O}_\mathrm{X}$-módulo coherente; para todo punto cerrado $x$ de $\mathrm{X}$ y todo entero $n\geq 0$, $\mathscr{F}/\mathfrak{m}_x^n\mathscr{F}$ es un $\sh{O}_\mathrm{X}$-módulo coherente cuyo soporte está contenido en $\{x\}$.
La familia de homomorfismos canónicos $\mathscr{F}\to\mathscr{F}/\mathfrak{m}_x^n\mathscr{F}$ (donde $x$ recorre el conjunto $\mathrm{X}_0$ de los puntos cerrados de $\mathrm{X}$ y $n$ el conjunto de los enteros $\geq 0$) es \emph{separante}: en efecto, si $t$ es una sección de $\mathscr{F}$ sobre un abierto $\mathrm{U}$ de $\mathrm{X}$ cuyas imágenes en todos los $\Gamma(\mathrm{U},\mathscr{F}/\mathfrak{m}_x^n\mathscr{F})$ son nulas, se tiene inmediatamente $t_x=0$ para todo punto cerrado $x\in\mathrm{U}$; como el conjunto de los puntos cerrados contenidos en $\mathrm{U}$ es muy denso en $\mathrm{U}$, esto implica que $t=0$ \sref{IV.10.2.1}.
Sin embargo, si se toma $\mathscr{F}=\sh{O}_\mathrm{X}$ y si $y\in\mathrm{X}$ es un punto \emph{no cerrado}, se tiene $(\sh{O}_\mathrm{X}/\mathfrak{m}_x^n\sh{O}_\mathrm{X})_y=0$ para todo punto cerrado $x$ de $\mathrm{X}$, pero $\sh{O}_{\mathrm{X},y}\neq 0$, y la intersección de los núcleos de los homomorfismos $\sh{O}_{\mathrm{X},y}\to(\sh{O}_\mathrm{X}/\mathfrak{m}_x^n\sh{O}_\mathrm{X})_y$ es igual a $\sh{O}_{\mathrm{X},y}$.
\end{remark}

\begin{lemma}[11.9.12]
\label{IV.11.9.12}
Sean las notaciones las de \sref{IV.11.9.1} y \sref{IV.11.9.5}, y supongamos que la familia $(u_\lambda)$ es separante; supongamos además que $\mathrm{X}$ es un $\mathrm{S}$-preesquema, donde $\mathrm{S}=\operatorname{Spec}(\mathrm{A})$ es un esquema afín, y que $\mathrm{X}'=\mathrm{X}\times_\mathrm{S}\mathrm{S}'$, donde $\mathrm{S}'=\operatorname{Spec}(\mathrm{A}')$, siendo $\mathrm{A}'$ una $\mathrm{A}$-álgebra; supongamos, por último, que los $\sh{G}_\lambda$ son planos sobre $\mathrm{S}$.
Sea $(\mathrm{A}'_\alpha)_{\alpha\in\mathrm{I}}$ la familia filtrante de las sub-$\mathrm{A}$-álgebras de tipo finito de $\mathrm{A}'$ y, para todo $\alpha\in\mathrm{I}$, sean $\mathrm{S}''_\alpha=\operatorname{Spec}(\mathrm{A}'_\alpha)$, $\mathrm{X}''_\alpha=\mathrm{X}\times_\mathrm{S}\mathrm{S}''_\alpha$, $\mathrm{Z}''_{\alpha\lambda}=\mathrm{Z}_\lambda\times_\mathrm{X}\mathrm{X}''_\alpha$, y sean $f''_{\alpha\lambda}:\mathrm{Z}''_{\alpha\lambda}\to\mathrm{X}''_\alpha$, $\mathscr{F}''_\alpha$, $\sh{G}''_{\alpha\lambda}$ y $u''_{\alpha\lambda}$ los morfismos, módulos y homomorfismos de módulos deducidos de $f_\lambda$, $\mathscr{F}$, $\sh{G}_\lambda$ y $u_\lambda$ mediante el cambio de base $\mathrm{X}''_\alpha\to\mathrm{X}$.
Entonces, si para todo $\alpha\in\mathrm{I}$ la familia $(u''_{\alpha\lambda})_{\lambda\in\mathrm{L}}$ es separante, también lo es $(u'_\lambda)_{\lambda\in\mathrm{L}}$.
\end{lemma}

\begin{proof}
Basta probar que, si $t'$ es una sección de $\mathscr{F}'$ sobre un abierto afín $\mathrm{U}'$ de $\mathrm{X}'$ tal que $u'_\lambda(t')=0$ para todo $\lambda\in\mathrm{L}$, entonces $t'=0$.
Si $h_\alpha:\mathrm{X}'\to\mathrm{X}''_\alpha$ es la proyección canónica,
\oldpage[IV-3]{166}
se deduce de \sref{IV.8.2.11} que existen un índice $\alpha$ y un abierto cuasicompacto $\mathrm{U}''_\alpha\subset\mathrm{X}''_\alpha$ tales que $\mathrm{U}'=h_\alpha^{-1}(\mathrm{U}''_\alpha)$; además, en virtud de \sref{IV.8.5.2},~(i), se puede suponer que $t'$ es la imagen canónica de una sección $t''_\alpha$ de $\mathscr{F}''_\alpha$ sobre $\mathrm{U}''_\alpha$.
Reemplazando, si es necesario, $\mathrm{S}$, $\mathrm{X}$, $\mathrm{Z}_\lambda$, $f_\lambda$, $\mathscr{F}$, $\sh{G}_\lambda$ y $u_\lambda$ por $\mathrm{S}''_\alpha$, $\mathrm{U}''_\alpha$, $(f''_{\alpha\lambda})^{-1}(\mathrm{U}''_\alpha)$ y las restricciones correspondientes de $\mathscr{F}''_\alpha$, $\sh{G}''_{\alpha\lambda}$ y $u''_{\alpha\lambda}$, podemos suponer, por tanto, que $\mathrm{U}'=\mathrm{X}'$, que $t'$ es la imagen canónica de una sección $t$ de $\mathscr{F}$ sobre $\mathrm{X}$ y que el homomorfismo $\mathrm{A}\to\mathrm{A}'$ es inyectivo; de manera equivalente, si $p:\mathrm{S}'\to\mathrm{S}$ es el morfismo correspondiente, el homomorfismo $\sh{O}_\mathrm{S}\to p_*(\sh{O}_{\mathrm{S}'})$ que interviene en la definición de $p$ es \emph{inyectivo}.
De la planitud de $\sh{G}_\lambda$ sobre $\mathrm{S}$ se deduce inmediatamente (reduciendo, por ejemplo, al caso en que $\mathrm{Z}_\lambda$ es afín sobre $\mathrm{S}$ \sref[I]{I.1.6.3} y \sref[I]{I.1.6.5}) que el homomorfismo canónico $\rho:\sh{G}_\lambda\to(g'_\lambda)_*(\sh{G}'_\lambda)$ es también \emph{inyectivo}.
Pero el homomorfismo compuesto
\[
  \Gamma(\mathrm{X},\mathscr{F})\xrightarrow{\Gamma(u_\lambda)}\Gamma(\mathrm{Z}_\lambda,\sh{G}_\lambda)\xrightarrow{\Gamma(\rho)}\Gamma(\mathrm{Z}'_\lambda,\sh{G}'_\lambda)
\]
es igual, por definición, al homomorfismo compuesto
\[
  \Gamma(\mathrm{X},\mathscr{F})\xrightarrow{\Gamma(\rho)}\Gamma(\mathrm{X}',\mathscr{F}')\xrightarrow{\Gamma(u'_\lambda)}\Gamma(\mathrm{Z}'_\lambda,\sh{G}'_\lambda)
\]
Por consiguiente, la imagen de $t$ mediante estos homomorfismos compuestos es $u'_\lambda(t')=0$; en virtud de la inyectividad del homomorfismo $\Gamma(\mathrm{Z}_\lambda,\sh{G}_\lambda)\xrightarrow{\Gamma(\rho)}\Gamma(\mathrm{Z}'_\lambda,\sh{G}'_\lambda)$, se concluye que $u_\lambda(t)=0$ para todo $\lambda\in\mathrm{L}$, de donde $t=0$ por hipótesis y, finalmente, $t'=0$.
\end{proof}

\begin{proposition}[11.9.13]
\label{IV.11.9.13}
Sean las notaciones las de \sref{IV.11.9.1} y \sref{IV.11.9.5}, y supongamos que $\mathrm{X}$ es un preesquema sobre un cuerpo $k$ y que, poniendo $\mathrm{S}=\operatorname{Spec}(k)$, se tiene $\mathrm{X}'=\mathrm{X}\times_\mathrm{S}\mathrm{S}'$, donde $\mathrm{S}'$ es un $k$-preesquema arbitrario.
Entonces, si la familia $(u_\lambda)_{\lambda\in\mathrm{L}}$ es separante, también lo es $(u'_\lambda)$.
\end{proposition}

\begin{proof}
Podemos limitarnos al caso en que $\mathrm{S}'=\operatorname{Spec}(\mathrm{A}')$ es afín.
Si $\mathrm{A}'$ es una $k$-álgebra de tipo finito, el morfismo $g:\mathrm{X}'\to\mathrm{X}$ es plano y de presentación finita; estamos, por tanto, bajo las hipótesis de \sref{IV.11.9.10},~(ii),~$b$).
En el caso general, se considera $\mathrm{A}'$ como límite inductivo de sus sub-$k$-álgebras $\mathrm{A}'_\alpha$ de tipo finito y se aplica a cada $\mathrm{A}'_\alpha$ el resultado de \sref{IV.11.9.10},~(ii),~$b$); se concluye entonces mediante el lema \sref{IV.11.9.12}, pues los $\sh{G}_\lambda$ son planos sobre $\mathrm{S}$.
\end{proof}

\begin{env}[11.9.14]
\label{IV.11.9.14}
Conservemos las notaciones de \sref{IV.11.9.1} y \sref{IV.11.9.5}, y supongamos que $\mathrm{X}$ es un $\mathrm{S}$-preesquema.
Si, para todo cambio de base $g:\mathrm{X}\times_\mathrm{S}\mathrm{S}'\to\mathrm{X}$, donde $\mathrm{S}'\to\mathrm{S}$ es un morfismo arbitrario, la familia correspondiente $(u'_\lambda)$ es separante, se dice que la familia $(u_\lambda)$ es \emph{universalmente separante con respecto a} $\mathrm{S}$.
Cuando la familia $(u_\lambda)$ se reduce a un solo elemento $u$, se dice asimismo que $u$ es \emph{universalmente inyectivo} con respecto a $\mathrm{S}$.
Está claro entonces que, para todo morfismo $h:\mathrm{S}'\to\mathrm{S}$, la familia correspondiente $(u'_\lambda)$ es universalmente separante con respecto a $\mathrm{S}'$; recíprocamente, si $h$ es \emph{fielmente plano} y si $(u'_\lambda)$ es universalmente separante con respecto a $\mathrm{S}'$, entonces $(u_\lambda)$ es universalmente separante con respecto a $\mathrm{S}$, como se deduce inmediatamente de \sref{IV.11.9.10},~(i) y del hecho de que, para todo morfismo $\mathrm{S}''\to\mathrm{S}$, el morfismo correspondiente $\mathrm{S}''\times_\mathrm{S}\mathrm{S}'\to\mathrm{S}''$ es fielmente plano.
\end{env}

\begin{proposition}[11.9.15]
\label{IV.11.9.15}
Sean las notaciones las de \sref{IV.11.9.1}, y supongamos que $\mathrm{X}$ es un $\mathrm{S}$-preesquema, siendo los $\sh{G}_\lambda$ planos sobre $\mathrm{S}$.
Sea $\mathrm{S}_0$ un subpreesquema cerrado de $\mathrm{S}$ definido por un ideal cuasicoherente nilpotente $\sh{J}$ de $\sh{O}_\mathrm{S}$, tal que los $(\sh{O}_\mathrm{S}/\sh{J})$-módulos $\sh{J}^k/\sh{J}^{k+1}$ sean todos localmente libres.
\oldpage[IV-3]{167}
Sea $(u_{\lambda 0})$ la familia de homomorfismos obtenida de $(u_\lambda)$ por el cambio de base $\mathrm{X}\leftarrow\mathrm{X}_0=\mathrm{X}\times_\mathrm{S}\mathrm{S}_0$.
Entonces, si la familia $(u_{\lambda 0})$ es separante (resp.\ universalmente separante con respecto a $\mathrm{S}_0$), la familia $(u_\lambda)$ es separante (resp.\ universalmente separante con respecto a $\mathrm{S}$).
\end{proposition}

\begin{proof}
Observemos que, si $\mathrm{S}'\to\mathrm{S}$ es un cambio de base arbitrario y $\mathrm{S}'_0=\mathrm{S}'\times_\mathrm{S}\mathrm{S}_0$, entonces $\mathrm{S}'_0$ es un subpreesquema cerrado de $\mathrm{S}'$ definido por un ideal cuasicoherente nilpotente $\sh{J}'$ de $\sh{O}_{\mathrm{S}'}$, tal que $\sh{J}'^k/\sh{J}'^{k+1}$ es un $(\sh{O}_{\mathrm{S}'}/\sh{J}')$-módulo localmente libre para todo $k$ \sref{IV.2.1.8},~(i); como, además, los $\sh{G}'_\lambda=\sh{G}_\lambda\otimes_{\sh{O}_\mathrm{S}}\sh{O}_{\mathrm{S}'}$ son planos sobre $\mathrm{S}'$, se ve que la afirmación relativa a las familias universalmente separantes es consecuencia de la relativa a las familias separantes.
Para probar esta última, se puede reducir, mediante \sref{IV.11.9.3}, al caso en que $\mathrm{S}=\operatorname{Spec}(\mathrm{A})$ y $\mathrm{X}=\operatorname{Spec}(\mathrm{B})$ son afines, $\mathrm{Z}_\lambda=\mathrm{X}$ para todo $\lambda$, $\mathscr{F}=\widetilde{\mathrm{M}}$ y $\sh{G}_\lambda=\widetilde{\mathrm{N}}_\lambda$, donde $\mathrm{M}$ y los $\mathrm{N}_\lambda$ son $\mathrm{B}$-módulos, siendo los $\mathrm{N}_\lambda$ planos como $\mathrm{A}$-módulos.
Además, como la cuestión es local sobre $\mathrm{X}$ y $\mathrm{S}$, basta probar que, si $t\in\mathrm{M}$ es tal que $u_\lambda(t)=0$ para todo $\lambda$, entonces $t=0$.
Se tiene $\sh{J}=\widetilde{\mathfrak{J}}$, donde $\mathfrak{J}$ es un ideal nilpotente de $\mathrm{A}$ tal que los $\mathfrak{J}^k/\mathfrak{J}^{k+1}$ son $(\mathrm{A}/\mathfrak{J})$-módulos libres y, por hipótesis, los homomorfismos $u_{\lambda 0}:\mathrm{M}/\mathfrak{J}\mathrm{M}\to\mathrm{N}_\lambda/\mathfrak{J}\mathrm{N}_\lambda$ forman una familia separante.
Supongamos que $\mathfrak{J}^{n+1}=0$ ($n$ un entero $\geq 0$) y razonemos por inducción sobre $n$, siendo la afirmación trivial para $n=0$.
Si $\bar{t}$ es la clase de $t$ en $\mathrm{M}/\mathfrak{J}^n\mathrm{M}$, la clase de $u_\lambda(t)$ en $\mathrm{N}_\lambda/\mathfrak{J}^n\mathrm{N}_\lambda$ es nula para todo $\lambda$; de la hipótesis de inducción se deduce que $\bar{t}=0$, es decir, que $t\in\mathfrak{J}^n\mathrm{M}$.
Ahora bien, $\mathfrak{J}^n=\mathfrak{J}^n/\mathfrak{J}^{n+1}$ es un $(\mathrm{A}/\mathfrak{J})$-módulo libre; si $(e_\alpha)$ es una base de este módulo, se puede escribir, por tanto, $t=\sum_\alpha e_\alpha t_\alpha^0$, con $t_\alpha^0\in\mathrm{M}/\mathfrak{J}\mathrm{M}$, siendo nulos todos salvo para un número finito de índices.
Por otra parte, como $\mathrm{N}_\lambda$ es un $\mathrm{A}$-módulo plano, $\mathfrak{J}^n\otimes_{(\mathrm{A}/\mathfrak{J})}(\mathrm{N}_\lambda/\mathfrak{J}\mathrm{N}_\lambda)$ se identifica con $\mathfrak{J}^n\mathrm{N}_\lambda$ y se puede escribir, por tanto, $u_\lambda(t)=\sum_\alpha e_\alpha u_{\lambda 0}(t_\alpha^0)=\sum_\alpha e_\alpha\otimes u_{\lambda 0}(t_\alpha^0)$.
Como, por hipótesis, $u_\lambda(t)=0$, se deduce que $u_{\lambda 0}(t_\alpha^0)=0$ para todo $\alpha$ y todo $\lambda$; de donde $t_\alpha^0=0$ para todo $\alpha$ y, por tanto, $t=0$.
\end{proof}

\begin{theorem}[11.9.16]
\label{IV.11.9.16}
Sean las notaciones las de \sref{IV.11.9.1}, y supongamos que $\mathrm{X}$ es un $\mathrm{S}$-preesquema localmente noetheriano, que $\mathscr{F}$ es un $\sh{O}_\mathrm{X}$-módulo coherente y que los $\sh{G}_\lambda$ son planos sobre $\mathrm{S}$.
Para todo $s\in\mathrm{S}$, sea $((u_\lambda)_s)_{\lambda\in\mathrm{L}}$ la familia obtenida de $(u_\lambda)$ por el cambio de base $\mathrm{X}_s=\mathrm{X}\times_\mathrm{S}\operatorname{Spec}(k(s))\to\mathrm{X}$.
Para que la familia $(u_\lambda)$ sea universalmente separante con respecto a $\mathrm{S}$, es necesario y suficiente que, para todo $s\in\mathrm{S}$, la familia $((u_\lambda)_s)$ sea separante.
\end{theorem}

\begin{proof}
La necesidad de la condición se deduce inmediatamente de las definiciones.
Recíprocamente, supongamos satisfecha la condición y probemos primero que la familia $(u_\lambda)$ es separante.
Mediante \sref{IV.11.9.3}, podemos reducirnos al caso en que $\mathrm{S}=\operatorname{Spec}(\mathrm{A})$ y $\mathrm{X}=\operatorname{Spec}(\mathrm{B})$ son afines, $\mathrm{Z}_\lambda=\mathrm{X}$ para todo $\lambda$, $\mathscr{F}=\widetilde{\mathrm{M}}$ y $\sh{G}_\lambda=\widetilde{\mathrm{N}}_\lambda$, donde $\mathrm{B}$ es noetheriano, $\mathrm{M}$ es un $\mathrm{B}$-módulo de tipo finito y los $\mathrm{N}_\lambda$ son $\mathrm{A}$-módulos planos, y limitarnos a probar que, si $t\in\mathrm{M}$ es tal que $u_\lambda(t)=0$ para todo $\lambda$, entonces $t=0$.
Para demostrar que $t=0$, basta probar que la imagen $t_\mathfrak{p}$ de $t$ en $\mathrm{M}_\mathfrak{p}$ es nula para todo ideal maximal $\mathfrak{p}$ de $\mathrm{B}$ (Bourbaki, \emph{Alg.\ comm.}, cap.~II, \S~3, n.º~3, cor.~1 del teor.~1).
Podemos, pues, limitarnos a probar que la intersección
\oldpage[IV-3]{168}
de los núcleos de los $(u_\lambda)_\mathfrak{p}:\mathrm{M}_\mathfrak{p}\to(\mathrm{N}_\lambda)_\mathfrak{p}$ deducidos de $(u_\lambda)$ por el cambio de base $\operatorname{Spec}(\mathrm{B}_\mathfrak{p})\to\operatorname{Spec}(\mathrm{B})$ se reduce a $0$.
Dicho de otro modo, nos reducimos al caso en que $\mathrm{B}$ es un anillo local noetheriano; considerando además el ideal primo de $\mathrm{A}$ que es la imagen inversa del ideal maximal de $\mathrm{B}$, podemos suponer también que $\mathrm{A}$ es un anillo local con ideal maximal $\mathfrak{m}$.
Entonces, como $\mathfrak{m}\mathrm{B}$ está contenido en el ideal maximal de $\mathrm{B}$ y $\mathrm{M}$ es un $\mathrm{B}$-módulo de tipo finito, la intersección de los $\mathfrak{m}^{n}\mathrm{M}$ se reduce a $0$ \sref[0]{0.7.3.5}; por tanto, basta probar que, para todo $n$, la imagen de $t$ en $\mathrm{M}/\mathfrak{m}^{n+1}\mathrm{M}$ es nula.
Basta, pues, demostrar que la familia deducida de $(u_\lambda)$ por el cambio de base $\operatorname{Spec}(\mathrm{B}/\mathfrak{m}^{n+1}\mathrm{B})\to\operatorname{Spec}(\mathrm{B})$ es separante, lo que significa que podemos limitarnos al caso en que $\mathrm{A}$ es un anillo local cuyo ideal maximal $\mathfrak{m}$ es nilpotente.
Pero entonces los $\mathfrak{m}^k/\mathfrak{m}^{k+1}$ son $(\mathrm{A}/\mathfrak{m})$-módulos libres y, en virtud de la hipótesis sobre los $(u_\lambda)_s$, estamos precisamente en las condiciones de aplicación de \sref{IV.11.9.15}, de donde se sigue la conclusión anunciada.

Sea ahora $h:\mathrm{S}'\to\mathrm{S}$ un morfismo de cambio de base, y sea $(u'_\lambda)$ la familia obtenida de $(u_\lambda)$ por el cambio de base $h':\mathrm{X}\times_\mathrm{S}\mathrm{S}'\to\mathrm{X}$; probemos que $(u'_\lambda)$ también es separante.
Supongamos primero que $h$ es \emph{localmente de tipo finito}; entonces lo mismo ocurre con $h'$ y, por consiguiente, $\mathrm{X}'=\mathrm{X}\times_\mathrm{S}\mathrm{S}'$ es localmente noetheriano. Además, si $s'\in\mathrm{S}'$ está sobre $s\in\mathrm{S}$, se deduce de \sref{IV.11.9.13}, aplicado a $\mathrm{X}_s$ y a $\mathrm{X}_{s'}=\mathrm{X}_s\otimes_{k(s)}k(s')$, que la familia $((u'_\lambda)_{s'})$ es separante para todo $s'\in\mathrm{S}'$; podemos concluir, por tanto, de la primera parte de la demostración que $(u'_\lambda)$ es separante en este caso.

Por último, si $h$ es arbitrario, podemos limitarnos evidentemente al caso en que $\mathrm{S}=\operatorname{Spec}(\mathrm{A})$ y $\mathrm{S}'=\operatorname{Spec}(\mathrm{A}')$ son afines, y considerar $\mathrm{A}'$ como límite inductivo de sus sub-$\mathrm{A}$-álgebras de tipo finito.
Como los $\sh{G}_\lambda$ son planos sobre $\mathrm{S}$, basta aplicar lo anterior y el lema \sref{IV.11.9.12} para concluir la demostración.
\end{proof}

\begin{proposition}[11.9.17]
\label{IV.11.9.17}
Sea $f:\mathrm{X}\to\mathrm{S}$ un morfismo localmente de presentación finita, sea $\mathscr{F}$ un $\sh{O}_\mathrm{X}$-módulo cuasicoherente de presentación finita y $f$-plano, sea $\mathrm{U}$ un abierto de $\mathrm{X}$, sea $j:\mathrm{U}\to\mathrm{X}$ la inyección canónica y sea $u:\mathscr{F}\to j_*(j^*(\mathscr{F}))$ el homomorfismo canónico \sref[0]{0.4.4.3.2}.
Para todo $s\in\mathrm{S}$, sean $\mathrm{X}_s$ la fibra $\mathrm{X}\times_\mathrm{S}\operatorname{Spec}(k(s))$, $\mathrm{U}_s$ el abierto $\mathrm{U}\cap\mathrm{X}_s$ de $\mathrm{X}_s$, $\mathscr{F}_s=\mathscr{F}\otimes_{\sh{O}_\mathrm{S}}k(s)$, $j_s:\mathrm{U}_s\to\mathrm{X}_s$ la inyección canónica y $u_s:\mathscr{F}_s\to(j_s)_*((j_s)^*(\mathscr{F}_s))$ el homomorfismo canónico.
Para que $u$ sea universalmente inyectivo con respecto a $\mathrm{S}$, es necesario y suficiente que $u_s$ sea inyectivo para todo $s\in\mathrm{S}$.
\end{proposition}

\begin{proof}
Solo queda probar la suficiencia de la condición.
Cuando $\mathrm{X}$ es localmente noetheriano, la proposición es un corolario inmediato de \sref{IV.11.9.16}.
Nos reduciremos ahora a este caso en dos etapas, limitándonos, como podemos hacer, al caso en que $\mathrm{S}=\operatorname{Spec}(\mathrm{A})$ y $\mathrm{X}$ son afines.
\medskip

\emph{A) Caso en que $\mathrm{U}$ es cuasicompacto.} --- Utilizaremos el lema siguiente:

\begin{lemma}[11.9.17.1]
\label{IV.11.9.17.1}
Bajo las hipótesis generales de \sref{IV.11.9.17}, y suponiendo además que $\mathrm{S}$ y $\mathrm{X}$ son afines y que $\mathrm{U}$ es cuasicompacto, el conjunto $\mathrm{E}$ de los $s\in\mathrm{S}$ tales que $u_s$ es inyectivo es constructible.
\end{lemma}

\begin{proof}
En efecto, las fibras $\mathrm{X}_s$ son preesquemas localmente noetherianos; de donde $\mathrm{E}$ también puede, en virtud de \sref{IV.5.10.2}, definirse como el conjunto de los $s\in\mathrm{S}$ tales que $\operatorname{Ass}(\mathscr{F}_s)\subset\mathrm{U}_s$.
Observemos además que $\mathrm{U}$, por ser cuasicompacto en un esquema afín, es constructible. Así, la verificación de la condición \sref{IV.9.2.1},~(i) se deduce inmediatamente de \sref{IV.4.2.7}; por otra parte, la
\oldpage[IV-3]{169}
verificación de \sref{IV.9.2.1},~(ii) se efectúa fácilmente utilizando el estudio de los ciclos primos asociados en un entorno de la fibra genérica \sref{IV.9.8.3}, así como \sref{IV.9.5.2} y \sref{IV.9.5.3}.
\end{proof}

Establecido este lema, podemos suponer, en virtud de \sref{IV.8.9.1} y \sref{IV.8.2.11}, que existen un subanillo noetheriano $\mathrm{A}_0$ de $\mathrm{A}$, un preesquema $\mathrm{X}_0$ de tipo finito sobre $\mathrm{S}_0=\operatorname{Spec}(\mathrm{A}_0)$, un abierto $\mathrm{U}_0$ de $\mathrm{X}_0$ y un $\sh{O}_{\mathrm{X}_0}$-módulo coherente $\mathscr{F}_0$, tales que $\mathrm{X}=\mathrm{X}_0\times_{\mathrm{S}_0}\mathrm{S}$ y, si $p_0:\mathrm{X}\to\mathrm{X}_0$ es la proyección canónica, $\mathrm{U}=p_0^{-1}(\mathrm{U}_0)$ y $\mathscr{F}=p_0^*(\mathscr{F}_0)$.
Sea $(\mathrm{A}_\alpha)$ la familia filtrante de los subanillos de $\mathrm{A}$ que son $\mathrm{A}_0$-álgebras de tipo finito, y pongamos $\mathrm{S}_\alpha=\operatorname{Spec}(\mathrm{A}_\alpha)$, $\mathrm{X}_\alpha=\mathrm{X}_0\times_{\mathrm{S}_0}\mathrm{S}_\alpha$, $\mathrm{U}_\alpha=\mathrm{U}_0\times_{\mathrm{S}_0}\mathrm{S}_\alpha$ y $\mathscr{F}_\alpha=\mathscr{F}_0\otimes_{\mathrm{S}_0}\mathrm{S}_\alpha$; sea $u_\alpha$ el homomorfismo canónico relativo a $\mathscr{F}_\alpha$ y $\mathrm{U}_\alpha$, definido como en \sref{IV.11.9.17}.
Para todo $s\in\mathrm{S}$, la hipótesis de que $u_s$ es inyectivo implica que $(u_\alpha)_{s_\alpha}$ es inyectivo \sref{IV.11.9.10},~(i), donde $s_\alpha=q_\alpha(s)$ es la imagen de $s$ por el morfismo $q_\alpha:\mathrm{S}\to\mathrm{S}_\alpha$.
Si $\mathrm{E}_\alpha$ es el conjunto de los $s_\alpha\in\mathrm{S}_\alpha$ tales que $(u_\alpha)_{s_\alpha}$ es inyectivo, se tiene, por tanto, $\mathrm{S}=q_\alpha^{-1}(\mathrm{E}_\alpha)$, y los $\mathrm{E}_\alpha$ forman un sistema proyectivo de conjuntos.
Pero el lema \sref{IV.11.9.17.1}, aplicado a $u_\alpha$, muestra que $\mathrm{E}_\alpha$ es constructible en $\mathrm{S}_\alpha$; se deduce entonces de \sref{IV.8.3.4} que existe un índice $\alpha$ tal que $\mathrm{E}_\alpha=\mathrm{S}_\alpha$.
Pero entonces, como $\mathrm{X}_\alpha$ es noetheriano, $u_\alpha$ es universalmente inyectivo en virtud de \sref{IV.11.9.16}; de donde $u$ también es universalmente inyectivo.
\medskip

\emph{B) Caso general.} --- El abierto $\mathrm{U}$ es una unión filtrante creciente de abiertos cuasicompactos $\mathrm{U}_\lambda$; si $j_\lambda:\mathrm{U}_\lambda\to\mathrm{X}$ es la inyección canónica y $u_\lambda:\mathscr{F}\to(j_\lambda)_*(j_\lambda^*(\mathscr{F}))$ el homomorfismo canónico correspondiente, se deduce del lema \sref{IV.11.9.17.1}, aplicado a $u_\lambda$, que el conjunto $\mathrm{E}_\lambda$ de los $s\in\mathrm{S}$ tales que $(u_\lambda)_s$ es inyectivo es constructible.
Por otra parte, para todo $s\in\mathrm{S}$, decir que $u_s$ (resp.\ $(u_\lambda)_s$) es inyectivo significa que $\operatorname{Ass}(\mathscr{F}_s)\subset\mathrm{U}\cap\mathrm{X}_s$ (resp.\ $\operatorname{Ass}(\mathscr{F}_s)\subset\mathrm{U}_\lambda\cap\mathrm{X}_s$) \sref{IV.5.10.2}.
Como $\operatorname{Ass}(\mathscr{F}_s)$ es finito \sref{IV.3.1.6}, decir que $u_s$ es inyectivo significa, por tanto, que existe $\lambda$ tal que $s\in\mathrm{E}_\lambda$.
La hipótesis significa, pues, que $\mathrm{S}=\bigcup_\lambda\mathrm{E}_\lambda$.
En virtud de \sref{IV.1.9.9}, existe un índice $\lambda$ tal que $\mathrm{S}=\mathrm{E}_\lambda$; de la primera parte del razonamiento se deduce entonces que $u_\lambda$ es universalmente inyectivo.
Se sigue de la factorización de $u_\lambda$:
\[
  \mathscr{F}\xrightarrow{u}\,j_*(\mathscr{F}|\mathrm{U})\to(j_\lambda)_*(\mathscr{F}|\mathrm{U}_\lambda)
\]
que $u$ también es universalmente inyectivo.
\end{proof}

\begin{remarks}[11.9.18]
\label{IV.11.9.18}
Sea $\mathrm{X}$ un preesquema, sean $\mathscr{F}$ y $\sh{G}$ dos $\sh{O}_\mathrm{X}$-módulos cuasicoherentes de presentación finita, suponiéndose además que $\sh{G}$ es \emph{localmente libre}, y sea $u:\mathscr{F}\to\sh{G}$ un homomorfismo.
Las condiciones siguientes son equivalentes:
\begin{enumerate}[label=\emph{\alph*)}]
  \item para todo morfismo $g:\mathrm{X}'\to\mathrm{X}$, el homomorfismo $g^*(u):g^*(\mathscr{F})\to g^*(\sh{G})$ es inyectivo;
  \item para todo $x\in\mathrm{X}$, el homomorfismo $u\otimes 1_{k(x)}:\mathscr{F}\otimes_{\sh{O}_{\mathrm{X}}}k(x)\to\sh{G}\otimes_{\sh{O}_\mathrm{X}}k(x)$ es inyectivo;
  \item para todo $x\in\mathrm{X}$, existe un entorno $\mathrm{U}$ de $x$ tal que $u|\mathrm{U}:\mathscr{F}|\mathrm{U}\to\sh{G}|\mathrm{U}$ es un isomorfismo de $\mathscr{F}|\mathrm{U}$ sobre un factor directo de $\sh{G}|\mathrm{U}$.
\end{enumerate}

Está claro que $c)$ implica $a)$ y que $a)$ implica $b)$.
El hecho de que $b)$ implica $c)$ se deduce de \sref[0]{0.19.1.12}, \sref[0]{0.5.2.5} y \sref[0]{0.5.5.5}.

Cuando se cumplen las condiciones equivalentes anteriores, se dice que $u$ es \emph{universalmente inyectivo}.
\end{remarks}

\oldpage[IV-3]{170}
\subsection{Familias esquemáticamente dominantes de morfismos y familias esquemáticamente densas de subpreesquemas}
\label{subsection:IV.11.10}

\begin{proposition}[11.10.1]
\label{IV.11.10.1}
Sea $\mathrm{X}$ un preesquema, $(Z_\lambda)_{\lambda\in\mathrm{L}}$ una familia de preesquemas y, para cada $\lambda\in\mathrm{L}$, sea $f_\lambda=(\psi_\lambda,\theta_\lambda):\mathrm{Z}_\lambda\to\mathrm{X}$ un morfismo.
Las condiciones siguientes son equivalentes:
\begin{enumerate}[label=\emph{\alph*)}]
  \item La familia de homomorfismos $\theta_\lambda:\sh{O}_\mathrm{X}\to(f_\lambda)_*(\sh{O}_{\mathrm{Z}_\lambda})$ es separante (en otras palabras \sref{IV.11.9.1}, la intersección de los núcleos de los $\theta_\lambda$ es \emph{cero}).
  \item Para todo abierto $\mathrm{U}$ de $\mathrm{X}$, toda sección $t$ de $\sh{O}_\mathrm{X}$ sobre $\mathrm{U}$ cuyas imágenes por todos los homomorfismos canónicos
\end{enumerate}
\[
  (\theta_\lambda)_\mathrm{U}:\Gamma(\mathrm{U},\sh{O}_\mathrm{X})\to\Gamma(f_\lambda^{-1}(\mathrm{U}),\sh{O}_{\mathrm{Z}_\lambda})
\label{IV.11.10.1.1}
\]
son cero es ella misma cero.
\begin{enumerate}[label=\emph{\alph*)}]
\setcounter{enumi}{2}
  \item Para todo abierto $\mathrm{U}$ de $\mathrm{X}$ y todo subpreesquema cerrado $\mathrm{Y}$ de $\mathrm{U}$ tal que, para cada $\lambda\in\mathrm{L}$, existe una factorización
\end{enumerate}
\[
  f_\lambda^{-1}(\mathrm{U})\xrightarrow{g_\lambda}\mathrm{Y}\xrightarrow{j}\mathrm{U}
\label{IV.11.10.1.2}
\]
de la restricción $f_\lambda^{-1}(\mathrm{U})\to\mathrm{U}$ de $f_\lambda$ (donde $j$ es la inmersión canónica), se tiene $\mathrm{Y}=\mathrm{U}$.

\emph{Si además $\mathrm{X}$ es un $\mathrm{S}$-preesquema, estas condiciones también son equivalentes a las siguientes:}
\begin{enumerate}[label=\emph{\alph*)}]
\setcounter{enumi}{3}
  \item Para todo morfismo separado $g:\mathrm{X}'\to\mathrm{S}$ y todo par $u$, $v$ de $\mathrm{S}$-morfismos de un abierto $\mathrm{U}$ de $\mathrm{X}$ a $\mathrm{X}'$, tales que, para cada $\lambda$, los compuestos de $u$ y $v$ con la restricción de $f_\lambda$ a $f_\lambda^{-1}(\mathrm{U})$ sean iguales, se tiene $u=v$.
\end{enumerate}
\end{proposition}

\begin{proof}
La equivalencia de $a)$ y $b)$ se deduce de las definiciones.
Para ver que $b)$ implica $c)$, basta considerar el ideal cuasicoherente $\sh{J}$ de $\sh{O}_\mathrm{U}$ que define $\mathrm{Y}$ y observar que, para todo abierto $\mathrm{V}\subset\mathrm{U}$, la hipótesis implica que el homomorfismo $(\theta_\lambda)_\mathrm{V}$ se factoriza como
\[
  \Gamma(\mathrm{V},\sh{O}_\mathrm{U})\to\Gamma(\mathrm{V},\sh{O}_\mathrm{U}/\sh{J})\to\Gamma(f_\lambda^{-1}(\mathrm{V}),\sh{O}_{\mathrm{Z}_\lambda}).
\]
Se concluye que toda sección $t$ de $\sh{J}$ sobre $\mathrm{V}$ tiene imagen $0$ en todos los $\Gamma(f_\lambda^{-1}(\mathrm{V}),\sh{O}_{\mathrm{Z}_\lambda})$, de donde, en virtud de $b)$, $\sh{J}=0$ y $\mathrm{Y}=\mathrm{U}$.
Recíprocamente, si se satisface $c)$, para demostrar $b)$ basta aplicar $c)$ al subpreesquema cerrado $\mathrm{Y}$ de $\mathrm{U}$ definido por el ideal $t\sh{O}_\mathrm{U}$: la hipótesis de que las imágenes de $t$ por los $(\theta_\lambda)_\mathrm{U}$ son todas cero implica que se tienen las factorizaciones \sref{IV.11.10.1.2} para cada $\lambda$ \sref[I]{I.4.1.9}.
Para demostrar que $c)$ implica $d)$, basta aplicar $c)$ al subpreesquema cerrado $\mathrm{Y}$ de $\mathrm{U}$ que es la imagen inversa, por $(u,v)_\mathrm{S}$, de la diagonal de $\mathrm{X}'$ en $\mathrm{X}'\times_\mathrm{S}\mathrm{X}'$, y utilizar \sref[I]{I.4.4.1}.
Recíprocamente, se deduce $b)$ de $d)$ considerando el $\mathrm{S}$-esquema $\mathrm{X}'=\mathrm{S}[\mathrm{T}]=\mathrm{S}\otimes_\mathbf{Z}\mathbf{Z}[\mathrm{T}]$ ($\mathrm{T}$ indeterminada) y observando que las secciones de $\sh{O}_\mathrm{U}$ sobre $\mathrm{U}$ corresponden biyectivamente a los $\mathrm{S}$-morfismos $\mathrm{U}\to\mathrm{S}[\mathrm{T}]$ \sref[I]{I.3.3.15}; decir que dos secciones de $\sh{O}_\mathrm{U}$ sobre $\mathrm{U}$ tienen las mismas imágenes por todos los $(\theta_\lambda)_\mathrm{U}$ equivale a decir que los compuestos de los dos morfismos correspondientes con todos los morfismos $f_\lambda^{-1}(\mathrm{U})\to\mathrm{U}$ son iguales.
\end{proof}

\oldpage[IV-3]{171}
\begin{definition}[11.10.2]
\label{IV.11.10.2}
Cuando se satisfacen las condiciones equivalentes de \sref{IV.11.10.1}, se dice que la familia $(f_\lambda)$ es \emph{esquemáticamente dominante}.
Cuando los $f_\lambda$ son las inmersiones canónicas en $\mathrm{X}$ de una familia $(\mathrm{Z}_\lambda)$ de subpreesquemas de $\mathrm{X}$, se dice también que la familia $(\mathrm{Z}_\lambda)$ es \emph{esquemáticamente densa}.
\end{definition}

\begin{remarks}[11.10.3]
\label{IV.11.10.3}
\medskip\noindent
\begin{enumerate}[label=(\roman*)]
  \item La noción de familia esquemáticamente dominante sobre $\mathrm{X}$, como resulta, por ejemplo, de una condición de la forma $b)$ de \sref{IV.11.10.1}, es \emph{local} sobre $\mathrm{X}$: si $(\mathrm{W}_\alpha)$ es un recubrimiento abierto de $\mathrm{X}$, la familia $(f_\lambda)$ es esquemáticamente dominante si y solo si lo es cada una de las familias formadas por los morfismos $f_\lambda^{-1}(\mathrm{W}_\alpha)\to\mathrm{W}_\alpha$, restricciones de los $f_\lambda$.
  \item Si $\mathrm{Z}$ es la suma de los preesquemas $\mathrm{Z}_\lambda$ y $f:\mathrm{Z}\to\mathrm{X}$ el morfismo que coincide con $f_\lambda$ sobre cada $\mathrm{Z}_\lambda$, equivale a decir que la familia $(f_\lambda)$ es esquemáticamente dominante o que lo es la familia reducida al único elemento $f$.
  \item Sea $\mathrm{M}$ un segundo conjunto de índices y, para cada $\lambda\in\mathrm{L}$, sea $(g_{\lambda\mu})_{\mu\in\mathrm{M}}$ una familia de morfismos $g_{\lambda\mu}:\mathrm{Z}'_{\lambda\mu}\to\mathrm{Z}_\lambda$; si, para cada $\lambda\in\mathrm{L}$, la familia $(g_{\lambda\mu})_{\mu\in\mathrm{M}}$ es esquemáticamente dominante, entonces equivale a decir que la familia $(f_\lambda)_{\lambda\in\mathrm{L}}$ es esquemáticamente dominante o que lo es la familia $(f_\lambda\circ g_{\lambda\mu})_{(\lambda,\mu)\in\mathrm{L}\times\mathrm{M}}$ \sref{IV.11.9.2}.
  \item Sea $f:\mathrm{Z}\to\mathrm{X}$ un morfismo tal que $f_*(\sh{O}_\mathrm{Z})$ sea un $\sh{O}_\mathrm{X}$-módulo cuasicoherente (por ejemplo, un morfismo cuasicompacto y cuasiseparado \sref{IV.1.7.4}).
  Entonces, decir que $f$ es esquemáticamente dominante significa que la \emph{imagen cerrada} de $\mathrm{Z}$ \emph{por} $f$ \sref[I]{I.9.5.3} \emph{es idéntica a} $\mathrm{X}$.
\end{enumerate}
\end{remarks}

\begin{proposition}[11.10.4]
\label{IV.11.10.4}
Si la familia de morfismos $f_\lambda:\mathrm{Z}_\lambda\to\mathrm{X}$ es esquemáticamente dominante, la unión de los $f_\lambda(\mathrm{Z}_\lambda)$ es densa en $\mathrm{X}$.
Por el contrario, si esta unión es densa en $\mathrm{X}$ y si además $\mathrm{X}$ es \emph{reducido}, la familia $(f_\lambda)$ es esquemáticamente dominante.
\end{proposition}

\begin{proof}
La primera afirmación se deduce inmediatamente de \sref{IV.11.10.1},~$b$)).
Por otra parte, si $\mathrm{X}$ es reducido y la unión de los $f_\lambda(\mathrm{Z}_\lambda)$ es densa en $\mathrm{X}$, entonces, si se tienen las factorizaciones \sref{IV.11.10.1.2} para cada $\lambda\in\mathrm{L}$, se tienen también las factorizaciones $(f_\lambda^{-1}(\mathrm{U}))_\mathrm{red}\xrightarrow{(g_\lambda)_\mathrm{red}}\mathrm{Y}_\mathrm{red}\xrightarrow{j_\mathrm{red}}\mathrm{U}$, y la hipótesis implica que el espacio subyacente de $\mathrm{Y}$ es idéntico a $\mathrm{U}$; de donde $\mathrm{Y}=\mathrm{Y}_\mathrm{red}=\mathrm{U}$, pues $\mathrm{U}$ es reducido.
\end{proof}

Los resultados de {\hyperlink{ega4.c04.subsec.11-9}{\textsection~11.9}} sobre las familias separantes se traducen en resultados sobre las familias esquemáticamente dominantes:

\begin{theorem}[11.10.5]
\label{IV.11.10.5}
Sea $(f_\lambda)$ una familia de morfismos $f_\lambda:\mathrm{Z}_\lambda\to\mathrm{X}$, sea $g:\mathrm{X}'\to\mathrm{X}$ un morfismo, y pongamos $\mathrm{Z}'_\lambda=\mathrm{Z}_\lambda\times_\mathrm{X}\mathrm{X}'$, $f'_\lambda=(f_\lambda)_{(\mathrm{X}')}:\mathrm{Z}'_\lambda\to\mathrm{X}'$.
\begin{enumerate}[label=(\roman*)]
  \item Si $g$ es fielmente plano y la familia $(f'_\lambda)$ es esquemáticamente dominante, entonces la familia $(f_\lambda)$ es esquemáticamente dominante.
  \item Supongamos que $g$ es un morfismo plano y que, además, se satisface una de las dos condiciones siguientes:
  \begin{enumerate}[label=\emph{\alph*)}]
    \item $\mathrm{L}$ es finito y los $f_\lambda$ son cuasicompactos.
    \item El morfismo $g$ es localmente de presentación finita.
  \end{enumerate}
  Entonces, si la familia $(f_\lambda)$ es esquemáticamente dominante, lo mismo ocurre con la familia $(f'_\lambda)$.
\end{enumerate}
\end{theorem}

\oldpage[IV-3]{172}
\begin{proposition}[11.10.6]
\label{IV.11.10.6}
Sean las notaciones las de \sref{IV.11.10.5}, y supongamos que $\mathrm{X}$ es un preesquema sobre un cuerpo $k$ y que, poniendo $\mathrm{S}=\operatorname{Spec}(k)$, se tiene $\mathrm{X}'=\mathrm{X}\times_\mathrm{S}\mathrm{S}'$, donde $\mathrm{S}'$ es un $k$-preesquema arbitrario.
Entonces, si la familia $(f_\lambda)$ es esquemáticamente dominante, lo mismo ocurre con $(f'_\lambda)$.
\end{proposition}

\begin{corollary}[11.10.7]
\label{IV.11.10.7}
Sea $\mathrm{X}$ un preesquema sobre un cuerpo $k$, sea $(Z_\lambda)_{\lambda\in\mathrm{L}}$ una familia de $k$-preesquemas geométricamente reducidos sobre $k$ y, para cada $\lambda\in\mathrm{L}$, sea $f_\lambda:\mathrm{Z}_\lambda\to\mathrm{X}$ un $k$-morfismo.
Sea $\mathrm{Y}$ el subpreesquema reducido de $\mathrm{X}$ cuyo espacio subyacente es la adherencia de $\bigcup_{\lambda\in\mathrm{L}}f_\lambda(\mathrm{Z}_\lambda)$.
Sea $k'$ una extensión de $k$, y sean $\mathrm{X}'$, $\mathrm{Z}'_\lambda$, $f'_\lambda:\mathrm{Z}'_\lambda\to\mathrm{X}'$ los preesquemas y morfismos deducidos de $\mathrm{X}$, $\mathrm{Z}_\lambda$ y $f_\lambda$ por extensión de la base a $k'$.
Entonces, si $\mathrm{Y}'$ es el subpreesquema reducido de $\mathrm{X}'$ cuyo espacio subyacente es la adherencia de $\bigcup_{\lambda\in\mathrm{L}}f'_\lambda(\mathrm{Z}'_\lambda)$, se tiene $\mathrm{Y}'=\mathrm{Y}\otimes_k k'$.
En particular, $\mathrm{Y}$ es geométricamente reducido sobre $k$.
\end{corollary}

\begin{proof}
Como los $\mathrm{Z}_\lambda$ son reducidos, los morfismos $f_\lambda$ se factorizan como $\mathrm{Z}_\lambda\xrightarrow{g_\lambda}\mathrm{Y}\xrightarrow{j}\mathrm{X}$, donde $j$ es la inmersión canónica \sref[I]{I.5.2.2}.
Se sigue entonces de \sref{IV.11.10.4} que $(g_\lambda)$ es una familia esquemáticamente dominante.
Sea $\mathrm{Y}'_1=\mathrm{Y}\otimes_k k'$, un subpreesquema cerrado de $\mathrm{X}'$, y sea $g'_\lambda$ el morfismo deducido de $g_\lambda$ por extensión de la base a $k'$, de modo que $f'_\lambda$ se factoriza como $\mathrm{Z}'_\lambda\xrightarrow{g'_\lambda}\mathrm{Y}'_1\xrightarrow{j'}\mathrm{X}'$, donde $j'$ es la inmersión canónica.
Se deduce de \sref{IV.11.10.6} que la familia $(g'_\lambda)$ es esquemáticamente dominante.
Pero, por hipótesis, los $\mathrm{Z}'_\lambda$ son \emph{reducidos}, de donde, por \sref[I]{I.5.2.2}, los $g'_\lambda$ se factorizan como $\mathrm{Z}'_\lambda\xrightarrow{h'_\lambda}\mathrm{Y}'\to\mathrm{Y}'_1$; se concluye, por tanto, de \sref{IV.11.10.2} que $\mathrm{Y}'=\mathrm{Y}'_1$ y $h'_\lambda=g'_\lambda$, lo que demuestra el corolario.
\end{proof}

\begin{definition}[11.10.8]
\label{IV.11.10.8}
Sean las notaciones las de \sref{IV.11.10.5}, y supongamos que $\mathrm{X}$ es un $\mathrm{S}$-preesquema.
Se dice que la familia $(f_\lambda)$ es \emph{universalmente esquemáticamente dominante} (relativamente a $\mathrm{S}$) si, para todo cambio de base $\mathrm{S}'\to\mathrm{S}$, la familia $(f'_\lambda)$ correspondiente al cambio de base $\mathrm{X}'=\mathrm{X}\times_\mathrm{S}\mathrm{S}'\to\mathrm{X}$ es esquemáticamente dominante.
\end{definition}

Cuando $\mathrm{S}$ es el espectro de un cuerpo, \sref{IV.11.10.6} muestra que toda familia esquemáticamente dominante es universalmente esquemáticamente dominante (relativamente a $\mathrm{S}$).

Cuando los $f_\lambda$ son inmersiones canónicas $\mathrm{Z}_\lambda\to\mathrm{X}$, se dice también que la familia de subpreesquemas $\mathrm{Z}_\lambda$ es \emph{universalmente esquemáticamente densa en} $\mathrm{X}$ (relativamente a $\mathrm{S}$), en vez de decir que la familia $(f_\lambda)$ es universalmente esquemáticamente dominante (relativamente a $\mathrm{S}$).

\begin{theorem}[11.10.9]
\label{IV.11.10.9}
Sean las notaciones las de \sref{IV.11.10.5}, y supongamos que $\mathrm{X}$ es un $\mathrm{S}$-preesquema localmente noetheriano y que todos los $\mathrm{Z}_\lambda$ son $\mathrm{S}$-planos.
Para cada $s\in\mathrm{S}$, pongamos $\mathrm{X}_s=\mathrm{X}\times_\mathrm{S}\operatorname{Spec}(k(s))$, $(\mathrm{Z}_\lambda)_s=\mathrm{Z}_\lambda\times_\mathrm{S}\operatorname{Spec}(k(s))$, $(f_\lambda)_s=f_\lambda\times 1:(\mathrm{Z}_\lambda)_s\to\mathrm{X}_s$.
Para que la familia $(f_\lambda)$ sea universalmente esquemáticamente dominante relativamente a $\mathrm{S}$, es necesario y suficiente que, para todo $s\in\mathrm{S}$, la familia $((f_\lambda)_s)$ sea esquemáticamente dominante.
\end{theorem}

\begin{proposition}[11.10.10]
\label{IV.11.10.10}
Sea $f:\mathrm{X}\to\mathrm{S}$ un morfismo plano localmente de presentación finita, y sea $\mathrm{U}$ un abierto de $\mathrm{X}$.
Para que $\mathrm{U}$ sea universalmente esquemáticamente denso en $\mathrm{X}$ relativamente a $\mathrm{S}$, es necesario y suficiente que, para todo $s\in\mathrm{S}$, $\mathrm{U}_s=\mathrm{U}\cap\mathrm{X}_s$ sea esquemáticamente denso en $\mathrm{X}_s$ (o, lo que equivale a lo mismo, que se tenga $\operatorname{Ass}(\sh{O}_{\mathrm{X}_s})\subset\mathrm{U}_s$).
\end{proposition}
